% SPDX-FileCopyrightText: 2026 Bart van de Lint
% SPDX-License-Identifier: AGPL-3.0-or-later
% The pictures every schematic figure shares: the paper's 140 m ship from the side, from
% ahead and from above, drawn from the model's hull in figure_ship.tex, and its rigs. They
% are drawn in metres, the origin at midship on the waterline (from the side and from
% above, bow to the right) or on the centreline at the waterline (from ahead), so that a
% figure scales them with `scale` [mm per m] and `transform shape`. The icons that stand
% for a concept come from the figure style.
% Written by the ship_pictures step of scripts/paper.jl.
\def\shipfoot{6.59}
\def\shipbow{73.50}
\def\hullface#1{\pgfsetfillcolor{#1}\pgfsetstrokecolor{#1}}
\def\hullmove#1#2{\pgfpathmoveto{\pgfqpointxy{#1}{#2}}}
\def\hullline#1#2{\pgfpathlineto{\pgfqpointxy{#1}{#2}}}
\def\hullfill{\pgfpathclose\pgfusepath{fill,stroke}}
\tikzset{
  ship side/.pic={\begin{scope}[hull face]\hullface{keel!57!black}\hullmove{67.73}{-7.45}\hullline{68.56}{-7.32}\hullline{68.56}{-7.32}\hullline{67.73}{-7.36}\hullfill
\hullface{topsides!78!black}\hullmove{72.68}{5.04}\hullline{73.50}{6.59}\hullline{73.50}{6.79}\hullline{72.68}{5.55}\hullfill
\hullface{keel!83!white}\hullmove{72.68}{-4.61}\hullline{73.50}{-4.54}\hullline{73.50}{-4.06}\hullline{72.68}{-4.48}\hullfill
\hullface{keel!57!black}\hullmove{72.68}{-6.39}\hullline{73.50}{-5.44}\hullline{73.50}{-5.44}\hullline{72.68}{-6.26}\hullfill
\hullface{keel!95!white}\hullmove{70.20}{-3.82}\hullline{71.03}{-3.90}\hullline{71.03}{-3.73}\hullline{70.20}{-3.74}\hullfill
\hullface{keel!57!black}\hullmove{66.91}{-8.15}\hullline{67.73}{-7.45}\hullline{67.73}{-7.36}\hullline{66.91}{-7.89}\hullfill
\hullface{keel!98!white}\hullmove{71.85}{-4.31}\hullline{72.68}{-4.48}\hullline{72.68}{-4.13}\hullline{71.85}{-4.13}\hullfill
\hullface{keel!57!black}\hullmove{66.08}{-8.65}\hullline{66.91}{-8.15}\hullline{66.91}{-7.89}\hullline{66.08}{-8.38}\hullfill
\hullface{topsides!66!black}\hullmove{72.68}{5.55}\hullline{73.50}{6.79}\hullline{73.50}{6.99}\hullline{72.68}{6.06}\hullfill
\hullface{keel!57!black}\hullmove{65.26}{-8.95}\hullline{66.08}{-8.65}\hullline{66.08}{-8.38}\hullline{65.26}{-8.67}\hullfill
\hullface{keel!57!black}\hullmove{64.44}{-9.06}\hullline{65.26}{-8.95}\hullline{65.26}{-8.67}\hullline{64.44}{-8.77}\hullfill
\hullface{keel!57!black}\hullmove{-74.90}{0.00}\hullline{-73.33}{-0.47}\hullline{-73.33}{-0.45}\hullline{-74.90}{0.00}\hullfill
\hullface{keel!57!black}\hullmove{63.61}{-9.06}\hullline{64.44}{-9.06}\hullline{64.44}{-8.77}\hullline{63.61}{-8.77}\hullfill
\hullface{keel!57!black}\hullmove{71.85}{-6.74}\hullline{72.68}{-6.39}\hullline{72.68}{-6.26}\hullline{71.85}{-6.56}\hullfill
\hullface{keel!90!black}\hullmove{72.68}{-4.77}\hullline{73.50}{-4.95}\hullline{73.50}{-4.54}\hullline{72.68}{-4.61}\hullfill
\hullface{keel!57!black}\hullmove{62.79}{-9.06}\hullline{63.61}{-9.06}\hullline{63.61}{-8.77}\hullline{62.79}{-8.77}\hullfill
\hullface{keel!57!black}\hullmove{72.68}{-6.19}\hullline{73.50}{-5.44}\hullline{73.50}{-5.44}\hullline{72.68}{-5.91}\hullfill
\hullface{keel!57!black}\hullmove{71.03}{-6.97}\hullline{71.85}{-6.74}\hullline{71.85}{-6.56}\hullline{71.03}{-6.76}\hullfill
\hullface{keel!93!black}\hullmove{71.03}{-4.11}\hullline{71.85}{-4.31}\hullline{71.85}{-4.13}\hullline{71.03}{-3.90}\hullfill
\hullface{topsides!57!black}\hullmove{72.68}{6.06}\hullline{73.50}{6.99}\hullline{73.50}{7.19}\hullline{72.68}{6.57}\hullfill
\hullface{keel!57!black}\hullmove{61.97}{-9.06}\hullline{62.79}{-9.06}\hullline{62.79}{-8.77}\hullline{61.97}{-8.77}\hullfill
\hullface{keel!81!black}\hullmove{72.68}{-4.96}\hullline{73.50}{-5.26}\hullline{73.50}{-4.95}\hullline{72.68}{-4.77}\hullfill
\hullface{keel!75!black}\hullmove{69.38}{-3.62}\hullline{70.20}{-3.74}\hullline{70.20}{-2.95}\hullline{69.38}{-3.00}\hullfill
\hullface{keel!57!black}\hullmove{70.20}{-7.13}\hullline{71.03}{-6.97}\hullline{71.03}{-6.76}\hullline{70.20}{-6.91}\hullfill
\hullface{keel!75!black}\hullmove{69.38}{-3.00}\hullline{70.20}{-2.95}\hullline{70.20}{-2.03}\hullline{70.20}{-1.04}\hullline{69.38}{-1.06}\hullline{69.38}{-2.06}\hullfill
\hullface{keel!75!black}\hullmove{69.38}{-1.06}\hullline{70.20}{-1.04}\hullline{70.20}{0.00}\hullline{70.20}{0.38}\hullline{69.38}{0.00}\hullline{69.38}{0.00}\hullfill
\hullface{keel!60!black}\hullmove{72.68}{-5.91}\hullline{73.50}{-5.44}\hullline{73.50}{-5.44}\hullline{72.68}{-5.82}\hullfill
\hullface{keel!57!black}\hullmove{69.38}{-7.25}\hullline{70.20}{-7.13}\hullline{70.20}{-6.91}\hullline{69.38}{-7.02}\hullfill
\hullface{keel!89!white}\hullmove{-74.90}{0.00}\hullline{-73.33}{-0.45}\hullline{-73.33}{-0.42}\hullline{-73.33}{-0.39}\hullline{-73.33}{-0.35}\hullline{-73.33}{-0.30}\hullline{-73.33}{-0.25}\hullline{-73.33}{-0.19}\hullline{-73.33}{-0.13}\hullline{-73.33}{-0.07}\hullline{-73.33}{-0.00}\hullline{-73.33}{0.00}\hullline{-74.90}{0.87}\hullline{-74.90}{0.00}\hullline{-74.90}{0.00}\hullline{-74.90}{0.00}\hullline{-74.90}{0.00}\hullline{-74.90}{0.00}\hullline{-74.90}{0.00}\hullline{-74.90}{0.00}\hullline{-74.90}{0.00}\hullline{-74.90}{0.00}\hullfill
\hullface{keel!57!black}\hullmove{61.14}{-9.06}\hullline{61.97}{-9.06}\hullline{61.97}{-8.77}\hullline{61.14}{-8.77}\hullfill
\hullface{keel!78!black}\hullmove{72.68}{-5.34}\hullline{73.50}{-5.44}\hullline{73.50}{-5.26}\hullline{72.68}{-4.96}\hullfill
\hullface{keel!57!black}\hullmove{68.56}{-7.32}\hullline{69.38}{-7.25}\hullline{69.38}{-7.02}\hullline{68.56}{-7.09}\hullfill
\hullface{keel!57!black}\hullmove{67.73}{-7.36}\hullline{68.56}{-7.32}\hullline{68.56}{-7.09}\hullline{67.73}{-7.22}\hullfill
\hullface{keel!57!black}\hullmove{-63.94}{-8.72}\hullline{-62.38}{-9.06}\hullline{-62.38}{-8.77}\hullline{-63.94}{-8.45}\hullfill
\hullface{keel!60!black}\hullmove{-67.07}{-4.83}\hullline{-65.51}{-7.07}\hullline{-65.51}{-6.84}\hullline{-67.07}{-4.83}\hullfill
\hullface{keel!69!black}\hullmove{72.68}{-5.82}\hullline{73.50}{-5.44}\hullline{73.50}{-5.44}\hullline{72.68}{-5.44}\hullfill
\hullface{keel!90!black}\hullmove{71.85}{-4.35}\hullline{72.68}{-4.61}\hullline{72.68}{-4.48}\hullline{71.85}{-4.31}\hullfill
\hullface{keel!72!black}\hullmove{72.68}{-5.44}\hullline{73.50}{-5.44}\hullline{73.50}{-5.44}\hullline{72.68}{-5.34}\hullfill
\hullface{keel!57!black}\hullmove{66.91}{-7.89}\hullline{67.73}{-7.36}\hullline{67.73}{-7.22}\hullline{66.91}{-7.42}\hullfill
\hullface{keel!57!black}\hullmove{60.32}{-9.06}\hullline{61.14}{-9.06}\hullline{61.14}{-8.77}\hullline{60.32}{-8.77}\hullfill
\hullface{keel!57!black}\hullmove{-62.38}{-9.06}\hullline{-60.81}{-9.06}\hullline{-60.81}{-8.77}\hullline{-62.38}{-8.77}\hullfill
\hullface{keel!96!black}\hullmove{70.20}{-3.97}\hullline{71.03}{-4.11}\hullline{71.03}{-3.90}\hullline{70.20}{-3.82}\hullfill
\hullface{keel!96!black}\hullmove{69.38}{-3.87}\hullline{70.20}{-3.82}\hullline{70.20}{-3.74}\hullline{69.38}{-3.62}\hullfill
\hullface{keel!63!black}\hullmove{66.08}{-8.38}\hullline{66.91}{-7.89}\hullline{66.91}{-7.42}\hullline{66.08}{-7.88}\hullfill
\hullface{topsides!57!black}\hullmove{72.68}{6.57}\hullline{73.50}{7.19}\hullline{73.50}{7.39}\hullline{72.68}{7.08}\hullfill
\hullface{keel!57!black}\hullmove{59.49}{-9.06}\hullline{60.32}{-9.06}\hullline{60.32}{-8.77}\hullline{59.49}{-8.77}\hullfill
\hullface{keel!57!black}\hullmove{71.85}{-6.56}\hullline{72.68}{-6.26}\hullline{72.68}{-6.19}\hullline{71.85}{-6.52}\hullfill
\hullface{keel!63!black}\hullmove{65.26}{-8.67}\hullline{66.08}{-8.38}\hullline{66.08}{-7.88}\hullline{65.26}{-8.15}\hullfill
\hullface{keel!57!black}\hullmove{-60.81}{-9.06}\hullline{-59.25}{-9.06}\hullline{-59.25}{-8.77}\hullline{-60.81}{-8.77}\hullfill
\hullface{keel!57!black}\hullmove{-65.51}{-7.07}\hullline{-63.94}{-8.72}\hullline{-63.94}{-8.45}\hullline{-65.51}{-6.84}\hullfill
\hullface{keel!57!black}\hullmove{58.67}{-9.06}\hullline{59.49}{-9.06}\hullline{59.49}{-8.77}\hullline{58.67}{-8.77}\hullfill
\hullface{keel!99!black}\hullmove{-72.10}{-8.61}\hullline{-67.90}{-8.61}\hullline{-67.90}{-2.26}\hullline{-72.10}{-2.26}\hullfill
\hullface{keel!63!black}\hullmove{64.44}{-8.77}\hullline{65.26}{-8.67}\hullline{65.26}{-8.15}\hullline{64.44}{-8.25}\hullfill
\hullface{keel!57!black}\hullmove{71.03}{-6.76}\hullline{71.85}{-6.56}\hullline{71.85}{-6.52}\hullline{71.03}{-6.75}\hullfill
\hullface{keel!99!black}\hullmove{71.85}{-4.78}\hullline{72.68}{-4.77}\hullline{72.68}{-4.61}\hullline{71.85}{-4.35}\hullfill
\hullface{keel!57!black}\hullmove{-59.25}{-9.06}\hullline{-57.68}{-9.06}\hullline{-57.68}{-8.77}\hullline{-59.25}{-8.77}\hullfill
\hullface{keel!57!black}\hullmove{57.85}{-9.06}\hullline{58.67}{-9.06}\hullline{58.67}{-8.77}\hullline{57.85}{-8.77}\hullfill
\hullface{keel!95!white}\hullmove{69.38}{-3.88}\hullline{70.20}{-3.97}\hullline{70.20}{-3.82}\hullline{69.38}{-3.87}\hullfill
\hullface{keel!57!black}\hullmove{-73.33}{-0.47}\hullline{-71.77}{-1.31}\hullline{-71.77}{-1.27}\hullline{-73.33}{-0.45}\hullfill
\hullface{keel!57!black}\hullmove{66.91}{-7.42}\hullline{67.73}{-7.22}\hullline{67.73}{-7.10}\hullline{66.91}{-7.36}\hullfill
\hullface{keel!63!black}\hullmove{71.85}{-6.52}\hullline{72.68}{-6.19}\hullline{72.68}{-5.91}\hullline{71.85}{-6.13}\hullfill
\hullface{keel!57!black}\hullmove{70.20}{-6.91}\hullline{71.03}{-6.76}\hullline{71.03}{-6.75}\hullline{70.20}{-6.90}\hullfill
\hullface{topsides!57!black}\hullmove{72.68}{7.08}\hullline{73.50}{7.39}\hullline{73.50}{7.59}\hullline{72.68}{7.59}\hullfill
\hullface{keel!60!black}\hullmove{63.61}{-8.77}\hullline{64.44}{-8.77}\hullline{64.44}{-8.25}\hullline{63.61}{-8.25}\hullfill
\hullface{keel!92!white}\hullmove{71.03}{-4.50}\hullline{71.85}{-4.35}\hullline{71.85}{-4.31}\hullline{71.03}{-4.11}\hullfill
\hullface{keel!57!black}\hullmove{69.38}{-7.02}\hullline{70.20}{-6.91}\hullline{70.20}{-6.90}\hullline{69.38}{-7.00}\hullfill
\hullface{keel!57!black}\hullmove{67.73}{-7.22}\hullline{68.56}{-7.09}\hullline{68.56}{-7.07}\hullline{67.73}{-7.10}\hullfill
\hullface{keel!66!black}\hullmove{66.08}{-7.88}\hullline{66.91}{-7.42}\hullline{66.91}{-7.36}\hullline{66.08}{-7.33}\hullfill
\hullface{topsides!57!black}\hullmove{71.85}{3.49}\hullline{72.68}{5.04}\hullline{72.68}{5.04}\hullline{71.85}{3.90}\hullfill
\hullface{keel!57!black}\hullmove{57.02}{-9.06}\hullline{57.85}{-9.06}\hullline{57.85}{-8.77}\hullline{57.02}{-8.77}\hullfill
\hullface{keel!57!black}\hullmove{68.56}{-7.09}\hullline{69.38}{-7.02}\hullline{69.38}{-7.00}\hullline{68.56}{-7.07}\hullfill
\hullface{keel!57!black}\hullmove{-57.68}{-9.06}\hullline{-56.12}{-9.06}\hullline{-56.12}{-8.77}\hullline{-57.68}{-8.77}\hullfill
\hullface{keel!63!black}\hullmove{71.03}{-6.75}\hullline{71.85}{-6.52}\hullline{71.85}{-6.13}\hullline{71.03}{-6.34}\hullfill
\hullface{topsides!72!black}\hullmove{-74.90}{0.87}\hullline{-73.33}{0.00}\hullline{-73.33}{0.78}\hullline{-74.90}{1.48}\hullfill
\hullface{keel!99!black}\hullmove{71.85}{-5.03}\hullline{72.68}{-4.96}\hullline{72.68}{-4.77}\hullline{71.85}{-4.78}\hullfill
\hullface{keel!75!black}\hullmove{71.85}{-6.13}\hullline{72.68}{-5.91}\hullline{72.68}{-5.82}\hullline{71.85}{-6.09}\hullfill
\hullface{keel!57!black}\hullmove{62.79}{-8.77}\hullline{63.61}{-8.77}\hullline{63.61}{-8.25}\hullline{62.79}{-8.25}\hullfill
\hullface{keel!69!black}\hullmove{66.91}{-7.36}\hullline{67.73}{-7.10}\hullline{67.73}{-6.79}\hullline{66.91}{-7.10}\hullfill
\hullface{keel!57!black}\hullmove{56.20}{-9.06}\hullline{57.02}{-9.06}\hullline{57.02}{-8.77}\hullline{56.20}{-8.77}\hullfill
\hullface{topsides!57!black}\hullmove{70.20}{0.38}\hullline{71.03}{1.94}\hullline{71.03}{3.07}\hullline{70.20}{1.83}\hullfill
\hullface{keel!63!black}\hullmove{70.20}{-6.90}\hullline{71.03}{-6.75}\hullline{71.03}{-6.34}\hullline{70.20}{-6.49}\hullfill
\hullface{keel!81!black}\hullmove{71.85}{-6.09}\hullline{72.68}{-5.82}\hullline{72.68}{-5.44}\hullline{71.85}{-5.63}\hullfill
\hullface{keel!69!black}\hullmove{66.08}{-7.33}\hullline{66.91}{-7.36}\hullline{66.91}{-7.10}\hullline{66.08}{-7.23}\hullfill
\hullface{keel!86!white}\hullmove{70.20}{-4.59}\hullline{71.03}{-4.50}\hullline{71.03}{-4.11}\hullline{70.20}{-3.97}\hullfill
\hullface{keel!96!black}\hullmove{71.85}{-5.44}\hullline{72.68}{-5.34}\hullline{72.68}{-4.96}\hullline{71.85}{-5.03}\hullfill
\hullface{keel!95!white}\hullmove{71.03}{-4.67}\hullline{71.85}{-4.78}\hullline{71.85}{-4.35}\hullline{71.03}{-4.50}\hullfill
\hullface{keel!69!black}\hullmove{65.26}{-8.15}\hullline{66.08}{-7.88}\hullline{66.08}{-7.33}\hullline{65.26}{-7.48}\hullfill
\hullface{keel!57!black}\hullmove{56.00}{-9.06}\hullline{56.20}{-9.06}\hullline{56.20}{-8.77}\hullline{56.00}{-8.77}\hullfill
\hullface{keel!63!black}\hullmove{69.38}{-7.00}\hullline{70.20}{-6.90}\hullline{70.20}{-6.49}\hullline{69.38}{-6.60}\hullfill
\hullface{keel!57!black}\hullmove{56.00}{-9.06}\hullline{56.00}{-9.06}\hullline{56.00}{-8.77}\hullline{56.00}{-8.77}\hullfill
\hullface{keel!90!black}\hullmove{71.85}{-5.63}\hullline{72.68}{-5.44}\hullline{72.68}{-5.34}\hullline{71.85}{-5.44}\hullfill
\hullface{keel!57!black}\hullmove{-56.12}{-9.06}\hullline{-54.55}{-9.06}\hullline{-54.55}{-8.77}\hullline{-56.12}{-8.77}\hullfill
\hullface{keel!84!black}\hullmove{68.56}{-3.80}\hullline{69.38}{-3.87}\hullline{69.38}{-3.62}\hullline{68.56}{-3.55}\hullfill
\hullface{keel!75!black}\hullmove{68.56}{-3.55}\hullline{69.38}{-3.62}\hullline{69.38}{-3.00}\hullline{68.56}{-3.03}\hullfill
\hullface{keel!63!black}\hullmove{68.56}{-7.07}\hullline{69.38}{-7.00}\hullline{69.38}{-6.60}\hullline{68.56}{-6.67}\hullfill
\hullface{keel!57!black}\hullmove{61.97}{-8.77}\hullline{62.79}{-8.77}\hullline{62.79}{-8.25}\hullline{61.97}{-8.25}\hullfill
\hullface{keel!63!black}\hullmove{67.73}{-7.10}\hullline{68.56}{-7.07}\hullline{68.56}{-6.67}\hullline{67.73}{-6.79}\hullfill
\hullface{keel!57!black}\hullmove{55.37}{-9.06}\hullline{56.00}{-9.06}\hullline{56.00}{-8.77}\hullline{55.37}{-8.77}\hullfill
\hullface{keel!84!black}\hullmove{68.56}{-3.92}\hullline{69.38}{-3.88}\hullline{69.38}{-3.87}\hullline{68.56}{-3.80}\hullfill
\hullface{keel!75!black}\hullmove{71.03}{-6.34}\hullline{71.85}{-6.13}\hullline{71.85}{-6.09}\hullline{71.03}{-6.20}\hullfill
\hullface{keel!86!white}\hullmove{69.38}{-4.53}\hullline{70.20}{-4.59}\hullline{70.20}{-3.97}\hullline{69.38}{-3.88}\hullfill
\hullface{keel!69!black}\hullmove{-67.07}{-4.83}\hullline{-65.51}{-6.84}\hullline{-65.51}{-6.44}\hullline{-67.07}{-4.83}\hullfill
\hullface{keel!72!black}\hullmove{68.56}{-3.03}\hullline{69.38}{-3.00}\hullline{69.38}{-2.06}\hullline{68.56}{-2.08}\hullfill
\hullface{keel!60!black}\hullmove{66.08}{-7.23}\hullline{66.91}{-7.10}\hullline{66.91}{-6.81}\hullline{66.08}{-7.08}\hullfill
\hullface{keel!66!black}\hullmove{64.44}{-8.25}\hullline{65.26}{-8.15}\hullline{65.26}{-7.48}\hullline{64.44}{-7.57}\hullfill
\hullface{keel!98!white}\hullmove{71.03}{-5.21}\hullline{71.85}{-5.03}\hullline{71.85}{-4.78}\hullline{71.03}{-4.67}\hullfill
\hullface{keel!81!black}\hullmove{71.03}{-6.20}\hullline{71.85}{-6.09}\hullline{71.85}{-5.63}\hullline{71.03}{-5.82}\hullfill
\hullface{topsides!57!black}\hullmove{69.38}{0.00}\hullline{70.20}{0.38}\hullline{70.20}{1.83}\hullline{69.38}{1.52}\hullfill
\hullface{keel!69!black}\hullmove{68.56}{-2.08}\hullline{69.38}{-2.06}\hullline{69.38}{-1.06}\hullline{68.56}{-1.07}\hullfill
\hullface{keel!57!black}\hullmove{54.55}{-9.06}\hullline{55.37}{-9.06}\hullline{55.37}{-8.77}\hullline{54.55}{-8.77}\hullfill
\hullface{keel!75!black}\hullmove{70.20}{-6.49}\hullline{71.03}{-6.34}\hullline{71.03}{-6.20}\hullline{70.20}{-6.28}\hullfill
\hullface{keel!69!black}\hullmove{66.91}{-7.10}\hullline{67.73}{-6.79}\hullline{67.73}{-6.40}\hullline{66.91}{-6.81}\hullfill
\hullface{keel!69!black}\hullmove{68.56}{-1.07}\hullline{69.38}{-1.06}\hullline{69.38}{0.00}\hullline{68.56}{0.00}\hullfill
\hullface{keel!57!black}\hullmove{-54.55}{-9.06}\hullline{-52.99}{-9.06}\hullline{-52.99}{-8.77}\hullline{-54.55}{-8.77}\hullfill
\hullface{keel!57!black}\hullmove{61.14}{-8.77}\hullline{61.97}{-8.77}\hullline{61.97}{-8.25}\hullline{61.14}{-8.25}\hullfill
\hullface{keel!89!white}\hullmove{70.20}{-4.61}\hullline{71.03}{-4.67}\hullline{71.03}{-4.50}\hullline{70.20}{-4.59}\hullfill
\hullface{keel!95!white}\hullmove{68.56}{-4.49}\hullline{69.38}{-4.53}\hullline{69.38}{-3.88}\hullline{68.56}{-3.92}\hullfill
\hullface{keel!96!black}\hullmove{71.03}{-5.44}\hullline{71.85}{-5.44}\hullline{71.85}{-5.03}\hullline{71.03}{-5.21}\hullfill
\hullface{keel!90!black}\hullmove{71.03}{-5.82}\hullline{71.85}{-5.63}\hullline{71.85}{-5.44}\hullline{71.03}{-5.44}\hullfill
\hullface{keel!57!black}\hullmove{-63.94}{-8.45}\hullline{-62.38}{-8.77}\hullline{-62.38}{-8.25}\hullline{-63.94}{-7.94}\hullfill
\hullface{keel!75!black}\hullmove{69.38}{-6.60}\hullline{70.20}{-6.49}\hullline{70.20}{-6.28}\hullline{69.38}{-6.34}\hullfill
\hullface{topsides!57!black}\hullmove{71.03}{1.94}\hullline{71.85}{3.49}\hullline{71.85}{3.90}\hullline{71.03}{3.07}\hullfill
\hullface{keel!84!black}\hullmove{70.20}{-6.28}\hullline{71.03}{-6.20}\hullline{71.03}{-5.82}\hullline{70.20}{-5.96}\hullfill
\hullface{keel!69!black}\hullmove{65.26}{-7.48}\hullline{66.08}{-7.33}\hullline{66.08}{-7.23}\hullline{66.08}{-7.08}\hullline{65.26}{-7.03}\hullline{65.26}{-7.27}\hullfill
\hullface{keel!57!black}\hullmove{53.73}{-9.06}\hullline{54.55}{-9.06}\hullline{54.55}{-8.77}\hullline{53.73}{-8.77}\hullfill
\hullface{keel!75!black}\hullmove{68.56}{-6.67}\hullline{69.38}{-6.60}\hullline{69.38}{-6.34}\hullline{68.56}{-6.38}\hullfill
\hullface{keel!95!white}\hullmove{70.20}{-5.33}\hullline{71.03}{-5.21}\hullline{71.03}{-4.67}\hullline{70.20}{-4.61}\hullfill
\hullface{keel!57!black}\hullmove{-71.77}{-1.31}\hullline{-70.20}{-2.15}\hullline{-70.20}{-2.09}\hullline{-71.77}{-1.27}\hullfill
\hullface{keel!63!black}\hullmove{63.61}{-8.25}\hullline{64.44}{-8.25}\hullline{64.44}{-7.57}\hullline{63.61}{-7.57}\hullfill
\hullface{keel!75!black}\hullmove{67.73}{-6.79}\hullline{68.56}{-6.67}\hullline{68.56}{-6.38}\hullline{67.73}{-6.40}\hullfill
\hullface{keel!89!white}\hullmove{69.38}{-4.68}\hullline{70.20}{-4.61}\hullline{70.20}{-4.59}\hullline{69.38}{-4.53}\hullfill
\hullface{keel!72!black}\hullmove{66.08}{-7.08}\hullline{66.91}{-6.81}\hullline{66.91}{-6.40}\hullline{66.08}{-6.46}\hullfill
\hullface{keel!57!black}\hullmove{60.32}{-8.77}\hullline{61.14}{-8.77}\hullline{61.14}{-8.25}\hullline{60.32}{-8.25}\hullfill
\hullface{keel!93!black}\hullmove{70.20}{-5.96}\hullline{71.03}{-5.82}\hullline{71.03}{-5.44}\hullline{70.20}{-5.44}\hullfill
\hullface{keel!84!black}\hullmove{69.38}{-6.34}\hullline{70.20}{-6.28}\hullline{70.20}{-5.96}\hullline{69.38}{-6.06}\hullfill
\hullface{keel!57!black}\hullmove{-52.99}{-9.06}\hullline{-51.42}{-9.06}\hullline{-51.42}{-8.77}\hullline{-52.99}{-8.77}\hullfill
\hullface{keel!57!black}\hullmove{-62.38}{-8.77}\hullline{-60.81}{-8.77}\hullline{-60.81}{-8.25}\hullline{-62.38}{-8.25}\hullfill
\hullface{keel!89!white}\hullmove{-73.33}{-0.45}\hullline{-71.77}{-1.27}\hullline{-71.77}{-1.19}\hullline{-71.77}{-1.10}\hullline{-71.77}{-0.98}\hullline{-71.77}{-0.85}\hullline{-71.77}{-0.70}\hullline{-71.77}{-0.54}\hullline{-71.77}{-0.37}\hullline{-71.77}{-0.19}\hullline{-71.77}{-0.00}\hullline{-73.33}{-0.00}\hullline{-73.33}{-0.07}\hullline{-73.33}{-0.13}\hullline{-73.33}{-0.19}\hullline{-73.33}{-0.25}\hullline{-73.33}{-0.30}\hullline{-73.33}{-0.35}\hullline{-73.33}{-0.39}\hullline{-73.33}{-0.42}\hullfill
\hullface{keel!78!black}\hullmove{66.91}{-6.81}\hullline{67.73}{-6.40}\hullline{67.73}{-6.23}\hullline{66.91}{-6.40}\hullfill
\hullface{keel!96!black}\hullmove{70.20}{-5.44}\hullline{71.03}{-5.44}\hullline{71.03}{-5.21}\hullline{70.20}{-5.33}\hullfill
\hullface{keel!86!white}\hullmove{68.56}{-4.73}\hullline{69.38}{-4.68}\hullline{69.38}{-4.53}\hullline{68.56}{-4.49}\hullfill
\hullface{keel!95!white}\hullmove{69.38}{-5.41}\hullline{70.20}{-5.33}\hullline{70.20}{-4.61}\hullline{69.38}{-4.68}\hullfill
\hullface{keel!57!black}\hullmove{52.90}{-9.06}\hullline{53.73}{-9.06}\hullline{53.73}{-8.77}\hullline{52.90}{-8.77}\hullfill
\hullface{keel!84!black}\hullmove{68.56}{-6.38}\hullline{69.38}{-6.34}\hullline{69.38}{-6.06}\hullline{68.56}{-6.12}\hullfill
\hullface{keel!84!black}\hullmove{67.73}{-6.40}\hullline{68.56}{-6.38}\hullline{68.56}{-6.12}\hullline{67.73}{-6.23}\hullfill
\hullface{keel!93!black}\hullmove{69.38}{-6.06}\hullline{70.20}{-5.96}\hullline{70.20}{-5.44}\hullline{69.38}{-5.44}\hullfill
\hullface{keel!66!black}\hullmove{64.44}{-7.57}\hullline{65.26}{-7.48}\hullline{65.26}{-7.27}\hullline{64.44}{-7.17}\hullfill
\hullface{keel!87!black}\hullmove{67.73}{-4.47}\hullline{68.56}{-4.49}\hullline{68.56}{-3.92}\hullline{67.73}{-3.99}\hullfill
\hullface{keel!75!black}\hullmove{-67.07}{-4.83}\hullline{-65.51}{-6.44}\hullline{-65.51}{-5.91}\hullline{-67.07}{-4.83}\hullfill
\hullface{keel!99!black}\hullmove{69.38}{-5.44}\hullline{70.20}{-5.44}\hullline{70.20}{-5.33}\hullline{69.38}{-5.41}\hullfill
\hullface{keel!75!black}\hullmove{67.73}{-3.99}\hullline{68.56}{-3.92}\hullline{68.56}{-3.80}\hullline{67.73}{-3.77}\hullfill
\hullface{keel!99!black}\hullmove{67.73}{-4.82}\hullline{68.56}{-4.73}\hullline{68.56}{-4.49}\hullline{67.73}{-4.47}\hullfill
\hullface{keel!95!white}\hullmove{68.56}{-5.44}\hullline{69.38}{-5.41}\hullline{69.38}{-4.68}\hullline{68.56}{-4.73}\hullfill
\hullface{keel!60!black}\hullmove{62.79}{-8.25}\hullline{63.61}{-8.25}\hullline{63.61}{-7.57}\hullline{62.79}{-7.57}\hullfill
\hullface{keel!90!black}\hullmove{66.91}{-6.40}\hullline{67.73}{-6.23}\hullline{67.73}{-5.57}\hullline{66.91}{-6.09}\hullfill
\hullface{keel!93!black}\hullmove{68.56}{-6.12}\hullline{69.38}{-6.06}\hullline{69.38}{-5.44}\hullline{68.56}{-5.47}\hullfill
\hullface{keel!57!black}\hullmove{59.49}{-8.77}\hullline{60.32}{-8.77}\hullline{60.32}{-8.25}\hullline{59.49}{-8.25}\hullfill
\hullface{keel!81!black}\hullmove{66.08}{-6.46}\hullline{66.91}{-6.40}\hullline{66.91}{-6.09}\hullline{66.08}{-6.38}\hullfill
\hullface{keel!75!black}\hullmove{67.73}{-3.77}\hullline{68.56}{-3.80}\hullline{68.56}{-3.55}\hullline{67.73}{-3.51}\hullfill
\hullface{keel!93!black}\hullmove{67.73}{-6.23}\hullline{68.56}{-6.12}\hullline{68.56}{-5.47}\hullline{67.73}{-5.57}\hullfill
\hullface{topsides!57!black}\hullmove{71.85}{3.90}\hullline{72.68}{5.04}\hullline{72.68}{5.55}\hullline{71.85}{4.31}\hullfill
\hullface{keel!95!white}\hullmove{67.73}{-5.44}\hullline{68.56}{-5.44}\hullline{68.56}{-4.73}\hullline{67.73}{-4.82}\hullfill
\hullface{keel!57!black}\hullmove{52.08}{-9.06}\hullline{52.90}{-9.06}\hullline{52.90}{-8.77}\hullline{52.08}{-8.77}\hullfill
\hullface{keel!57!black}\hullmove{-51.42}{-9.06}\hullline{-49.85}{-9.06}\hullline{-49.85}{-8.77}\hullline{-51.42}{-8.77}\hullfill
\hullface{keel!69!black}\hullmove{65.26}{-7.03}\hullline{66.08}{-7.08}\hullline{66.08}{-6.46}\hullline{66.08}{-6.38}\hullline{65.26}{-6.35}\hullline{65.26}{-6.69}\hullfill
\hullface{keel!63!black}\hullmove{-65.51}{-6.84}\hullline{-63.94}{-8.45}\hullline{-63.94}{-7.94}\hullline{-65.51}{-6.44}\hullfill
\hullface{keel!99!black}\hullmove{67.73}{-5.57}\hullline{68.56}{-5.47}\hullline{68.56}{-5.44}\hullline{67.73}{-5.44}\hullfill
\hullface{keel!93!black}\hullmove{66.91}{-6.09}\hullline{67.73}{-5.57}\hullline{67.73}{-5.44}\hullline{66.91}{-5.44}\hullfill
\hullface{keel!75!black}\hullmove{67.73}{-3.51}\hullline{68.56}{-3.55}\hullline{68.56}{-3.03}\hullline{67.73}{-3.09}\hullfill
\hullface{keel!57!black}\hullmove{-60.81}{-8.77}\hullline{-59.25}{-8.77}\hullline{-59.25}{-8.25}\hullline{-60.81}{-8.25}\hullfill
\hullface{keel!96!black}\hullmove{66.91}{-5.44}\hullline{67.73}{-5.44}\hullline{67.73}{-4.82}\hullline{66.91}{-5.27}\hullfill
\hullface{keel!69!black}\hullmove{64.44}{-7.17}\hullline{65.26}{-7.27}\hullline{65.26}{-7.03}\hullline{64.44}{-6.94}\hullfill
\hullface{topsides!57!black}\hullmove{71.85}{4.31}\hullline{72.68}{5.55}\hullline{72.68}{6.06}\hullline{71.85}{5.13}\hullfill
\hullface{keel!78!black}\hullmove{66.08}{-6.38}\hullline{66.91}{-6.09}\hullline{66.91}{-5.44}\hullline{66.08}{-5.59}\hullfill
\hullface{keel!57!black}\hullmove{58.67}{-8.77}\hullline{59.49}{-8.77}\hullline{59.49}{-8.25}\hullline{58.67}{-8.25}\hullfill
\hullface{keel!78!black}\hullmove{66.91}{-5.27}\hullline{67.73}{-4.82}\hullline{67.73}{-4.47}\hullline{66.91}{-4.47}\hullfill
\hullface{keel!57!black}\hullmove{61.97}{-8.25}\hullline{62.79}{-8.25}\hullline{62.79}{-7.57}\hullline{61.97}{-7.57}\hullfill
\hullface{keel!72!black}\hullmove{67.73}{-3.09}\hullline{68.56}{-3.03}\hullline{68.56}{-2.08}\hullline{67.73}{-2.12}\hullfill
\hullface{keel!57!black}\hullmove{51.25}{-9.06}\hullline{52.08}{-9.06}\hullline{52.08}{-8.77}\hullline{51.25}{-8.77}\hullfill
\hullface{keel!66!black}\hullmove{63.61}{-7.57}\hullline{64.44}{-7.57}\hullline{64.44}{-7.17}\hullline{63.61}{-7.02}\hullfill
\hullface{topsides!72!black}\hullmove{-73.33}{0.00}\hullline{-71.77}{0.00}\hullline{-71.77}{0.78}\hullline{-73.33}{0.78}\hullfill
\hullface{keel!81!black}\hullmove{-67.07}{-4.83}\hullline{-65.51}{-5.91}\hullline{-65.51}{-5.44}\hullline{-67.07}{-4.83}\hullfill
\hullface{keel!69!black}\hullmove{64.44}{-6.94}\hullline{65.26}{-7.03}\hullline{65.26}{-6.69}\hullline{64.44}{-6.77}\hullfill
\hullface{topsides!57!black}\hullmove{71.85}{5.13}\hullline{72.68}{6.06}\hullline{72.68}{6.57}\hullline{71.85}{5.95}\hullfill
\hullface{keel!57!black}\hullmove{-70.20}{-2.15}\hullline{-68.64}{-2.94}\hullline{-68.64}{-2.84}\hullline{-70.20}{-2.09}\hullfill
\hullface{keel!57!black}\hullmove{-49.85}{-9.06}\hullline{-48.29}{-9.06}\hullline{-48.29}{-8.77}\hullline{-49.85}{-8.77}\hullfill
\hullface{keel!72!black}\hullmove{66.08}{-5.59}\hullline{66.91}{-5.44}\hullline{66.91}{-5.27}\hullline{66.08}{-5.44}\hullfill
\hullface{keel!75!black}\hullmove{66.91}{-4.47}\hullline{67.73}{-4.47}\hullline{67.73}{-3.99}\hullline{66.91}{-4.36}\hullfill
\hullface{keel!72!black}\hullmove{67.73}{-2.12}\hullline{68.56}{-2.08}\hullline{68.56}{-1.07}\hullline{67.73}{-1.09}\hullfill
\hullface{topsides!57!black}\hullmove{68.56}{0.00}\hullline{69.38}{0.00}\hullline{69.38}{1.52}\hullline{68.56}{1.52}\hullfill
\hullface{keel!57!black}\hullmove{-59.25}{-8.77}\hullline{-57.68}{-8.77}\hullline{-57.68}{-8.25}\hullline{-59.25}{-8.25}\hullfill
\hullface{keel!57!black}\hullmove{57.85}{-8.77}\hullline{58.67}{-8.77}\hullline{58.67}{-8.25}\hullline{57.85}{-8.25}\hullfill
\hullface{keel!69!black}\hullmove{65.26}{-6.35}\hullline{66.08}{-6.38}\hullline{66.08}{-5.59}\hullline{65.26}{-5.79}\hullfill
\hullface{keel!57!black}\hullmove{50.43}{-9.06}\hullline{51.25}{-9.06}\hullline{51.25}{-8.77}\hullline{50.43}{-8.77}\hullfill
\hullface{keel!72!black}\hullmove{66.91}{-4.36}\hullline{67.73}{-3.99}\hullline{67.73}{-3.77}\hullline{66.91}{-3.77}\hullfill
\hullface{keel!69!black}\hullmove{67.73}{-1.09}\hullline{68.56}{-1.07}\hullline{68.56}{0.00}\hullline{67.73}{0.00}\hullfill
\hullface{keel!57!black}\hullmove{61.14}{-8.25}\hullline{61.97}{-8.25}\hullline{61.97}{-7.57}\hullline{61.14}{-7.57}\hullfill
\hullface{topsides!57!black}\hullmove{71.85}{5.95}\hullline{72.68}{6.57}\hullline{72.68}{7.08}\hullline{71.85}{6.77}\hullfill
\hullface{topsides!60!black}\hullmove{-74.90}{1.48}\hullline{-73.33}{0.78}\hullline{-73.33}{1.56}\hullline{-74.90}{2.08}\hullfill
\hullface{keel!63!black}\hullmove{-63.94}{-7.94}\hullline{-62.38}{-8.25}\hullline{-62.38}{-7.57}\hullline{-63.94}{-7.29}\hullfill
\hullface{keel!69!black}\hullmove{64.44}{-6.77}\hullline{65.26}{-6.69}\hullline{65.26}{-6.35}\hullline{64.44}{-6.30}\hullfill
\hullface{keel!66!black}\hullmove{63.61}{-7.02}\hullline{64.44}{-7.17}\hullline{64.44}{-6.94}\hullline{63.61}{-6.81}\hullfill
\hullface{keel!72!black}\hullmove{66.08}{-5.44}\hullline{66.91}{-5.27}\hullline{66.91}{-4.47}\hullline{66.08}{-4.63}\hullfill
\hullface{keel!72!black}\hullmove{66.91}{-3.77}\hullline{67.73}{-3.77}\hullline{67.73}{-3.51}\hullline{66.91}{-3.51}\hullfill
\hullface{keel!63!black}\hullmove{62.79}{-7.57}\hullline{63.61}{-7.57}\hullline{63.61}{-7.02}\hullline{62.79}{-6.81}\hullfill
\hullface{keel!72!black}\hullmove{66.91}{-3.51}\hullline{67.73}{-3.51}\hullline{67.73}{-3.09}\hullline{66.91}{-3.38}\hullfill
\hullface{keel!57!black}\hullmove{-48.29}{-9.06}\hullline{-46.72}{-9.06}\hullline{-46.72}{-8.77}\hullline{-48.29}{-8.77}\hullfill
\hullface{keel!66!black}\hullmove{63.61}{-6.81}\hullline{64.44}{-6.94}\hullline{64.44}{-6.77}\hullline{63.61}{-6.77}\hullfill
\hullface{keel!57!black}\hullmove{57.02}{-8.77}\hullline{57.85}{-8.77}\hullline{57.85}{-8.25}\hullline{57.02}{-8.25}\hullfill
\hullface{keel!57!black}\hullmove{49.61}{-9.06}\hullline{50.43}{-9.06}\hullline{50.43}{-8.77}\hullline{49.61}{-8.77}\hullfill
\hullface{keel!69!black}\hullmove{65.26}{-5.79}\hullline{66.08}{-5.59}\hullline{66.08}{-5.44}\hullline{65.26}{-5.44}\hullfill
\hullface{topsides!57!black}\hullmove{71.85}{6.77}\hullline{72.68}{7.08}\hullline{72.68}{7.59}\hullline{71.85}{7.59}\hullfill
\hullface{keel!57!black}\hullmove{60.32}{-8.25}\hullline{61.14}{-8.25}\hullline{61.14}{-7.57}\hullline{60.32}{-7.57}\hullfill
\hullface{keel!72!black}\hullmove{66.08}{-4.63}\hullline{66.91}{-4.47}\hullline{66.91}{-4.36}\hullline{66.08}{-4.49}\hullfill
\hullface{keel!57!black}\hullmove{-57.68}{-8.77}\hullline{-56.12}{-8.77}\hullline{-56.12}{-8.25}\hullline{-57.68}{-8.25}\hullfill
\hullface{keel!60!black}\hullmove{-62.38}{-8.25}\hullline{-60.81}{-8.25}\hullline{-60.81}{-7.57}\hullline{-62.38}{-7.57}\hullfill
\hullface{keel!69!black}\hullmove{64.44}{-6.30}\hullline{65.26}{-6.35}\hullline{65.26}{-5.79}\hullline{64.44}{-5.86}\hullfill
\hullface{keel!72!black}\hullmove{66.91}{-3.38}\hullline{67.73}{-3.09}\hullline{67.73}{-2.12}\hullline{66.91}{-2.32}\hullfill
\hullface{keel!57!black}\hullmove{-68.64}{-2.94}\hullline{-67.07}{-4.83}\hullline{-67.07}{-4.68}\hullline{-68.64}{-2.84}\hullfill
\hullface{keel!66!black}\hullmove{63.61}{-6.77}\hullline{64.44}{-6.77}\hullline{64.44}{-6.30}\hullline{63.61}{-6.23}\hullfill
\hullface{keel!72!black}\hullmove{66.08}{-4.49}\hullline{66.91}{-4.36}\hullline{66.91}{-3.77}\hullline{66.08}{-3.79}\hullfill
\hullface{keel!57!black}\hullmove{48.78}{-9.06}\hullline{49.61}{-9.06}\hullline{49.61}{-8.77}\hullline{48.78}{-8.77}\hullfill
\hullface{keel!57!black}\hullmove{56.20}{-8.77}\hullline{57.02}{-8.77}\hullline{57.02}{-8.25}\hullline{56.20}{-8.25}\hullfill
\hullface{keel!69!black}\hullmove{65.26}{-5.44}\hullline{66.08}{-5.44}\hullline{66.08}{-4.63}\hullline{65.26}{-4.79}\hullfill
\hullface{keel!66!black}\hullmove{-65.51}{-6.44}\hullline{-63.94}{-7.94}\hullline{-63.94}{-7.29}\hullline{-65.51}{-5.91}\hullfill
\hullface{keel!60!black}\hullmove{61.97}{-7.57}\hullline{62.79}{-7.57}\hullline{62.79}{-6.81}\hullline{61.97}{-6.77}\hullfill
\hullface{keel!57!black}\hullmove{-46.72}{-9.06}\hullline{-45.16}{-9.06}\hullline{-45.16}{-8.77}\hullline{-46.72}{-8.77}\hullfill
\hullface{keel!63!black}\hullmove{62.79}{-6.81}\hullline{63.61}{-7.02}\hullline{63.61}{-6.81}\hullline{63.61}{-6.77}\hullline{62.79}{-6.63}\hullline{62.79}{-6.77}\hullfill
\hullface{keel!57!black}\hullmove{59.49}{-8.25}\hullline{60.32}{-8.25}\hullline{60.32}{-7.57}\hullline{59.49}{-7.57}\hullfill
\hullface{keel!89!white}\hullmove{-71.77}{-1.27}\hullline{-70.20}{-2.09}\hullline{-70.20}{-1.96}\hullline{-70.20}{-1.80}\hullline{-70.20}{-1.61}\hullline{-70.20}{-1.39}\hullline{-70.20}{-1.15}\hullline{-70.20}{-0.89}\hullline{-70.20}{-0.61}\hullline{-70.20}{-0.31}\hullline{-70.20}{-0.00}\hullline{-71.77}{-0.00}\hullline{-71.77}{-0.19}\hullline{-71.77}{-0.37}\hullline{-71.77}{-0.54}\hullline{-71.77}{-0.70}\hullline{-71.77}{-0.85}\hullline{-71.77}{-0.98}\hullline{-71.77}{-1.10}\hullline{-71.77}{-1.19}\hullfill
\hullface{keel!69!black}\hullmove{64.44}{-5.86}\hullline{65.26}{-5.79}\hullline{65.26}{-5.44}\hullline{64.44}{-5.44}\hullfill
\hullface{keel!72!black}\hullmove{66.91}{-2.32}\hullline{67.73}{-2.12}\hullline{67.73}{-1.09}\hullline{66.91}{-1.19}\hullfill
\hullface{keel!57!black}\hullmove{56.00}{-8.77}\hullline{56.20}{-8.77}\hullline{56.20}{-8.25}\hullline{56.00}{-8.25}\hullfill
\hullface{keel!72!black}\hullmove{66.08}{-3.79}\hullline{66.91}{-3.77}\hullline{66.91}{-3.51}\hullline{66.91}{-3.38}\hullline{66.08}{-3.54}\hullline{66.08}{-3.58}\hullfill
\hullface{keel!57!black}\hullmove{56.00}{-8.77}\hullline{56.00}{-8.77}\hullline{56.00}{-8.25}\hullline{56.00}{-8.25}\hullfill
\hullface{keel!57!black}\hullmove{-56.12}{-8.77}\hullline{-54.55}{-8.77}\hullline{-54.55}{-8.25}\hullline{-56.12}{-8.25}\hullfill
\hullface{keel!69!black}\hullmove{65.26}{-4.79}\hullline{66.08}{-4.63}\hullline{66.08}{-4.49}\hullline{65.26}{-4.52}\hullfill
\hullface{keel!57!black}\hullmove{47.96}{-9.06}\hullline{48.78}{-9.06}\hullline{48.78}{-8.77}\hullline{47.96}{-8.77}\hullfill
\hullface{keel!57!black}\hullmove{-60.81}{-8.25}\hullline{-59.25}{-8.25}\hullline{-59.25}{-7.57}\hullline{-60.81}{-7.57}\hullfill
\hullface{keel!66!black}\hullmove{63.61}{-6.23}\hullline{64.44}{-6.30}\hullline{64.44}{-5.86}\hullline{63.61}{-5.86}\hullfill
\hullface{topsides!57!black}\hullmove{67.73}{0.00}\hullline{68.56}{0.00}\hullline{68.56}{1.52}\hullline{67.73}{1.52}\hullfill
\hullface{keel!57!black}\hullmove{55.37}{-8.77}\hullline{56.00}{-8.77}\hullline{56.00}{-8.25}\hullline{55.37}{-8.25}\hullfill
\hullface{keel!69!black}\hullmove{66.91}{-1.19}\hullline{67.73}{-1.09}\hullline{67.73}{0.00}\hullline{66.91}{0.00}\hullfill
\hullface{keel!63!black}\hullmove{62.79}{-6.63}\hullline{63.61}{-6.77}\hullline{63.61}{-6.23}\hullline{62.79}{-6.12}\hullfill
\hullface{keel!57!black}\hullmove{58.67}{-8.25}\hullline{59.49}{-8.25}\hullline{59.49}{-7.57}\hullline{58.67}{-7.57}\hullfill
\hullface{keel!57!black}\hullmove{61.14}{-7.57}\hullline{61.97}{-7.57}\hullline{61.97}{-6.77}\hullline{61.14}{-6.77}\hullfill
\hullface{keel!72!black}\hullmove{66.08}{-3.54}\hullline{66.91}{-3.38}\hullline{66.91}{-2.32}\hullline{66.08}{-2.46}\hullfill
\hullface{keel!69!black}\hullmove{64.44}{-5.44}\hullline{65.26}{-5.44}\hullline{65.26}{-4.79}\hullline{64.44}{-4.85}\hullfill
\hullface{keel!69!black}\hullmove{-63.94}{-7.29}\hullline{-62.38}{-7.57}\hullline{-62.38}{-6.77}\hullline{-63.94}{-6.52}\hullfill
\hullface{keel!69!black}\hullmove{65.26}{-4.52}\hullline{66.08}{-4.49}\hullline{66.08}{-3.79}\hullline{65.26}{-3.85}\hullfill
\hullface{keel!57!black}\hullmove{-45.16}{-9.06}\hullline{-43.59}{-9.06}\hullline{-43.59}{-8.77}\hullline{-45.16}{-8.77}\hullfill
\hullface{topsides!57!black}\hullmove{-73.33}{0.78}\hullline{-71.77}{0.78}\hullline{-71.77}{1.56}\hullline{-73.33}{1.56}\hullfill
\hullface{topsides!75!black}\hullmove{-71.77}{0.00}\hullline{-70.20}{0.00}\hullline{-70.20}{0.78}\hullline{-71.77}{0.78}\hullfill
\hullface{keel!60!black}\hullmove{61.97}{-6.77}\hullline{62.79}{-6.81}\hullline{62.79}{-6.77}\hullline{61.97}{-6.51}\hullfill
\hullface{topsides!57!black}\hullmove{70.20}{1.83}\hullline{71.03}{3.07}\hullline{71.03}{3.90}\hullline{70.20}{3.27}\hullfill
\hullface{keel!57!black}\hullmove{47.13}{-9.06}\hullline{47.96}{-9.06}\hullline{47.96}{-8.77}\hullline{47.13}{-8.77}\hullfill
\hullface{keel!66!black}\hullmove{63.61}{-5.86}\hullline{64.44}{-5.86}\hullline{64.44}{-5.44}\hullline{63.61}{-5.44}\hullfill
\hullface{keel!57!black}\hullmove{54.55}{-8.77}\hullline{55.37}{-8.77}\hullline{55.37}{-8.25}\hullline{54.55}{-8.25}\hullfill
\hullface{topsides!57!black}\hullmove{71.03}{3.07}\hullline{71.85}{3.90}\hullline{71.85}{4.31}\hullline{71.03}{3.90}\hullfill
\hullface{keel!60!black}\hullmove{61.97}{-6.51}\hullline{62.79}{-6.77}\hullline{62.79}{-6.63}\hullline{61.97}{-6.37}\hullfill
\hullface{keel!57!black}\hullmove{-54.55}{-8.77}\hullline{-52.99}{-8.77}\hullline{-52.99}{-8.25}\hullline{-54.55}{-8.25}\hullfill
\hullface{keel!69!black}\hullmove{65.26}{-3.85}\hullline{66.08}{-3.79}\hullline{66.08}{-3.58}\hullline{66.08}{-3.54}\hullline{65.26}{-3.60}\hullline{65.26}{-3.71}\hullfill
\hullface{keel!69!black}\hullmove{64.44}{-4.85}\hullline{65.26}{-4.79}\hullline{65.26}{-4.52}\hullline{64.44}{-4.57}\hullfill
\hullface{topsides!57!black}\hullmove{69.38}{1.52}\hullline{70.20}{1.83}\hullline{70.20}{3.27}\hullline{69.38}{3.04}\hullfill
\hullface{keel!57!black}\hullmove{-59.25}{-8.25}\hullline{-57.68}{-8.25}\hullline{-57.68}{-7.57}\hullline{-59.25}{-7.57}\hullfill
\hullface{keel!57!black}\hullmove{57.85}{-8.25}\hullline{58.67}{-8.25}\hullline{58.67}{-7.57}\hullline{57.85}{-7.57}\hullfill
\hullface{keel!63!black}\hullmove{62.79}{-6.12}\hullline{63.61}{-6.23}\hullline{63.61}{-5.86}\hullline{62.79}{-5.86}\hullfill
\hullface{keel!57!black}\hullmove{60.32}{-7.57}\hullline{61.14}{-7.57}\hullline{61.14}{-6.77}\hullline{60.32}{-6.77}\hullfill
\hullface{keel!69!black}\hullmove{-65.51}{-5.91}\hullline{-63.94}{-7.29}\hullline{-63.94}{-6.52}\hullline{-65.51}{-5.44}\hullfill
\hullface{keel!69!black}\hullmove{66.08}{-2.46}\hullline{66.91}{-2.32}\hullline{66.91}{-1.19}\hullline{66.08}{-1.26}\hullfill
\hullface{keel!66!black}\hullmove{-62.38}{-7.57}\hullline{-60.81}{-7.57}\hullline{-60.81}{-6.77}\hullline{-62.38}{-6.77}\hullfill
\hullface{keel!57!black}\hullmove{-67.07}{-4.83}\hullline{-65.51}{-5.44}\hullline{-65.51}{-5.44}\hullline{-67.07}{-4.68}\hullfill
\hullface{keel!57!black}\hullmove{46.31}{-9.06}\hullline{47.13}{-9.06}\hullline{47.13}{-8.77}\hullline{46.31}{-8.77}\hullfill
\hullface{keel!66!black}\hullmove{63.61}{-5.44}\hullline{64.44}{-5.44}\hullline{64.44}{-4.85}\hullline{63.61}{-4.85}\hullfill
\hullface{keel!57!black}\hullmove{-43.59}{-9.06}\hullline{-42.03}{-9.06}\hullline{-42.03}{-8.77}\hullline{-43.59}{-8.77}\hullfill
\hullface{keel!60!black}\hullmove{61.97}{-6.37}\hullline{62.79}{-6.63}\hullline{62.79}{-6.12}\hullline{61.97}{-5.97}\hullfill
\hullface{keel!66!black}\hullmove{64.44}{-4.57}\hullline{65.26}{-4.52}\hullline{65.26}{-3.85}\hullline{64.44}{-3.93}\hullfill
\hullface{keel!57!black}\hullmove{53.73}{-8.77}\hullline{54.55}{-8.77}\hullline{54.55}{-8.25}\hullline{53.73}{-8.25}\hullfill
\hullface{keel!69!black}\hullmove{65.26}{-3.60}\hullline{66.08}{-3.54}\hullline{66.08}{-2.46}\hullline{65.26}{-2.55}\hullfill
\hullface{keel!63!black}\hullmove{62.79}{-5.86}\hullline{63.61}{-5.86}\hullline{63.61}{-5.44}\hullline{62.79}{-5.44}\hullfill
\hullface{keel!57!black}\hullmove{57.02}{-8.25}\hullline{57.85}{-8.25}\hullline{57.85}{-7.57}\hullline{57.02}{-7.57}\hullfill
\hullface{keel!69!black}\hullmove{66.08}{-1.26}\hullline{66.91}{-1.19}\hullline{66.91}{0.00}\hullline{66.08}{0.00}\hullfill
\hullface{keel!60!black}\hullmove{61.14}{-6.77}\hullline{61.97}{-6.77}\hullline{61.97}{-6.51}\hullline{61.14}{-5.96}\hullfill
\hullface{keel!57!black}\hullmove{59.49}{-7.57}\hullline{60.32}{-7.57}\hullline{60.32}{-6.77}\hullline{59.49}{-6.77}\hullfill
\hullface{topsides!57!black}\hullmove{66.91}{0.00}\hullline{67.73}{0.00}\hullline{67.73}{1.52}\hullline{66.91}{1.52}\hullfill
\hullface{keel!66!black}\hullmove{63.61}{-4.85}\hullline{64.44}{-4.85}\hullline{64.44}{-4.57}\hullline{63.61}{-4.64}\hullfill
\hullface{keel!57!black}\hullmove{-52.99}{-8.77}\hullline{-51.42}{-8.77}\hullline{-51.42}{-8.25}\hullline{-52.99}{-8.25}\hullfill
\hullface{keel!66!black}\hullmove{64.44}{-3.93}\hullline{65.26}{-3.85}\hullline{65.26}{-3.71}\hullline{65.26}{-3.60}\hullline{64.44}{-3.70}\hullline{64.44}{-3.75}\hullfill
\hullface{keel!57!black}\hullmove{-57.68}{-8.25}\hullline{-56.12}{-8.25}\hullline{-56.12}{-7.57}\hullline{-57.68}{-7.57}\hullfill
\hullface{keel!57!black}\hullmove{45.49}{-9.06}\hullline{46.31}{-9.06}\hullline{46.31}{-8.77}\hullline{45.49}{-8.77}\hullfill
\hullface{keel!63!black}\hullmove{61.97}{-5.97}\hullline{62.79}{-6.12}\hullline{62.79}{-5.86}\hullline{61.97}{-5.86}\hullfill
\hullface{topsides!63!black}\hullmove{71.03}{3.90}\hullline{71.85}{4.31}\hullline{71.85}{5.13}\hullline{71.85}{5.95}\hullline{71.85}{6.77}\hullline{71.85}{7.59}\hullline{71.03}{7.59}\hullline{71.03}{6.46}\hullline{71.03}{5.33}\hullline{71.03}{4.20}\hullfill
\hullface{topsides!57!black}\hullmove{-74.90}{2.08}\hullline{-73.33}{1.56}\hullline{-73.33}{2.34}\hullline{-74.90}{2.69}\hullfill
\hullface{keel!89!white}\hullmove{-70.20}{-2.09}\hullline{-68.64}{-2.84}\hullline{-68.64}{-2.67}\hullline{-68.64}{-2.45}\hullline{-68.64}{-2.19}\hullline{-68.64}{-1.90}\hullline{-68.64}{-1.57}\hullline{-68.64}{-1.22}\hullline{-68.64}{-0.84}\hullline{-68.64}{-0.43}\hullline{-68.64}{-0.00}\hullline{-70.20}{-0.00}\hullline{-70.20}{-0.31}\hullline{-70.20}{-0.61}\hullline{-70.20}{-0.89}\hullline{-70.20}{-1.15}\hullline{-70.20}{-1.39}\hullline{-70.20}{-1.61}\hullline{-70.20}{-1.80}\hullline{-70.20}{-1.96}\hullfill
\hullface{keel!72!black}\hullmove{-63.94}{-6.52}\hullline{-62.38}{-6.77}\hullline{-62.38}{-5.86}\hullline{-63.94}{-5.64}\hullfill
\hullface{keel!57!black}\hullmove{52.90}{-8.77}\hullline{53.73}{-8.77}\hullline{53.73}{-8.25}\hullline{52.90}{-8.25}\hullfill
\hullface{keel!63!black}\hullmove{-60.81}{-7.57}\hullline{-59.25}{-7.57}\hullline{-59.25}{-6.77}\hullline{-60.81}{-6.77}\hullfill
\hullface{keel!72!black}\hullmove{-65.51}{-5.44}\hullline{-63.94}{-6.52}\hullline{-63.94}{-5.64}\hullline{-65.51}{-5.44}\hullfill
\hullface{keel!63!black}\hullmove{62.79}{-5.44}\hullline{63.61}{-5.44}\hullline{63.61}{-4.85}\hullline{62.79}{-4.85}\hullfill
\hullface{keel!57!black}\hullmove{-42.03}{-9.06}\hullline{-40.46}{-9.06}\hullline{-40.46}{-8.77}\hullline{-42.03}{-8.77}\hullfill
\hullface{keel!57!black}\hullmove{56.20}{-8.25}\hullline{57.02}{-8.25}\hullline{57.02}{-7.57}\hullline{56.20}{-7.57}\hullfill
\hullface{keel!69!black}\hullmove{65.26}{-2.55}\hullline{66.08}{-2.46}\hullline{66.08}{-1.26}\hullline{65.26}{-1.31}\hullfill
\hullface{keel!66!black}\hullmove{63.61}{-4.64}\hullline{64.44}{-4.57}\hullline{64.44}{-3.93}\hullline{63.61}{-4.06}\hullfill
\hullface{keel!60!black}\hullmove{61.14}{-5.96}\hullline{61.97}{-6.51}\hullline{61.97}{-6.37}\hullline{61.14}{-5.89}\hullfill
\hullface{keel!63!black}\hullmove{61.97}{-5.86}\hullline{62.79}{-5.86}\hullline{62.79}{-5.44}\hullline{61.97}{-5.44}\hullfill
\hullface{topsides!57!black}\hullmove{68.56}{1.52}\hullline{69.38}{1.52}\hullline{69.38}{3.04}\hullline{68.56}{3.04}\hullfill
\hullface{keel!57!black}\hullmove{58.67}{-7.57}\hullline{59.49}{-7.57}\hullline{59.49}{-6.77}\hullline{58.67}{-6.77}\hullfill
\hullface{keel!66!black}\hullmove{64.44}{-3.70}\hullline{65.26}{-3.60}\hullline{65.26}{-2.55}\hullline{64.44}{-2.58}\hullfill
\hullface{keel!57!black}\hullmove{44.66}{-9.06}\hullline{45.49}{-9.06}\hullline{45.49}{-8.77}\hullline{44.66}{-8.77}\hullfill
\hullface{keel!72!black}\hullmove{-65.51}{-5.44}\hullline{-63.94}{-5.64}\hullline{-63.94}{-5.44}\hullline{-65.51}{-5.44}\hullfill
\hullface{topsides!60!black}\hullmove{-71.77}{0.78}\hullline{-70.20}{0.78}\hullline{-70.20}{1.56}\hullline{-71.77}{1.56}\hullfill
\hullface{keel!57!black}\hullmove{56.00}{-8.25}\hullline{56.20}{-8.25}\hullline{56.20}{-7.57}\hullline{56.00}{-7.57}\hullfill
\hullface{topsides!78!black}\hullmove{-70.20}{0.00}\hullline{-68.64}{0.00}\hullline{-68.64}{0.78}\hullline{-70.20}{0.78}\hullfill
\hullface{keel!63!black}\hullmove{62.79}{-4.85}\hullline{63.61}{-4.85}\hullline{63.61}{-4.64}\hullline{62.79}{-4.75}\hullfill
\hullface{keel!60!black}\hullmove{61.14}{-5.89}\hullline{61.97}{-6.37}\hullline{61.97}{-5.97}\hullline{61.97}{-5.86}\hullline{61.14}{-5.70}\hullline{61.14}{-5.86}\hullfill
\hullface{keel!69!black}\hullmove{-65.51}{-5.44}\hullline{-63.94}{-5.44}\hullline{-63.94}{-5.44}\hullline{-65.51}{-5.28}\hullfill
\hullface{keel!57!black}\hullmove{56.00}{-8.25}\hullline{56.00}{-8.25}\hullline{56.00}{-7.57}\hullline{56.00}{-7.57}\hullfill
\hullface{keel!57!black}\hullmove{60.32}{-6.77}\hullline{61.14}{-6.77}\hullline{61.14}{-5.96}\hullline{60.32}{-5.86}\hullfill
\hullface{keel!66!black}\hullmove{63.61}{-4.06}\hullline{64.44}{-3.93}\hullline{64.44}{-3.75}\hullline{64.44}{-3.70}\hullline{63.61}{-3.75}\hullline{63.61}{-3.85}\hullfill
\hullface{keel!57!black}\hullmove{52.08}{-8.77}\hullline{52.90}{-8.77}\hullline{52.90}{-8.25}\hullline{52.08}{-8.25}\hullfill
\hullface{keel!57!black}\hullmove{-56.12}{-8.25}\hullline{-54.55}{-8.25}\hullline{-54.55}{-7.57}\hullline{-56.12}{-7.57}\hullfill
\hullface{keel!57!black}\hullmove{-51.42}{-8.77}\hullline{-49.85}{-8.77}\hullline{-49.85}{-8.25}\hullline{-51.42}{-8.25}\hullfill
\hullface{keel!72!black}\hullmove{-62.38}{-6.77}\hullline{-60.81}{-6.77}\hullline{-60.81}{-5.86}\hullline{-62.38}{-5.86}\hullfill
\hullface{keel!69!black}\hullmove{65.26}{-1.31}\hullline{66.08}{-1.26}\hullline{66.08}{0.00}\hullline{65.26}{0.00}\hullfill
\hullface{keel!57!black}\hullmove{55.37}{-8.25}\hullline{56.00}{-8.25}\hullline{56.00}{-7.57}\hullline{55.37}{-7.57}\hullfill
\hullface{keel!63!black}\hullmove{61.97}{-5.44}\hullline{62.79}{-5.44}\hullline{62.79}{-4.85}\hullline{61.97}{-4.90}\hullfill
\hullface{keel!72!black}\hullmove{-65.51}{-5.28}\hullline{-63.94}{-5.44}\hullline{-63.94}{-5.44}\hullline{-65.51}{-4.57}\hullfill
\hullface{keel!63!black}\hullmove{62.79}{-4.75}\hullline{63.61}{-4.64}\hullline{63.61}{-4.06}\hullline{62.79}{-4.24}\hullfill
\hullface{keel!75!black}\hullmove{-63.94}{-5.64}\hullline{-62.38}{-5.86}\hullline{-62.38}{-5.44}\hullline{-63.94}{-5.44}\hullfill
\hullface{keel!63!black}\hullmove{-59.25}{-7.57}\hullline{-57.68}{-7.57}\hullline{-57.68}{-6.77}\hullline{-59.25}{-6.77}\hullfill
\hullface{keel!57!black}\hullmove{57.85}{-7.57}\hullline{58.67}{-7.57}\hullline{58.67}{-6.77}\hullline{57.85}{-6.77}\hullfill
\hullface{keel!57!black}\hullmove{-40.46}{-9.06}\hullline{-38.90}{-9.06}\hullline{-38.90}{-8.77}\hullline{-40.46}{-8.77}\hullfill
\hullface{keel!57!black}\hullmove{43.84}{-9.06}\hullline{44.66}{-9.06}\hullline{44.66}{-8.77}\hullline{43.84}{-8.77}\hullfill
\hullface{topsides!57!black}\hullmove{66.08}{0.00}\hullline{66.91}{0.00}\hullline{66.91}{1.52}\hullline{66.08}{1.52}\hullfill
\hullface{keel!60!black}\hullmove{61.14}{-5.70}\hullline{61.97}{-5.86}\hullline{61.97}{-5.44}\hullline{61.14}{-5.44}\hullfill
\hullface{keel!69!black}\hullmove{64.44}{-2.58}\hullline{65.26}{-2.55}\hullline{65.26}{-1.31}\hullline{64.44}{-1.32}\hullfill
\hullface{keel!66!black}\hullmove{63.61}{-3.75}\hullline{64.44}{-3.70}\hullline{64.44}{-2.58}\hullline{63.61}{-2.58}\hullfill
\hullface{keel!57!black}\hullmove{59.49}{-6.77}\hullline{60.32}{-6.77}\hullline{60.32}{-5.86}\hullline{59.49}{-5.86}\hullfill
\hullface{topsides!57!black}\hullmove{-73.33}{1.56}\hullline{-71.77}{1.56}\hullline{-71.77}{2.34}\hullline{-73.33}{2.34}\hullfill
\hullface{keel!57!black}\hullmove{51.25}{-8.77}\hullline{52.08}{-8.77}\hullline{52.08}{-8.25}\hullline{51.25}{-8.25}\hullfill
\hullface{topsides!57!black}\hullmove{70.20}{3.27}\hullline{71.03}{3.90}\hullline{71.03}{4.20}\hullline{70.20}{3.90}\hullfill
\hullface{keel!57!black}\hullmove{54.55}{-8.25}\hullline{55.37}{-8.25}\hullline{55.37}{-7.57}\hullline{54.55}{-7.57}\hullfill
\hullface{keel!66!black}\hullmove{62.79}{-4.24}\hullline{63.61}{-4.06}\hullline{63.61}{-3.85}\hullline{63.61}{-3.75}\hullline{62.79}{-3.75}\hullline{62.79}{-4.06}\hullfill
\hullface{keel!78!black}\hullmove{-67.07}{-4.68}\hullline{-65.51}{-5.44}\hullline{-65.51}{-5.44}\hullline{-65.51}{-5.44}\hullline{-67.07}{-4.04}\hullline{-67.07}{-4.40}\hullfill
\hullface{keel!63!black}\hullmove{61.97}{-4.90}\hullline{62.79}{-4.85}\hullline{62.79}{-4.75}\hullline{62.79}{-4.24}\hullline{61.97}{-4.51}\hullline{61.97}{-4.85}\hullfill
\hullface{keel!78!black}\hullmove{-67.07}{-4.04}\hullline{-65.51}{-5.44}\hullline{-65.51}{-5.28}\hullline{-67.07}{-3.61}\hullfill
\hullface{keel!57!black}\hullmove{-54.55}{-8.25}\hullline{-52.99}{-8.25}\hullline{-52.99}{-7.57}\hullline{-54.55}{-7.57}\hullfill
\hullface{keel!57!black}\hullmove{43.02}{-9.06}\hullline{43.84}{-9.06}\hullline{43.84}{-8.77}\hullline{43.02}{-8.77}\hullfill
\hullface{keel!57!black}\hullmove{57.02}{-7.57}\hullline{57.85}{-7.57}\hullline{57.85}{-6.77}\hullline{57.02}{-6.77}\hullfill
\hullface{keel!78!black}\hullmove{-63.94}{-5.44}\hullline{-62.38}{-5.44}\hullline{-62.38}{-4.85}\hullline{-63.94}{-4.67}\hullfill
\hullface{keel!75!black}\hullmove{-65.51}{-4.57}\hullline{-63.94}{-5.44}\hullline{-63.94}{-4.67}\hullline{-65.51}{-3.78}\hullfill
\hullface{keel!60!black}\hullmove{60.32}{-5.86}\hullline{61.14}{-5.96}\hullline{61.14}{-5.89}\hullline{61.14}{-5.86}\hullline{61.14}{-5.70}\hullline{60.32}{-5.44}\hullline{60.32}{-5.44}\hullline{60.32}{-5.44}\hullfill
\hullface{keel!89!white}\hullmove{-68.64}{-2.84}\hullline{-67.07}{-4.68}\hullline{-67.07}{-4.40}\hullline{-67.07}{-4.04}\hullline{-67.07}{-3.61}\hullline{-67.07}{-3.13}\hullline{-67.07}{-2.59}\hullline{-67.07}{-2.00}\hullline{-67.07}{-1.38}\hullline{-67.07}{-0.71}\hullline{-67.07}{-0.00}\hullline{-68.64}{-0.00}\hullline{-68.64}{-0.43}\hullline{-68.64}{-0.84}\hullline{-68.64}{-1.22}\hullline{-68.64}{-1.57}\hullline{-68.64}{-1.90}\hullline{-68.64}{-2.19}\hullline{-68.64}{-2.45}\hullline{-68.64}{-2.67}\hullfill
\hullface{keel!57!black}\hullmove{-49.85}{-8.77}\hullline{-48.29}{-8.77}\hullline{-48.29}{-8.25}\hullline{-49.85}{-8.25}\hullfill
\hullface{topsides!57!black}\hullmove{67.73}{1.52}\hullline{68.56}{1.52}\hullline{68.56}{3.04}\hullline{67.73}{3.04}\hullfill
\hullface{keel!63!black}\hullmove{61.14}{-5.44}\hullline{61.97}{-5.44}\hullline{61.97}{-4.90}\hullline{61.97}{-4.85}\hullline{61.14}{-4.98}\hullline{61.14}{-5.17}\hullfill
\hullface{keel!69!black}\hullmove{-60.81}{-6.77}\hullline{-59.25}{-6.77}\hullline{-59.25}{-5.86}\hullline{-60.81}{-5.86}\hullfill
\hullface{keel!57!black}\hullmove{-38.90}{-9.06}\hullline{-37.33}{-9.06}\hullline{-37.33}{-8.77}\hullline{-38.90}{-8.77}\hullfill
\hullface{keel!81!black}\hullmove{-67.07}{-3.61}\hullline{-65.51}{-5.28}\hullline{-65.51}{-4.57}\hullline{-67.07}{-3.13}\hullfill
\hullface{keel!75!black}\hullmove{-62.38}{-5.86}\hullline{-60.81}{-5.86}\hullline{-60.81}{-5.44}\hullline{-62.38}{-5.44}\hullfill
\hullface{keel!60!black}\hullmove{-57.68}{-7.57}\hullline{-56.12}{-7.57}\hullline{-56.12}{-6.77}\hullline{-57.68}{-6.77}\hullfill
\hullface{topsides!66!black}\hullmove{70.20}{3.90}\hullline{71.03}{4.20}\hullline{71.03}{5.33}\hullline{71.03}{6.46}\hullline{71.03}{7.59}\hullline{70.20}{7.59}\hullline{70.20}{6.15}\hullline{70.20}{4.71}\hullfill
\hullface{keel!60!black}\hullmove{60.32}{-5.44}\hullline{61.14}{-5.70}\hullline{61.14}{-5.44}\hullline{60.32}{-5.44}\hullfill
\hullface{keel!69!black}\hullmove{64.44}{-1.32}\hullline{65.26}{-1.31}\hullline{65.26}{0.00}\hullline{64.44}{0.00}\hullfill
\hullface{keel!57!black}\hullmove{58.67}{-6.77}\hullline{59.49}{-6.77}\hullline{59.49}{-5.86}\hullline{58.67}{-5.86}\hullfill
\hullface{keel!57!black}\hullmove{53.73}{-8.25}\hullline{54.55}{-8.25}\hullline{54.55}{-7.57}\hullline{53.73}{-7.57}\hullfill
\hullface{keel!57!black}\hullmove{50.43}{-8.77}\hullline{51.25}{-8.77}\hullline{51.25}{-8.25}\hullline{50.43}{-8.25}\hullfill
\hullface{keel!57!black}\hullmove{42.19}{-9.06}\hullline{43.02}{-9.06}\hullline{43.02}{-8.77}\hullline{42.19}{-8.77}\hullfill
\hullface{keel!63!black}\hullmove{61.97}{-4.51}\hullline{62.79}{-4.24}\hullline{62.79}{-4.06}\hullline{62.79}{-3.75}\hullline{61.97}{-3.75}\hullline{61.97}{-4.36}\hullfill
\hullface{topsides!63!black}\hullmove{-70.20}{0.78}\hullline{-68.64}{0.78}\hullline{-68.64}{1.56}\hullline{-70.20}{1.56}\hullfill
\hullface{keel!66!black}\hullmove{62.79}{-3.75}\hullline{63.61}{-3.75}\hullline{63.61}{-2.58}\hullline{62.79}{-2.58}\hullfill
\hullface{keel!63!black}\hullmove{61.14}{-4.98}\hullline{61.97}{-4.85}\hullline{61.97}{-4.51}\hullline{61.97}{-4.36}\hullline{61.14}{-4.85}\hullline{61.14}{-4.91}\hullfill
\hullface{keel!60!black}\hullmove{60.32}{-5.44}\hullline{61.14}{-5.44}\hullline{61.14}{-5.17}\hullline{61.14}{-4.98}\hullline{60.32}{-5.44}\hullline{60.32}{-5.44}\hullfill
\hullface{keel!66!black}\hullmove{63.61}{-2.58}\hullline{64.44}{-2.58}\hullline{64.44}{-1.32}\hullline{63.61}{-1.32}\hullfill
\hullface{topsides!81!black}\hullmove{-68.64}{0.00}\hullline{-67.07}{0.00}\hullline{-67.07}{0.78}\hullline{-68.64}{0.78}\hullfill
\hullface{topsides!57!black}\hullmove{65.26}{0.00}\hullline{66.08}{0.00}\hullline{66.08}{1.52}\hullline{65.26}{1.52}\hullfill
\hullface{keel!57!black}\hullmove{56.20}{-7.57}\hullline{57.02}{-7.57}\hullline{57.02}{-6.77}\hullline{56.20}{-6.77}\hullfill
\hullface{keel!90!black}\hullmove{-67.07}{-3.13}\hullline{-65.51}{-4.57}\hullline{-65.51}{-3.78}\hullline{-67.07}{-2.59}\hullfill
\hullface{keel!63!black}\hullmove{60.32}{-5.44}\hullline{61.14}{-4.98}\hullline{61.14}{-4.91}\hullline{60.32}{-5.44}\hullfill
\hullface{topsides!57!black}\hullmove{69.38}{3.04}\hullline{70.20}{3.27}\hullline{70.20}{3.90}\hullline{69.38}{3.90}\hullfill
\hullface{keel!57!black}\hullmove{-37.33}{-9.06}\hullline{-35.77}{-9.06}\hullline{-35.77}{-8.77}\hullline{-37.33}{-8.77}\hullfill
\hullface{keel!60!black}\hullmove{59.49}{-5.86}\hullline{60.32}{-5.86}\hullline{60.32}{-5.44}\hullline{59.49}{-5.44}\hullfill
\hullface{keel!57!black}\hullmove{-52.99}{-8.25}\hullline{-51.42}{-8.25}\hullline{-51.42}{-7.57}\hullline{-52.99}{-7.57}\hullfill
\hullface{keel!57!black}\hullmove{41.37}{-9.06}\hullline{42.19}{-9.06}\hullline{42.19}{-8.77}\hullline{41.37}{-8.77}\hullfill
\hullface{keel!81!black}\hullmove{-65.51}{-3.78}\hullline{-63.94}{-4.67}\hullline{-63.94}{-3.61}\hullline{-65.51}{-2.93}\hullfill
\hullface{keel!78!black}\hullmove{-62.38}{-5.44}\hullline{-60.81}{-5.44}\hullline{-60.81}{-4.85}\hullline{-62.38}{-4.85}\hullfill
\hullface{keel!60!black}\hullmove{60.32}{-5.44}\hullline{61.14}{-4.91}\hullline{61.14}{-4.85}\hullline{60.32}{-4.85}\hullfill
\hullface{keel!96!black}\hullmove{-67.07}{-2.59}\hullline{-65.51}{-3.78}\hullline{-65.51}{-2.93}\hullline{-67.07}{-2.00}\hullfill
\hullface{keel!57!black}\hullmove{56.00}{-7.57}\hullline{56.20}{-7.57}\hullline{56.20}{-6.77}\hullline{56.00}{-6.77}\hullfill
\hullface{keel!81!black}\hullmove{-63.94}{-4.67}\hullline{-62.38}{-4.85}\hullline{-62.38}{-3.75}\hullline{-63.94}{-3.61}\hullfill
\hullface{keel!57!black}\hullmove{-35.77}{-9.06}\hullline{41.37}{-9.06}\hullline{41.37}{-8.77}\hullline{-35.77}{-8.77}\hullfill
\hullface{keel!69!black}\hullmove{-59.25}{-6.77}\hullline{-57.68}{-6.77}\hullline{-57.68}{-5.86}\hullline{-59.25}{-5.86}\hullfill
\hullface{keel!57!black}\hullmove{57.85}{-6.77}\hullline{58.67}{-6.77}\hullline{58.67}{-5.86}\hullline{57.85}{-5.86}\hullfill
\hullface{keel!57!black}\hullmove{56.00}{-7.57}\hullline{56.00}{-7.57}\hullline{56.00}{-6.77}\hullline{56.00}{-6.77}\hullfill
\hullface{keel!57!black}\hullmove{52.90}{-8.25}\hullline{53.73}{-8.25}\hullline{53.73}{-7.57}\hullline{52.90}{-7.57}\hullfill
\hullface{keel!57!black}\hullmove{-48.29}{-8.77}\hullline{-46.72}{-8.77}\hullline{-46.72}{-8.25}\hullline{-48.29}{-8.25}\hullfill
\hullface{keel!60!black}\hullmove{-56.12}{-7.57}\hullline{-54.55}{-7.57}\hullline{-54.55}{-6.77}\hullline{-56.12}{-6.77}\hullfill
\hullface{keel!63!black}\hullmove{61.14}{-4.85}\hullline{61.97}{-4.36}\hullline{61.97}{-3.75}\hullline{61.14}{-3.75}\hullfill
\hullface{keel!57!black}\hullmove{49.61}{-8.77}\hullline{50.43}{-8.77}\hullline{50.43}{-8.25}\hullline{49.61}{-8.25}\hullfill
\hullface{keel!72!black}\hullmove{-60.81}{-5.86}\hullline{-59.25}{-5.86}\hullline{-59.25}{-5.44}\hullline{-60.81}{-5.44}\hullfill
\hullface{topsides!57!black}\hullmove{-71.77}{1.56}\hullline{-70.20}{1.56}\hullline{-70.20}{2.34}\hullline{-71.77}{2.34}\hullfill
\hullface{keel!99!black}\hullmove{-67.07}{-2.00}\hullline{-65.51}{-2.93}\hullline{-65.51}{-2.01}\hullline{-67.07}{-1.38}\hullfill
\hullface{topsides!57!black}\hullmove{66.91}{1.52}\hullline{67.73}{1.52}\hullline{67.73}{3.04}\hullline{66.91}{3.04}\hullfill
\hullface{keel!57!black}\hullmove{55.37}{-7.57}\hullline{56.00}{-7.57}\hullline{56.00}{-6.77}\hullline{55.37}{-6.77}\hullfill
\hullface{keel!69!black}\hullmove{63.61}{-1.32}\hullline{64.44}{-1.32}\hullline{64.44}{0.00}\hullline{63.61}{0.00}\hullfill
\hullface{keel!66!black}\hullmove{61.97}{-3.75}\hullline{62.79}{-3.75}\hullline{62.79}{-2.58}\hullline{61.97}{-2.58}\hullfill
\hullface{keel!95!white}\hullmove{-67.07}{-1.38}\hullline{-65.51}{-2.01}\hullline{-65.51}{-1.03}\hullline{-67.07}{-0.71}\hullfill
\hullface{keel!57!black}\hullmove{58.67}{-5.86}\hullline{59.49}{-5.86}\hullline{59.49}{-5.44}\hullline{58.67}{-5.44}\hullfill
\hullface{keel!66!black}\hullmove{62.79}{-2.58}\hullline{63.61}{-2.58}\hullline{63.61}{-1.32}\hullline{62.79}{-1.32}\hullfill
\hullface{keel!60!black}\hullmove{59.49}{-5.44}\hullline{60.32}{-5.44}\hullline{60.32}{-4.85}\hullline{59.49}{-4.85}\hullfill
\hullface{keel!92!white}\hullmove{-67.07}{-0.71}\hullline{-65.51}{-1.03}\hullline{-65.51}{-0.00}\hullline{-67.07}{-0.00}\hullfill
\hullface{topsides!57!black}\hullmove{-74.90}{2.69}\hullline{-73.33}{2.34}\hullline{-73.33}{3.12}\hullline{-74.90}{3.29}\hullfill
\hullface{topsides!57!black}\hullmove{64.44}{0.00}\hullline{65.26}{0.00}\hullline{65.26}{1.52}\hullline{64.44}{1.52}\hullfill
\hullface{keel!87!black}\hullmove{-65.51}{-2.93}\hullline{-63.94}{-3.61}\hullline{-63.94}{-2.48}\hullline{-65.51}{-2.01}\hullfill
\hullface{keel!57!black}\hullmove{57.02}{-6.77}\hullline{57.85}{-6.77}\hullline{57.85}{-5.86}\hullline{57.02}{-5.86}\hullfill
\hullface{keel!57!black}\hullmove{52.08}{-8.25}\hullline{52.90}{-8.25}\hullline{52.90}{-7.57}\hullline{52.08}{-7.57}\hullfill
\hullface{keel!57!black}\hullmove{-51.42}{-8.25}\hullline{-49.85}{-8.25}\hullline{-49.85}{-7.57}\hullline{-51.42}{-7.57}\hullfill
\hullface{keel!63!black}\hullmove{60.32}{-4.85}\hullline{61.14}{-4.85}\hullline{61.14}{-3.75}\hullline{60.32}{-3.75}\hullfill
\hullface{keel!57!black}\hullmove{54.55}{-7.57}\hullline{55.37}{-7.57}\hullline{55.37}{-6.77}\hullline{54.55}{-6.77}\hullfill
\hullface{topsides!66!black}\hullmove{-68.64}{0.78}\hullline{-67.07}{0.78}\hullline{-67.07}{1.56}\hullline{-68.64}{1.56}\hullfill
\hullface{keel!75!black}\hullmove{-60.81}{-5.44}\hullline{-59.25}{-5.44}\hullline{-59.25}{-4.85}\hullline{-60.81}{-4.85}\hullfill
\hullface{keel!81!black}\hullmove{-62.38}{-4.85}\hullline{-60.81}{-4.85}\hullline{-60.81}{-3.75}\hullline{-62.38}{-3.75}\hullfill
\hullface{topsides!69!black}\hullmove{69.38}{3.90}\hullline{70.20}{3.90}\hullline{70.20}{4.71}\hullline{70.20}{6.15}\hullline{70.20}{7.59}\hullline{69.38}{7.59}\hullline{69.38}{6.07}\hullline{69.38}{4.55}\hullfill
\hullface{keel!57!black}\hullmove{48.78}{-8.77}\hullline{49.61}{-8.77}\hullline{49.61}{-8.25}\hullline{48.78}{-8.25}\hullfill
\hullface{keel!66!black}\hullmove{-57.68}{-6.77}\hullline{-56.12}{-6.77}\hullline{-56.12}{-5.86}\hullline{-57.68}{-5.86}\hullfill
\hullface{topsides!81!black}\hullmove{-67.07}{0.00}\hullline{-65.51}{0.00}\hullline{-65.51}{0.78}\hullline{-67.07}{0.78}\hullfill
\hullface{keel!60!black}\hullmove{-54.55}{-7.57}\hullline{-52.99}{-7.57}\hullline{-52.99}{-6.77}\hullline{-54.55}{-6.77}\hullfill
\hullface{topsides!57!black}\hullmove{68.56}{3.04}\hullline{69.38}{3.04}\hullline{69.38}{3.90}\hullline{68.56}{3.90}\hullfill
\hullface{keel!84!black}\hullmove{-63.94}{-3.61}\hullline{-62.38}{-3.75}\hullline{-62.38}{-2.58}\hullline{-63.94}{-2.48}\hullfill
\hullface{keel!72!black}\hullmove{-59.25}{-5.86}\hullline{-57.68}{-5.86}\hullline{-57.68}{-5.44}\hullline{-59.25}{-5.44}\hullfill
\hullface{keel!57!black}\hullmove{57.85}{-5.86}\hullline{58.67}{-5.86}\hullline{58.67}{-5.44}\hullline{57.85}{-5.44}\hullfill
\hullface{keel!66!black}\hullmove{61.14}{-3.75}\hullline{61.97}{-3.75}\hullline{61.97}{-2.58}\hullline{61.14}{-2.58}\hullfill
\hullface{keel!57!black}\hullmove{-46.72}{-8.77}\hullline{-45.16}{-8.77}\hullline{-45.16}{-8.25}\hullline{-46.72}{-8.25}\hullfill
\hullface{keel!60!black}\hullmove{58.67}{-5.44}\hullline{59.49}{-5.44}\hullline{59.49}{-4.85}\hullline{58.67}{-4.85}\hullfill
\hullface{keel!72!black}\hullmove{62.79}{-1.32}\hullline{63.61}{-1.32}\hullline{63.61}{0.00}\hullline{62.79}{0.00}\hullfill
\hullface{keel!93!black}\hullmove{-65.51}{-2.01}\hullline{-63.94}{-2.48}\hullline{-63.94}{-1.28}\hullline{-65.51}{-1.03}\hullfill
\hullface{topsides!57!black}\hullmove{66.08}{1.52}\hullline{66.91}{1.52}\hullline{66.91}{3.04}\hullline{66.08}{3.04}\hullfill
\hullface{keel!57!black}\hullmove{56.20}{-6.77}\hullline{57.02}{-6.77}\hullline{57.02}{-5.86}\hullline{56.20}{-5.86}\hullfill
\hullface{keel!69!black}\hullmove{61.97}{-2.58}\hullline{62.79}{-2.58}\hullline{62.79}{-1.32}\hullline{61.97}{-1.32}\hullfill
\hullface{keel!57!black}\hullmove{51.25}{-8.25}\hullline{52.08}{-8.25}\hullline{52.08}{-7.57}\hullline{51.25}{-7.57}\hullfill
\hullface{topsides!57!black}\hullmove{-70.20}{1.56}\hullline{-68.64}{1.56}\hullline{-68.64}{2.34}\hullline{-70.20}{2.34}\hullfill
\hullface{keel!57!black}\hullmove{53.73}{-7.57}\hullline{54.55}{-7.57}\hullline{54.55}{-6.77}\hullline{53.73}{-6.77}\hullfill
\hullface{keel!63!black}\hullmove{59.49}{-4.85}\hullline{60.32}{-4.85}\hullline{60.32}{-3.75}\hullline{59.49}{-3.75}\hullfill
\hullface{topsides!60!black}\hullmove{63.61}{0.00}\hullline{64.44}{0.00}\hullline{64.44}{1.52}\hullline{63.61}{1.52}\hullfill
\hullface{topsides!57!black}\hullmove{-73.33}{2.34}\hullline{-71.77}{2.34}\hullline{-71.77}{3.12}\hullline{-73.33}{3.12}\hullfill
\hullface{keel!98!white}\hullmove{-65.51}{-1.03}\hullline{-63.94}{-1.28}\hullline{-63.94}{-0.00}\hullline{-65.51}{-0.00}\hullfill
\hullface{keel!57!black}\hullmove{56.00}{-6.77}\hullline{56.20}{-6.77}\hullline{56.20}{-5.86}\hullline{56.00}{-5.86}\hullfill
\hullface{keel!57!black}\hullmove{47.96}{-8.77}\hullline{48.78}{-8.77}\hullline{48.78}{-8.25}\hullline{47.96}{-8.25}\hullfill
\hullface{keel!57!black}\hullmove{-49.85}{-8.25}\hullline{-48.29}{-8.25}\hullline{-48.29}{-7.57}\hullline{-49.85}{-7.57}\hullfill
\hullface{keel!57!black}\hullmove{56.00}{-6.77}\hullline{56.00}{-6.77}\hullline{56.00}{-5.86}\hullline{56.00}{-5.86}\hullfill
\hullface{keel!57!black}\hullmove{57.02}{-5.86}\hullline{57.85}{-5.86}\hullline{57.85}{-5.44}\hullline{57.02}{-5.44}\hullfill
\hullface{keel!75!black}\hullmove{-59.25}{-5.44}\hullline{-57.68}{-5.44}\hullline{-57.68}{-4.85}\hullline{-59.25}{-4.85}\hullfill
\hullface{keel!66!black}\hullmove{-56.12}{-6.77}\hullline{-54.55}{-6.77}\hullline{-54.55}{-5.86}\hullline{-56.12}{-5.86}\hullfill
\hullface{keel!60!black}\hullmove{57.85}{-5.44}\hullline{58.67}{-5.44}\hullline{58.67}{-4.85}\hullline{57.85}{-4.85}\hullfill
\hullface{keel!81!black}\hullmove{-60.81}{-4.85}\hullline{-59.25}{-4.85}\hullline{-59.25}{-3.75}\hullline{-60.81}{-3.75}\hullfill
\hullface{keel!60!black}\hullmove{-52.99}{-7.57}\hullline{-51.42}{-7.57}\hullline{-51.42}{-6.77}\hullline{-52.99}{-6.77}\hullfill
\hullface{keel!66!black}\hullmove{60.32}{-3.75}\hullline{61.14}{-3.75}\hullline{61.14}{-2.58}\hullline{60.32}{-2.58}\hullfill
\hullface{keel!90!black}\hullmove{-63.94}{-2.48}\hullline{-62.38}{-2.58}\hullline{-62.38}{-1.32}\hullline{-63.94}{-1.28}\hullfill
\hullface{keel!84!black}\hullmove{-62.38}{-3.75}\hullline{-60.81}{-3.75}\hullline{-60.81}{-2.58}\hullline{-62.38}{-2.58}\hullfill
\hullface{keel!57!black}\hullmove{55.37}{-6.77}\hullline{56.00}{-6.77}\hullline{56.00}{-5.86}\hullline{55.37}{-5.86}\hullfill
\hullface{keel!72!black}\hullmove{-57.68}{-5.86}\hullline{-56.12}{-5.86}\hullline{-56.12}{-5.44}\hullline{-57.68}{-5.44}\hullfill
\hullface{topsides!69!black}\hullmove{-67.07}{0.78}\hullline{-65.51}{0.78}\hullline{-65.51}{1.56}\hullline{-67.07}{1.56}\hullfill
\hullface{keel!57!black}\hullmove{50.43}{-8.25}\hullline{51.25}{-8.25}\hullline{51.25}{-7.57}\hullline{50.43}{-7.57}\hullfill
\hullface{keel!72!black}\hullmove{61.97}{-1.32}\hullline{62.79}{-1.32}\hullline{62.79}{0.00}\hullline{61.97}{0.00}\hullfill
\hullface{keel!57!black}\hullmove{52.90}{-7.57}\hullline{53.73}{-7.57}\hullline{53.73}{-6.77}\hullline{52.90}{-6.77}\hullfill
\hullface{keel!69!black}\hullmove{61.14}{-2.58}\hullline{61.97}{-2.58}\hullline{61.97}{-1.32}\hullline{61.14}{-1.32}\hullfill
\hullface{keel!63!black}\hullmove{58.67}{-4.85}\hullline{59.49}{-4.85}\hullline{59.49}{-3.75}\hullline{58.67}{-3.75}\hullfill
\hullface{topsides!84!black}\hullmove{-65.51}{0.00}\hullline{-63.94}{0.00}\hullline{-63.94}{0.78}\hullline{-65.51}{0.78}\hullfill
\hullface{topsides!57!black}\hullmove{67.73}{3.04}\hullline{68.56}{3.04}\hullline{68.56}{3.90}\hullline{67.73}{3.90}\hullfill
\hullface{topsides!57!black}\hullmove{65.26}{1.52}\hullline{66.08}{1.52}\hullline{66.08}{3.04}\hullline{65.26}{3.04}\hullfill
\hullface{topsides!72!black}\hullmove{68.56}{3.90}\hullline{69.38}{3.90}\hullline{69.38}{4.55}\hullline{69.38}{6.07}\hullline{69.38}{7.59}\hullline{68.56}{7.59}\hullline{68.56}{6.07}\hullline{68.56}{4.55}\hullfill
\hullface{keel!57!black}\hullmove{-45.16}{-8.77}\hullline{-43.59}{-8.77}\hullline{-43.59}{-8.25}\hullline{-45.16}{-8.25}\hullfill
\hullface{topsides!60!black}\hullmove{62.79}{0.00}\hullline{63.61}{0.00}\hullline{63.61}{1.52}\hullline{62.79}{1.52}\hullfill
\hullface{keel!57!black}\hullmove{56.20}{-5.86}\hullline{57.02}{-5.86}\hullline{57.02}{-5.44}\hullline{56.20}{-5.44}\hullfill
\hullface{keel!60!black}\hullmove{57.02}{-5.44}\hullline{57.85}{-5.44}\hullline{57.85}{-4.85}\hullline{57.02}{-4.85}\hullfill
\hullface{keel!57!black}\hullmove{47.13}{-8.77}\hullline{47.96}{-8.77}\hullline{47.96}{-8.25}\hullline{47.13}{-8.25}\hullfill
\hullface{keel!57!black}\hullmove{54.55}{-6.77}\hullline{55.37}{-6.77}\hullline{55.37}{-5.86}\hullline{54.55}{-5.86}\hullfill
\hullface{topsides!57!black}\hullmove{-68.64}{1.56}\hullline{-67.07}{1.56}\hullline{-67.07}{2.34}\hullline{-68.64}{2.34}\hullfill
\hullface{keel!66!black}\hullmove{59.49}{-3.75}\hullline{60.32}{-3.75}\hullline{60.32}{-2.58}\hullline{59.49}{-2.58}\hullfill
\hullface{keel!57!black}\hullmove{-48.29}{-8.25}\hullline{-46.72}{-8.25}\hullline{-46.72}{-7.57}\hullline{-48.29}{-7.57}\hullfill
\hullface{keel!96!black}\hullmove{-63.94}{-1.28}\hullline{-62.38}{-1.32}\hullline{-62.38}{-0.00}\hullline{-63.94}{-0.00}\hullfill
\hullface{keel!66!black}\hullmove{-54.55}{-6.77}\hullline{-52.99}{-6.77}\hullline{-52.99}{-5.86}\hullline{-54.55}{-5.86}\hullfill
\hullface{keel!75!black}\hullmove{-57.68}{-5.44}\hullline{-56.12}{-5.44}\hullline{-56.12}{-4.85}\hullline{-57.68}{-4.85}\hullfill
\hullface{keel!57!black}\hullmove{52.08}{-7.57}\hullline{52.90}{-7.57}\hullline{52.90}{-6.77}\hullline{52.08}{-6.77}\hullfill
\hullface{keel!57!black}\hullmove{49.61}{-8.25}\hullline{50.43}{-8.25}\hullline{50.43}{-7.57}\hullline{49.61}{-7.57}\hullfill
\hullface{topsides!57!black}\hullmove{-71.77}{2.34}\hullline{-70.20}{2.34}\hullline{-70.20}{3.12}\hullline{-71.77}{3.12}\hullfill
\hullface{keel!60!black}\hullmove{-51.42}{-7.57}\hullline{-49.85}{-7.57}\hullline{-49.85}{-6.77}\hullline{-51.42}{-6.77}\hullfill
\hullface{keel!57!black}\hullmove{56.00}{-5.86}\hullline{56.20}{-5.86}\hullline{56.20}{-5.44}\hullline{56.00}{-5.44}\hullfill
\hullface{keel!78!black}\hullmove{-59.25}{-4.85}\hullline{-57.68}{-4.85}\hullline{-57.68}{-3.75}\hullline{-59.25}{-3.75}\hullfill
\hullface{keel!63!black}\hullmove{57.85}{-4.85}\hullline{58.67}{-4.85}\hullline{58.67}{-3.75}\hullline{57.85}{-3.75}\hullfill
\hullface{keel!75!black}\hullmove{61.14}{-1.32}\hullline{61.97}{-1.32}\hullline{61.97}{0.00}\hullline{61.14}{0.00}\hullfill
\hullface{keel!57!black}\hullmove{56.00}{-5.86}\hullline{56.00}{-5.86}\hullline{56.00}{-5.44}\hullline{56.00}{-5.44}\hullfill
\hullface{keel!69!black}\hullmove{60.32}{-2.58}\hullline{61.14}{-2.58}\hullline{61.14}{-1.32}\hullline{60.32}{-1.32}\hullfill
\hullface{keel!72!black}\hullmove{-56.12}{-5.86}\hullline{-54.55}{-5.86}\hullline{-54.55}{-5.44}\hullline{-56.12}{-5.44}\hullfill
\hullface{keel!90!black}\hullmove{-62.38}{-2.58}\hullline{-60.81}{-2.58}\hullline{-60.81}{-1.32}\hullline{-62.38}{-1.32}\hullfill
\hullface{keel!84!black}\hullmove{-60.81}{-3.75}\hullline{-59.25}{-3.75}\hullline{-59.25}{-2.58}\hullline{-60.81}{-2.58}\hullfill
\hullface{topsides!57!black}\hullmove{64.44}{1.52}\hullline{65.26}{1.52}\hullline{65.26}{3.04}\hullline{64.44}{3.04}\hullfill
\hullface{keel!60!black}\hullmove{56.20}{-5.44}\hullline{57.02}{-5.44}\hullline{57.02}{-4.85}\hullline{56.20}{-4.85}\hullfill
\hullface{keel!57!black}\hullmove{55.37}{-5.86}\hullline{56.00}{-5.86}\hullline{56.00}{-5.44}\hullline{55.37}{-5.44}\hullfill
\hullface{topsides!72!black}\hullmove{-65.51}{0.78}\hullline{-63.94}{0.78}\hullline{-63.94}{1.56}\hullline{-65.51}{1.56}\hullfill
\hullface{topsides!63!black}\hullmove{61.97}{0.00}\hullline{62.79}{0.00}\hullline{62.79}{1.52}\hullline{61.97}{1.52}\hullfill
\hullface{keel!57!black}\hullmove{53.73}{-6.77}\hullline{54.55}{-6.77}\hullline{54.55}{-5.86}\hullline{53.73}{-5.86}\hullfill
\hullface{topsides!57!black}\hullmove{66.91}{3.04}\hullline{67.73}{3.04}\hullline{67.73}{3.90}\hullline{66.91}{3.90}\hullfill
\hullface{topsides!87!black}\hullmove{-63.94}{0.00}\hullline{-62.38}{0.00}\hullline{-62.38}{0.78}\hullline{-63.94}{0.78}\hullfill
\hullface{keel!66!black}\hullmove{58.67}{-3.75}\hullline{59.49}{-3.75}\hullline{59.49}{-2.58}\hullline{58.67}{-2.58}\hullfill
\hullface{keel!57!black}\hullmove{46.31}{-8.77}\hullline{47.13}{-8.77}\hullline{47.13}{-8.25}\hullline{46.31}{-8.25}\hullfill
\hullface{keel!57!black}\hullmove{51.25}{-7.57}\hullline{52.08}{-7.57}\hullline{52.08}{-6.77}\hullline{51.25}{-6.77}\hullfill
\hullface{topsides!75!black}\hullmove{67.73}{3.90}\hullline{68.56}{3.90}\hullline{68.56}{4.55}\hullline{68.56}{6.07}\hullline{68.56}{7.59}\hullline{67.73}{7.59}\hullline{67.73}{6.07}\hullline{67.73}{4.55}\hullfill
\hullface{keel!57!black}\hullmove{48.78}{-8.25}\hullline{49.61}{-8.25}\hullline{49.61}{-7.57}\hullline{48.78}{-7.57}\hullfill
\hullface{keel!63!black}\hullmove{57.02}{-4.85}\hullline{57.85}{-4.85}\hullline{57.85}{-3.75}\hullline{57.02}{-3.75}\hullfill
\hullface{keel!57!black}\hullmove{-43.59}{-8.77}\hullline{-42.03}{-8.77}\hullline{-42.03}{-8.25}\hullline{-43.59}{-8.25}\hullfill
\hullface{keel!60!black}\hullmove{56.00}{-5.44}\hullline{56.20}{-5.44}\hullline{56.20}{-4.85}\hullline{56.00}{-4.85}\hullfill
\hullface{keel!69!black}\hullmove{-52.99}{-6.77}\hullline{-51.42}{-6.77}\hullline{-51.42}{-5.86}\hullline{-52.99}{-5.86}\hullfill
\hullface{keel!60!black}\hullmove{56.00}{-5.44}\hullline{56.00}{-5.44}\hullline{56.00}{-4.85}\hullline{56.00}{-4.85}\hullfill
\hullface{topsides!60!black}\hullmove{-67.07}{1.56}\hullline{-65.51}{1.56}\hullline{-65.51}{2.34}\hullline{-67.07}{2.34}\hullfill
\hullface{keel!69!black}\hullmove{59.49}{-2.58}\hullline{60.32}{-2.58}\hullline{60.32}{-1.32}\hullline{59.49}{-1.32}\hullfill
\hullface{keel!75!black}\hullmove{60.32}{-1.32}\hullline{61.14}{-1.32}\hullline{61.14}{0.00}\hullline{60.32}{0.00}\hullfill
\hullface{keel!75!black}\hullmove{-56.12}{-5.44}\hullline{-54.55}{-5.44}\hullline{-54.55}{-4.85}\hullline{-56.12}{-4.85}\hullfill
\hullface{keel!57!black}\hullmove{54.55}{-5.86}\hullline{55.37}{-5.86}\hullline{55.37}{-5.44}\hullline{54.55}{-5.44}\hullfill
\hullface{keel!57!black}\hullmove{-46.72}{-8.25}\hullline{-45.16}{-8.25}\hullline{-45.16}{-7.57}\hullline{-46.72}{-7.57}\hullfill
\hullface{keel!60!black}\hullmove{-49.85}{-7.57}\hullline{-48.29}{-7.57}\hullline{-48.29}{-6.77}\hullline{-49.85}{-6.77}\hullfill
\hullface{keel!96!black}\hullmove{-62.38}{-1.32}\hullline{-60.81}{-1.32}\hullline{-60.81}{-0.00}\hullline{-62.38}{-0.00}\hullfill
\hullface{keel!78!black}\hullmove{-57.68}{-4.85}\hullline{-56.12}{-4.85}\hullline{-56.12}{-3.75}\hullline{-57.68}{-3.75}\hullfill
\hullface{topsides!57!black}\hullmove{-70.20}{2.34}\hullline{-68.64}{2.34}\hullline{-68.64}{3.12}\hullline{-70.20}{3.12}\hullfill
\hullface{keel!72!black}\hullmove{-54.55}{-5.86}\hullline{-52.99}{-5.86}\hullline{-52.99}{-5.44}\hullline{-54.55}{-5.44}\hullfill
\hullface{keel!57!black}\hullmove{52.90}{-6.77}\hullline{53.73}{-6.77}\hullline{53.73}{-5.86}\hullline{52.90}{-5.86}\hullfill
\hullface{keel!60!black}\hullmove{55.37}{-5.44}\hullline{56.00}{-5.44}\hullline{56.00}{-4.85}\hullline{55.37}{-4.85}\hullfill
\hullface{topsides!57!black}\hullmove{63.61}{1.52}\hullline{64.44}{1.52}\hullline{64.44}{3.04}\hullline{63.61}{3.04}\hullfill
\hullface{topsides!66!black}\hullmove{61.14}{0.00}\hullline{61.97}{0.00}\hullline{61.97}{1.52}\hullline{61.14}{1.52}\hullfill
\hullface{topsides!57!black}\hullmove{-74.90}{3.29}\hullline{-73.33}{3.12}\hullline{-73.33}{3.90}\hullline{-74.90}{3.90}\hullfill
\hullface{keel!84!black}\hullmove{-59.25}{-3.75}\hullline{-57.68}{-3.75}\hullline{-57.68}{-2.58}\hullline{-59.25}{-2.58}\hullfill
\hullface{keel!66!black}\hullmove{57.85}{-3.75}\hullline{58.67}{-3.75}\hullline{58.67}{-2.58}\hullline{57.85}{-2.58}\hullfill
\hullface{keel!90!black}\hullmove{-60.81}{-2.58}\hullline{-59.25}{-2.58}\hullline{-59.25}{-1.32}\hullline{-60.81}{-1.32}\hullfill
\hullface{keel!57!black}\hullmove{50.43}{-7.57}\hullline{51.25}{-7.57}\hullline{51.25}{-6.77}\hullline{50.43}{-6.77}\hullfill
\hullface{keel!63!black}\hullmove{56.20}{-4.85}\hullline{57.02}{-4.85}\hullline{57.02}{-3.75}\hullline{56.20}{-3.75}\hullfill
\hullface{topsides!75!black}\hullmove{-63.94}{0.78}\hullline{-62.38}{0.78}\hullline{-62.38}{1.56}\hullline{-63.94}{1.56}\hullfill
\hullface{topsides!57!black}\hullmove{66.08}{3.04}\hullline{66.91}{3.04}\hullline{66.91}{3.90}\hullline{66.08}{3.90}\hullfill
\hullface{keel!57!black}\hullmove{47.96}{-8.25}\hullline{48.78}{-8.25}\hullline{48.78}{-7.57}\hullline{47.96}{-7.57}\hullfill
\hullface{topsides!90!black}\hullmove{-62.38}{0.00}\hullline{-60.81}{0.00}\hullline{-60.81}{0.78}\hullline{-62.38}{0.78}\hullfill
\hullface{keel!57!black}\hullmove{45.49}{-8.77}\hullline{46.31}{-8.77}\hullline{46.31}{-8.25}\hullline{45.49}{-8.25}\hullfill
\hullface{keel!60!black}\hullmove{53.73}{-5.86}\hullline{54.55}{-5.86}\hullline{54.55}{-5.44}\hullline{53.73}{-5.44}\hullfill
\hullface{keel!72!black}\hullmove{58.67}{-2.58}\hullline{59.49}{-2.58}\hullline{59.49}{-1.32}\hullline{58.67}{-1.32}\hullfill
\hullface{keel!78!black}\hullmove{59.49}{-1.32}\hullline{60.32}{-1.32}\hullline{60.32}{0.00}\hullline{59.49}{0.00}\hullfill
\hullface{keel!60!black}\hullmove{54.55}{-5.44}\hullline{55.37}{-5.44}\hullline{55.37}{-4.85}\hullline{54.55}{-4.85}\hullfill
\hullface{keel!57!black}\hullmove{52.08}{-6.77}\hullline{52.90}{-6.77}\hullline{52.90}{-5.86}\hullline{52.08}{-5.86}\hullfill
\hullface{keel!69!black}\hullmove{-51.42}{-6.77}\hullline{-49.85}{-6.77}\hullline{-49.85}{-5.86}\hullline{-51.42}{-5.86}\hullfill
\hullface{topsides!78!black}\hullmove{66.91}{3.90}\hullline{67.73}{3.90}\hullline{67.73}{4.55}\hullline{67.73}{6.07}\hullline{67.73}{7.59}\hullline{66.91}{7.59}\hullline{66.91}{6.07}\hullline{66.91}{4.55}\hullfill
\hullface{keel!75!black}\hullmove{-54.55}{-5.44}\hullline{-52.99}{-5.44}\hullline{-52.99}{-4.85}\hullline{-54.55}{-4.85}\hullfill
\hullface{topsides!63!black}\hullmove{-65.51}{1.56}\hullline{-63.94}{1.56}\hullline{-63.94}{2.34}\hullline{-65.51}{2.34}\hullfill
\hullface{keel!66!black}\hullmove{56.00}{-4.85}\hullline{56.20}{-4.85}\hullline{56.20}{-3.75}\hullline{56.00}{-3.75}\hullfill
\hullface{topsides!57!black}\hullmove{62.79}{1.52}\hullline{63.61}{1.52}\hullline{63.61}{3.04}\hullline{62.79}{3.04}\hullfill
\hullface{topsides!69!black}\hullmove{60.32}{0.00}\hullline{61.14}{0.00}\hullline{61.14}{1.52}\hullline{60.32}{1.52}\hullfill
\hullface{keel!66!black}\hullmove{56.00}{-4.85}\hullline{56.00}{-4.85}\hullline{56.00}{-3.75}\hullline{56.00}{-3.75}\hullfill
\hullface{keel!81!black}\hullmove{-56.12}{-4.85}\hullline{-54.55}{-4.85}\hullline{-54.55}{-3.75}\hullline{-56.12}{-3.75}\hullfill
\hullface{keel!69!black}\hullmove{57.02}{-3.75}\hullline{57.85}{-3.75}\hullline{57.85}{-2.58}\hullline{57.02}{-2.58}\hullfill
\hullface{keel!63!black}\hullmove{-48.29}{-7.57}\hullline{-46.72}{-7.57}\hullline{-46.72}{-6.77}\hullline{-48.29}{-6.77}\hullfill
\hullface{topsides!57!black}\hullmove{-68.64}{2.34}\hullline{-67.07}{2.34}\hullline{-67.07}{3.12}\hullline{-68.64}{3.12}\hullfill
\hullface{keel!72!black}\hullmove{-52.99}{-5.86}\hullline{-51.42}{-5.86}\hullline{-51.42}{-5.44}\hullline{-52.99}{-5.44}\hullfill
\hullface{keel!96!black}\hullmove{-60.81}{-1.32}\hullline{-59.25}{-1.32}\hullline{-59.25}{-0.00}\hullline{-60.81}{-0.00}\hullfill
\hullface{keel!57!black}\hullmove{49.61}{-7.57}\hullline{50.43}{-7.57}\hullline{50.43}{-6.77}\hullline{49.61}{-6.77}\hullfill
\hullface{keel!57!black}\hullmove{-45.16}{-8.25}\hullline{-43.59}{-8.25}\hullline{-43.59}{-7.57}\hullline{-45.16}{-7.57}\hullfill
\hullface{keel!57!black}\hullmove{-42.03}{-8.77}\hullline{-40.46}{-8.77}\hullline{-40.46}{-8.25}\hullline{-42.03}{-8.25}\hullfill
\hullface{keel!66!black}\hullmove{55.37}{-4.85}\hullline{56.00}{-4.85}\hullline{56.00}{-3.75}\hullline{55.37}{-3.75}\hullfill
\hullface{topsides!57!black}\hullmove{-73.33}{3.12}\hullline{-71.77}{3.12}\hullline{-71.77}{3.90}\hullline{-73.33}{3.90}\hullfill
\hullface{keel!84!black}\hullmove{-57.68}{-3.75}\hullline{-56.12}{-3.75}\hullline{-56.12}{-2.58}\hullline{-57.68}{-2.58}\hullfill
\hullface{keel!60!black}\hullmove{52.90}{-5.86}\hullline{53.73}{-5.86}\hullline{53.73}{-5.44}\hullline{52.90}{-5.44}\hullfill
\hullface{keel!90!black}\hullmove{-59.25}{-2.58}\hullline{-57.68}{-2.58}\hullline{-57.68}{-1.32}\hullline{-59.25}{-1.32}\hullfill
\hullface{keel!72!black}\hullmove{57.85}{-2.58}\hullline{58.67}{-2.58}\hullline{58.67}{-1.32}\hullline{57.85}{-1.32}\hullfill
\hullface{keel!63!black}\hullmove{53.73}{-5.44}\hullline{54.55}{-5.44}\hullline{54.55}{-4.85}\hullline{53.73}{-4.85}\hullfill
\hullface{keel!57!black}\hullmove{47.13}{-8.25}\hullline{47.96}{-8.25}\hullline{47.96}{-7.57}\hullline{47.13}{-7.57}\hullfill
\hullface{keel!78!black}\hullmove{58.67}{-1.32}\hullline{59.49}{-1.32}\hullline{59.49}{0.00}\hullline{58.67}{0.00}\hullfill
\hullface{topsides!57!black}\hullmove{65.26}{3.04}\hullline{66.08}{3.04}\hullline{66.08}{3.90}\hullline{65.26}{3.90}\hullfill
\hullface{topsides!78!black}\hullmove{-62.38}{0.78}\hullline{-60.81}{0.78}\hullline{-60.81}{1.56}\hullline{-62.38}{1.56}\hullfill
\hullface{keel!57!black}\hullmove{51.25}{-6.77}\hullline{52.08}{-6.77}\hullline{52.08}{-5.86}\hullline{51.25}{-5.86}\hullfill
\hullface{topsides!72!black}\hullmove{59.49}{0.00}\hullline{60.32}{0.00}\hullline{60.32}{1.52}\hullline{59.49}{1.52}\hullfill
\hullface{topsides!57!black}\hullmove{61.97}{1.52}\hullline{62.79}{1.52}\hullline{62.79}{3.04}\hullline{61.97}{3.04}\hullfill
\hullface{topsides!93!black}\hullmove{-60.81}{0.00}\hullline{-59.25}{0.00}\hullline{-59.25}{0.78}\hullline{-60.81}{0.78}\hullfill
\hullface{keel!57!black}\hullmove{44.66}{-8.77}\hullline{45.49}{-8.77}\hullline{45.49}{-8.25}\hullline{44.66}{-8.25}\hullfill
\hullface{ship!99!black}\hullmove{-42.00}{2.90}\hullline{-26.93}{2.90}\hullline{-26.93}{6.59}\hullline{-42.00}{6.59}\hullfill
\hullface{ship!99!black}\hullmove{-26.37}{2.90}\hullline{-11.29}{2.90}\hullline{-11.29}{6.59}\hullline{-26.37}{6.59}\hullfill
\hullface{ship!99!black}\hullmove{-10.73}{2.90}\hullline{4.34}{2.90}\hullline{4.34}{6.59}\hullline{-10.73}{6.59}\hullfill
\hullface{ship!99!black}\hullmove{4.90}{2.90}\hullline{19.97}{2.90}\hullline{19.97}{6.59}\hullline{4.90}{6.59}\hullfill
\hullface{ship!99!black}\hullmove{20.53}{2.90}\hullline{35.61}{2.90}\hullline{35.61}{6.59}\hullline{20.53}{6.59}\hullfill
\hullface{ship!99!black}\hullmove{36.17}{2.90}\hullline{51.24}{2.90}\hullline{51.24}{6.59}\hullline{36.17}{6.59}\hullfill
\hullface{keel!69!black}\hullmove{56.20}{-3.75}\hullline{57.02}{-3.75}\hullline{57.02}{-2.58}\hullline{56.20}{-2.58}\hullfill
\hullface{topsides!69!black}\hullmove{-63.94}{1.56}\hullline{-62.38}{1.56}\hullline{-62.38}{2.34}\hullline{-63.94}{2.34}\hullfill
\hullface{keel!69!black}\hullmove{-49.85}{-6.77}\hullline{-48.29}{-6.77}\hullline{-48.29}{-5.86}\hullline{-49.85}{-5.86}\hullfill
\hullface{keel!66!black}\hullmove{54.55}{-4.85}\hullline{55.37}{-4.85}\hullline{55.37}{-3.75}\hullline{54.55}{-3.75}\hullfill
\hullface{keel!78!black}\hullmove{-52.99}{-5.44}\hullline{-51.42}{-5.44}\hullline{-51.42}{-4.85}\hullline{-52.99}{-4.85}\hullfill
\hullface{topsides!81!black}\hullmove{66.08}{3.90}\hullline{66.91}{3.90}\hullline{66.91}{4.55}\hullline{66.91}{6.07}\hullline{66.91}{7.59}\hullline{66.08}{7.59}\hullline{66.08}{6.07}\hullline{66.08}{4.55}\hullfill
\hullface{keel!57!black}\hullmove{48.78}{-7.57}\hullline{49.61}{-7.57}\hullline{49.61}{-6.77}\hullline{48.78}{-6.77}\hullfill
\hullface{keel!81!black}\hullmove{-54.55}{-4.85}\hullline{-52.99}{-4.85}\hullline{-52.99}{-3.75}\hullline{-54.55}{-3.75}\hullfill
\hullface{topsides!57!black}\hullmove{-67.07}{2.34}\hullline{-65.51}{2.34}\hullline{-65.51}{3.12}\hullline{-67.07}{3.12}\hullfill
\hullface{keel!60!black}\hullmove{52.08}{-5.86}\hullline{52.90}{-5.86}\hullline{52.90}{-5.44}\hullline{52.08}{-5.44}\hullfill
\hullface{keel!75!black}\hullmove{-51.42}{-5.86}\hullline{-49.85}{-5.86}\hullline{-49.85}{-5.44}\hullline{-51.42}{-5.44}\hullfill
\hullface{keel!63!black}\hullmove{52.90}{-5.44}\hullline{53.73}{-5.44}\hullline{53.73}{-4.85}\hullline{52.90}{-4.85}\hullfill
\hullface{keel!75!black}\hullmove{57.02}{-2.58}\hullline{57.85}{-2.58}\hullline{57.85}{-1.32}\hullline{57.02}{-1.32}\hullfill
\hullface{keel!63!black}\hullmove{-46.72}{-7.57}\hullline{-45.16}{-7.57}\hullline{-45.16}{-6.77}\hullline{-46.72}{-6.77}\hullfill
\hullface{keel!99!black}\hullmove{-59.25}{-1.32}\hullline{-57.68}{-1.32}\hullline{-57.68}{-0.00}\hullline{-59.25}{-0.00}\hullfill
\hullface{keel!81!black}\hullmove{57.85}{-1.32}\hullline{58.67}{-1.32}\hullline{58.67}{0.00}\hullline{57.85}{0.00}\hullfill
\hullface{keel!72!black}\hullmove{56.00}{-3.75}\hullline{56.20}{-3.75}\hullline{56.20}{-2.58}\hullline{56.00}{-2.58}\hullfill
\hullface{keel!57!black}\hullmove{50.43}{-6.77}\hullline{51.25}{-6.77}\hullline{51.25}{-5.86}\hullline{50.43}{-5.86}\hullfill
\hullface{keel!72!black}\hullmove{56.00}{-3.75}\hullline{56.00}{-3.75}\hullline{56.00}{-2.58}\hullline{56.00}{-2.58}\hullfill
\hullface{keel!87!black}\hullmove{-56.12}{-3.75}\hullline{-54.55}{-3.75}\hullline{-54.55}{-2.58}\hullline{-56.12}{-2.58}\hullfill
\hullface{topsides!57!black}\hullmove{64.44}{3.04}\hullline{65.26}{3.04}\hullline{65.26}{3.90}\hullline{64.44}{3.90}\hullfill
\hullface{topsides!57!black}\hullmove{-71.77}{3.12}\hullline{-70.20}{3.12}\hullline{-70.20}{3.90}\hullline{-71.77}{3.90}\hullfill
\hullface{keel!57!black}\hullmove{46.31}{-8.25}\hullline{47.13}{-8.25}\hullline{47.13}{-7.57}\hullline{46.31}{-7.57}\hullfill
\hullface{topsides!75!black}\hullmove{58.67}{0.00}\hullline{59.49}{0.00}\hullline{59.49}{1.52}\hullline{58.67}{1.52}\hullfill
\hullface{topsides!60!black}\hullmove{61.14}{1.52}\hullline{61.97}{1.52}\hullline{61.97}{3.04}\hullline{61.14}{3.04}\hullfill
\hullface{keel!90!black}\hullmove{-57.68}{-2.58}\hullline{-56.12}{-2.58}\hullline{-56.12}{-1.32}\hullline{-57.68}{-1.32}\hullfill
\hullface{topsides!81!black}\hullmove{-60.81}{0.78}\hullline{-59.25}{0.78}\hullline{-59.25}{1.56}\hullline{-60.81}{1.56}\hullfill
\hullface{keel!57!black}\hullmove{-43.59}{-8.25}\hullline{-42.03}{-8.25}\hullline{-42.03}{-7.57}\hullline{-43.59}{-7.57}\hullfill
\hullface{keel!72!black}\hullmove{55.37}{-3.75}\hullline{56.00}{-3.75}\hullline{56.00}{-2.58}\hullline{55.37}{-2.58}\hullfill
\hullface{keel!69!black}\hullmove{53.73}{-4.85}\hullline{54.55}{-4.85}\hullline{54.55}{-3.75}\hullline{53.73}{-3.75}\hullfill
\hullface{topsides!93!black}\hullmove{-59.25}{0.00}\hullline{-57.68}{0.00}\hullline{-57.68}{0.78}\hullline{-59.25}{0.78}\hullfill
\hullface{keel!57!black}\hullmove{47.96}{-7.57}\hullline{48.78}{-7.57}\hullline{48.78}{-6.77}\hullline{47.96}{-6.77}\hullfill
\hullface{topsides!72!black}\hullmove{-62.38}{1.56}\hullline{-60.81}{1.56}\hullline{-60.81}{2.34}\hullline{-62.38}{2.34}\hullfill
\hullface{keel!63!black}\hullmove{51.25}{-5.86}\hullline{52.08}{-5.86}\hullline{52.08}{-5.44}\hullline{51.25}{-5.44}\hullfill
\hullface{keel!57!black}\hullmove{-40.46}{-8.77}\hullline{-38.90}{-8.77}\hullline{-38.90}{-8.25}\hullline{-40.46}{-8.25}\hullfill
\hullface{keel!66!black}\hullmove{52.08}{-5.44}\hullline{52.90}{-5.44}\hullline{52.90}{-4.85}\hullline{52.08}{-4.85}\hullfill
\hullface{keel!75!black}\hullmove{56.20}{-2.58}\hullline{57.02}{-2.58}\hullline{57.02}{-1.32}\hullline{56.20}{-1.32}\hullfill
\hullface{topsides!84!black}\hullmove{65.26}{3.90}\hullline{66.08}{3.90}\hullline{66.08}{4.55}\hullline{66.08}{6.07}\hullline{66.08}{7.59}\hullline{65.26}{7.59}\hullline{65.26}{6.07}\hullline{65.26}{4.55}\hullfill
\hullface{keel!78!black}\hullmove{-51.42}{-5.44}\hullline{-49.85}{-5.44}\hullline{-49.85}{-4.85}\hullline{-51.42}{-4.85}\hullfill
\hullface{topsides!60!black}\hullmove{-65.51}{2.34}\hullline{-63.94}{2.34}\hullline{-63.94}{3.12}\hullline{-65.51}{3.12}\hullfill
\hullface{keel!84!black}\hullmove{57.02}{-1.32}\hullline{57.85}{-1.32}\hullline{57.85}{0.00}\hullline{57.02}{0.00}\hullfill
\hullface{keel!57!black}\hullmove{43.84}{-8.77}\hullline{44.66}{-8.77}\hullline{44.66}{-8.25}\hullline{43.84}{-8.25}\hullfill
\hullface{keel!72!black}\hullmove{-48.29}{-6.77}\hullline{-46.72}{-6.77}\hullline{-46.72}{-5.86}\hullline{-48.29}{-5.86}\hullfill
\hullface{keel!60!black}\hullmove{49.61}{-6.77}\hullline{50.43}{-6.77}\hullline{50.43}{-5.86}\hullline{49.61}{-5.86}\hullfill
\hullface{keel!84!black}\hullmove{-52.99}{-4.85}\hullline{-51.42}{-4.85}\hullline{-51.42}{-3.75}\hullline{-52.99}{-3.75}\hullfill
\hullface{topsides!63!black}\hullmove{60.32}{1.52}\hullline{61.14}{1.52}\hullline{61.14}{3.04}\hullline{60.32}{3.04}\hullfill
\hullface{topsides!75!black}\hullmove{57.85}{0.00}\hullline{58.67}{0.00}\hullline{58.67}{1.52}\hullline{57.85}{1.52}\hullfill
\hullface{keel!78!black}\hullmove{-49.85}{-5.86}\hullline{-48.29}{-5.86}\hullline{-48.29}{-5.44}\hullline{-49.85}{-5.44}\hullfill
\hullface{topsides!57!black}\hullmove{63.61}{3.04}\hullline{64.44}{3.04}\hullline{64.44}{3.90}\hullline{63.61}{3.90}\hullfill
\hullface{keel!99!black}\hullmove{-57.68}{-1.32}\hullline{-56.12}{-1.32}\hullline{-56.12}{-0.00}\hullline{-57.68}{-0.00}\hullfill
\hullface{keel!72!black}\hullmove{54.55}{-3.75}\hullline{55.37}{-3.75}\hullline{55.37}{-2.58}\hullline{54.55}{-2.58}\hullfill
\hullface{keel!69!black}\hullmove{52.90}{-4.85}\hullline{53.73}{-4.85}\hullline{53.73}{-3.75}\hullline{52.90}{-3.75}\hullfill
\hullface{topsides!57!black}\hullmove{-70.20}{3.12}\hullline{-68.64}{3.12}\hullline{-68.64}{3.90}\hullline{-70.20}{3.90}\hullfill
\hullface{keel!78!black}\hullmove{56.00}{-2.58}\hullline{56.20}{-2.58}\hullline{56.20}{-1.32}\hullline{56.00}{-1.32}\hullfill
\hullface{keel!57!black}\hullmove{45.49}{-8.25}\hullline{46.31}{-8.25}\hullline{46.31}{-7.57}\hullline{45.49}{-7.57}\hullfill
\hullface{keel!69!black}\hullmove{-45.16}{-7.57}\hullline{-43.59}{-7.57}\hullline{-43.59}{-6.77}\hullline{-45.16}{-6.77}\hullfill
\hullface{keel!87!black}\hullmove{-54.55}{-3.75}\hullline{-52.99}{-3.75}\hullline{-52.99}{-2.58}\hullline{-54.55}{-2.58}\hullfill
\hullface{topsides!84!black}\hullmove{-59.25}{0.78}\hullline{-57.68}{0.78}\hullline{-57.68}{1.56}\hullline{-59.25}{1.56}\hullfill
\hullface{keel!78!black}\hullmove{56.00}{-2.58}\hullline{56.00}{-2.58}\hullline{56.00}{-1.32}\hullline{56.00}{-1.32}\hullfill
\hullface{keel!93!black}\hullmove{-56.12}{-2.58}\hullline{-54.55}{-2.58}\hullline{-54.55}{-1.32}\hullline{-56.12}{-1.32}\hullfill
\hullface{keel!66!black}\hullmove{50.43}{-5.86}\hullline{51.25}{-5.86}\hullline{51.25}{-5.44}\hullline{50.43}{-5.44}\hullfill
\hullface{keel!66!black}\hullmove{51.25}{-5.44}\hullline{52.08}{-5.44}\hullline{52.08}{-4.85}\hullline{51.25}{-4.85}\hullfill
\hullface{keel!57!black}\hullmove{47.13}{-7.57}\hullline{47.96}{-7.57}\hullline{47.96}{-6.77}\hullline{47.13}{-6.77}\hullfill
\hullface{keel!84!black}\hullmove{56.20}{-1.32}\hullline{57.02}{-1.32}\hullline{57.02}{0.00}\hullline{56.20}{0.00}\hullfill
\hullface{topsides!75!black}\hullmove{-60.81}{1.56}\hullline{-59.25}{1.56}\hullline{-59.25}{2.34}\hullline{-60.81}{2.34}\hullfill
\hullface{topsides!96!black}\hullmove{-57.68}{0.00}\hullline{-56.12}{0.00}\hullline{-56.12}{0.78}\hullline{-57.68}{0.78}\hullfill
\hullface{keel!78!black}\hullmove{55.37}{-2.58}\hullline{56.00}{-2.58}\hullline{56.00}{-1.32}\hullline{55.37}{-1.32}\hullfill
\hullface{topsides!66!black}\hullmove{59.49}{1.52}\hullline{60.32}{1.52}\hullline{60.32}{3.04}\hullline{59.49}{3.04}\hullfill
\hullface{topsides!63!black}\hullmove{-63.94}{2.34}\hullline{-62.38}{2.34}\hullline{-62.38}{3.12}\hullline{-63.94}{3.12}\hullfill
\hullface{topsides!78!black}\hullmove{57.02}{0.00}\hullline{57.85}{0.00}\hullline{57.85}{1.52}\hullline{57.02}{1.52}\hullfill
\hullface{topsides!87!black}\hullmove{64.44}{3.90}\hullline{65.26}{3.90}\hullline{65.26}{4.55}\hullline{65.26}{6.07}\hullline{65.26}{7.59}\hullline{64.44}{7.59}\hullline{64.44}{6.07}\hullline{64.44}{4.55}\hullfill
\hullface{keel!63!black}\hullmove{48.78}{-6.77}\hullline{49.61}{-6.77}\hullline{49.61}{-5.86}\hullline{48.78}{-5.86}\hullfill
\hullface{ship!99!black}\hullmove{-67.20}{2.90}\hullline{-56.00}{2.90}\hullline{-56.00}{16.57}\hullline{-67.20}{16.57}\hullfill
\hullface{keel!60!black}\hullmove{-42.03}{-8.25}\hullline{-40.46}{-8.25}\hullline{-40.46}{-7.57}\hullline{-42.03}{-7.57}\hullfill
\hullface{keel!75!black}\hullmove{53.73}{-3.75}\hullline{54.55}{-3.75}\hullline{54.55}{-2.58}\hullline{53.73}{-2.58}\hullfill
\hullface{topsides!57!black}\hullmove{62.79}{3.04}\hullline{63.61}{3.04}\hullline{63.61}{3.90}\hullline{62.79}{3.90}\hullfill
\hullface{keel!81!black}\hullmove{-49.85}{-5.44}\hullline{-48.29}{-5.44}\hullline{-48.29}{-4.85}\hullline{-49.85}{-4.85}\hullfill
\hullface{keel!72!black}\hullmove{52.08}{-4.85}\hullline{52.90}{-4.85}\hullline{52.90}{-3.75}\hullline{52.08}{-3.75}\hullfill
\hullface{keel!84!black}\hullmove{-51.42}{-4.85}\hullline{-49.85}{-4.85}\hullline{-49.85}{-3.75}\hullline{-51.42}{-3.75}\hullfill
\hullface{keel!75!black}\hullmove{-46.72}{-6.77}\hullline{-45.16}{-6.77}\hullline{-45.16}{-5.86}\hullline{-46.72}{-5.86}\hullfill
\hullface{keel!87!black}\hullmove{56.00}{-1.32}\hullline{56.20}{-1.32}\hullline{56.20}{0.00}\hullline{56.00}{0.00}\hullfill
\hullface{keel!87!black}\hullmove{56.00}{-1.32}\hullline{56.00}{-1.32}\hullline{56.00}{0.00}\hullline{56.00}{0.00}\hullfill
\hullface{keel!99!black}\hullmove{-56.12}{-1.32}\hullline{-54.55}{-1.32}\hullline{-54.55}{-0.00}\hullline{-56.12}{-0.00}\hullfill
\hullface{keel!57!black}\hullmove{43.02}{-8.77}\hullline{43.84}{-8.77}\hullline{43.84}{-8.25}\hullline{43.02}{-8.25}\hullfill
\hullface{keel!81!black}\hullmove{-48.29}{-5.86}\hullline{-46.72}{-5.86}\hullline{-46.72}{-5.44}\hullline{-48.29}{-5.44}\hullfill
\hullface{topsides!57!black}\hullmove{-68.64}{3.12}\hullline{-67.07}{3.12}\hullline{-67.07}{3.90}\hullline{-68.64}{3.90}\hullfill
\hullface{keel!81!black}\hullmove{54.55}{-2.58}\hullline{55.37}{-2.58}\hullline{55.37}{-1.32}\hullline{54.55}{-1.32}\hullfill
\hullface{keel!66!black}\hullmove{49.61}{-5.86}\hullline{50.43}{-5.86}\hullline{50.43}{-5.44}\hullline{49.61}{-5.44}\hullfill
\hullface{topsides!87!black}\hullmove{-57.68}{0.78}\hullline{-56.12}{0.78}\hullline{-56.12}{1.56}\hullline{-57.68}{1.56}\hullfill
\hullface{keel!69!black}\hullmove{50.43}{-5.44}\hullline{51.25}{-5.44}\hullline{51.25}{-4.85}\hullline{50.43}{-4.85}\hullfill
\hullface{topsides!69!black}\hullmove{58.67}{1.52}\hullline{59.49}{1.52}\hullline{59.49}{3.04}\hullline{58.67}{3.04}\hullfill
\hullface{keel!90!black}\hullmove{-52.99}{-3.75}\hullline{-51.42}{-3.75}\hullline{-51.42}{-2.58}\hullline{-52.99}{-2.58}\hullfill
\hullface{keel!87!black}\hullmove{55.37}{-1.32}\hullline{56.00}{-1.32}\hullline{56.00}{0.00}\hullline{55.37}{0.00}\hullfill
\hullface{keel!57!black}\hullmove{44.66}{-8.25}\hullline{45.49}{-8.25}\hullline{45.49}{-7.57}\hullline{44.66}{-7.57}\hullfill
\hullface{topsides!81!black}\hullmove{56.20}{0.00}\hullline{57.02}{0.00}\hullline{57.02}{1.52}\hullline{56.20}{1.52}\hullfill
\hullface{keel!93!black}\hullmove{-54.55}{-2.58}\hullline{-52.99}{-2.58}\hullline{-52.99}{-1.32}\hullline{-54.55}{-1.32}\hullfill
\hullface{topsides!81!black}\hullmove{-59.25}{1.56}\hullline{-57.68}{1.56}\hullline{-57.68}{2.34}\hullline{-59.25}{2.34}\hullfill
\hullface{keel!60!black}\hullmove{46.31}{-7.57}\hullline{47.13}{-7.57}\hullline{47.13}{-6.77}\hullline{46.31}{-6.77}\hullfill
\hullface{topsides!66!black}\hullmove{-62.38}{2.34}\hullline{-60.81}{2.34}\hullline{-60.81}{3.12}\hullline{-62.38}{3.12}\hullfill
\hullface{topsides!96!black}\hullmove{-56.12}{0.00}\hullline{-54.55}{0.00}\hullline{-54.55}{0.78}\hullline{-56.12}{0.78}\hullfill
\hullface{keel!66!black}\hullmove{47.96}{-6.77}\hullline{48.78}{-6.77}\hullline{48.78}{-5.86}\hullline{47.96}{-5.86}\hullfill
\hullface{keel!78!black}\hullmove{52.90}{-3.75}\hullline{53.73}{-3.75}\hullline{53.73}{-2.58}\hullline{52.90}{-2.58}\hullfill
\hullface{keel!72!black}\hullmove{-43.59}{-7.57}\hullline{-42.03}{-7.57}\hullline{-42.03}{-6.77}\hullline{-43.59}{-6.77}\hullfill
\hullface{topsides!57!black}\hullmove{61.97}{3.04}\hullline{62.79}{3.04}\hullline{62.79}{3.90}\hullline{61.97}{3.90}\hullfill
\hullface{topsides!90!black}\hullmove{63.61}{3.90}\hullline{64.44}{3.90}\hullline{64.44}{4.55}\hullline{64.44}{6.07}\hullline{64.44}{7.59}\hullline{63.61}{7.59}\hullline{63.61}{6.07}\hullline{63.61}{4.55}\hullfill
\hullface{keel!57!black}\hullmove{-38.90}{-8.77}\hullline{-37.33}{-8.77}\hullline{-37.33}{-8.25}\hullline{-38.90}{-8.25}\hullfill
\hullface{keel!75!black}\hullmove{51.25}{-4.85}\hullline{52.08}{-4.85}\hullline{52.08}{-3.75}\hullline{51.25}{-3.75}\hullfill
\hullface{topsides!84!black}\hullmove{55.37}{0.00}\hullline{56.00}{0.00}\hullline{56.00}{0.78}\hullline{55.37}{0.78}\hullfill
\hullface{topsides!84!black}\hullmove{56.00}{0.00}\hullline{56.20}{0.00}\hullline{56.20}{1.52}\hullline{56.00}{1.52}\hullfill
\hullface{topsides!72!black}\hullmove{57.85}{1.52}\hullline{58.67}{1.52}\hullline{58.67}{3.04}\hullline{57.85}{3.04}\hullfill
\hullface{keel!87!black}\hullmove{54.55}{-1.32}\hullline{55.37}{-1.32}\hullline{55.37}{0.00}\hullline{54.55}{0.00}\hullfill
\hullface{keel!84!black}\hullmove{53.73}{-2.58}\hullline{54.55}{-2.58}\hullline{54.55}{-1.32}\hullline{53.73}{-1.32}\hullfill
\hullface{keel!84!black}\hullmove{-48.29}{-5.44}\hullline{-46.72}{-5.44}\hullline{-46.72}{-4.85}\hullline{-48.29}{-4.85}\hullfill
\hullface{topsides!57!black}\hullmove{-67.07}{3.12}\hullline{-65.51}{3.12}\hullline{-65.51}{3.90}\hullline{-67.07}{3.90}\hullfill
\hullface{keel!87!black}\hullmove{-49.85}{-4.85}\hullline{-48.29}{-4.85}\hullline{-48.29}{-3.75}\hullline{-49.85}{-3.75}\hullfill
\hullface{keel!69!black}\hullmove{48.78}{-5.86}\hullline{49.61}{-5.86}\hullline{49.61}{-5.44}\hullline{48.78}{-5.44}\hullfill
\hullface{keel!72!black}\hullmove{49.61}{-5.44}\hullline{50.43}{-5.44}\hullline{50.43}{-4.85}\hullline{49.61}{-4.85}\hullfill
\hullface{keel!98!white}\hullmove{-54.55}{-1.32}\hullline{-52.99}{-1.32}\hullline{-52.99}{-0.00}\hullline{-54.55}{-0.00}\hullfill
\hullface{topsides!90!black}\hullmove{-56.12}{0.78}\hullline{-54.55}{0.78}\hullline{-54.55}{1.56}\hullline{-56.12}{1.56}\hullfill
\hullface{keel!81!black}\hullmove{-45.16}{-6.77}\hullline{-43.59}{-6.77}\hullline{-43.59}{-5.86}\hullline{-45.16}{-5.86}\hullfill
\hullface{topsides!72!black}\hullmove{-60.81}{2.34}\hullline{-59.25}{2.34}\hullline{-59.25}{3.12}\hullline{-60.81}{3.12}\hullfill
\hullface{topsides!84!black}\hullmove{-57.68}{1.56}\hullline{-56.12}{1.56}\hullline{-56.12}{2.34}\hullline{-57.68}{2.34}\hullfill
\hullface{topsides!57!black}\hullmove{61.14}{3.04}\hullline{61.97}{3.04}\hullline{61.97}{3.90}\hullline{61.14}{3.90}\hullfill
\hullface{keel!84!black}\hullmove{-46.72}{-5.86}\hullline{-45.16}{-5.86}\hullline{-45.16}{-5.44}\hullline{-46.72}{-5.44}\hullfill
\hullface{keel!81!black}\hullmove{52.08}{-3.75}\hullline{52.90}{-3.75}\hullline{52.90}{-2.58}\hullline{52.08}{-2.58}\hullfill
\hullface{topsides!81!black}\hullmove{55.37}{0.78}\hullline{56.00}{0.78}\hullline{56.00}{1.56}\hullline{55.37}{1.56}\hullfill
\hullface{topsides!87!black}\hullmove{54.55}{0.00}\hullline{55.37}{0.00}\hullline{55.37}{0.78}\hullline{54.55}{0.78}\hullfill
\hullface{keel!90!black}\hullmove{-51.42}{-3.75}\hullline{-49.85}{-3.75}\hullline{-49.85}{-2.58}\hullline{-51.42}{-2.58}\hullfill
\hullface{keel!96!black}\hullmove{-52.99}{-2.58}\hullline{-51.42}{-2.58}\hullline{-51.42}{-1.32}\hullline{-52.99}{-1.32}\hullfill
\hullface{topsides!90!black}\hullmove{62.79}{3.90}\hullline{63.61}{3.90}\hullline{63.61}{4.55}\hullline{63.61}{6.07}\hullline{63.61}{7.59}\hullline{62.79}{7.59}\hullline{62.79}{6.07}\hullline{62.79}{4.55}\hullfill
\hullface{keel!69!black}\hullmove{47.13}{-6.77}\hullline{47.96}{-6.77}\hullline{47.96}{-5.86}\hullline{47.13}{-5.86}\hullfill
\hullface{keel!63!black}\hullmove{45.49}{-7.57}\hullline{46.31}{-7.57}\hullline{46.31}{-6.77}\hullline{45.49}{-6.77}\hullfill
\hullface{keel!75!black}\hullmove{50.43}{-4.85}\hullline{51.25}{-4.85}\hullline{51.25}{-3.75}\hullline{50.43}{-3.75}\hullfill
\hullface{keel!69!black}\hullmove{-40.46}{-8.25}\hullline{-38.90}{-8.25}\hullline{-38.90}{-7.57}\hullline{-40.46}{-7.57}\hullfill
\hullface{topsides!99!black}\hullmove{-54.55}{0.00}\hullline{-52.99}{0.00}\hullline{-52.99}{0.78}\hullline{-54.55}{0.78}\hullfill
\hullface{topsides!75!black}\hullmove{57.02}{1.52}\hullline{57.85}{1.52}\hullline{57.85}{3.04}\hullline{57.02}{3.04}\hullfill
\hullface{keel!57!black}\hullmove{43.84}{-8.25}\hullline{44.66}{-8.25}\hullline{44.66}{-7.57}\hullline{43.84}{-7.57}\hullfill
\hullface{keel!57!black}\hullmove{42.19}{-8.77}\hullline{43.02}{-8.77}\hullline{43.02}{-8.25}\hullline{42.19}{-8.25}\hullfill
\hullface{keel!90!black}\hullmove{53.73}{-1.32}\hullline{54.55}{-1.32}\hullline{54.55}{0.00}\hullline{53.73}{0.00}\hullfill
\hullface{keel!84!black}\hullmove{52.90}{-2.58}\hullline{53.73}{-2.58}\hullline{53.73}{-1.32}\hullline{52.90}{-1.32}\hullfill
\hullface{topsides!57!black}\hullmove{-65.51}{3.12}\hullline{-63.94}{3.12}\hullline{-63.94}{3.90}\hullline{-65.51}{3.90}\hullfill
\hullface{topsides!84!black}\hullmove{54.55}{0.78}\hullline{55.37}{0.78}\hullline{55.37}{1.56}\hullline{54.55}{1.56}\hullfill
\hullface{keel!75!black}\hullmove{48.78}{-5.44}\hullline{49.61}{-5.44}\hullline{49.61}{-4.85}\hullline{48.78}{-4.85}\hullfill
\hullface{keel!72!black}\hullmove{47.96}{-5.86}\hullline{48.78}{-5.86}\hullline{48.78}{-5.44}\hullline{47.96}{-5.44}\hullfill
\hullface{topsides!60!black}\hullmove{60.32}{3.04}\hullline{61.14}{3.04}\hullline{61.14}{3.90}\hullline{60.32}{3.90}\hullfill
\hullface{topsides!90!black}\hullmove{53.73}{0.00}\hullline{54.55}{0.00}\hullline{54.55}{0.78}\hullline{53.73}{0.78}\hullfill
\hullface{topsides!75!black}\hullmove{-59.25}{2.34}\hullline{-57.68}{2.34}\hullline{-57.68}{3.12}\hullline{-59.25}{3.12}\hullfill
\hullface{keel!78!black}\hullmove{-42.03}{-7.57}\hullline{-40.46}{-7.57}\hullline{-40.46}{-6.77}\hullline{-42.03}{-6.77}\hullfill
\hullface{topsides!93!black}\hullmove{-54.55}{0.78}\hullline{-52.99}{0.78}\hullline{-52.99}{1.56}\hullline{-54.55}{1.56}\hullfill
\hullface{keel!98!white}\hullmove{-52.99}{-1.32}\hullline{-51.42}{-1.32}\hullline{-51.42}{-0.00}\hullline{-52.99}{-0.00}\hullfill
\hullface{keel!81!black}\hullmove{51.25}{-3.75}\hullline{52.08}{-3.75}\hullline{52.08}{-2.58}\hullline{51.25}{-2.58}\hullfill
\hullface{topsides!78!black}\hullmove{56.20}{1.52}\hullline{57.02}{1.52}\hullline{57.02}{3.04}\hullline{56.20}{3.04}\hullfill
\hullface{topsides!87!black}\hullmove{-56.12}{1.56}\hullline{-54.55}{1.56}\hullline{-54.55}{2.34}\hullline{-56.12}{2.34}\hullfill
\hullface{keel!87!black}\hullmove{-46.72}{-5.44}\hullline{-45.16}{-5.44}\hullline{-45.16}{-4.85}\hullline{-46.72}{-4.85}\hullfill
\hullface{keel!90!black}\hullmove{-48.29}{-4.85}\hullline{-46.72}{-4.85}\hullline{-46.72}{-3.75}\hullline{-48.29}{-3.75}\hullfill
\hullface{topsides!81!black}\hullmove{55.37}{1.56}\hullline{56.00}{1.56}\hullline{56.00}{2.34}\hullline{55.37}{2.34}\hullfill
\hullface{topsides!93!black}\hullmove{61.97}{3.90}\hullline{62.79}{3.90}\hullline{62.79}{4.55}\hullline{62.79}{6.07}\hullline{62.79}{7.59}\hullline{61.97}{7.59}\hullline{61.97}{6.07}\hullline{61.97}{4.55}\hullfill
\hullface{keel!78!black}\hullmove{49.61}{-4.85}\hullline{50.43}{-4.85}\hullline{50.43}{-3.75}\hullline{49.61}{-3.75}\hullfill
\hullface{keel!90!black}\hullmove{52.90}{-1.32}\hullline{53.73}{-1.32}\hullline{53.73}{0.00}\hullline{52.90}{0.00}\hullfill
\hullface{topsides!81!black}\hullmove{56.00}{0.00}\hullline{56.00}{0.00}\hullline{56.00}{1.52}\hullline{56.00}{3.04}\hullline{56.00}{3.90}\hullline{56.00}{4.55}\hullline{56.00}{3.12}\hullline{56.00}{2.34}\hullline{56.00}{1.56}\hullline{56.00}{0.78}\hullfill
\hullface{keel!72!black}\hullmove{46.31}{-6.77}\hullline{47.13}{-6.77}\hullline{47.13}{-5.86}\hullline{46.31}{-5.86}\hullfill
\hullface{keel!87!black}\hullmove{52.08}{-2.58}\hullline{52.90}{-2.58}\hullline{52.90}{-1.32}\hullline{52.08}{-1.32}\hullfill
\hullface{keel!93!black}\hullmove{-49.85}{-3.75}\hullline{-48.29}{-3.75}\hullline{-48.29}{-2.58}\hullline{-49.85}{-2.58}\hullfill
\hullface{topsides!99!black}\hullmove{-52.99}{0.00}\hullline{-51.42}{0.00}\hullline{-51.42}{0.78}\hullline{-52.99}{0.78}\hullfill
\hullface{keel!96!black}\hullmove{-51.42}{-2.58}\hullline{-49.85}{-2.58}\hullline{-49.85}{-1.32}\hullline{-51.42}{-1.32}\hullfill
\hullface{topsides!78!black}\hullmove{56.00}{1.52}\hullline{56.20}{1.52}\hullline{56.20}{3.04}\hullline{56.00}{3.04}\hullfill
\hullface{keel!57!black}\hullmove{-37.33}{-8.77}\hullline{-35.77}{-8.77}\hullline{-35.77}{-8.25}\hullline{-37.33}{-8.25}\hullfill
\hullface{keel!84!black}\hullmove{-43.59}{-6.77}\hullline{-42.03}{-6.77}\hullline{-42.03}{-5.86}\hullline{-43.59}{-5.86}\hullfill
\hullface{topsides!87!black}\hullmove{53.73}{0.78}\hullline{54.55}{0.78}\hullline{54.55}{1.56}\hullline{53.73}{1.56}\hullfill
\hullface{keel!87!black}\hullmove{-45.16}{-5.86}\hullline{-43.59}{-5.86}\hullline{-43.59}{-5.44}\hullline{-45.16}{-5.44}\hullfill
\hullface{keel!66!black}\hullmove{44.66}{-7.57}\hullline{45.49}{-7.57}\hullline{45.49}{-6.77}\hullline{44.66}{-6.77}\hullfill
\hullface{topsides!60!black}\hullmove{-63.94}{3.12}\hullline{-62.38}{3.12}\hullline{-62.38}{3.90}\hullline{-63.94}{3.90}\hullfill
\hullface{topsides!63!black}\hullmove{59.49}{3.04}\hullline{60.32}{3.04}\hullline{60.32}{3.90}\hullline{59.49}{3.90}\hullfill
\hullface{topsides!90!black}\hullmove{52.90}{0.00}\hullline{53.73}{0.00}\hullline{53.73}{0.78}\hullline{52.90}{0.78}\hullfill
\hullface{topsides!84!black}\hullmove{54.55}{1.56}\hullline{55.37}{1.56}\hullline{55.37}{2.34}\hullline{54.55}{2.34}\hullfill
\hullface{keel!78!black}\hullmove{47.96}{-5.44}\hullline{48.78}{-5.44}\hullline{48.78}{-4.85}\hullline{47.96}{-4.85}\hullfill
\hullface{topsides!81!black}\hullmove{-57.68}{2.34}\hullline{-56.12}{2.34}\hullline{-56.12}{3.12}\hullline{-57.68}{3.12}\hullfill
\hullface{keel!78!black}\hullmove{47.13}{-5.86}\hullline{47.96}{-5.86}\hullline{47.96}{-5.44}\hullline{47.13}{-5.44}\hullfill
\hullface{keel!84!black}\hullmove{50.43}{-3.75}\hullline{51.25}{-3.75}\hullline{51.25}{-2.58}\hullline{50.43}{-2.58}\hullfill
\hullface{keel!60!black}\hullmove{43.02}{-8.25}\hullline{43.84}{-8.25}\hullline{43.84}{-7.57}\hullline{43.02}{-7.57}\hullfill
\hullface{keel!57!black}\hullmove{41.37}{-8.77}\hullline{42.19}{-8.77}\hullline{42.19}{-8.25}\hullline{41.37}{-8.25}\hullfill
\hullface{topsides!90!black}\hullmove{-54.55}{1.56}\hullline{-52.99}{1.56}\hullline{-52.99}{2.34}\hullline{-54.55}{2.34}\hullfill
\hullface{topsides!96!black}\hullmove{-52.99}{0.78}\hullline{-51.42}{0.78}\hullline{-51.42}{1.56}\hullline{-52.99}{1.56}\hullfill
\hullface{keel!93!black}\hullmove{52.08}{-1.32}\hullline{52.90}{-1.32}\hullline{52.90}{0.00}\hullline{52.08}{0.00}\hullfill
\hullface{topsides!93!black}\hullmove{61.14}{3.90}\hullline{61.97}{3.90}\hullline{61.97}{4.55}\hullline{61.97}{6.07}\hullline{61.97}{7.59}\hullline{61.14}{7.59}\hullline{61.14}{6.07}\hullline{61.14}{4.55}\hullfill
\hullface{keel!98!white}\hullmove{-51.42}{-1.32}\hullline{-49.85}{-1.32}\hullline{-49.85}{-0.00}\hullline{-51.42}{-0.00}\hullfill
\hullface{keel!81!black}\hullmove{48.78}{-4.85}\hullline{49.61}{-4.85}\hullline{49.61}{-3.75}\hullline{48.78}{-3.75}\hullfill
\hullface{keel!90!black}\hullmove{51.25}{-2.58}\hullline{52.08}{-2.58}\hullline{52.08}{-1.32}\hullline{51.25}{-1.32}\hullfill
\hullface{topsides!90!black}\hullmove{52.90}{0.78}\hullline{53.73}{0.78}\hullline{53.73}{1.56}\hullline{52.90}{1.56}\hullfill
\hullface{topsides!66!black}\hullmove{58.67}{3.04}\hullline{59.49}{3.04}\hullline{59.49}{3.90}\hullline{58.67}{3.90}\hullfill
\hullface{topsides!63!black}\hullmove{-62.38}{3.12}\hullline{-60.81}{3.12}\hullline{-60.81}{3.90}\hullline{-62.38}{3.90}\hullfill
\hullface{keel!93!black}\hullmove{-46.72}{-4.85}\hullline{-45.16}{-4.85}\hullline{-45.16}{-3.75}\hullline{-46.72}{-3.75}\hullfill
\hullface{topsides!93!black}\hullmove{52.08}{0.00}\hullline{52.90}{0.00}\hullline{52.90}{0.78}\hullline{52.08}{0.78}\hullfill
\hullface{topsides!87!black}\hullmove{53.73}{1.56}\hullline{54.55}{1.56}\hullline{54.55}{2.34}\hullline{53.73}{2.34}\hullfill
\hullface{keel!78!black}\hullmove{-38.90}{-8.25}\hullline{-37.33}{-8.25}\hullline{-37.33}{-7.57}\hullline{-38.90}{-7.57}\hullfill
\hullface{topsides!99!black}\hullmove{-51.42}{0.00}\hullline{-49.85}{0.00}\hullline{-49.85}{0.78}\hullline{-51.42}{0.78}\hullfill
\hullface{keel!90!black}\hullmove{-45.16}{-5.44}\hullline{-43.59}{-5.44}\hullline{-43.59}{-4.85}\hullline{-45.16}{-4.85}\hullfill
\hullface{keel!78!black}\hullmove{45.49}{-6.77}\hullline{46.31}{-6.77}\hullline{46.31}{-5.86}\hullline{45.49}{-5.86}\hullfill
\hullface{keel!60!black}\hullmove{-35.77}{-8.77}\hullline{41.37}{-8.77}\hullline{41.37}{-8.25}\hullline{-35.77}{-8.25}\hullfill
\hullface{keel!99!black}\hullmove{-49.85}{-2.58}\hullline{-48.29}{-2.58}\hullline{-48.29}{-1.32}\hullline{-49.85}{-1.32}\hullfill
\hullface{topsides!84!black}\hullmove{-56.12}{2.34}\hullline{-54.55}{2.34}\hullline{-54.55}{3.12}\hullline{-56.12}{3.12}\hullfill
\hullface{topsides!81!black}\hullmove{55.37}{2.34}\hullline{56.00}{2.34}\hullline{56.00}{3.12}\hullline{55.37}{3.12}\hullfill
\hullface{keel!96!black}\hullmove{-48.29}{-3.75}\hullline{-46.72}{-3.75}\hullline{-46.72}{-2.58}\hullline{-48.29}{-2.58}\hullfill
\hullface{keel!87!black}\hullmove{49.61}{-3.75}\hullline{50.43}{-3.75}\hullline{50.43}{-2.58}\hullline{49.61}{-2.58}\hullfill
\hullface{keel!81!black}\hullmove{47.13}{-5.44}\hullline{47.96}{-5.44}\hullline{47.96}{-4.85}\hullline{47.13}{-4.85}\hullfill
\hullface{keel!84!black}\hullmove{-40.46}{-7.57}\hullline{-38.90}{-7.57}\hullline{-38.90}{-6.77}\hullline{-40.46}{-6.77}\hullfill
\hullface{keel!93!black}\hullmove{51.25}{-1.32}\hullline{52.08}{-1.32}\hullline{52.08}{0.00}\hullline{51.25}{0.00}\hullfill
\hullface{topsides!93!black}\hullmove{-52.99}{1.56}\hullline{-51.42}{1.56}\hullline{-51.42}{2.34}\hullline{-52.99}{2.34}\hullfill
\hullface{topsides!96!black}\hullmove{60.32}{3.90}\hullline{61.14}{3.90}\hullline{61.14}{4.55}\hullline{61.14}{6.07}\hullline{61.14}{7.59}\hullline{60.32}{7.59}\hullline{60.32}{6.07}\hullline{60.32}{4.55}\hullfill
\hullface{keel!72!black}\hullmove{43.84}{-7.57}\hullline{44.66}{-7.57}\hullline{44.66}{-6.77}\hullline{43.84}{-6.77}\hullfill
\hullface{keel!81!black}\hullmove{46.31}{-5.86}\hullline{47.13}{-5.86}\hullline{47.13}{-5.44}\hullline{46.31}{-5.44}\hullfill
\hullface{topsides!90!black}\hullmove{52.08}{0.78}\hullline{52.90}{0.78}\hullline{52.90}{1.56}\hullline{52.08}{1.56}\hullfill
\hullface{topsides!69!black}\hullmove{57.85}{3.04}\hullline{58.67}{3.04}\hullline{58.67}{3.90}\hullline{57.85}{3.90}\hullfill
\hullface{topsides!96!black}\hullmove{-51.42}{0.78}\hullline{-49.85}{0.78}\hullline{-49.85}{1.56}\hullline{-51.42}{1.56}\hullfill
\hullface{keel!90!black}\hullmove{50.43}{-2.58}\hullline{51.25}{-2.58}\hullline{51.25}{-1.32}\hullline{50.43}{-1.32}\hullfill
\hullface{keel!90!black}\hullmove{-43.59}{-5.86}\hullline{-42.03}{-5.86}\hullline{-42.03}{-5.44}\hullline{-43.59}{-5.44}\hullfill
\hullface{topsides!87!black}\hullmove{52.90}{1.56}\hullline{53.73}{1.56}\hullline{53.73}{2.34}\hullline{52.90}{2.34}\hullfill
\hullface{keel!84!black}\hullmove{47.96}{-4.85}\hullline{48.78}{-4.85}\hullline{48.78}{-3.75}\hullline{47.96}{-3.75}\hullfill
\hullface{keel!90!black}\hullmove{-42.03}{-6.77}\hullline{-40.46}{-6.77}\hullline{-40.46}{-5.86}\hullline{-42.03}{-5.86}\hullfill
\hullface{topsides!84!black}\hullmove{54.55}{2.34}\hullline{55.37}{2.34}\hullline{55.37}{3.12}\hullline{54.55}{3.12}\hullfill
\hullface{topsides!93!black}\hullmove{51.25}{0.00}\hullline{52.08}{0.00}\hullline{52.08}{0.78}\hullline{51.25}{0.78}\hullfill
\hullface{topsides!66!black}\hullmove{-60.81}{3.12}\hullline{-59.25}{3.12}\hullline{-59.25}{3.90}\hullline{-60.81}{3.90}\hullfill
\hullface{keel!98!white}\hullmove{-49.85}{-1.32}\hullline{-48.29}{-1.32}\hullline{-48.29}{-0.00}\hullline{-49.85}{-0.00}\hullfill
\hullface{topsides!87!black}\hullmove{-54.55}{2.34}\hullline{-52.99}{2.34}\hullline{-52.99}{3.12}\hullline{-54.55}{3.12}\hullfill
\hullface{topsides!99!black}\hullmove{-49.85}{0.00}\hullline{-48.29}{0.00}\hullline{-48.29}{0.78}\hullline{-49.85}{0.78}\hullfill
\hullface{keel!69!black}\hullmove{42.19}{-8.25}\hullline{43.02}{-8.25}\hullline{43.02}{-7.57}\hullline{42.19}{-7.57}\hullfill
\hullface{topsides!72!black}\hullmove{57.02}{3.04}\hullline{57.85}{3.04}\hullline{57.85}{3.90}\hullline{57.02}{3.90}\hullfill
\hullface{keel!96!black}\hullmove{50.43}{-1.32}\hullline{51.25}{-1.32}\hullline{51.25}{0.00}\hullline{50.43}{0.00}\hullfill
\hullface{topsides!93!black}\hullmove{51.25}{0.78}\hullline{52.08}{0.78}\hullline{52.08}{1.56}\hullline{51.25}{1.56}\hullfill
\hullface{keel!90!black}\hullmove{48.78}{-3.75}\hullline{49.61}{-3.75}\hullline{49.61}{-2.58}\hullline{48.78}{-2.58}\hullfill
\hullface{topsides!96!black}\hullmove{59.49}{3.90}\hullline{60.32}{3.90}\hullline{60.32}{4.55}\hullline{60.32}{6.07}\hullline{60.32}{7.59}\hullline{59.49}{7.59}\hullline{59.49}{6.07}\hullline{59.49}{4.55}\hullfill
\hullface{topsides!90!black}\hullmove{52.08}{1.56}\hullline{52.90}{1.56}\hullline{52.90}{2.34}\hullline{52.08}{2.34}\hullfill
\hullface{keel!81!black}\hullmove{44.66}{-6.77}\hullline{45.49}{-6.77}\hullline{45.49}{-5.86}\hullline{44.66}{-5.86}\hullfill
\hullface{topsides!87!black}\hullmove{53.73}{2.34}\hullline{54.55}{2.34}\hullline{54.55}{3.12}\hullline{53.73}{3.12}\hullfill
\hullface{keel!99!black}\hullmove{-48.29}{-2.58}\hullline{-46.72}{-2.58}\hullline{-46.72}{-1.32}\hullline{-48.29}{-1.32}\hullfill
\hullface{keel!96!black}\hullmove{-45.16}{-4.85}\hullline{-43.59}{-4.85}\hullline{-43.59}{-3.75}\hullline{-45.16}{-3.75}\hullfill
\hullface{topsides!96!black}\hullmove{-51.42}{1.56}\hullline{-49.85}{1.56}\hullline{-49.85}{2.34}\hullline{-51.42}{2.34}\hullfill
\hullface{keel!84!black}\hullmove{46.31}{-5.44}\hullline{47.13}{-5.44}\hullline{47.13}{-4.85}\hullline{46.31}{-4.85}\hullfill
\hullface{keel!93!black}\hullmove{49.61}{-2.58}\hullline{50.43}{-2.58}\hullline{50.43}{-1.32}\hullline{49.61}{-1.32}\hullfill
\hullface{topsides!96!black}\hullmove{50.43}{0.00}\hullline{51.25}{0.00}\hullline{51.25}{0.78}\hullline{50.43}{0.78}\hullfill
\hullface{topsides!72!black}\hullmove{-59.25}{3.12}\hullline{-57.68}{3.12}\hullline{-57.68}{3.90}\hullline{-59.25}{3.90}\hullfill
\hullface{keel!96!black}\hullmove{-46.72}{-3.75}\hullline{-45.16}{-3.75}\hullline{-45.16}{-2.58}\hullline{-46.72}{-2.58}\hullfill
\hullface{topsides!99!black}\hullmove{-49.85}{0.78}\hullline{-48.29}{0.78}\hullline{-48.29}{1.56}\hullline{-49.85}{1.56}\hullfill
\hullface{keel!93!black}\hullmove{-43.59}{-5.44}\hullline{-42.03}{-5.44}\hullline{-42.03}{-4.85}\hullline{-43.59}{-4.85}\hullfill
\hullface{keel!84!black}\hullmove{45.49}{-5.86}\hullline{46.31}{-5.86}\hullline{46.31}{-5.44}\hullline{45.49}{-5.44}\hullfill
\hullface{keel!87!black}\hullmove{47.13}{-4.85}\hullline{47.96}{-4.85}\hullline{47.96}{-3.75}\hullline{47.13}{-3.75}\hullfill
\hullface{topsides!75!black}\hullmove{56.20}{3.04}\hullline{57.02}{3.04}\hullline{57.02}{3.90}\hullline{56.20}{3.90}\hullfill
\hullface{topsides!90!black}\hullmove{-52.99}{2.34}\hullline{-51.42}{2.34}\hullline{-51.42}{3.12}\hullline{-52.99}{3.12}\hullfill
\hullface{keel!98!white}\hullmove{-48.29}{-1.32}\hullline{-46.72}{-1.32}\hullline{-46.72}{-0.00}\hullline{-48.29}{-0.00}\hullfill
\hullface{topsides!93!black}\hullmove{50.43}{0.78}\hullline{51.25}{0.78}\hullline{51.25}{1.56}\hullline{50.43}{1.56}\hullfill
\hullface{topsides!87!black}\hullmove{52.90}{2.34}\hullline{53.73}{2.34}\hullline{53.73}{3.12}\hullline{52.90}{3.12}\hullfill
\hullface{keel!81!black}\hullmove{43.02}{-7.57}\hullline{43.84}{-7.57}\hullline{43.84}{-6.77}\hullline{43.02}{-6.77}\hullfill
\hullface{topsides!93!black}\hullmove{51.25}{1.56}\hullline{52.08}{1.56}\hullline{52.08}{2.34}\hullline{51.25}{2.34}\hullfill
\hullface{keel!96!black}\hullmove{49.61}{-1.32}\hullline{50.43}{-1.32}\hullline{50.43}{0.00}\hullline{49.61}{0.00}\hullfill
\hullface{topsides!96!black}\hullmove{58.67}{3.90}\hullline{59.49}{3.90}\hullline{59.49}{4.55}\hullline{59.49}{6.07}\hullline{59.49}{7.59}\hullline{58.67}{7.59}\hullline{58.67}{6.07}\hullline{58.67}{4.55}\hullfill
\hullface{topsides!99!black}\hullmove{-48.29}{0.00}\hullline{-46.72}{0.00}\hullline{-46.72}{0.78}\hullline{-48.29}{0.78}\hullfill
\hullface{topsides!78!black}\hullmove{56.00}{3.04}\hullline{56.20}{3.04}\hullline{56.20}{3.90}\hullline{56.00}{3.90}\hullfill
\hullface{keel!84!black}\hullmove{-37.33}{-8.25}\hullline{-35.77}{-8.25}\hullline{-35.77}{-7.57}\hullline{-37.33}{-7.57}\hullfill
\hullface{keel!93!black}\hullmove{47.96}{-3.75}\hullline{48.78}{-3.75}\hullline{48.78}{-2.58}\hullline{47.96}{-2.58}\hullfill
\hullface{topsides!75!black}\hullmove{-57.68}{3.12}\hullline{-56.12}{3.12}\hullline{-56.12}{3.90}\hullline{-57.68}{3.90}\hullfill
\hullface{topsides!96!black}\hullmove{49.61}{0.00}\hullline{50.43}{0.00}\hullline{50.43}{0.78}\hullline{49.61}{0.78}\hullfill
\hullface{keel!93!black}\hullmove{-42.03}{-5.86}\hullline{-40.46}{-5.86}\hullline{-40.46}{-5.44}\hullline{-42.03}{-5.44}\hullfill
\hullface{topsides!96!black}\hullmove{-49.85}{1.56}\hullline{-48.29}{1.56}\hullline{-48.29}{2.34}\hullline{-49.85}{2.34}\hullfill
\hullface{keel!93!black}\hullmove{48.78}{-2.58}\hullline{49.61}{-2.58}\hullline{49.61}{-1.32}\hullline{48.78}{-1.32}\hullfill
\hullface{keel!93!black}\hullmove{-38.90}{-7.57}\hullline{-37.33}{-7.57}\hullline{-37.33}{-6.77}\hullline{-38.90}{-6.77}\hullfill
\hullface{topsides!90!black}\hullmove{52.08}{2.34}\hullline{52.90}{2.34}\hullline{52.90}{3.12}\hullline{52.08}{3.12}\hullfill
\hullface{topsides!99!black}\hullmove{-48.29}{0.78}\hullline{-46.72}{0.78}\hullline{-46.72}{1.56}\hullline{-48.29}{1.56}\hullfill
\hullface{keel!81!black}\hullmove{41.37}{-8.25}\hullline{42.19}{-8.25}\hullline{42.19}{-7.57}\hullline{41.37}{-7.57}\hullfill
\hullface{keel!99!black}\hullmove{-46.72}{-2.58}\hullline{-45.16}{-2.58}\hullline{-45.16}{-1.32}\hullline{-46.72}{-1.32}\hullfill
\hullface{keel!93!black}\hullmove{-40.46}{-6.77}\hullline{-38.90}{-6.77}\hullline{-38.90}{-5.86}\hullline{-40.46}{-5.86}\hullfill
\hullface{topsides!81!black}\hullmove{55.37}{3.12}\hullline{56.00}{3.12}\hullline{56.00}{3.90}\hullline{55.37}{3.90}\hullfill
\hullface{keel!90!black}\hullmove{45.49}{-5.44}\hullline{46.31}{-5.44}\hullline{46.31}{-4.85}\hullline{45.49}{-4.85}\hullfill
\hullface{topsides!93!black}\hullmove{-51.42}{2.34}\hullline{-49.85}{2.34}\hullline{-49.85}{3.12}\hullline{-51.42}{3.12}\hullfill
\hullface{topsides!93!black}\hullmove{50.43}{1.56}\hullline{51.25}{1.56}\hullline{51.25}{2.34}\hullline{50.43}{2.34}\hullfill
\hullface{topsides!96!black}\hullmove{49.61}{0.78}\hullline{50.43}{0.78}\hullline{50.43}{1.56}\hullline{49.61}{1.56}\hullfill
\hullface{keel!87!black}\hullmove{43.84}{-6.77}\hullline{44.66}{-6.77}\hullline{44.66}{-5.86}\hullline{43.84}{-5.86}\hullfill
\hullface{keel!90!black}\hullmove{46.31}{-4.85}\hullline{47.13}{-4.85}\hullline{47.13}{-3.75}\hullline{46.31}{-3.75}\hullfill
\hullface{topsides!99!black}\hullmove{57.85}{3.90}\hullline{58.67}{3.90}\hullline{58.67}{4.55}\hullline{58.67}{6.07}\hullline{58.67}{7.59}\hullline{57.85}{7.59}\hullline{57.85}{6.07}\hullline{57.85}{4.55}\hullfill
\hullface{keel!96!black}\hullmove{48.78}{-1.32}\hullline{49.61}{-1.32}\hullline{49.61}{0.00}\hullline{48.78}{0.00}\hullfill
\hullface{keel!99!black}\hullmove{-45.16}{-3.75}\hullline{-43.59}{-3.75}\hullline{-43.59}{-2.58}\hullline{-45.16}{-2.58}\hullfill
\hullface{keel!96!black}\hullmove{-43.59}{-4.85}\hullline{-42.03}{-4.85}\hullline{-42.03}{-3.75}\hullline{-43.59}{-3.75}\hullfill
\hullface{topsides!81!black}\hullmove{-56.12}{3.12}\hullline{-54.55}{3.12}\hullline{-54.55}{3.90}\hullline{-56.12}{3.90}\hullfill
\hullface{keel!90!black}\hullmove{44.66}{-5.86}\hullline{45.49}{-5.86}\hullline{45.49}{-5.44}\hullline{44.66}{-5.44}\hullfill
\hullface{topsides!96!black}\hullmove{48.78}{0.00}\hullline{49.61}{0.00}\hullline{49.61}{0.78}\hullline{48.78}{0.78}\hullfill
\hullface{keel!98!white}\hullmove{-46.72}{-1.32}\hullline{-45.16}{-1.32}\hullline{-45.16}{-0.00}\hullline{-46.72}{-0.00}\hullfill
\hullface{topsides!84!black}\hullmove{54.55}{3.12}\hullline{55.37}{3.12}\hullline{55.37}{3.90}\hullline{54.55}{3.90}\hullfill
\hullface{topsides!93!black}\hullmove{51.25}{2.34}\hullline{52.08}{2.34}\hullline{52.08}{3.12}\hullline{51.25}{3.12}\hullfill
\hullface{keel!93!black}\hullmove{47.13}{-3.75}\hullline{47.96}{-3.75}\hullline{47.96}{-2.58}\hullline{47.13}{-2.58}\hullfill
\hullface{keel!96!black}\hullmove{47.96}{-2.58}\hullline{48.78}{-2.58}\hullline{48.78}{-1.32}\hullline{47.96}{-1.32}\hullfill
\hullface{topsides!98!white}\hullmove{-46.72}{0.00}\hullline{-45.16}{0.00}\hullline{-45.16}{0.78}\hullline{-46.72}{0.78}\hullfill
\hullface{keel!87!black}\hullmove{-35.77}{-8.25}\hullline{41.37}{-8.25}\hullline{41.37}{-7.57}\hullline{-35.77}{-7.57}\hullfill
\hullface{keel!96!black}\hullmove{-42.03}{-5.44}\hullline{-40.46}{-5.44}\hullline{-40.46}{-4.85}\hullline{-42.03}{-4.85}\hullfill
\hullface{topsides!99!black}\hullmove{-48.29}{1.56}\hullline{-46.72}{1.56}\hullline{-46.72}{2.34}\hullline{-48.29}{2.34}\hullfill
\hullface{topsides!96!black}\hullmove{49.61}{1.56}\hullline{50.43}{1.56}\hullline{50.43}{2.34}\hullline{49.61}{2.34}\hullfill
\hullface{topsides!57!black}\hullmove{56.00}{3.12}\hullline{56.00}{4.55}\hullline{56.00}{6.07}\hullline{56.00}{3.90}\hullfill
\hullface{topsides!96!black}\hullmove{48.78}{0.78}\hullline{49.61}{0.78}\hullline{49.61}{1.56}\hullline{48.78}{1.56}\hullfill
\hullface{topsides!96!black}\hullmove{-49.85}{2.34}\hullline{-48.29}{2.34}\hullline{-48.29}{3.12}\hullline{-49.85}{3.12}\hullfill
\hullface{topsides!99!black}\hullmove{57.02}{3.90}\hullline{57.85}{3.90}\hullline{57.85}{4.55}\hullline{57.85}{6.07}\hullline{57.85}{7.59}\hullline{57.02}{7.59}\hullline{57.02}{6.07}\hullline{57.02}{4.55}\hullfill
\hullface{keel!87!black}\hullmove{42.19}{-7.57}\hullline{43.02}{-7.57}\hullline{43.02}{-6.77}\hullline{42.19}{-6.77}\hullfill
\hullface{keel!99!black}\hullmove{47.96}{-1.32}\hullline{48.78}{-1.32}\hullline{48.78}{0.00}\hullline{47.96}{0.00}\hullfill
\hullface{topsides!84!black}\hullmove{-54.55}{3.12}\hullline{-52.99}{3.12}\hullline{-52.99}{3.90}\hullline{-54.55}{3.90}\hullfill
\hullface{topsides!87!black}\hullmove{53.73}{3.12}\hullline{54.55}{3.12}\hullline{54.55}{3.90}\hullline{53.73}{3.90}\hullfill
\hullface{topsides!93!black}\hullmove{50.43}{2.34}\hullline{51.25}{2.34}\hullline{51.25}{3.12}\hullline{50.43}{3.12}\hullfill
\hullface{topsides!96!black}\hullmove{47.96}{0.00}\hullline{48.78}{0.00}\hullline{48.78}{0.78}\hullline{47.96}{0.78}\hullfill
\hullface{keel!99!black}\hullmove{-45.16}{-2.58}\hullline{-43.59}{-2.58}\hullline{-43.59}{-1.32}\hullline{-45.16}{-1.32}\hullfill
\hullface{keel!93!black}\hullmove{44.66}{-5.44}\hullline{45.49}{-5.44}\hullline{45.49}{-4.85}\hullline{44.66}{-4.85}\hullfill
\hullface{topsides!99!black}\hullmove{-46.72}{0.78}\hullline{-45.16}{0.78}\hullline{-45.16}{1.56}\hullline{-45.16}{2.34}\hullline{-46.72}{2.34}\hullline{-46.72}{1.56}\hullfill
\hullface{keel!93!black}\hullmove{45.49}{-4.85}\hullline{46.31}{-4.85}\hullline{46.31}{-3.75}\hullline{45.49}{-3.75}\hullfill
\hullface{topsides!96!black}\hullmove{48.78}{1.56}\hullline{49.61}{1.56}\hullline{49.61}{2.34}\hullline{48.78}{2.34}\hullfill
\hullface{keel!96!black}\hullmove{47.13}{-2.58}\hullline{47.96}{-2.58}\hullline{47.96}{-1.32}\hullline{47.13}{-1.32}\hullfill
\hullface{keel!98!white}\hullmove{-45.16}{-1.32}\hullline{-43.59}{-1.32}\hullline{-43.59}{-0.00}\hullline{-45.16}{-0.00}\hullfill
\hullface{topsides!87!black}\hullmove{52.90}{3.12}\hullline{53.73}{3.12}\hullline{53.73}{3.90}\hullline{52.90}{3.90}\hullfill
\hullface{keel!96!black}\hullmove{-40.46}{-5.86}\hullline{-38.90}{-5.86}\hullline{-38.90}{-5.44}\hullline{-40.46}{-5.44}\hullfill
\hullface{keel!90!black}\hullmove{43.02}{-6.77}\hullline{43.84}{-6.77}\hullline{43.84}{-5.86}\hullline{43.02}{-5.86}\hullfill
\hullface{topsides!99!black}\hullmove{56.20}{3.90}\hullline{57.02}{3.90}\hullline{57.02}{4.55}\hullline{57.02}{6.07}\hullline{57.02}{7.59}\hullline{56.20}{7.59}\hullline{56.20}{6.07}\hullline{56.20}{4.55}\hullfill
\hullface{keel!96!black}\hullmove{46.31}{-3.75}\hullline{47.13}{-3.75}\hullline{47.13}{-2.58}\hullline{46.31}{-2.58}\hullfill
\hullface{topsides!90!black}\hullmove{-52.99}{3.12}\hullline{-51.42}{3.12}\hullline{-51.42}{3.90}\hullline{-52.99}{3.90}\hullfill
\hullface{topsides!96!black}\hullmove{-48.29}{2.34}\hullline{-46.72}{2.34}\hullline{-46.72}{3.12}\hullline{-48.29}{3.12}\hullfill
\hullface{keel!99!black}\hullmove{-43.59}{-3.75}\hullline{-42.03}{-3.75}\hullline{-42.03}{-2.58}\hullline{-43.59}{-2.58}\hullfill
\hullface{keel!93!black}\hullmove{43.84}{-5.86}\hullline{44.66}{-5.86}\hullline{44.66}{-5.44}\hullline{43.84}{-5.44}\hullfill
\hullface{topsides!96!black}\hullmove{49.61}{2.34}\hullline{50.43}{2.34}\hullline{50.43}{3.12}\hullline{49.61}{3.12}\hullfill
\hullface{topsides!96!black}\hullmove{47.96}{0.78}\hullline{48.78}{0.78}\hullline{48.78}{1.56}\hullline{48.78}{2.34}\hullline{47.96}{2.34}\hullline{47.96}{1.56}\hullfill
\hullface{keel!99!black}\hullmove{47.13}{-1.32}\hullline{47.96}{-1.32}\hullline{47.96}{0.00}\hullline{47.13}{0.00}\hullfill
\hullface{topsides!99!black}\hullmove{-45.16}{0.00}\hullline{-43.59}{0.00}\hullline{-43.59}{0.78}\hullline{-43.59}{1.56}\hullline{-45.16}{1.56}\hullline{-45.16}{0.78}\hullfill
\hullface{keel!99!black}\hullmove{-42.03}{-4.85}\hullline{-40.46}{-4.85}\hullline{-40.46}{-3.75}\hullline{-42.03}{-3.75}\hullfill
\hullface{topsides!99!black}\hullmove{56.00}{3.90}\hullline{56.20}{3.90}\hullline{56.20}{4.55}\hullline{56.20}{6.07}\hullline{56.20}{7.59}\hullline{56.00}{7.59}\hullline{56.00}{6.07}\hullline{56.00}{4.55}\hullfill
\hullface{keel!96!black}\hullmove{-37.33}{-7.57}\hullline{-35.77}{-7.57}\hullline{-35.77}{-6.77}\hullline{-37.33}{-6.77}\hullfill
\hullface{keel!96!black}\hullmove{-38.90}{-6.77}\hullline{-37.33}{-6.77}\hullline{-37.33}{-5.86}\hullline{-38.90}{-5.86}\hullfill
\hullface{topsides!90!black}\hullmove{52.08}{3.12}\hullline{52.90}{3.12}\hullline{52.90}{3.90}\hullline{52.08}{3.90}\hullfill
\hullface{topsides!99!black}\hullmove{56.00}{3.90}\hullline{56.00}{6.07}\hullline{56.00}{7.59}\hullline{56.00}{3.90}\hullfill
\hullface{topsides!99!black}\hullmove{47.13}{0.00}\hullline{47.96}{0.00}\hullline{47.96}{0.78}\hullline{47.96}{1.56}\hullline{47.13}{1.56}\hullline{47.13}{0.78}\hullfill
\hullface{topsides!93!black}\hullmove{-51.42}{3.12}\hullline{-49.85}{3.12}\hullline{-49.85}{3.90}\hullline{-51.42}{3.90}\hullfill
\hullface{topsides!96!black}\hullmove{48.78}{2.34}\hullline{49.61}{2.34}\hullline{49.61}{3.12}\hullline{48.78}{3.12}\hullfill
\hullface{keel!99!black}\hullmove{46.31}{-2.58}\hullline{47.13}{-2.58}\hullline{47.13}{-1.32}\hullline{46.31}{-1.32}\hullfill
\hullface{keel!93!black}\hullmove{41.37}{-7.57}\hullline{42.19}{-7.57}\hullline{42.19}{-6.77}\hullline{41.37}{-6.77}\hullfill
\hullface{keel!99!black}\hullmove{-40.46}{-5.44}\hullline{-38.90}{-5.44}\hullline{-38.90}{-4.85}\hullline{-40.46}{-4.85}\hullfill
\hullface{keel!96!black}\hullmove{44.66}{-4.85}\hullline{45.49}{-4.85}\hullline{45.49}{-3.75}\hullline{44.66}{-3.75}\hullfill
\hullface{topsides!99!black}\hullmove{-46.72}{2.34}\hullline{-45.16}{2.34}\hullline{-45.16}{3.12}\hullline{-46.72}{3.12}\hullfill
\hullface{keel!98!white}\hullmove{-43.59}{-2.58}\hullline{-42.03}{-2.58}\hullline{-42.03}{-1.32}\hullline{-42.03}{-0.00}\hullline{-43.59}{-0.00}\hullline{-43.59}{-1.32}\hullfill
\hullface{topsides!93!black}\hullmove{51.25}{3.12}\hullline{52.08}{3.12}\hullline{52.08}{3.90}\hullline{51.25}{3.90}\hullfill
\hullface{keel!96!black}\hullmove{43.84}{-5.44}\hullline{44.66}{-5.44}\hullline{44.66}{-4.85}\hullline{43.84}{-4.85}\hullfill
\hullface{keel!96!black}\hullmove{45.49}{-3.75}\hullline{46.31}{-3.75}\hullline{46.31}{-2.58}\hullline{45.49}{-2.58}\hullfill
\hullface{keel!99!black}\hullmove{46.31}{-1.32}\hullline{47.13}{-1.32}\hullline{47.13}{0.00}\hullline{46.31}{0.00}\hullfill
\hullface{topsides!99!black}\hullmove{46.31}{0.00}\hullline{47.13}{0.00}\hullline{47.13}{0.78}\hullline{47.13}{1.56}\hullline{46.31}{1.56}\hullline{46.31}{0.78}\hullfill
\hullface{topsides!96!black}\hullmove{47.96}{2.34}\hullline{48.78}{2.34}\hullline{48.78}{3.12}\hullline{47.96}{3.12}\hullfill
\hullface{topsides!99!black}\hullmove{-45.16}{1.56}\hullline{-43.59}{1.56}\hullline{-43.59}{2.34}\hullline{-43.59}{3.12}\hullline{-45.16}{3.12}\hullline{-45.16}{2.34}\hullfill
\hullface{topsides!93!black}\hullmove{-49.85}{3.12}\hullline{-48.29}{3.12}\hullline{-48.29}{3.90}\hullline{-49.85}{3.90}\hullfill
\hullface{topsides!93!black}\hullmove{50.43}{3.12}\hullline{51.25}{3.12}\hullline{51.25}{3.90}\hullline{50.43}{3.90}\hullfill
\hullface{keel!96!black}\hullmove{42.19}{-6.77}\hullline{43.02}{-6.77}\hullline{43.02}{-5.86}\hullline{42.19}{-5.86}\hullfill
\hullface{topsides!99!black}\hullmove{47.13}{1.56}\hullline{47.96}{1.56}\hullline{47.96}{2.34}\hullline{47.96}{3.12}\hullline{47.13}{3.12}\hullline{47.13}{2.34}\hullfill
\hullface{keel!96!black}\hullmove{-35.77}{-7.57}\hullline{41.37}{-7.57}\hullline{41.37}{-6.77}\hullline{-35.77}{-6.77}\hullfill
\hullface{keel!99!black}\hullmove{-42.03}{-3.75}\hullline{-40.46}{-3.75}\hullline{-40.46}{-2.58}\hullline{-42.03}{-2.58}\hullfill
\hullface{keel!96!black}\hullmove{43.02}{-5.86}\hullline{43.84}{-5.86}\hullline{43.84}{-5.44}\hullline{43.02}{-5.44}\hullfill
\hullface{topsides!99!black}\hullmove{-43.59}{0.00}\hullline{-42.03}{0.00}\hullline{-42.03}{0.78}\hullline{-42.03}{1.56}\hullline{-42.03}{2.34}\hullline{-43.59}{2.34}\hullline{-43.59}{1.56}\hullline{-43.59}{0.78}\hullfill
\hullface{keel!99!black}\hullmove{45.49}{-2.58}\hullline{46.31}{-2.58}\hullline{46.31}{-1.32}\hullline{46.31}{0.00}\hullline{45.49}{0.00}\hullline{45.49}{-1.32}\hullfill
\hullface{topsides!96!black}\hullmove{49.61}{3.12}\hullline{50.43}{3.12}\hullline{50.43}{3.90}\hullline{49.61}{3.90}\hullfill
\hullface{topsides!96!black}\hullmove{-48.29}{3.12}\hullline{-46.72}{3.12}\hullline{-46.72}{3.90}\hullline{-48.29}{3.90}\hullfill
\hullface{keel!99!black}\hullmove{-40.46}{-4.85}\hullline{-38.90}{-4.85}\hullline{-38.90}{-3.75}\hullline{-40.46}{-3.75}\hullfill
\hullface{keel!99!black}\hullmove{-38.90}{-5.86}\hullline{-37.33}{-5.86}\hullline{-37.33}{-5.44}\hullline{-38.90}{-5.44}\hullfill
\hullface{keel!99!black}\hullmove{44.66}{-3.75}\hullline{45.49}{-3.75}\hullline{45.49}{-2.58}\hullline{44.66}{-2.58}\hullfill
\hullface{keel!99!black}\hullmove{43.84}{-4.85}\hullline{44.66}{-4.85}\hullline{44.66}{-3.75}\hullline{43.84}{-3.75}\hullfill
\hullface{keel!98!white}\hullmove{-42.03}{-2.58}\hullline{-40.46}{-2.58}\hullline{-40.46}{-1.32}\hullline{-40.46}{-0.00}\hullline{-42.03}{-0.00}\hullline{-42.03}{-1.32}\hullfill
\hullface{topsides!99!black}\hullmove{46.31}{1.56}\hullline{47.13}{1.56}\hullline{47.13}{2.34}\hullline{47.13}{3.12}\hullline{46.31}{3.12}\hullline{46.31}{2.34}\hullfill
\hullface{topsides!99!black}\hullmove{45.49}{0.00}\hullline{46.31}{0.00}\hullline{46.31}{0.78}\hullline{46.31}{1.56}\hullline{46.31}{2.34}\hullline{45.49}{2.34}\hullline{45.49}{1.56}\hullline{45.49}{0.78}\hullfill
\hullface{keel!96!black}\hullmove{43.02}{-5.44}\hullline{43.84}{-5.44}\hullline{43.84}{-4.85}\hullline{43.02}{-4.85}\hullfill
\hullface{topsides!96!black}\hullmove{48.78}{3.12}\hullline{49.61}{3.12}\hullline{49.61}{3.90}\hullline{48.78}{3.90}\hullfill
\hullface{keel!99!black}\hullmove{-37.33}{-6.77}\hullline{-35.77}{-6.77}\hullline{-35.77}{-5.86}\hullline{-37.33}{-5.86}\hullfill
\hullface{topsides!96!black}\hullmove{-46.72}{3.12}\hullline{-45.16}{3.12}\hullline{-45.16}{3.90}\hullline{-46.72}{3.90}\hullfill
\hullface{keel!99!black}\hullmove{44.66}{-2.58}\hullline{45.49}{-2.58}\hullline{45.49}{-1.32}\hullline{45.49}{0.00}\hullline{44.66}{0.00}\hullline{44.66}{-1.32}\hullfill
\hullface{keel!99!black}\hullmove{-38.90}{-5.44}\hullline{-37.33}{-5.44}\hullline{-37.33}{-4.85}\hullline{-38.90}{-4.85}\hullfill
\hullface{topsides!99!black}\hullmove{-42.03}{0.00}\hullline{-40.46}{0.00}\hullline{-40.46}{0.78}\hullline{-40.46}{1.56}\hullline{-40.46}{2.34}\hullline{-40.46}{3.12}\hullline{-42.03}{3.12}\hullline{-42.03}{2.34}\hullline{-42.03}{1.56}\hullline{-42.03}{0.78}\hullfill
\hullface{topsides!99!black}\hullmove{47.96}{3.12}\hullline{48.78}{3.12}\hullline{48.78}{3.90}\hullline{47.96}{3.90}\hullfill
\hullface{keel!99!black}\hullmove{41.37}{-6.77}\hullline{42.19}{-6.77}\hullline{42.19}{-5.86}\hullline{41.37}{-5.86}\hullfill
\hullface{keel!99!black}\hullmove{42.19}{-5.86}\hullline{43.02}{-5.86}\hullline{43.02}{-5.44}\hullline{42.19}{-5.44}\hullfill
\hullface{topsides!99!black}\hullmove{-43.59}{2.34}\hullline{-42.03}{2.34}\hullline{-42.03}{3.12}\hullline{-42.03}{3.90}\hullline{-43.59}{3.90}\hullline{-43.59}{3.12}\hullfill
\hullface{topsides!99!black}\hullmove{44.66}{0.00}\hullline{45.49}{0.00}\hullline{45.49}{0.78}\hullline{45.49}{1.56}\hullline{45.49}{2.34}\hullline{45.49}{3.12}\hullline{44.66}{3.12}\hullline{44.66}{2.34}\hullline{44.66}{1.56}\hullline{44.66}{0.78}\hullfill
\hullface{topsides!99!black}\hullmove{-45.16}{3.12}\hullline{-43.59}{3.12}\hullline{-43.59}{3.90}\hullline{-45.16}{3.90}\hullfill
\hullface{keel!99!black}\hullmove{-40.46}{-3.75}\hullline{-38.90}{-3.75}\hullline{-38.90}{-2.58}\hullline{-38.90}{-1.32}\hullline{-38.90}{-0.00}\hullline{-40.46}{-0.00}\hullline{-40.46}{-1.32}\hullline{-40.46}{-2.58}\hullfill
\hullface{keel!99!black}\hullmove{43.84}{-3.75}\hullline{44.66}{-3.75}\hullline{44.66}{-2.58}\hullline{44.66}{-1.32}\hullline{44.66}{0.00}\hullline{43.84}{0.00}\hullline{43.84}{-1.32}\hullline{43.84}{-2.58}\hullfill
\hullface{topsides!99!black}\hullmove{47.13}{3.12}\hullline{47.96}{3.12}\hullline{47.96}{3.90}\hullline{47.13}{3.90}\hullfill
\hullface{keel!99!black}\hullmove{43.02}{-4.85}\hullline{43.84}{-4.85}\hullline{43.84}{-3.75}\hullline{43.02}{-3.75}\hullfill
\hullface{topsides!99!black}\hullmove{45.49}{2.34}\hullline{46.31}{2.34}\hullline{46.31}{3.12}\hullline{46.31}{3.90}\hullline{45.49}{3.90}\hullline{45.49}{3.12}\hullfill
\hullface{topsides!99!black}\hullmove{46.31}{3.12}\hullline{47.13}{3.12}\hullline{47.13}{3.90}\hullline{46.31}{3.90}\hullfill
\hullface{keel!99!black}\hullmove{-35.77}{-6.77}\hullline{41.37}{-6.77}\hullline{41.37}{-5.86}\hullline{41.37}{-5.44}\hullline{-35.77}{-5.44}\hullline{-35.77}{-5.86}\hullfill
\hullface{keel!99!black}\hullmove{-37.33}{-5.86}\hullline{-35.77}{-5.86}\hullline{-35.77}{-5.44}\hullline{-37.33}{-5.44}\hullfill
\hullface{topsides!99!black}\hullmove{-40.46}{0.00}\hullline{-38.90}{0.00}\hullline{-38.90}{0.78}\hullline{-38.90}{1.56}\hullline{-38.90}{2.34}\hullline{-38.90}{3.12}\hullline{-38.90}{3.90}\hullline{-40.46}{3.90}\hullline{-40.46}{3.12}\hullline{-40.46}{2.34}\hullline{-40.46}{1.56}\hullline{-40.46}{0.78}\hullfill
\hullface{keel!99!black}\hullmove{42.19}{-5.44}\hullline{43.02}{-5.44}\hullline{43.02}{-4.85}\hullline{43.02}{-3.75}\hullline{42.19}{-3.75}\hullline{42.19}{-4.85}\hullfill
\hullface{topsides!99!black}\hullmove{43.84}{0.00}\hullline{44.66}{0.00}\hullline{44.66}{0.78}\hullline{44.66}{1.56}\hullline{44.66}{2.34}\hullline{44.66}{3.12}\hullline{44.66}{3.90}\hullline{43.84}{3.90}\hullline{43.84}{3.12}\hullline{43.84}{2.34}\hullline{43.84}{1.56}\hullline{43.84}{0.78}\hullfill
\hullface{keel!99!black}\hullmove{-38.90}{-4.85}\hullline{-37.33}{-4.85}\hullline{-37.33}{-3.75}\hullline{-37.33}{-2.58}\hullline{-37.33}{-1.32}\hullline{-37.33}{-0.00}\hullline{-38.90}{-0.00}\hullline{-38.90}{-1.32}\hullline{-38.90}{-2.58}\hullline{-38.90}{-3.75}\hullfill
\hullface{keel!99!black}\hullmove{43.02}{-3.75}\hullline{43.84}{-3.75}\hullline{43.84}{-2.58}\hullline{43.84}{-1.32}\hullline{43.84}{0.00}\hullline{43.02}{0.00}\hullline{43.02}{-1.32}\hullline{43.02}{-2.58}\hullfill
\hullface{keel!99!black}\hullmove{41.37}{-5.86}\hullline{42.19}{-5.86}\hullline{42.19}{-5.44}\hullline{41.37}{-5.44}\hullfill
\hullface{topsides!99!black}\hullmove{-42.03}{3.12}\hullline{-40.46}{3.12}\hullline{-40.46}{3.90}\hullline{-42.03}{3.90}\hullfill
\hullface{topsides!99!black}\hullmove{44.66}{3.12}\hullline{45.49}{3.12}\hullline{45.49}{3.90}\hullline{44.66}{3.90}\hullfill
\hullface{topsides!99!black}\hullmove{43.02}{0.00}\hullline{43.84}{0.00}\hullline{43.84}{0.78}\hullline{43.84}{1.56}\hullline{43.84}{2.34}\hullline{43.84}{3.12}\hullline{43.84}{3.90}\hullline{43.02}{3.90}\hullline{43.02}{3.12}\hullline{43.02}{2.34}\hullline{43.02}{1.56}\hullline{43.02}{0.78}\hullfill
\hullface{topsides!99!black}\hullmove{-38.90}{0.00}\hullline{-37.33}{0.00}\hullline{-37.33}{0.78}\hullline{-37.33}{1.56}\hullline{-37.33}{2.34}\hullline{-37.33}{3.12}\hullline{-37.33}{3.90}\hullline{-38.90}{3.90}\hullline{-38.90}{3.12}\hullline{-38.90}{2.34}\hullline{-38.90}{1.56}\hullline{-38.90}{0.78}\hullfill
\hullface{keel!99!black}\hullmove{-37.33}{-5.44}\hullline{-35.77}{-5.44}\hullline{-35.77}{-4.85}\hullline{-35.77}{-3.75}\hullline{-35.77}{-2.58}\hullline{-35.77}{-1.32}\hullline{-35.77}{-0.00}\hullline{-37.33}{-0.00}\hullline{-37.33}{-1.32}\hullline{-37.33}{-2.58}\hullline{-37.33}{-3.75}\hullline{-37.33}{-4.85}\hullfill
\hullface{keel!99!black}\hullmove{42.19}{-3.75}\hullline{43.02}{-3.75}\hullline{43.02}{-2.58}\hullline{43.02}{-1.32}\hullline{43.02}{0.00}\hullline{42.19}{0.00}\hullline{42.19}{-1.32}\hullline{42.19}{-2.58}\hullfill
\hullface{keel!99!black}\hullmove{41.37}{-5.44}\hullline{42.19}{-5.44}\hullline{42.19}{-4.85}\hullline{42.19}{-3.75}\hullline{42.19}{-2.58}\hullline{42.19}{-1.32}\hullline{42.19}{0.00}\hullline{41.37}{0.00}\hullline{41.37}{-1.32}\hullline{41.37}{-2.58}\hullline{41.37}{-3.75}\hullline{41.37}{-4.85}\hullfill
\hullface{topsides!99!black}\hullmove{42.19}{0.00}\hullline{43.02}{0.00}\hullline{43.02}{0.78}\hullline{43.02}{1.56}\hullline{43.02}{2.34}\hullline{43.02}{3.12}\hullline{43.02}{3.90}\hullline{42.19}{3.90}\hullline{42.19}{3.12}\hullline{42.19}{2.34}\hullline{42.19}{1.56}\hullline{42.19}{0.78}\hullfill
\hullface{keel!99!black}\hullmove{-35.77}{-5.44}\hullline{41.37}{-5.44}\hullline{41.37}{-4.85}\hullline{41.37}{-3.75}\hullline{41.37}{-2.58}\hullline{41.37}{-1.32}\hullline{41.37}{0.00}\hullline{-35.77}{-0.00}\hullline{-35.77}{-1.32}\hullline{-35.77}{-2.58}\hullline{-35.77}{-3.75}\hullline{-35.77}{-4.85}\hullfill
\hullface{topsides!99!black}\hullmove{-37.33}{0.00}\hullline{-35.77}{0.00}\hullline{-35.77}{0.78}\hullline{-35.77}{1.56}\hullline{-35.77}{2.34}\hullline{-35.77}{3.12}\hullline{-35.77}{3.90}\hullline{-37.33}{3.90}\hullline{-37.33}{3.12}\hullline{-37.33}{2.34}\hullline{-37.33}{1.56}\hullline{-37.33}{0.78}\hullfill
\hullface{topsides!99!black}\hullmove{41.37}{0.00}\hullline{42.19}{0.00}\hullline{42.19}{0.78}\hullline{42.19}{1.56}\hullline{42.19}{2.34}\hullline{42.19}{3.12}\hullline{42.19}{3.90}\hullline{41.37}{3.90}\hullline{41.37}{3.12}\hullline{41.37}{2.34}\hullline{41.37}{1.56}\hullline{41.37}{0.78}\hullfill
\hullface{ship!99!black}\hullmove{-60.20}{16.57}\hullline{-56.00}{16.57}\hullline{-56.00}{19.37}\hullline{-60.20}{19.37}\hullfill
\hullface{topsides!99!black}\hullmove{-35.77}{0.00}\hullline{41.37}{0.00}\hullline{41.37}{0.78}\hullline{41.37}{1.56}\hullline{41.37}{2.34}\hullline{41.37}{3.12}\hullline{41.37}{3.90}\hullline{-35.77}{3.90}\hullline{-35.77}{3.12}\hullline{-35.77}{2.34}\hullline{-35.77}{1.56}\hullline{-35.77}{0.78}\hullfill
\end{scope}},
  ship above/.pic={\begin{scope}[hull face]\hullface{deck!99!black}\hullmove{-74.90}{-8.15}\hullline{-73.33}{-8.72}\hullline{-71.77}{-9.19}\hullline{-70.20}{-9.57}\hullline{-68.64}{-9.88}\hullline{-67.07}{-10.12}\hullline{-65.51}{-10.31}\hullline{-63.94}{-10.46}\hullline{-62.38}{-10.58}\hullline{-60.81}{-10.66}\hullline{-59.25}{-10.73}\hullline{-57.68}{-10.78}\hullline{-56.12}{-10.81}\hullline{-54.55}{-10.83}\hullline{-52.99}{-10.85}\hullline{-51.42}{-10.86}\hullline{-49.85}{-10.86}\hullline{-48.29}{-10.87}\hullline{-46.72}{-10.87}\hullline{-45.16}{-10.87}\hullline{-43.59}{-10.87}\hullline{-42.03}{-10.87}\hullline{-40.46}{-10.87}\hullline{-38.90}{-10.87}\hullline{-37.33}{-10.87}\hullline{-35.77}{-10.87}\hullline{41.37}{-10.87}\hullline{42.19}{-10.87}\hullline{43.02}{-10.87}\hullline{43.84}{-10.87}\hullline{44.66}{-10.87}\hullline{45.49}{-10.87}\hullline{46.31}{-10.87}\hullline{47.13}{-10.87}\hullline{47.96}{-10.87}\hullline{48.78}{-10.87}\hullline{49.61}{-10.87}\hullline{50.43}{-10.86}\hullline{51.25}{-10.86}\hullline{52.08}{-10.85}\hullline{52.90}{-10.84}\hullline{53.73}{-10.83}\hullline{54.55}{-10.81}\hullline{55.37}{-10.78}\hullline{56.00}{-10.76}\hullline{56.00}{10.76}\hullline{55.37}{10.78}\hullline{54.55}{10.81}\hullline{53.73}{10.83}\hullline{52.90}{10.84}\hullline{52.08}{10.85}\hullline{51.25}{10.86}\hullline{50.43}{10.86}\hullline{49.61}{10.87}\hullline{48.78}{10.87}\hullline{47.96}{10.87}\hullline{47.13}{10.87}\hullline{46.31}{10.87}\hullline{45.49}{10.87}\hullline{44.66}{10.87}\hullline{43.84}{10.87}\hullline{43.02}{10.87}\hullline{42.19}{10.87}\hullline{41.37}{10.87}\hullline{-35.77}{10.87}\hullline{-37.33}{10.87}\hullline{-38.90}{10.87}\hullline{-40.46}{10.87}\hullline{-42.03}{10.87}\hullline{-43.59}{10.87}\hullline{-45.16}{10.87}\hullline{-46.72}{10.87}\hullline{-48.29}{10.87}\hullline{-49.85}{10.86}\hullline{-51.42}{10.86}\hullline{-52.99}{10.85}\hullline{-54.55}{10.83}\hullline{-56.12}{10.81}\hullline{-57.68}{10.78}\hullline{-59.25}{10.73}\hullline{-60.81}{10.66}\hullline{-62.38}{10.58}\hullline{-63.94}{10.46}\hullline{-65.51}{10.31}\hullline{-67.07}{10.12}\hullline{-68.64}{9.88}\hullline{-70.20}{9.57}\hullline{-71.77}{9.19}\hullline{-73.33}{8.72}\hullline{-74.90}{8.15}\hullfill
\hullface{topsides!99!black}\hullmove{-74.90}{7.65}\hullline{-73.33}{8.22}\hullline{-73.33}{8.72}\hullline{-74.90}{8.15}\hullfill
\hullface{topsides!99!black}\hullmove{-74.90}{-8.15}\hullline{-73.33}{-8.72}\hullline{-73.33}{-8.22}\hullline{-74.90}{-7.65}\hullfill
\hullface{topsides!99!black}\hullmove{-73.33}{8.22}\hullline{-71.77}{8.69}\hullline{-71.77}{9.19}\hullline{-73.33}{8.72}\hullfill
\hullface{topsides!99!black}\hullmove{-73.33}{-8.72}\hullline{-71.77}{-9.19}\hullline{-71.77}{-8.69}\hullline{-73.33}{-8.22}\hullfill
\hullface{topsides!99!black}\hullmove{-71.77}{8.69}\hullline{-70.20}{9.07}\hullline{-70.20}{9.57}\hullline{-71.77}{9.19}\hullfill
\hullface{topsides!99!black}\hullmove{-71.77}{-9.19}\hullline{-70.20}{-9.57}\hullline{-70.20}{-9.07}\hullline{-71.77}{-8.69}\hullfill
\hullface{topsides!99!black}\hullmove{-70.20}{9.07}\hullline{-68.64}{9.38}\hullline{-68.64}{9.88}\hullline{-70.20}{9.57}\hullfill
\hullface{topsides!99!black}\hullmove{-70.20}{-9.57}\hullline{-68.64}{-9.88}\hullline{-68.64}{-9.38}\hullline{-70.20}{-9.07}\hullfill
\hullface{topsides!99!black}\hullmove{-68.64}{9.38}\hullline{-67.07}{9.62}\hullline{-67.07}{10.12}\hullline{-68.64}{9.88}\hullfill
\hullface{topsides!99!black}\hullmove{-68.64}{-9.88}\hullline{-67.07}{-10.12}\hullline{-67.07}{-9.62}\hullline{-68.64}{-9.38}\hullfill
\hullface{topsides!99!black}\hullmove{-67.07}{9.62}\hullline{-65.51}{9.81}\hullline{-65.51}{10.31}\hullline{-67.07}{10.12}\hullfill
\hullface{topsides!99!black}\hullmove{-67.07}{-10.12}\hullline{-65.51}{-10.31}\hullline{-65.51}{-9.81}\hullline{-67.07}{-9.62}\hullfill
\hullface{topsides!99!black}\hullmove{-65.51}{9.81}\hullline{-63.94}{9.96}\hullline{-63.94}{10.46}\hullline{-65.51}{10.31}\hullfill
\hullface{topsides!99!black}\hullmove{-65.51}{-10.31}\hullline{-63.94}{-10.46}\hullline{-63.94}{-9.96}\hullline{-65.51}{-9.81}\hullfill
\hullface{topsides!99!black}\hullmove{-63.94}{9.96}\hullline{-62.38}{10.08}\hullline{-62.38}{10.58}\hullline{-63.94}{10.46}\hullfill
\hullface{topsides!99!black}\hullmove{-63.94}{-10.46}\hullline{-62.38}{-10.58}\hullline{-62.38}{-10.08}\hullline{-63.94}{-9.96}\hullfill
\hullface{topsides!99!black}\hullmove{-62.38}{10.08}\hullline{-60.81}{10.16}\hullline{-60.81}{10.66}\hullline{-62.38}{10.58}\hullfill
\hullface{topsides!99!black}\hullmove{-62.38}{-10.58}\hullline{-60.81}{-10.66}\hullline{-60.81}{-10.16}\hullline{-62.38}{-10.08}\hullfill
\hullface{topsides!99!black}\hullmove{-60.81}{10.16}\hullline{-59.25}{10.23}\hullline{-59.25}{10.73}\hullline{-60.81}{10.66}\hullfill
\hullface{topsides!99!black}\hullmove{-60.81}{-10.66}\hullline{-59.25}{-10.73}\hullline{-59.25}{-10.23}\hullline{-60.81}{-10.16}\hullfill
\hullface{topsides!99!black}\hullmove{-56.12}{10.31}\hullline{-54.55}{10.33}\hullline{-54.55}{10.83}\hullline{-56.12}{10.81}\hullfill
\hullface{topsides!99!black}\hullmove{-56.12}{-10.81}\hullline{-54.55}{-10.83}\hullline{-54.55}{-10.33}\hullline{-56.12}{-10.31}\hullfill
\hullface{topsides!99!black}\hullmove{-54.55}{10.33}\hullline{-52.99}{10.35}\hullline{-52.99}{10.85}\hullline{-54.55}{10.83}\hullfill
\hullface{topsides!99!black}\hullmove{-54.55}{-10.83}\hullline{-52.99}{-10.85}\hullline{-52.99}{-10.35}\hullline{-54.55}{-10.33}\hullfill
\hullface{topsides!99!black}\hullmove{-52.99}{10.35}\hullline{-51.42}{10.36}\hullline{-51.42}{10.86}\hullline{-52.99}{10.85}\hullfill
\hullface{topsides!99!black}\hullmove{-52.99}{-10.85}\hullline{-51.42}{-10.86}\hullline{-51.42}{-10.36}\hullline{-52.99}{-10.35}\hullfill
\hullface{topsides!99!black}\hullmove{-51.42}{10.36}\hullline{-49.85}{10.36}\hullline{-49.85}{10.86}\hullline{-51.42}{10.86}\hullfill
\hullface{topsides!99!black}\hullmove{-51.42}{-10.86}\hullline{-49.85}{-10.86}\hullline{-49.85}{-10.36}\hullline{-51.42}{-10.36}\hullfill
\hullface{topsides!99!black}\hullmove{-49.85}{10.36}\hullline{-48.29}{10.37}\hullline{-48.29}{10.87}\hullline{-49.85}{10.86}\hullfill
\hullface{topsides!99!black}\hullmove{-49.85}{-10.86}\hullline{-48.29}{-10.87}\hullline{-48.29}{-10.37}\hullline{-49.85}{-10.36}\hullfill
\hullface{topsides!99!black}\hullmove{-48.29}{10.37}\hullline{-46.72}{10.37}\hullline{-46.72}{10.87}\hullline{-48.29}{10.87}\hullfill
\hullface{topsides!99!black}\hullmove{-48.29}{-10.87}\hullline{-46.72}{-10.87}\hullline{-46.72}{-10.37}\hullline{-48.29}{-10.37}\hullfill
\hullface{topsides!99!black}\hullmove{-46.72}{10.37}\hullline{-45.16}{10.37}\hullline{-45.16}{10.87}\hullline{-46.72}{10.87}\hullfill
\hullface{topsides!99!black}\hullmove{-46.72}{-10.87}\hullline{-45.16}{-10.87}\hullline{-45.16}{-10.37}\hullline{-46.72}{-10.37}\hullfill
\hullface{topsides!99!black}\hullmove{-45.16}{10.37}\hullline{-43.59}{10.37}\hullline{-43.59}{10.87}\hullline{-45.16}{10.87}\hullfill
\hullface{topsides!99!black}\hullmove{-45.16}{-10.87}\hullline{-43.59}{-10.87}\hullline{-43.59}{-10.37}\hullline{-45.16}{-10.37}\hullfill
\hullface{topsides!99!black}\hullmove{-43.59}{10.37}\hullline{-42.03}{10.37}\hullline{-42.03}{10.87}\hullline{-43.59}{10.87}\hullfill
\hullface{topsides!99!black}\hullmove{-43.59}{-10.87}\hullline{-42.03}{-10.87}\hullline{-42.03}{-10.37}\hullline{-43.59}{-10.37}\hullfill
\hullface{topsides!99!black}\hullmove{-42.03}{10.37}\hullline{-40.46}{10.37}\hullline{-40.46}{10.87}\hullline{-42.03}{10.87}\hullfill
\hullface{topsides!99!black}\hullmove{-42.03}{-10.87}\hullline{-40.46}{-10.87}\hullline{-40.46}{-10.37}\hullline{-42.03}{-10.37}\hullfill
\hullface{topsides!99!black}\hullmove{-40.46}{10.37}\hullline{-38.90}{10.37}\hullline{-38.90}{10.87}\hullline{-40.46}{10.87}\hullfill
\hullface{topsides!99!black}\hullmove{-40.46}{-10.87}\hullline{-38.90}{-10.87}\hullline{-38.90}{-10.37}\hullline{-40.46}{-10.37}\hullfill
\hullface{topsides!99!black}\hullmove{-38.90}{10.37}\hullline{-37.33}{10.37}\hullline{-37.33}{10.87}\hullline{-38.90}{10.87}\hullfill
\hullface{topsides!99!black}\hullmove{-38.90}{-10.87}\hullline{-37.33}{-10.87}\hullline{-37.33}{-10.37}\hullline{-38.90}{-10.37}\hullfill
\hullface{topsides!99!black}\hullmove{-37.33}{10.37}\hullline{-35.77}{10.37}\hullline{-35.77}{10.87}\hullline{-37.33}{10.87}\hullfill
\hullface{topsides!99!black}\hullmove{-37.33}{-10.87}\hullline{-35.77}{-10.87}\hullline{-35.77}{-10.37}\hullline{-37.33}{-10.37}\hullfill
\hullface{topsides!99!black}\hullmove{-35.77}{10.37}\hullline{41.37}{10.37}\hullline{41.37}{10.87}\hullline{-35.77}{10.87}\hullfill
\hullface{topsides!99!black}\hullmove{-35.77}{-10.87}\hullline{41.37}{-10.87}\hullline{41.37}{-10.37}\hullline{-35.77}{-10.37}\hullfill
\hullface{topsides!99!black}\hullmove{41.37}{10.37}\hullline{42.19}{10.37}\hullline{42.19}{10.87}\hullline{41.37}{10.87}\hullfill
\hullface{topsides!99!black}\hullmove{41.37}{-10.87}\hullline{42.19}{-10.87}\hullline{42.19}{-10.37}\hullline{41.37}{-10.37}\hullfill
\hullface{topsides!99!black}\hullmove{42.19}{10.37}\hullline{43.02}{10.37}\hullline{43.02}{10.87}\hullline{42.19}{10.87}\hullfill
\hullface{topsides!99!black}\hullmove{42.19}{-10.87}\hullline{43.02}{-10.87}\hullline{43.02}{-10.37}\hullline{42.19}{-10.37}\hullfill
\hullface{topsides!99!black}\hullmove{43.02}{10.37}\hullline{43.84}{10.37}\hullline{43.84}{10.87}\hullline{43.02}{10.87}\hullfill
\hullface{topsides!99!black}\hullmove{43.02}{-10.87}\hullline{43.84}{-10.87}\hullline{43.84}{-10.37}\hullline{43.02}{-10.37}\hullfill
\hullface{topsides!99!black}\hullmove{43.84}{10.37}\hullline{44.66}{10.37}\hullline{44.66}{10.87}\hullline{43.84}{10.87}\hullfill
\hullface{topsides!99!black}\hullmove{43.84}{-10.87}\hullline{44.66}{-10.87}\hullline{44.66}{-10.37}\hullline{43.84}{-10.37}\hullfill
\hullface{topsides!99!black}\hullmove{44.66}{10.37}\hullline{45.49}{10.37}\hullline{45.49}{10.87}\hullline{44.66}{10.87}\hullfill
\hullface{topsides!99!black}\hullmove{44.66}{-10.87}\hullline{45.49}{-10.87}\hullline{45.49}{-10.37}\hullline{44.66}{-10.37}\hullfill
\hullface{topsides!99!black}\hullmove{45.49}{10.37}\hullline{46.31}{10.37}\hullline{46.31}{10.87}\hullline{45.49}{10.87}\hullfill
\hullface{topsides!99!black}\hullmove{45.49}{-10.87}\hullline{46.31}{-10.87}\hullline{46.31}{-10.37}\hullline{45.49}{-10.37}\hullfill
\hullface{topsides!99!black}\hullmove{46.31}{10.37}\hullline{47.13}{10.37}\hullline{47.13}{10.87}\hullline{46.31}{10.87}\hullfill
\hullface{topsides!99!black}\hullmove{46.31}{-10.87}\hullline{47.13}{-10.87}\hullline{47.13}{-10.37}\hullline{46.31}{-10.37}\hullfill
\hullface{topsides!99!black}\hullmove{47.13}{10.37}\hullline{47.96}{10.37}\hullline{47.96}{10.87}\hullline{47.13}{10.87}\hullfill
\hullface{topsides!99!black}\hullmove{47.13}{-10.87}\hullline{47.96}{-10.87}\hullline{47.96}{-10.37}\hullline{47.13}{-10.37}\hullfill
\hullface{topsides!99!black}\hullmove{47.96}{10.37}\hullline{48.78}{10.37}\hullline{48.78}{10.87}\hullline{47.96}{10.87}\hullfill
\hullface{topsides!99!black}\hullmove{47.96}{-10.87}\hullline{48.78}{-10.87}\hullline{48.78}{-10.37}\hullline{47.96}{-10.37}\hullfill
\hullface{topsides!99!black}\hullmove{48.78}{10.37}\hullline{49.61}{10.37}\hullline{49.61}{10.87}\hullline{48.78}{10.87}\hullfill
\hullface{topsides!99!black}\hullmove{48.78}{-10.87}\hullline{49.61}{-10.87}\hullline{49.61}{-10.37}\hullline{48.78}{-10.37}\hullfill
\hullface{topsides!99!black}\hullmove{49.61}{10.37}\hullline{50.43}{10.36}\hullline{50.43}{10.86}\hullline{49.61}{10.87}\hullfill
\hullface{topsides!99!black}\hullmove{49.61}{-10.87}\hullline{50.43}{-10.86}\hullline{50.43}{-10.36}\hullline{49.61}{-10.37}\hullfill
\hullface{topsides!99!black}\hullmove{50.43}{10.36}\hullline{51.25}{10.36}\hullline{51.25}{10.86}\hullline{50.43}{10.86}\hullfill
\hullface{topsides!99!black}\hullmove{50.43}{-10.86}\hullline{51.25}{-10.86}\hullline{51.25}{-10.36}\hullline{50.43}{-10.36}\hullfill
\hullface{topsides!99!black}\hullmove{51.25}{10.36}\hullline{52.08}{10.35}\hullline{52.08}{10.85}\hullline{51.25}{10.86}\hullfill
\hullface{topsides!99!black}\hullmove{51.25}{-10.86}\hullline{52.08}{-10.85}\hullline{52.08}{-10.35}\hullline{51.25}{-10.36}\hullfill
\hullface{topsides!99!black}\hullmove{52.08}{10.35}\hullline{52.90}{10.34}\hullline{52.90}{10.84}\hullline{52.08}{10.85}\hullfill
\hullface{topsides!99!black}\hullmove{52.08}{-10.85}\hullline{52.90}{-10.84}\hullline{52.90}{-10.34}\hullline{52.08}{-10.35}\hullfill
\hullface{topsides!99!black}\hullmove{52.90}{10.34}\hullline{53.73}{10.33}\hullline{53.73}{10.83}\hullline{52.90}{10.84}\hullfill
\hullface{topsides!99!black}\hullmove{52.90}{-10.84}\hullline{53.73}{-10.83}\hullline{53.73}{-10.33}\hullline{52.90}{-10.34}\hullfill
\hullface{topsides!99!black}\hullmove{53.73}{10.33}\hullline{54.55}{10.31}\hullline{54.55}{10.81}\hullline{53.73}{10.83}\hullfill
\hullface{topsides!99!black}\hullmove{53.73}{-10.83}\hullline{54.55}{-10.81}\hullline{54.55}{-10.31}\hullline{53.73}{-10.33}\hullfill
\hullface{topsides!99!black}\hullmove{54.55}{10.31}\hullline{55.37}{10.28}\hullline{55.37}{10.78}\hullline{54.55}{10.81}\hullfill
\hullface{topsides!99!black}\hullmove{54.55}{-10.81}\hullline{55.37}{-10.78}\hullline{55.37}{-10.28}\hullline{54.55}{-10.31}\hullfill
\hullface{topsides!99!black}\hullmove{55.37}{10.28}\hullline{56.00}{10.26}\hullline{56.00}{10.76}\hullline{55.37}{10.78}\hullfill
\hullface{topsides!99!black}\hullmove{55.37}{-10.78}\hullline{56.00}{-10.76}\hullline{56.00}{-10.26}\hullline{55.37}{-10.28}\hullfill
\hullface{deck!99!black}\hullmove{56.00}{-10.76}\hullline{56.20}{-10.75}\hullline{57.02}{-10.70}\hullline{57.85}{-10.64}\hullline{58.67}{-10.57}\hullline{59.49}{-10.47}\hullline{60.32}{-10.36}\hullline{61.14}{-10.21}\hullline{61.97}{-10.04}\hullline{62.79}{-9.83}\hullline{63.61}{-9.57}\hullline{64.44}{-9.27}\hullline{65.26}{-8.91}\hullline{66.08}{-8.49}\hullline{66.91}{-7.99}\hullline{67.73}{-7.41}\hullline{68.56}{-6.73}\hullline{69.38}{-5.95}\hullline{70.20}{-5.06}\hullline{71.03}{-4.03}\hullline{71.85}{-2.85}\hullline{72.68}{-1.52}\hullline{73.50}{0.00}\hullline{73.50}{0.00}\hullline{72.68}{1.52}\hullline{71.85}{2.85}\hullline{71.03}{4.03}\hullline{70.20}{5.06}\hullline{69.38}{5.95}\hullline{68.56}{6.73}\hullline{67.73}{7.41}\hullline{66.91}{7.99}\hullline{66.08}{8.49}\hullline{65.26}{8.91}\hullline{64.44}{9.27}\hullline{63.61}{9.57}\hullline{62.79}{9.83}\hullline{61.97}{10.04}\hullline{61.14}{10.21}\hullline{60.32}{10.36}\hullline{59.49}{10.47}\hullline{58.67}{10.57}\hullline{57.85}{10.64}\hullline{57.02}{10.70}\hullline{56.20}{10.75}\hullline{56.00}{10.76}\hullfill
\hullface{ship!99!black}\hullmove{-42.00}{-8.15}\hullline{-26.93}{-8.15}\hullline{-26.93}{8.15}\hullline{-42.00}{8.15}\hullfill
\hullface{ship!99!black}\hullmove{-26.37}{-8.15}\hullline{-11.29}{-8.15}\hullline{-11.29}{8.15}\hullline{-26.37}{8.15}\hullfill
\hullface{ship!99!black}\hullmove{-10.73}{-8.15}\hullline{4.34}{-8.15}\hullline{4.34}{8.15}\hullline{-10.73}{8.15}\hullfill
\hullface{ship!99!black}\hullmove{4.90}{-8.15}\hullline{19.97}{-8.15}\hullline{19.97}{8.15}\hullline{4.90}{8.15}\hullfill
\hullface{ship!99!black}\hullmove{20.53}{-8.15}\hullline{35.61}{-8.15}\hullline{35.61}{8.15}\hullline{20.53}{8.15}\hullfill
\hullface{ship!99!black}\hullmove{36.17}{-8.15}\hullline{51.24}{-8.15}\hullline{51.24}{8.15}\hullline{36.17}{8.15}\hullfill
\hullface{topsides!99!black}\hullmove{56.00}{10.26}\hullline{56.20}{10.25}\hullline{56.20}{10.75}\hullline{56.00}{10.76}\hullfill
\hullface{topsides!99!black}\hullmove{56.00}{-10.76}\hullline{56.20}{-10.75}\hullline{56.20}{-10.25}\hullline{56.00}{-10.26}\hullfill
\hullface{topsides!99!black}\hullmove{56.20}{10.25}\hullline{57.02}{10.20}\hullline{57.02}{10.70}\hullline{56.20}{10.75}\hullfill
\hullface{topsides!99!black}\hullmove{56.20}{-10.75}\hullline{57.02}{-10.70}\hullline{57.02}{-10.20}\hullline{56.20}{-10.25}\hullfill
\hullface{topsides!99!black}\hullmove{57.02}{10.20}\hullline{57.85}{10.14}\hullline{57.85}{10.64}\hullline{57.02}{10.70}\hullfill
\hullface{topsides!99!black}\hullmove{57.02}{-10.70}\hullline{57.85}{-10.64}\hullline{57.85}{-10.14}\hullline{57.02}{-10.20}\hullfill
\hullface{topsides!99!black}\hullmove{57.85}{10.14}\hullline{58.67}{10.07}\hullline{58.67}{10.57}\hullline{57.85}{10.64}\hullfill
\hullface{topsides!99!black}\hullmove{57.85}{-10.64}\hullline{58.67}{-10.57}\hullline{58.67}{-10.07}\hullline{57.85}{-10.14}\hullfill
\hullface{topsides!99!black}\hullmove{58.67}{10.07}\hullline{59.49}{9.97}\hullline{59.49}{10.47}\hullline{58.67}{10.57}\hullfill
\hullface{topsides!99!black}\hullmove{58.67}{-10.57}\hullline{59.49}{-10.47}\hullline{59.49}{-9.97}\hullline{58.67}{-10.07}\hullfill
\hullface{topsides!99!black}\hullmove{59.49}{9.97}\hullline{60.32}{9.86}\hullline{60.32}{10.36}\hullline{59.49}{10.47}\hullfill
\hullface{topsides!99!black}\hullmove{59.49}{-10.47}\hullline{60.32}{-10.36}\hullline{60.32}{-9.86}\hullline{59.49}{-9.97}\hullfill
\hullface{topsides!99!black}\hullmove{60.32}{9.86}\hullline{61.14}{9.71}\hullline{61.14}{10.21}\hullline{60.32}{10.36}\hullfill
\hullface{topsides!99!black}\hullmove{60.32}{-10.36}\hullline{61.14}{-10.21}\hullline{61.14}{-9.71}\hullline{60.32}{-9.86}\hullfill
\hullface{topsides!99!black}\hullmove{61.14}{9.71}\hullline{61.97}{9.54}\hullline{61.97}{10.04}\hullline{61.14}{10.21}\hullfill
\hullface{topsides!99!black}\hullmove{61.14}{-10.21}\hullline{61.97}{-10.04}\hullline{61.97}{-9.54}\hullline{61.14}{-9.71}\hullfill
\hullface{topsides!99!black}\hullmove{61.97}{9.54}\hullline{62.79}{9.33}\hullline{62.79}{9.83}\hullline{61.97}{10.04}\hullfill
\hullface{topsides!99!black}\hullmove{61.97}{-10.04}\hullline{62.79}{-9.83}\hullline{62.79}{-9.33}\hullline{61.97}{-9.54}\hullfill
\hullface{topsides!99!black}\hullmove{62.79}{9.33}\hullline{63.61}{9.07}\hullline{63.61}{9.57}\hullline{62.79}{9.83}\hullfill
\hullface{topsides!99!black}\hullmove{62.79}{-9.83}\hullline{63.61}{-9.57}\hullline{63.61}{-9.07}\hullline{62.79}{-9.33}\hullfill
\hullface{topsides!99!black}\hullmove{63.61}{9.07}\hullline{64.44}{8.77}\hullline{64.44}{9.27}\hullline{63.61}{9.57}\hullfill
\hullface{topsides!99!black}\hullmove{63.61}{-9.57}\hullline{64.44}{-9.27}\hullline{64.44}{-8.77}\hullline{63.61}{-9.07}\hullfill
\hullface{topsides!99!black}\hullmove{64.44}{8.77}\hullline{65.26}{8.41}\hullline{65.26}{8.91}\hullline{64.44}{9.27}\hullfill
\hullface{topsides!99!black}\hullmove{64.44}{-9.27}\hullline{65.26}{-8.91}\hullline{65.26}{-8.41}\hullline{64.44}{-8.77}\hullfill
\hullface{topsides!99!black}\hullmove{65.26}{8.41}\hullline{66.08}{7.99}\hullline{66.08}{8.49}\hullline{65.26}{8.91}\hullfill
\hullface{topsides!99!black}\hullmove{65.26}{-8.91}\hullline{66.08}{-8.49}\hullline{66.08}{-7.99}\hullline{65.26}{-8.41}\hullfill
\hullface{topsides!99!black}\hullmove{66.08}{7.99}\hullline{66.91}{7.49}\hullline{66.91}{7.99}\hullline{66.08}{8.49}\hullfill
\hullface{topsides!99!black}\hullmove{66.08}{-8.49}\hullline{66.91}{-7.99}\hullline{66.91}{-7.49}\hullline{66.08}{-7.99}\hullfill
\hullface{topsides!99!black}\hullmove{66.91}{7.49}\hullline{67.73}{6.91}\hullline{67.73}{7.41}\hullline{66.91}{7.99}\hullfill
\hullface{topsides!99!black}\hullmove{66.91}{-7.99}\hullline{67.73}{-7.41}\hullline{67.73}{-6.91}\hullline{66.91}{-7.49}\hullfill
\hullface{topsides!99!black}\hullmove{67.73}{6.91}\hullline{68.56}{6.23}\hullline{68.56}{6.73}\hullline{67.73}{7.41}\hullfill
\hullface{topsides!99!black}\hullmove{67.73}{-7.41}\hullline{68.56}{-6.73}\hullline{68.56}{-6.23}\hullline{67.73}{-6.91}\hullfill
\hullface{topsides!99!black}\hullmove{68.56}{6.23}\hullline{69.38}{5.45}\hullline{69.38}{5.95}\hullline{68.56}{6.73}\hullfill
\hullface{topsides!99!black}\hullmove{68.56}{-6.73}\hullline{69.38}{-5.95}\hullline{69.38}{-5.45}\hullline{68.56}{-6.23}\hullfill
\hullface{topsides!99!black}\hullmove{69.38}{5.45}\hullline{70.20}{4.56}\hullline{70.20}{5.06}\hullline{69.38}{5.95}\hullfill
\hullface{topsides!99!black}\hullmove{69.38}{-5.95}\hullline{70.20}{-5.06}\hullline{70.20}{-4.56}\hullline{69.38}{-5.45}\hullfill
\hullface{topsides!99!black}\hullmove{70.20}{4.56}\hullline{71.03}{3.53}\hullline{71.03}{4.03}\hullline{70.20}{5.06}\hullfill
\hullface{topsides!99!black}\hullmove{70.20}{-5.06}\hullline{71.03}{-4.03}\hullline{71.03}{-3.53}\hullline{70.20}{-4.56}\hullfill
\hullface{topsides!99!black}\hullmove{71.03}{3.53}\hullline{71.85}{2.35}\hullline{71.85}{2.85}\hullline{71.03}{4.03}\hullfill
\hullface{topsides!99!black}\hullmove{71.03}{-4.03}\hullline{71.85}{-2.85}\hullline{71.85}{-2.35}\hullline{71.03}{-3.53}\hullfill
\hullface{topsides!99!black}\hullmove{71.85}{2.35}\hullline{72.68}{1.02}\hullline{72.68}{1.52}\hullline{71.85}{2.85}\hullfill
\hullface{topsides!99!black}\hullmove{71.85}{-2.85}\hullline{72.68}{-1.52}\hullline{72.68}{-1.02}\hullline{71.85}{-2.35}\hullfill
\hullface{topsides!99!black}\hullmove{72.68}{1.02}\hullline{73.50}{0.00}\hullline{73.50}{0.00}\hullline{72.68}{1.52}\hullfill
\hullface{topsides!99!black}\hullmove{72.68}{-1.52}\hullline{73.50}{0.00}\hullline{73.50}{0.00}\hullline{72.68}{-1.02}\hullfill
\hullface{ship!99!black}\hullmove{-67.20}{-9.09}\hullline{-56.00}{-9.09}\hullline{-56.00}{9.09}\hullline{-67.20}{9.09}\hullfill
\hullface{ship!99!black}\hullmove{-60.20}{-10.87}\hullline{-56.00}{-10.87}\hullline{-56.00}{10.87}\hullline{-60.20}{10.87}\hullfill
\end{scope}},
  ship ahead/.pic={
    \begin{scope}[hull face]\colorlet{keel}{topsides}\hullface{ship!99!black}\hullmove{-9.09}{2.90}\hullline{9.09}{2.90}\hullline{9.09}{16.57}\hullline{-9.09}{16.57}\hullfill
\hullface{ship!99!black}\hullmove{-10.87}{16.57}\hullline{10.87}{16.57}\hullline{10.87}{19.37}\hullline{-10.87}{19.37}\hullfill
\hullface{keel!57!black}\hullmove{10.87}{-4.85}\hullline{10.87}{-3.75}\hullline{10.87}{-2.58}\hullline{10.87}{-1.32}\hullline{10.87}{0.00}\hullline{10.87}{0.00}\hullline{10.87}{-1.32}\hullline{10.87}{-2.58}\hullline{10.87}{-3.75}\hullline{10.87}{-4.85}\hullline{10.87}{-5.44}\hullline{10.87}{-5.44}\hullfill
\hullface{topsides!57!black}\hullmove{10.87}{0.78}\hullline{10.87}{1.56}\hullline{10.87}{2.34}\hullline{10.87}{3.12}\hullline{10.87}{3.90}\hullline{10.87}{3.90}\hullline{10.87}{3.12}\hullline{10.87}{2.34}\hullline{10.87}{1.56}\hullline{10.87}{0.78}\hullline{10.87}{0.00}\hullline{10.87}{0.00}\hullfill
\hullface{keel!57!black}\hullmove{10.42}{-8.77}\hullline{9.69}{-8.77}\hullline{0.00}{-9.06}\hullline{0.00}{-9.06}\hullfill
\hullface{keel!57!black}\hullmove{0.00}{-9.06}\hullline{0.00}{-9.06}\hullline{-9.69}{-8.77}\hullline{-10.42}{-8.77}\hullfill
\hullface{keel!57!black}\hullmove{10.87}{-8.25}\hullline{10.85}{-8.25}\hullline{9.69}{-8.77}\hullline{10.42}{-8.77}\hullfill
\hullface{keel!57!black}\hullmove{-10.42}{-8.77}\hullline{-9.69}{-8.77}\hullline{-10.85}{-8.25}\hullline{-10.87}{-8.25}\hullfill
\hullface{keel!57!black}\hullmove{10.87}{-7.57}\hullline{10.87}{-6.77}\hullline{10.87}{-5.86}\hullline{10.87}{-5.44}\hullline{10.87}{-5.44}\hullline{10.87}{-5.86}\hullline{10.87}{-6.77}\hullline{10.87}{-7.57}\hullline{10.85}{-8.25}\hullline{10.87}{-8.25}\hullfill
\hullface{keel!57!black}\hullmove{9.69}{-8.77}\hullline{8.24}{-8.77}\hullline{0.00}{-9.06}\hullline{0.00}{-9.06}\hullfill
\hullface{keel!57!black}\hullmove{0.00}{-9.06}\hullline{0.00}{-9.06}\hullline{-8.24}{-8.77}\hullline{-9.69}{-8.77}\hullfill
\hullface{keel!57!black}\hullmove{10.85}{-8.25}\hullline{10.69}{-8.25}\hullline{8.24}{-8.77}\hullline{9.69}{-8.77}\hullfill
\hullface{keel!57!black}\hullmove{-9.69}{-8.77}\hullline{-8.24}{-8.77}\hullline{-10.69}{-8.25}\hullline{-10.85}{-8.25}\hullfill
\hullface{keel!57!black}\hullmove{10.87}{-7.57}\hullline{10.86}{-7.57}\hullline{10.69}{-8.25}\hullline{10.85}{-8.25}\hullfill
\hullface{keel!81!black}\hullmove{-10.85}{-8.25}\hullline{-10.69}{-8.25}\hullline{-10.86}{-7.57}\hullline{-10.87}{-7.57}\hullfill
\hullface{keel!57!black}\hullmove{10.87}{-6.77}\hullline{10.87}{-5.86}\hullline{10.87}{-5.44}\hullline{10.87}{-5.44}\hullline{10.87}{-5.86}\hullline{10.87}{-6.77}\hullline{10.86}{-7.57}\hullline{10.87}{-7.57}\hullfill
\hullface{keel!57!black}\hullmove{10.87}{-4.85}\hullline{10.87}{-3.75}\hullline{10.87}{-2.58}\hullline{10.87}{-1.32}\hullline{10.87}{0.00}\hullline{10.87}{0.00}\hullline{10.87}{-1.32}\hullline{10.87}{-2.58}\hullline{10.87}{-3.75}\hullline{10.87}{-4.85}\hullline{10.87}{-5.44}\hullline{10.87}{-5.44}\hullfill
\hullface{topsides!57!black}\hullmove{10.87}{0.78}\hullline{10.87}{1.56}\hullline{10.87}{2.34}\hullline{10.87}{3.12}\hullline{10.87}{3.90}\hullline{10.87}{3.90}\hullline{10.87}{3.12}\hullline{10.87}{2.34}\hullline{10.87}{1.56}\hullline{10.87}{0.78}\hullline{10.87}{0.00}\hullline{10.87}{0.00}\hullfill
\hullface{keel!57!black}\hullmove{8.24}{-8.77}\hullline{6.72}{-8.77}\hullline{0.00}{-9.06}\hullline{0.00}{-9.06}\hullfill
\hullface{keel!57!black}\hullmove{0.00}{-9.06}\hullline{0.00}{-9.06}\hullline{-6.72}{-8.77}\hullline{-8.24}{-8.77}\hullfill
\hullface{keel!57!black}\hullmove{10.69}{-8.25}\hullline{10.21}{-8.25}\hullline{6.72}{-8.77}\hullline{8.24}{-8.77}\hullfill
\hullface{keel!57!black}\hullmove{-8.24}{-8.77}\hullline{-6.72}{-8.77}\hullline{-10.21}{-8.25}\hullline{-10.69}{-8.25}\hullfill
\hullface{keel!57!black}\hullmove{10.86}{-7.57}\hullline{10.81}{-7.57}\hullline{10.21}{-8.25}\hullline{10.69}{-8.25}\hullfill
\hullface{keel!72!black}\hullmove{-10.69}{-8.25}\hullline{-10.21}{-8.25}\hullline{-10.81}{-7.57}\hullline{-10.86}{-7.57}\hullfill
\hullface{keel!81!black}\hullmove{-10.86}{-7.57}\hullline{-10.81}{-7.57}\hullline{-10.86}{-6.77}\hullline{-10.87}{-6.77}\hullfill
\hullface{keel!57!black}\hullmove{10.87}{-5.86}\hullline{10.87}{-5.44}\hullline{10.86}{-5.44}\hullline{10.86}{-5.86}\hullline{10.86}{-6.77}\hullline{10.87}{-6.77}\hullfill
\hullface{topsides!57!black}\hullmove{10.87}{0.78}\hullline{10.87}{1.56}\hullline{10.87}{2.34}\hullline{10.87}{3.12}\hullline{10.87}{3.90}\hullline{10.87}{3.90}\hullline{10.87}{3.12}\hullline{10.87}{2.34}\hullline{10.86}{1.56}\hullline{10.86}{0.78}\hullline{10.86}{0.00}\hullline{10.87}{0.00}\hullfill
\hullface{keel!57!black}\hullmove{6.72}{-8.77}\hullline{5.45}{-8.77}\hullline{0.00}{-9.06}\hullline{0.00}{-9.06}\hullfill
\hullface{keel!57!black}\hullmove{0.00}{-9.06}\hullline{0.00}{-9.06}\hullline{-5.45}{-8.77}\hullline{-6.72}{-8.77}\hullfill
\hullface{keel!57!black}\hullmove{10.21}{-8.25}\hullline{9.43}{-8.25}\hullline{5.45}{-8.77}\hullline{6.72}{-8.77}\hullfill
\hullface{keel!57!black}\hullmove{-6.72}{-8.77}\hullline{-5.45}{-8.77}\hullline{-9.43}{-8.25}\hullline{-10.21}{-8.25}\hullfill
\hullface{keel!57!black}\hullmove{10.81}{-7.57}\hullline{10.63}{-7.57}\hullline{9.43}{-8.25}\hullline{10.21}{-8.25}\hullfill
\hullface{keel!60!black}\hullmove{-10.21}{-8.25}\hullline{-9.43}{-8.25}\hullline{-10.63}{-7.57}\hullline{-10.81}{-7.57}\hullfill
\hullface{keel!57!black}\hullmove{10.86}{-6.77}\hullline{10.83}{-6.77}\hullline{10.63}{-7.57}\hullline{10.81}{-7.57}\hullfill
\hullface{keel!81!black}\hullmove{-10.81}{-7.57}\hullline{-10.63}{-7.57}\hullline{-10.83}{-6.77}\hullline{-10.86}{-6.77}\hullfill
\hullface{keel!57!black}\hullmove{5.45}{-8.77}\hullline{4.46}{-8.77}\hullline{0.00}{-9.06}\hullline{0.00}{-9.06}\hullfill
\hullface{keel!57!black}\hullmove{0.00}{-9.06}\hullline{0.00}{-9.06}\hullline{-4.46}{-8.77}\hullline{-5.45}{-8.77}\hullfill
\hullface{keel!57!black}\hullmove{9.43}{-8.25}\hullline{8.53}{-8.25}\hullline{4.46}{-8.77}\hullline{5.45}{-8.77}\hullfill
\hullface{keel!57!black}\hullmove{-5.45}{-8.77}\hullline{-4.46}{-8.77}\hullline{-8.53}{-8.25}\hullline{-9.43}{-8.25}\hullfill
\hullface{keel!57!black}\hullmove{10.63}{-7.57}\hullline{10.27}{-7.57}\hullline{8.53}{-8.25}\hullline{9.43}{-8.25}\hullfill
\hullface{keel!57!black}\hullmove{-9.43}{-8.25}\hullline{-8.53}{-8.25}\hullline{-10.27}{-7.57}\hullline{-10.63}{-7.57}\hullfill
\hullface{keel!57!black}\hullmove{10.83}{-6.77}\hullline{10.74}{-6.77}\hullline{10.27}{-7.57}\hullline{10.63}{-7.57}\hullfill
\hullface{keel!75!black}\hullmove{-10.63}{-7.57}\hullline{-10.27}{-7.57}\hullline{-10.74}{-6.77}\hullline{-10.83}{-6.77}\hullfill
\hullface{keel!57!black}\hullmove{10.85}{-5.86}\hullline{10.82}{-5.86}\hullline{10.74}{-6.77}\hullline{10.83}{-6.77}\hullfill
\hullface{keel!81!black}\hullmove{-10.83}{-6.77}\hullline{-10.74}{-6.77}\hullline{-10.82}{-5.86}\hullline{-10.85}{-5.86}\hullfill
\hullface{keel!57!black}\hullmove{10.85}{-5.44}\hullline{10.83}{-5.44}\hullline{10.82}{-5.86}\hullline{10.85}{-5.86}\hullfill
\hullface{topsides!57!black}\hullmove{10.87}{3.90}\hullline{10.87}{3.90}\hullline{10.86}{3.12}\hullline{10.86}{3.12}\hullfill
\hullface{keel!57!black}\hullmove{4.46}{-8.77}\hullline{3.70}{-8.77}\hullline{0.00}{-9.06}\hullline{0.00}{-9.06}\hullfill
\hullface{keel!57!black}\hullmove{0.00}{-9.06}\hullline{0.00}{-9.06}\hullline{-3.70}{-8.77}\hullline{-4.46}{-8.77}\hullfill
\hullface{keel!57!black}\hullmove{8.53}{-8.25}\hullline{7.63}{-8.25}\hullline{3.70}{-8.77}\hullline{4.46}{-8.77}\hullfill
\hullface{keel!57!black}\hullmove{-4.46}{-8.77}\hullline{-3.70}{-8.77}\hullline{-7.63}{-8.25}\hullline{-8.53}{-8.25}\hullfill
\hullface{keel!57!black}\hullmove{10.27}{-7.57}\hullline{9.77}{-7.57}\hullline{7.63}{-8.25}\hullline{8.53}{-8.25}\hullfill
\hullface{keel!57!black}\hullmove{-8.53}{-8.25}\hullline{-7.63}{-8.25}\hullline{-9.77}{-7.57}\hullline{-10.27}{-7.57}\hullfill
\hullface{keel!57!black}\hullmove{10.74}{-6.77}\hullline{10.56}{-6.77}\hullline{9.77}{-7.57}\hullline{10.27}{-7.57}\hullfill
\hullface{keel!69!black}\hullmove{-10.27}{-7.57}\hullline{-9.77}{-7.57}\hullline{-10.56}{-6.77}\hullline{-10.74}{-6.77}\hullfill
\hullface{keel!57!black}\hullmove{10.82}{-5.86}\hullline{10.77}{-5.86}\hullline{10.56}{-6.77}\hullline{10.74}{-6.77}\hullfill
\hullface{keel!81!black}\hullmove{-10.74}{-6.77}\hullline{-10.56}{-6.77}\hullline{-10.77}{-5.86}\hullline{-10.82}{-5.86}\hullfill
\hullface{keel!57!black}\hullmove{10.83}{-5.44}\hullline{10.79}{-5.44}\hullline{10.77}{-5.86}\hullline{10.82}{-5.86}\hullfill
\hullface{keel!81!black}\hullmove{-10.82}{-5.86}\hullline{-10.77}{-5.86}\hullline{-10.79}{-5.44}\hullline{-10.83}{-5.44}\hullfill
\hullface{topsides!81!black}\hullmove{-10.83}{0.00}\hullline{-10.80}{0.00}\hullline{-10.81}{0.78}\hullline{-10.82}{1.56}\hullline{-10.84}{2.34}\hullline{-10.85}{2.34}\hullline{-10.84}{1.56}\hullline{-10.83}{0.78}\hullfill
\hullface{topsides!57!black}\hullmove{10.86}{3.12}\hullline{10.87}{3.90}\hullline{10.87}{3.90}\hullline{10.85}{3.12}\hullline{10.84}{2.34}\hullline{10.85}{2.34}\hullfill
\hullface{keel!81!black}\hullmove{-10.83}{-5.44}\hullline{-10.79}{-5.44}\hullline{-10.80}{-4.85}\hullline{-10.80}{-3.75}\hullline{-10.80}{-2.58}\hullline{-10.80}{-1.32}\hullline{-10.80}{0.00}\hullline{-10.83}{0.00}\hullline{-10.83}{-1.32}\hullline{-10.83}{-2.58}\hullline{-10.83}{-3.75}\hullline{-10.83}{-4.85}\hullfill
\hullface{topsides!57!black}\hullmove{10.81}{0.78}\hullline{10.82}{1.56}\hullline{10.79}{1.56}\hullline{10.78}{0.78}\hullline{10.76}{0.00}\hullline{10.80}{0.00}\hullfill
\hullface{keel!57!black}\hullmove{10.80}{-3.75}\hullline{10.80}{-2.58}\hullline{10.80}{-1.32}\hullline{10.80}{0.00}\hullline{10.76}{0.00}\hullline{10.76}{-1.32}\hullline{10.76}{-2.58}\hullline{10.76}{-3.75}\hullline{10.75}{-4.85}\hullline{10.80}{-4.85}\hullfill
\hullface{keel!57!black}\hullmove{3.70}{-8.77}\hullline{3.12}{-8.77}\hullline{0.00}{-9.06}\hullline{0.00}{-9.06}\hullfill
\hullface{keel!57!black}\hullmove{0.00}{-9.06}\hullline{0.00}{-9.06}\hullline{-3.12}{-8.77}\hullline{-3.70}{-8.77}\hullfill
\hullface{keel!57!black}\hullmove{7.63}{-8.25}\hullline{6.80}{-8.25}\hullline{3.12}{-8.77}\hullline{3.70}{-8.77}\hullfill
\hullface{keel!57!black}\hullmove{-3.70}{-8.77}\hullline{-3.12}{-8.77}\hullline{-6.80}{-8.25}\hullline{-7.63}{-8.25}\hullfill
\hullface{keel!57!black}\hullmove{9.77}{-7.57}\hullline{9.17}{-7.57}\hullline{6.80}{-8.25}\hullline{7.63}{-8.25}\hullfill
\hullface{keel!57!black}\hullmove{-7.63}{-8.25}\hullline{-6.80}{-8.25}\hullline{-9.17}{-7.57}\hullline{-9.77}{-7.57}\hullfill
\hullface{keel!57!black}\hullmove{10.56}{-6.77}\hullline{10.28}{-6.77}\hullline{9.17}{-7.57}\hullline{9.77}{-7.57}\hullfill
\hullface{keel!60!black}\hullmove{-9.77}{-7.57}\hullline{-9.17}{-7.57}\hullline{-10.28}{-6.77}\hullline{-10.56}{-6.77}\hullfill
\hullface{keel!57!black}\hullmove{10.77}{-5.86}\hullline{10.66}{-5.86}\hullline{10.28}{-6.77}\hullline{10.56}{-6.77}\hullfill
\hullface{keel!78!black}\hullmove{-10.56}{-6.77}\hullline{-10.28}{-6.77}\hullline{-10.66}{-5.86}\hullline{-10.77}{-5.86}\hullfill
\hullface{keel!57!black}\hullmove{10.79}{-5.44}\hullline{10.72}{-5.44}\hullline{10.66}{-5.86}\hullline{10.77}{-5.86}\hullfill
\hullface{keel!81!black}\hullmove{-10.77}{-5.86}\hullline{-10.66}{-5.86}\hullline{-10.72}{-5.44}\hullline{-10.79}{-5.44}\hullfill
\hullface{keel!57!black}\hullmove{10.80}{-4.85}\hullline{10.75}{-4.85}\hullline{10.72}{-5.44}\hullline{10.79}{-5.44}\hullfill
\hullface{keel!81!black}\hullmove{-10.79}{-5.44}\hullline{-10.72}{-5.44}\hullline{-10.75}{-4.85}\hullline{-10.80}{-4.85}\hullfill
\hullface{topsides!81!black}\hullmove{-10.80}{0.00}\hullline{-10.76}{0.00}\hullline{-10.78}{0.78}\hullline{-10.79}{1.56}\hullline{-10.82}{1.56}\hullline{-10.81}{0.78}\hullfill
\hullface{topsides!81!black}\hullmove{-10.82}{1.56}\hullline{-10.79}{1.56}\hullline{-10.82}{2.34}\hullline{-10.84}{3.12}\hullline{-10.85}{3.12}\hullline{-10.84}{2.34}\hullfill
\hullface{topsides!57!black}\hullmove{10.87}{3.90}\hullline{10.87}{3.90}\hullline{10.84}{3.12}\hullline{10.85}{3.12}\hullfill
\hullface{keel!84!black}\hullmove{-10.76}{-3.75}\hullline{-10.71}{-3.75}\hullline{-10.71}{-2.58}\hullline{-10.71}{-1.32}\hullline{-10.71}{0.00}\hullline{-10.76}{0.00}\hullline{-10.76}{-1.32}\hullline{-10.76}{-2.58}\hullfill
\hullface{topsides!57!black}\hullmove{10.78}{0.78}\hullline{10.79}{1.56}\hullline{10.76}{1.56}\hullline{10.73}{0.78}\hullline{10.71}{0.00}\hullline{10.76}{0.00}\hullfill
\hullface{topsides!81!black}\hullmove{-10.76}{0.00}\hullline{-10.71}{0.00}\hullline{-10.73}{0.78}\hullline{-10.76}{1.56}\hullline{-10.79}{1.56}\hullline{-10.78}{0.78}\hullfill
\hullface{topsides!57!black}\hullmove{10.82}{2.34}\hullline{10.84}{3.12}\hullline{10.83}{3.12}\hullline{10.79}{2.34}\hullline{10.76}{1.56}\hullline{10.79}{1.56}\hullfill
\hullface{topsides!81!black}\hullmove{-10.79}{1.56}\hullline{-10.76}{1.56}\hullline{-10.79}{2.34}\hullline{-10.83}{3.12}\hullline{-10.84}{3.12}\hullline{-10.82}{2.34}\hullfill
\hullface{keel!57!black}\hullmove{3.12}{-8.77}\hullline{2.66}{-8.77}\hullline{0.00}{-9.06}\hullline{0.00}{-9.06}\hullfill
\hullface{keel!57!black}\hullmove{0.00}{-9.06}\hullline{0.00}{-9.06}\hullline{-2.66}{-8.77}\hullline{-3.12}{-8.77}\hullfill
\hullface{keel!57!black}\hullmove{6.80}{-8.25}\hullline{6.06}{-8.25}\hullline{2.66}{-8.77}\hullline{3.12}{-8.77}\hullfill
\hullface{keel!57!black}\hullmove{-3.12}{-8.77}\hullline{-2.66}{-8.77}\hullline{-6.06}{-8.25}\hullline{-6.80}{-8.25}\hullfill
\hullface{keel!57!black}\hullmove{9.17}{-7.57}\hullline{8.54}{-7.57}\hullline{6.06}{-8.25}\hullline{6.80}{-8.25}\hullfill
\hullface{keel!57!black}\hullmove{-6.80}{-8.25}\hullline{-6.06}{-8.25}\hullline{-8.54}{-7.57}\hullline{-9.17}{-7.57}\hullfill
\hullface{keel!57!black}\hullmove{10.28}{-6.77}\hullline{9.91}{-6.77}\hullline{8.54}{-7.57}\hullline{9.17}{-7.57}\hullfill
\hullface{keel!57!black}\hullmove{-9.17}{-7.57}\hullline{-8.54}{-7.57}\hullline{-9.91}{-6.77}\hullline{-10.28}{-6.77}\hullfill
\hullface{keel!57!black}\hullmove{10.66}{-5.86}\hullline{10.49}{-5.86}\hullline{9.91}{-6.77}\hullline{10.28}{-6.77}\hullfill
\hullface{keel!75!black}\hullmove{-10.28}{-6.77}\hullline{-9.91}{-6.77}\hullline{-10.49}{-5.86}\hullline{-10.66}{-5.86}\hullfill
\hullface{keel!57!black}\hullmove{10.72}{-5.44}\hullline{10.60}{-5.44}\hullline{10.49}{-5.86}\hullline{10.66}{-5.86}\hullfill
\hullface{keel!81!black}\hullmove{-10.66}{-5.86}\hullline{-10.49}{-5.86}\hullline{-10.60}{-5.44}\hullline{-10.72}{-5.44}\hullfill
\hullface{keel!57!black}\hullmove{10.75}{-4.85}\hullline{10.67}{-4.85}\hullline{10.60}{-5.44}\hullline{10.72}{-5.44}\hullfill
\hullface{keel!81!black}\hullmove{-10.72}{-5.44}\hullline{-10.60}{-5.44}\hullline{-10.67}{-4.85}\hullline{-10.75}{-4.85}\hullfill
\hullface{keel!57!black}\hullmove{10.76}{-3.75}\hullline{10.71}{-3.75}\hullline{10.67}{-4.85}\hullline{10.75}{-4.85}\hullfill
\hullface{keel!84!black}\hullmove{-10.75}{-4.85}\hullline{-10.67}{-4.85}\hullline{-10.71}{-3.75}\hullline{-10.76}{-3.75}\hullfill
\hullface{topsides!57!black}\hullmove{10.87}{3.90}\hullline{10.87}{3.90}\hullline{10.83}{3.12}\hullline{10.84}{3.12}\hullfill
\hullface{keel!57!black}\hullmove{2.66}{-8.77}\hullline{2.30}{-8.77}\hullline{0.00}{-9.06}\hullline{0.00}{-9.06}\hullfill
\hullface{keel!57!black}\hullmove{0.00}{-9.06}\hullline{0.00}{-9.06}\hullline{-2.30}{-8.77}\hullline{-2.66}{-8.77}\hullfill
\hullface{keel!57!black}\hullmove{6.06}{-8.25}\hullline{5.41}{-8.25}\hullline{2.30}{-8.77}\hullline{2.66}{-8.77}\hullfill
\hullface{keel!57!black}\hullmove{-2.66}{-8.77}\hullline{-2.30}{-8.77}\hullline{-5.41}{-8.25}\hullline{-6.06}{-8.25}\hullfill
\hullface{keel!57!black}\hullmove{8.54}{-7.57}\hullline{7.91}{-7.57}\hullline{5.41}{-8.25}\hullline{6.06}{-8.25}\hullfill
\hullface{keel!57!black}\hullmove{-6.06}{-8.25}\hullline{-5.41}{-8.25}\hullline{-7.91}{-7.57}\hullline{-8.54}{-7.57}\hullfill
\hullface{keel!57!black}\hullmove{9.91}{-6.77}\hullline{9.48}{-6.77}\hullline{7.91}{-7.57}\hullline{8.54}{-7.57}\hullfill
\hullface{keel!57!black}\hullmove{-8.54}{-7.57}\hullline{-7.91}{-7.57}\hullline{-9.48}{-6.77}\hullline{-9.91}{-6.77}\hullfill
\hullface{keel!57!black}\hullmove{10.49}{-5.86}\hullline{10.26}{-5.86}\hullline{9.48}{-6.77}\hullline{9.91}{-6.77}\hullfill
\hullface{keel!69!black}\hullmove{-9.91}{-6.77}\hullline{-9.48}{-6.77}\hullline{-10.26}{-5.86}\hullline{-10.49}{-5.86}\hullfill
\hullface{keel!57!black}\hullmove{10.60}{-5.44}\hullline{10.43}{-5.44}\hullline{10.26}{-5.86}\hullline{10.49}{-5.86}\hullfill
\hullface{keel!78!black}\hullmove{-10.49}{-5.86}\hullline{-10.26}{-5.86}\hullline{-10.43}{-5.44}\hullline{-10.60}{-5.44}\hullfill
\hullface{keel!57!black}\hullmove{10.67}{-4.85}\hullline{10.56}{-4.85}\hullline{10.43}{-5.44}\hullline{10.60}{-5.44}\hullfill
\hullface{keel!81!black}\hullmove{-10.60}{-5.44}\hullline{-10.43}{-5.44}\hullline{-10.56}{-4.85}\hullline{-10.67}{-4.85}\hullfill
\hullface{keel!57!black}\hullmove{10.71}{-3.75}\hullline{10.64}{-3.75}\hullline{10.56}{-4.85}\hullline{10.67}{-4.85}\hullfill
\hullface{keel!84!black}\hullmove{-10.67}{-4.85}\hullline{-10.56}{-4.85}\hullline{-10.64}{-3.75}\hullline{-10.71}{-3.75}\hullfill
\hullface{keel!57!black}\hullmove{10.71}{-2.58}\hullline{10.71}{-1.32}\hullline{10.71}{0.00}\hullline{10.65}{0.00}\hullline{10.65}{-1.32}\hullline{10.65}{-2.58}\hullline{10.64}{-3.75}\hullline{10.71}{-3.75}\hullfill
\hullface{topsides!57!black}\hullmove{10.73}{0.78}\hullline{10.67}{0.78}\hullline{10.65}{0.00}\hullline{10.71}{0.00}\hullfill
\hullface{topsides!81!black}\hullmove{-10.71}{0.00}\hullline{-10.65}{0.00}\hullline{-10.67}{0.78}\hullline{-10.73}{0.78}\hullfill
\hullface{topsides!57!black}\hullmove{10.76}{1.56}\hullline{10.79}{2.34}\hullline{10.76}{2.34}\hullline{10.71}{1.56}\hullline{10.67}{0.78}\hullline{10.73}{0.78}\hullfill
\hullface{topsides!81!black}\hullmove{-10.73}{0.78}\hullline{-10.67}{0.78}\hullline{-10.71}{1.56}\hullline{-10.76}{2.34}\hullline{-10.79}{2.34}\hullline{-10.76}{1.56}\hullfill
\hullface{topsides!57!black}\hullmove{10.83}{3.12}\hullline{10.81}{3.12}\hullline{10.76}{2.34}\hullline{10.79}{2.34}\hullfill
\hullface{topsides!81!black}\hullmove{-10.79}{2.34}\hullline{-10.76}{2.34}\hullline{-10.81}{3.12}\hullline{-10.83}{3.12}\hullfill
\hullface{topsides!57!black}\hullmove{10.87}{3.90}\hullline{10.87}{3.90}\hullline{10.81}{3.12}\hullline{10.83}{3.12}\hullfill
\hullface{topsides!78!black}\hullmove{-10.83}{3.12}\hullline{-10.81}{3.12}\hullline{-10.87}{3.90}\hullline{-10.87}{3.90}\hullfill
\hullface{keel!57!black}\hullmove{2.30}{-8.77}\hullline{2.00}{-8.77}\hullline{0.00}{-9.06}\hullline{0.00}{-9.06}\hullfill
\hullface{keel!57!black}\hullmove{0.00}{-9.06}\hullline{0.00}{-9.06}\hullline{-2.00}{-8.77}\hullline{-2.30}{-8.77}\hullfill
\hullface{keel!57!black}\hullmove{5.41}{-8.25}\hullline{4.84}{-8.25}\hullline{2.00}{-8.77}\hullline{2.30}{-8.77}\hullfill
\hullface{keel!57!black}\hullmove{-2.30}{-8.77}\hullline{-2.00}{-8.77}\hullline{-4.84}{-8.25}\hullline{-5.41}{-8.25}\hullfill
\hullface{keel!57!black}\hullmove{7.91}{-7.57}\hullline{7.30}{-7.57}\hullline{4.84}{-8.25}\hullline{5.41}{-8.25}\hullfill
\hullface{keel!57!black}\hullmove{-5.41}{-8.25}\hullline{-4.84}{-8.25}\hullline{-7.30}{-7.57}\hullline{-7.91}{-7.57}\hullfill
\hullface{keel!57!black}\hullmove{9.48}{-6.77}\hullline{9.00}{-6.77}\hullline{7.30}{-7.57}\hullline{7.91}{-7.57}\hullfill
\hullface{keel!57!black}\hullmove{-7.91}{-7.57}\hullline{-7.30}{-7.57}\hullline{-9.00}{-6.77}\hullline{-9.48}{-6.77}\hullfill
\hullface{keel!57!black}\hullmove{10.26}{-5.86}\hullline{9.96}{-5.86}\hullline{9.00}{-6.77}\hullline{9.48}{-6.77}\hullfill
\hullface{keel!66!black}\hullmove{-9.48}{-6.77}\hullline{-9.00}{-6.77}\hullline{-9.96}{-5.86}\hullline{-10.26}{-5.86}\hullfill
\hullface{keel!57!black}\hullmove{10.43}{-5.44}\hullline{10.20}{-5.44}\hullline{9.96}{-5.86}\hullline{10.26}{-5.86}\hullfill
\hullface{keel!75!black}\hullmove{-10.26}{-5.86}\hullline{-9.96}{-5.86}\hullline{-10.20}{-5.44}\hullline{-10.43}{-5.44}\hullfill
\hullface{keel!57!black}\hullmove{10.56}{-4.85}\hullline{10.40}{-4.85}\hullline{10.20}{-5.44}\hullline{10.43}{-5.44}\hullfill
\hullface{keel!81!black}\hullmove{-10.43}{-5.44}\hullline{-10.20}{-5.44}\hullline{-10.40}{-4.85}\hullline{-10.56}{-4.85}\hullfill
\hullface{keel!57!black}\hullmove{10.64}{-3.75}\hullline{10.54}{-3.75}\hullline{10.40}{-4.85}\hullline{10.56}{-4.85}\hullfill
\hullface{keel!84!black}\hullmove{-10.56}{-4.85}\hullline{-10.40}{-4.85}\hullline{-10.54}{-3.75}\hullline{-10.64}{-3.75}\hullfill
\hullface{keel!57!black}\hullmove{10.65}{-2.58}\hullline{10.57}{-2.58}\hullline{10.54}{-3.75}\hullline{10.64}{-3.75}\hullfill
\hullface{keel!84!black}\hullmove{-10.64}{-3.75}\hullline{-10.54}{-3.75}\hullline{-10.57}{-2.58}\hullline{-10.65}{-2.58}\hullfill
\hullface{topsides!57!black}\hullmove{10.67}{0.78}\hullline{10.60}{0.78}\hullline{10.57}{0.00}\hullline{10.65}{0.00}\hullfill
\hullface{topsides!84!black}\hullmove{-10.65}{0.00}\hullline{-10.57}{0.00}\hullline{-10.60}{0.78}\hullline{-10.67}{0.78}\hullfill
\hullface{topsides!57!black}\hullmove{10.71}{1.56}\hullline{10.66}{1.56}\hullline{10.60}{0.78}\hullline{10.67}{0.78}\hullfill
\hullface{topsides!81!black}\hullmove{-10.67}{0.78}\hullline{-10.60}{0.78}\hullline{-10.66}{1.56}\hullline{-10.71}{1.56}\hullfill
\hullface{topsides!57!black}\hullmove{10.76}{2.34}\hullline{10.72}{2.34}\hullline{10.66}{1.56}\hullline{10.71}{1.56}\hullfill
\hullface{topsides!81!black}\hullmove{-10.71}{1.56}\hullline{-10.66}{1.56}\hullline{-10.72}{2.34}\hullline{-10.76}{2.34}\hullfill
\hullface{topsides!57!black}\hullmove{10.81}{3.12}\hullline{10.79}{3.12}\hullline{10.72}{2.34}\hullline{10.76}{2.34}\hullfill
\hullface{topsides!57!black}\hullmove{10.87}{3.90}\hullline{10.87}{3.90}\hullline{10.79}{3.12}\hullline{10.81}{3.12}\hullfill
\hullface{topsides!78!black}\hullmove{-10.81}{3.12}\hullline{-10.79}{3.12}\hullline{-10.87}{3.90}\hullline{-10.87}{3.90}\hullfill
\hullface{keel!57!black}\hullmove{10.65}{-1.32}\hullline{10.65}{0.00}\hullline{10.57}{0.00}\hullline{10.57}{-1.32}\hullline{10.57}{-2.58}\hullline{10.65}{-2.58}\hullfill
\hullface{keel!84!black}\hullmove{-10.65}{-2.58}\hullline{-10.57}{-2.58}\hullline{-10.57}{-1.32}\hullline{-10.57}{0.00}\hullline{-10.65}{0.00}\hullline{-10.65}{-1.32}\hullfill
\hullface{keel!57!black}\hullmove{10.57}{-1.32}\hullline{10.57}{0.00}\hullline{10.48}{0.00}\hullline{10.48}{-1.32}\hullline{10.47}{-2.58}\hullline{10.57}{-2.58}\hullfill
\hullface{keel!84!black}\hullmove{-10.57}{-2.58}\hullline{-10.47}{-2.58}\hullline{-10.48}{-1.32}\hullline{-10.48}{0.00}\hullline{-10.57}{0.00}\hullline{-10.57}{-1.32}\hullfill
\hullface{keel!57!black}\hullmove{2.00}{-8.77}\hullline{1.76}{-8.77}\hullline{0.00}{-9.06}\hullline{0.00}{-9.06}\hullfill
\hullface{keel!57!black}\hullmove{0.00}{-9.06}\hullline{0.00}{-9.06}\hullline{-1.76}{-8.77}\hullline{-2.00}{-8.77}\hullfill
\hullface{keel!57!black}\hullmove{4.84}{-8.25}\hullline{4.35}{-8.25}\hullline{1.76}{-8.77}\hullline{2.00}{-8.77}\hullfill
\hullface{keel!57!black}\hullmove{-2.00}{-8.77}\hullline{-1.76}{-8.77}\hullline{-4.35}{-8.25}\hullline{-4.84}{-8.25}\hullfill
\hullface{keel!57!black}\hullmove{7.30}{-7.57}\hullline{6.72}{-7.57}\hullline{4.35}{-8.25}\hullline{4.84}{-8.25}\hullfill
\hullface{keel!57!black}\hullmove{-4.84}{-8.25}\hullline{-4.35}{-8.25}\hullline{-6.72}{-7.57}\hullline{-7.30}{-7.57}\hullfill
\hullface{keel!57!black}\hullmove{9.00}{-6.77}\hullline{8.50}{-6.77}\hullline{6.72}{-7.57}\hullline{7.30}{-7.57}\hullfill
\hullface{keel!57!black}\hullmove{-7.30}{-7.57}\hullline{-6.72}{-7.57}\hullline{-8.50}{-6.77}\hullline{-9.00}{-6.77}\hullfill
\hullface{keel!57!black}\hullmove{9.96}{-5.86}\hullline{9.61}{-5.86}\hullline{8.50}{-6.77}\hullline{9.00}{-6.77}\hullfill
\hullface{keel!63!black}\hullmove{-9.00}{-6.77}\hullline{-8.50}{-6.77}\hullline{-9.61}{-5.86}\hullline{-9.96}{-5.86}\hullfill
\hullface{keel!57!black}\hullmove{10.20}{-5.44}\hullline{9.91}{-5.44}\hullline{9.61}{-5.86}\hullline{9.96}{-5.86}\hullfill
\hullface{keel!72!black}\hullmove{-9.96}{-5.86}\hullline{-9.61}{-5.86}\hullline{-9.91}{-5.44}\hullline{-10.20}{-5.44}\hullfill
\hullface{keel!57!black}\hullmove{10.40}{-4.85}\hullline{10.19}{-4.85}\hullline{9.91}{-5.44}\hullline{10.20}{-5.44}\hullfill
\hullface{keel!78!black}\hullmove{-10.20}{-5.44}\hullline{-9.91}{-5.44}\hullline{-10.19}{-4.85}\hullline{-10.40}{-4.85}\hullfill
\hullface{keel!57!black}\hullmove{10.54}{-3.75}\hullline{10.41}{-3.75}\hullline{10.19}{-4.85}\hullline{10.40}{-4.85}\hullfill
\hullface{keel!84!black}\hullmove{-10.40}{-4.85}\hullline{-10.19}{-4.85}\hullline{-10.41}{-3.75}\hullline{-10.54}{-3.75}\hullfill
\hullface{keel!57!black}\hullmove{10.57}{-2.58}\hullline{10.47}{-2.58}\hullline{10.41}{-3.75}\hullline{10.54}{-3.75}\hullfill
\hullface{keel!84!black}\hullmove{-10.54}{-3.75}\hullline{-10.41}{-3.75}\hullline{-10.47}{-2.58}\hullline{-10.57}{-2.58}\hullfill
\hullface{topsides!57!black}\hullmove{10.60}{0.78}\hullline{10.52}{0.78}\hullline{10.48}{0.00}\hullline{10.57}{0.00}\hullfill
\hullface{topsides!84!black}\hullmove{-10.57}{0.00}\hullline{-10.48}{0.00}\hullline{-10.52}{0.78}\hullline{-10.60}{0.78}\hullfill
\hullface{topsides!57!black}\hullmove{10.66}{1.56}\hullline{10.59}{1.56}\hullline{10.52}{0.78}\hullline{10.60}{0.78}\hullfill
\hullface{topsides!81!black}\hullmove{-10.60}{0.78}\hullline{-10.52}{0.78}\hullline{-10.59}{1.56}\hullline{-10.66}{1.56}\hullfill
\hullface{topsides!57!black}\hullmove{10.72}{2.34}\hullline{10.67}{2.34}\hullline{10.59}{1.56}\hullline{10.66}{1.56}\hullfill
\hullface{topsides!81!black}\hullmove{-10.66}{1.56}\hullline{-10.59}{1.56}\hullline{-10.67}{2.34}\hullline{-10.72}{2.34}\hullfill
\hullface{topsides!57!black}\hullmove{10.79}{3.12}\hullline{10.76}{3.12}\hullline{10.67}{2.34}\hullline{10.72}{2.34}\hullfill
\hullface{topsides!78!black}\hullmove{-10.72}{2.34}\hullline{-10.67}{2.34}\hullline{-10.76}{3.12}\hullline{-10.79}{3.12}\hullfill
\hullface{topsides!57!black}\hullmove{10.87}{3.90}\hullline{10.86}{3.90}\hullline{10.76}{3.12}\hullline{10.79}{3.12}\hullfill
\hullface{topsides!78!black}\hullmove{-10.79}{3.12}\hullline{-10.76}{3.12}\hullline{-10.86}{3.90}\hullline{-10.87}{3.90}\hullfill
\hullface{keel!57!black}\hullmove{1.76}{-8.77}\hullline{1.56}{-8.77}\hullline{0.00}{-9.06}\hullline{0.00}{-9.06}\hullfill
\hullface{keel!57!black}\hullmove{0.00}{-9.06}\hullline{0.00}{-9.06}\hullline{-1.56}{-8.77}\hullline{-1.76}{-8.77}\hullfill
\hullface{keel!57!black}\hullmove{4.35}{-8.25}\hullline{3.91}{-8.25}\hullline{1.56}{-8.77}\hullline{1.76}{-8.77}\hullfill
\hullface{keel!57!black}\hullmove{-1.76}{-8.77}\hullline{-1.56}{-8.77}\hullline{-3.91}{-8.25}\hullline{-4.35}{-8.25}\hullfill
\hullface{keel!57!black}\hullmove{6.72}{-7.57}\hullline{6.19}{-7.57}\hullline{3.91}{-8.25}\hullline{4.35}{-8.25}\hullfill
\hullface{keel!57!black}\hullmove{-4.35}{-8.25}\hullline{-3.91}{-8.25}\hullline{-6.19}{-7.57}\hullline{-6.72}{-7.57}\hullfill
\hullface{keel!57!black}\hullmove{8.50}{-6.77}\hullline{8.00}{-6.77}\hullline{6.19}{-7.57}\hullline{6.72}{-7.57}\hullfill
\hullface{keel!57!black}\hullmove{-6.72}{-7.57}\hullline{-6.19}{-7.57}\hullline{-8.00}{-6.77}\hullline{-8.50}{-6.77}\hullfill
\hullface{keel!57!black}\hullmove{9.61}{-5.86}\hullline{9.23}{-5.86}\hullline{8.00}{-6.77}\hullline{8.50}{-6.77}\hullfill
\hullface{keel!60!black}\hullmove{-8.50}{-6.77}\hullline{-8.00}{-6.77}\hullline{-9.23}{-5.86}\hullline{-9.61}{-5.86}\hullfill
\hullface{keel!57!black}\hullmove{9.91}{-5.44}\hullline{9.59}{-5.44}\hullline{9.23}{-5.86}\hullline{9.61}{-5.86}\hullfill
\hullface{keel!69!black}\hullmove{-9.61}{-5.86}\hullline{-9.23}{-5.86}\hullline{-9.59}{-5.44}\hullline{-9.91}{-5.44}\hullfill
\hullface{keel!57!black}\hullmove{10.19}{-4.85}\hullline{9.93}{-4.85}\hullline{9.59}{-5.44}\hullline{9.91}{-5.44}\hullfill
\hullface{keel!75!black}\hullmove{-9.91}{-5.44}\hullline{-9.59}{-5.44}\hullline{-9.93}{-4.85}\hullline{-10.19}{-4.85}\hullfill
\hullface{keel!57!black}\hullmove{10.41}{-3.75}\hullline{10.25}{-3.75}\hullline{9.93}{-4.85}\hullline{10.19}{-4.85}\hullfill
\hullface{keel!81!black}\hullmove{-10.19}{-4.85}\hullline{-9.93}{-4.85}\hullline{-10.25}{-3.75}\hullline{-10.41}{-3.75}\hullfill
\hullface{keel!57!black}\hullmove{10.47}{-2.58}\hullline{10.35}{-2.58}\hullline{10.25}{-3.75}\hullline{10.41}{-3.75}\hullfill
\hullface{keel!84!black}\hullmove{-10.41}{-3.75}\hullline{-10.25}{-3.75}\hullline{-10.35}{-2.58}\hullline{-10.47}{-2.58}\hullfill
\hullface{keel!57!black}\hullmove{10.48}{-1.32}\hullline{10.48}{0.00}\hullline{10.37}{0.00}\hullline{10.37}{-1.32}\hullline{10.35}{-2.58}\hullline{10.47}{-2.58}\hullfill
\hullface{keel!87!black}\hullmove{-10.47}{-2.58}\hullline{-10.35}{-2.58}\hullline{-10.37}{-1.32}\hullline{-10.37}{0.00}\hullline{-10.48}{0.00}\hullline{-10.48}{-1.32}\hullfill
\hullface{topsides!57!black}\hullmove{10.52}{0.78}\hullline{10.42}{0.78}\hullline{10.37}{0.00}\hullline{10.48}{0.00}\hullfill
\hullface{topsides!84!black}\hullmove{-10.48}{0.00}\hullline{-10.37}{0.00}\hullline{-10.42}{0.78}\hullline{-10.52}{0.78}\hullfill
\hullface{topsides!57!black}\hullmove{10.59}{1.56}\hullline{10.50}{1.56}\hullline{10.42}{0.78}\hullline{10.52}{0.78}\hullfill
\hullface{topsides!81!black}\hullmove{-10.52}{0.78}\hullline{-10.42}{0.78}\hullline{-10.50}{1.56}\hullline{-10.59}{1.56}\hullfill
\hullface{topsides!57!black}\hullmove{10.67}{2.34}\hullline{10.61}{2.34}\hullline{10.50}{1.56}\hullline{10.59}{1.56}\hullfill
\hullface{topsides!81!black}\hullmove{-10.59}{1.56}\hullline{-10.50}{1.56}\hullline{-10.61}{2.34}\hullline{-10.67}{2.34}\hullfill
\hullface{topsides!57!black}\hullmove{10.76}{3.12}\hullline{10.73}{3.12}\hullline{10.61}{2.34}\hullline{10.67}{2.34}\hullfill
\hullface{topsides!78!black}\hullmove{-10.67}{2.34}\hullline{-10.61}{2.34}\hullline{-10.73}{3.12}\hullline{-10.76}{3.12}\hullfill
\hullface{topsides!57!black}\hullmove{10.86}{3.90}\hullline{10.86}{3.90}\hullline{10.73}{3.12}\hullline{10.76}{3.12}\hullfill
\hullface{topsides!75!black}\hullmove{-10.76}{3.12}\hullline{-10.73}{3.12}\hullline{-10.86}{3.90}\hullline{-10.86}{3.90}\hullfill
\hullface{keel!57!black}\hullmove{1.56}{-8.77}\hullline{1.38}{-8.77}\hullline{0.00}{-9.06}\hullline{0.00}{-9.06}\hullfill
\hullface{keel!57!black}\hullmove{0.00}{-9.06}\hullline{0.00}{-9.06}\hullline{-1.38}{-8.77}\hullline{-1.56}{-8.77}\hullfill
\hullface{keel!57!black}\hullmove{3.91}{-8.25}\hullline{3.53}{-8.25}\hullline{1.38}{-8.77}\hullline{1.56}{-8.77}\hullfill
\hullface{keel!57!black}\hullmove{-1.56}{-8.77}\hullline{-1.38}{-8.77}\hullline{-3.53}{-8.25}\hullline{-3.91}{-8.25}\hullfill
\hullface{keel!57!black}\hullmove{6.19}{-7.57}\hullline{5.69}{-7.57}\hullline{3.53}{-8.25}\hullline{3.91}{-8.25}\hullfill
\hullface{keel!57!black}\hullmove{-3.91}{-8.25}\hullline{-3.53}{-8.25}\hullline{-5.69}{-7.57}\hullline{-6.19}{-7.57}\hullfill
\hullface{keel!57!black}\hullmove{8.00}{-6.77}\hullline{7.50}{-6.77}\hullline{5.69}{-7.57}\hullline{6.19}{-7.57}\hullfill
\hullface{keel!57!black}\hullmove{-6.19}{-7.57}\hullline{-5.69}{-7.57}\hullline{-7.50}{-6.77}\hullline{-8.00}{-6.77}\hullfill
\hullface{keel!57!black}\hullmove{9.23}{-5.86}\hullline{8.81}{-5.86}\hullline{7.50}{-6.77}\hullline{8.00}{-6.77}\hullfill
\hullface{keel!57!black}\hullmove{-8.00}{-6.77}\hullline{-7.50}{-6.77}\hullline{-8.81}{-5.86}\hullline{-9.23}{-5.86}\hullfill
\hullface{keel!57!black}\hullmove{9.59}{-5.44}\hullline{9.22}{-5.44}\hullline{8.81}{-5.86}\hullline{9.23}{-5.86}\hullfill
\hullface{keel!66!black}\hullmove{-9.23}{-5.86}\hullline{-8.81}{-5.86}\hullline{-9.22}{-5.44}\hullline{-9.59}{-5.44}\hullfill
\hullface{keel!57!black}\hullmove{9.93}{-4.85}\hullline{9.63}{-4.85}\hullline{9.22}{-5.44}\hullline{9.59}{-5.44}\hullfill
\hullface{keel!72!black}\hullmove{-9.59}{-5.44}\hullline{-9.22}{-5.44}\hullline{-9.63}{-4.85}\hullline{-9.93}{-4.85}\hullfill
\hullface{keel!57!black}\hullmove{10.25}{-3.75}\hullline{10.05}{-3.75}\hullline{9.63}{-4.85}\hullline{9.93}{-4.85}\hullfill
\hullface{keel!81!black}\hullmove{-9.93}{-4.85}\hullline{-9.63}{-4.85}\hullline{-10.05}{-3.75}\hullline{-10.25}{-3.75}\hullfill
\hullface{keel!57!black}\hullmove{10.35}{-2.58}\hullline{10.20}{-2.58}\hullline{10.05}{-3.75}\hullline{10.25}{-3.75}\hullfill
\hullface{keel!84!black}\hullmove{-10.25}{-3.75}\hullline{-10.05}{-3.75}\hullline{-10.20}{-2.58}\hullline{-10.35}{-2.58}\hullfill
\hullface{keel!57!black}\hullmove{10.37}{-1.32}\hullline{10.24}{-1.32}\hullline{10.20}{-2.58}\hullline{10.35}{-2.58}\hullfill
\hullface{keel!87!black}\hullmove{-10.35}{-2.58}\hullline{-10.20}{-2.58}\hullline{-10.24}{-1.32}\hullline{-10.37}{-1.32}\hullfill
\hullface{keel!57!black}\hullmove{10.37}{0.00}\hullline{10.24}{0.00}\hullline{10.24}{-1.32}\hullline{10.37}{-1.32}\hullfill
\hullface{keel!87!black}\hullmove{-10.37}{-1.32}\hullline{-10.24}{-1.32}\hullline{-10.24}{0.00}\hullline{-10.37}{0.00}\hullfill
\hullface{topsides!57!black}\hullmove{10.42}{0.78}\hullline{10.30}{0.78}\hullline{10.24}{0.00}\hullline{10.37}{0.00}\hullfill
\hullface{topsides!84!black}\hullmove{-10.37}{0.00}\hullline{-10.24}{0.00}\hullline{-10.30}{0.78}\hullline{-10.42}{0.78}\hullfill
\hullface{topsides!57!black}\hullmove{10.50}{1.56}\hullline{10.41}{1.56}\hullline{10.30}{0.78}\hullline{10.42}{0.78}\hullfill
\hullface{topsides!81!black}\hullmove{-10.42}{0.78}\hullline{-10.30}{0.78}\hullline{-10.41}{1.56}\hullline{-10.50}{1.56}\hullfill
\hullface{topsides!57!black}\hullmove{10.61}{2.34}\hullline{10.54}{2.34}\hullline{10.41}{1.56}\hullline{10.50}{1.56}\hullfill
\hullface{topsides!78!black}\hullmove{-10.50}{1.56}\hullline{-10.41}{1.56}\hullline{-10.54}{2.34}\hullline{-10.61}{2.34}\hullfill
\hullface{topsides!57!black}\hullmove{10.73}{3.12}\hullline{10.69}{3.12}\hullline{10.54}{2.34}\hullline{10.61}{2.34}\hullfill
\hullface{topsides!78!black}\hullmove{-10.61}{2.34}\hullline{-10.54}{2.34}\hullline{-10.69}{3.12}\hullline{-10.73}{3.12}\hullfill
\hullface{topsides!57!black}\hullmove{10.86}{3.90}\hullline{10.85}{3.90}\hullline{10.69}{3.12}\hullline{10.73}{3.12}\hullfill
\hullface{topsides!75!black}\hullmove{-10.73}{3.12}\hullline{-10.69}{3.12}\hullline{-10.85}{3.90}\hullline{-10.86}{3.90}\hullfill
\hullface{keel!57!black}\hullmove{1.38}{-8.77}\hullline{1.24}{-8.77}\hullline{0.00}{-9.06}\hullline{0.00}{-9.06}\hullfill
\hullface{keel!57!black}\hullmove{0.00}{-9.06}\hullline{0.00}{-9.06}\hullline{-1.24}{-8.77}\hullline{-1.38}{-8.77}\hullfill
\hullface{keel!57!black}\hullmove{3.53}{-8.25}\hullline{3.19}{-8.25}\hullline{1.24}{-8.77}\hullline{1.38}{-8.77}\hullfill
\hullface{keel!57!black}\hullmove{-1.38}{-8.77}\hullline{-1.24}{-8.77}\hullline{-3.19}{-8.25}\hullline{-3.53}{-8.25}\hullfill
\hullface{keel!57!black}\hullmove{5.69}{-7.57}\hullline{5.22}{-7.57}\hullline{3.19}{-8.25}\hullline{3.53}{-8.25}\hullfill
\hullface{keel!57!black}\hullmove{-3.53}{-8.25}\hullline{-3.19}{-8.25}\hullline{-5.22}{-7.57}\hullline{-5.69}{-7.57}\hullfill
\hullface{keel!57!black}\hullmove{7.50}{-6.77}\hullline{7.00}{-6.77}\hullline{5.22}{-7.57}\hullline{5.69}{-7.57}\hullfill
\hullface{keel!57!black}\hullmove{-5.69}{-7.57}\hullline{-5.22}{-7.57}\hullline{-7.00}{-6.77}\hullline{-7.50}{-6.77}\hullfill
\hullface{keel!57!black}\hullmove{8.81}{-5.86}\hullline{8.37}{-5.86}\hullline{7.00}{-6.77}\hullline{7.50}{-6.77}\hullfill
\hullface{keel!57!black}\hullmove{-7.50}{-6.77}\hullline{-7.00}{-6.77}\hullline{-8.37}{-5.86}\hullline{-8.81}{-5.86}\hullfill
\hullface{keel!57!black}\hullmove{9.22}{-5.44}\hullline{8.82}{-5.44}\hullline{8.37}{-5.86}\hullline{8.81}{-5.86}\hullfill
\hullface{keel!63!black}\hullmove{-8.81}{-5.86}\hullline{-8.37}{-5.86}\hullline{-8.82}{-5.44}\hullline{-9.22}{-5.44}\hullfill
\hullface{keel!57!black}\hullmove{9.63}{-4.85}\hullline{9.29}{-4.85}\hullline{8.82}{-5.44}\hullline{9.22}{-5.44}\hullfill
\hullface{keel!69!black}\hullmove{-9.22}{-5.44}\hullline{-8.82}{-5.44}\hullline{-9.29}{-4.85}\hullline{-9.63}{-4.85}\hullfill
\hullface{keel!57!black}\hullmove{10.05}{-3.75}\hullline{9.81}{-3.75}\hullline{9.29}{-4.85}\hullline{9.63}{-4.85}\hullfill
\hullface{keel!78!black}\hullmove{-9.63}{-4.85}\hullline{-9.29}{-4.85}\hullline{-9.81}{-3.75}\hullline{-10.05}{-3.75}\hullfill
\hullface{keel!57!black}\hullmove{10.20}{-2.58}\hullline{10.03}{-2.58}\hullline{9.81}{-3.75}\hullline{10.05}{-3.75}\hullfill
\hullface{keel!84!black}\hullmove{-10.05}{-3.75}\hullline{-9.81}{-3.75}\hullline{-10.03}{-2.58}\hullline{-10.20}{-2.58}\hullfill
\hullface{keel!57!black}\hullmove{10.24}{-1.32}\hullline{10.08}{-1.32}\hullline{10.03}{-2.58}\hullline{10.20}{-2.58}\hullfill
\hullface{keel!87!black}\hullmove{-10.20}{-2.58}\hullline{-10.03}{-2.58}\hullline{-10.08}{-1.32}\hullline{-10.24}{-1.32}\hullfill
\hullface{keel!57!black}\hullmove{10.24}{0.00}\hullline{10.09}{0.00}\hullline{10.08}{-1.32}\hullline{10.24}{-1.32}\hullfill
\hullface{keel!87!black}\hullmove{-10.24}{-1.32}\hullline{-10.08}{-1.32}\hullline{-10.09}{0.00}\hullline{-10.24}{0.00}\hullfill
\hullface{topsides!57!black}\hullmove{10.30}{0.78}\hullline{10.17}{0.78}\hullline{10.09}{0.00}\hullline{10.24}{0.00}\hullfill
\hullface{topsides!84!black}\hullmove{-10.24}{0.00}\hullline{-10.09}{0.00}\hullline{-10.17}{0.78}\hullline{-10.30}{0.78}\hullfill
\hullface{topsides!57!black}\hullmove{10.41}{1.56}\hullline{10.30}{1.56}\hullline{10.17}{0.78}\hullline{10.30}{0.78}\hullfill
\hullface{topsides!81!black}\hullmove{-10.30}{0.78}\hullline{-10.17}{0.78}\hullline{-10.30}{1.56}\hullline{-10.41}{1.56}\hullfill
\hullface{topsides!57!black}\hullmove{10.54}{2.34}\hullline{10.46}{2.34}\hullline{10.30}{1.56}\hullline{10.41}{1.56}\hullfill
\hullface{topsides!78!black}\hullmove{-10.41}{1.56}\hullline{-10.30}{1.56}\hullline{-10.46}{2.34}\hullline{-10.54}{2.34}\hullfill
\hullface{topsides!57!black}\hullmove{10.69}{3.12}\hullline{10.64}{3.12}\hullline{10.46}{2.34}\hullline{10.54}{2.34}\hullfill
\hullface{topsides!75!black}\hullmove{-10.54}{2.34}\hullline{-10.46}{2.34}\hullline{-10.64}{3.12}\hullline{-10.69}{3.12}\hullfill
\hullface{topsides!57!black}\hullmove{10.85}{3.90}\hullline{10.84}{3.90}\hullline{10.64}{3.12}\hullline{10.69}{3.12}\hullfill
\hullface{topsides!72!black}\hullmove{-10.69}{3.12}\hullline{-10.64}{3.12}\hullline{-10.84}{3.90}\hullline{-10.85}{3.90}\hullfill
\hullface{keel!57!black}\hullmove{1.24}{-8.77}\hullline{1.11}{-8.77}\hullline{0.00}{-9.06}\hullline{0.00}{-9.06}\hullfill
\hullface{keel!57!black}\hullmove{0.00}{-9.06}\hullline{0.00}{-9.06}\hullline{-1.11}{-8.77}\hullline{-1.24}{-8.77}\hullfill
\hullface{keel!57!black}\hullmove{3.19}{-8.25}\hullline{2.89}{-8.25}\hullline{1.11}{-8.77}\hullline{1.24}{-8.77}\hullfill
\hullface{keel!57!black}\hullmove{-1.24}{-8.77}\hullline{-1.11}{-8.77}\hullline{-2.89}{-8.25}\hullline{-3.19}{-8.25}\hullfill
\hullface{keel!57!black}\hullmove{5.22}{-7.57}\hullline{4.80}{-7.57}\hullline{2.89}{-8.25}\hullline{3.19}{-8.25}\hullfill
\hullface{keel!57!black}\hullmove{-3.19}{-8.25}\hullline{-2.89}{-8.25}\hullline{-4.80}{-7.57}\hullline{-5.22}{-7.57}\hullfill
\hullface{keel!57!black}\hullmove{7.00}{-6.77}\hullline{6.53}{-6.77}\hullline{4.80}{-7.57}\hullline{5.22}{-7.57}\hullfill
\hullface{keel!57!black}\hullmove{-5.22}{-7.57}\hullline{-4.80}{-7.57}\hullline{-6.53}{-6.77}\hullline{-7.00}{-6.77}\hullfill
\hullface{keel!57!black}\hullmove{8.37}{-5.86}\hullline{7.93}{-5.86}\hullline{6.53}{-6.77}\hullline{7.00}{-6.77}\hullfill
\hullface{keel!57!black}\hullmove{-7.00}{-6.77}\hullline{-6.53}{-6.77}\hullline{-7.93}{-5.86}\hullline{-8.37}{-5.86}\hullfill
\hullface{keel!57!black}\hullmove{8.82}{-5.44}\hullline{8.40}{-5.44}\hullline{7.93}{-5.86}\hullline{8.37}{-5.86}\hullfill
\hullface{keel!63!black}\hullmove{-8.37}{-5.86}\hullline{-7.93}{-5.86}\hullline{-8.40}{-5.44}\hullline{-8.82}{-5.44}\hullfill
\hullface{keel!57!black}\hullmove{9.29}{-4.85}\hullline{8.92}{-4.85}\hullline{8.40}{-5.44}\hullline{8.82}{-5.44}\hullfill
\hullface{keel!69!black}\hullmove{-8.82}{-5.44}\hullline{-8.40}{-5.44}\hullline{-8.92}{-4.85}\hullline{-9.29}{-4.85}\hullfill
\hullface{keel!57!black}\hullmove{9.81}{-3.75}\hullline{9.53}{-3.75}\hullline{8.92}{-4.85}\hullline{9.29}{-4.85}\hullfill
\hullface{keel!78!black}\hullmove{-9.29}{-4.85}\hullline{-8.92}{-4.85}\hullline{-9.53}{-3.75}\hullline{-9.81}{-3.75}\hullfill
\hullface{keel!57!black}\hullmove{10.03}{-2.58}\hullline{9.82}{-2.58}\hullline{9.53}{-3.75}\hullline{9.81}{-3.75}\hullfill
\hullface{keel!84!black}\hullmove{-9.81}{-3.75}\hullline{-9.53}{-3.75}\hullline{-9.82}{-2.58}\hullline{-10.03}{-2.58}\hullfill
\hullface{keel!57!black}\hullmove{10.08}{-1.32}\hullline{9.90}{-1.32}\hullline{9.82}{-2.58}\hullline{10.03}{-2.58}\hullfill
\hullface{keel!87!black}\hullmove{-10.03}{-2.58}\hullline{-9.82}{-2.58}\hullline{-9.90}{-1.32}\hullline{-10.08}{-1.32}\hullfill
\hullface{keel!57!black}\hullmove{10.09}{0.00}\hullline{9.91}{0.00}\hullline{9.90}{-1.32}\hullline{10.08}{-1.32}\hullfill
\hullface{keel!90!black}\hullmove{-10.08}{-1.32}\hullline{-9.90}{-1.32}\hullline{-9.91}{0.00}\hullline{-10.09}{0.00}\hullfill
\hullface{topsides!57!black}\hullmove{10.17}{0.78}\hullline{10.01}{0.78}\hullline{9.91}{0.00}\hullline{10.09}{0.00}\hullfill
\hullface{topsides!84!black}\hullmove{-10.09}{0.00}\hullline{-9.91}{0.00}\hullline{-10.01}{0.78}\hullline{-10.17}{0.78}\hullfill
\hullface{topsides!57!black}\hullmove{10.30}{1.56}\hullline{10.17}{1.56}\hullline{10.01}{0.78}\hullline{10.17}{0.78}\hullfill
\hullface{topsides!81!black}\hullmove{-10.17}{0.78}\hullline{-10.01}{0.78}\hullline{-10.17}{1.56}\hullline{-10.30}{1.56}\hullfill
\hullface{topsides!57!black}\hullmove{10.46}{2.34}\hullline{10.36}{2.34}\hullline{10.17}{1.56}\hullline{10.30}{1.56}\hullfill
\hullface{topsides!78!black}\hullmove{-10.30}{1.56}\hullline{-10.17}{1.56}\hullline{-10.36}{2.34}\hullline{-10.46}{2.34}\hullfill
\hullface{topsides!57!black}\hullmove{10.64}{3.12}\hullline{10.58}{3.12}\hullline{10.36}{2.34}\hullline{10.46}{2.34}\hullfill
\hullface{topsides!75!black}\hullmove{-10.46}{2.34}\hullline{-10.36}{2.34}\hullline{-10.58}{3.12}\hullline{-10.64}{3.12}\hullfill
\hullface{topsides!57!black}\hullmove{10.84}{3.90}\hullline{10.83}{3.90}\hullline{10.58}{3.12}\hullline{10.64}{3.12}\hullfill
\hullface{topsides!72!black}\hullmove{-10.64}{3.12}\hullline{-10.58}{3.12}\hullline{-10.83}{3.90}\hullline{-10.84}{3.90}\hullfill
\hullface{keel!57!black}\hullmove{1.11}{-8.77}\hullline{0.99}{-8.77}\hullline{0.00}{-9.06}\hullline{0.00}{-9.06}\hullfill
\hullface{keel!57!black}\hullmove{0.00}{-9.06}\hullline{0.00}{-9.06}\hullline{-0.99}{-8.77}\hullline{-1.11}{-8.77}\hullfill
\hullface{keel!57!black}\hullmove{2.89}{-8.25}\hullline{2.63}{-8.25}\hullline{0.99}{-8.77}\hullline{1.11}{-8.77}\hullfill
\hullface{keel!57!black}\hullmove{-1.11}{-8.77}\hullline{-0.99}{-8.77}\hullline{-2.63}{-8.25}\hullline{-2.89}{-8.25}\hullfill
\hullface{keel!57!black}\hullmove{4.80}{-7.57}\hullline{4.40}{-7.57}\hullline{2.63}{-8.25}\hullline{2.89}{-8.25}\hullfill
\hullface{keel!57!black}\hullmove{-2.89}{-8.25}\hullline{-2.63}{-8.25}\hullline{-4.40}{-7.57}\hullline{-4.80}{-7.57}\hullfill
\hullface{keel!57!black}\hullmove{6.53}{-6.77}\hullline{6.07}{-6.77}\hullline{4.40}{-7.57}\hullline{4.80}{-7.57}\hullfill
\hullface{keel!57!black}\hullmove{-4.80}{-7.57}\hullline{-4.40}{-7.57}\hullline{-6.07}{-6.77}\hullline{-6.53}{-6.77}\hullfill
\hullface{keel!57!black}\hullmove{7.93}{-5.86}\hullline{7.47}{-5.86}\hullline{6.07}{-6.77}\hullline{6.53}{-6.77}\hullfill
\hullface{keel!57!black}\hullmove{-6.53}{-6.77}\hullline{-6.07}{-6.77}\hullline{-7.47}{-5.86}\hullline{-7.93}{-5.86}\hullfill
\hullface{keel!57!black}\hullmove{8.40}{-5.44}\hullline{7.97}{-5.44}\hullline{7.47}{-5.86}\hullline{7.93}{-5.86}\hullfill
\hullface{keel!63!black}\hullmove{-7.93}{-5.86}\hullline{-7.47}{-5.86}\hullline{-7.97}{-5.44}\hullline{-8.40}{-5.44}\hullfill
\hullface{keel!57!black}\hullmove{8.92}{-4.85}\hullline{8.53}{-4.85}\hullline{7.97}{-5.44}\hullline{8.40}{-5.44}\hullfill
\hullface{keel!66!black}\hullmove{-8.40}{-5.44}\hullline{-7.97}{-5.44}\hullline{-8.53}{-4.85}\hullline{-8.92}{-4.85}\hullfill
\hullface{keel!57!black}\hullmove{9.53}{-3.75}\hullline{9.21}{-3.75}\hullline{8.53}{-4.85}\hullline{8.92}{-4.85}\hullfill
\hullface{keel!75!black}\hullmove{-8.92}{-4.85}\hullline{-8.53}{-4.85}\hullline{-9.21}{-3.75}\hullline{-9.53}{-3.75}\hullfill
\hullface{keel!57!black}\hullmove{9.82}{-2.58}\hullline{9.57}{-2.58}\hullline{9.21}{-3.75}\hullline{9.53}{-3.75}\hullfill
\hullface{keel!84!black}\hullmove{-9.53}{-3.75}\hullline{-9.21}{-3.75}\hullline{-9.57}{-2.58}\hullline{-9.82}{-2.58}\hullfill
\hullface{keel!57!black}\hullmove{9.90}{-1.32}\hullline{9.70}{-1.32}\hullline{9.57}{-2.58}\hullline{9.82}{-2.58}\hullfill
\hullface{keel!87!black}\hullmove{-9.82}{-2.58}\hullline{-9.57}{-2.58}\hullline{-9.70}{-1.32}\hullline{-9.90}{-1.32}\hullfill
\hullface{keel!57!black}\hullmove{9.91}{0.00}\hullline{9.72}{0.00}\hullline{9.70}{-1.32}\hullline{9.90}{-1.32}\hullfill
\hullface{keel!90!black}\hullmove{-9.90}{-1.32}\hullline{-9.70}{-1.32}\hullline{-9.72}{0.00}\hullline{-9.91}{0.00}\hullfill
\hullface{topsides!57!black}\hullmove{10.01}{0.78}\hullline{9.83}{0.78}\hullline{9.72}{0.00}\hullline{9.91}{0.00}\hullfill
\hullface{topsides!84!black}\hullmove{-9.91}{0.00}\hullline{-9.72}{0.00}\hullline{-9.83}{0.78}\hullline{-10.01}{0.78}\hullfill
\hullface{topsides!57!black}\hullmove{10.17}{1.56}\hullline{10.02}{1.56}\hullline{9.83}{0.78}\hullline{10.01}{0.78}\hullfill
\hullface{topsides!81!black}\hullmove{-10.01}{0.78}\hullline{-9.83}{0.78}\hullline{-10.02}{1.56}\hullline{-10.17}{1.56}\hullfill
\hullface{topsides!57!black}\hullmove{10.36}{2.34}\hullline{10.25}{2.34}\hullline{10.02}{1.56}\hullline{10.17}{1.56}\hullfill
\hullface{topsides!78!black}\hullmove{-10.17}{1.56}\hullline{-10.02}{1.56}\hullline{-10.25}{2.34}\hullline{-10.36}{2.34}\hullfill
\hullface{topsides!57!black}\hullmove{10.58}{3.12}\hullline{10.51}{3.12}\hullline{10.25}{2.34}\hullline{10.36}{2.34}\hullfill
\hullface{topsides!75!black}\hullmove{-10.36}{2.34}\hullline{-10.25}{2.34}\hullline{-10.51}{3.12}\hullline{-10.58}{3.12}\hullfill
\hullface{topsides!57!black}\hullmove{10.83}{3.90}\hullline{10.81}{3.90}\hullline{10.51}{3.12}\hullline{10.58}{3.12}\hullfill
\hullface{topsides!69!black}\hullmove{-10.58}{3.12}\hullline{-10.51}{3.12}\hullline{-10.81}{3.90}\hullline{-10.83}{3.90}\hullfill
\hullface{keel!57!black}\hullmove{0.99}{-8.77}\hullline{0.90}{-8.77}\hullline{0.00}{-9.06}\hullline{0.00}{-9.06}\hullfill
\hullface{keel!57!black}\hullmove{0.00}{-9.06}\hullline{0.00}{-9.06}\hullline{-0.90}{-8.77}\hullline{-0.99}{-8.77}\hullfill
\hullface{keel!57!black}\hullmove{2.63}{-8.25}\hullline{2.38}{-8.25}\hullline{0.90}{-8.77}\hullline{0.99}{-8.77}\hullfill
\hullface{keel!57!black}\hullmove{-0.99}{-8.77}\hullline{-0.90}{-8.77}\hullline{-2.38}{-8.25}\hullline{-2.63}{-8.25}\hullfill
\hullface{keel!57!black}\hullmove{4.40}{-7.57}\hullline{4.03}{-7.57}\hullline{2.38}{-8.25}\hullline{2.63}{-8.25}\hullfill
\hullface{keel!57!black}\hullmove{-2.63}{-8.25}\hullline{-2.38}{-8.25}\hullline{-4.03}{-7.57}\hullline{-4.40}{-7.57}\hullfill
\hullface{keel!57!black}\hullmove{6.07}{-6.77}\hullline{5.63}{-6.77}\hullline{4.03}{-7.57}\hullline{4.40}{-7.57}\hullfill
\hullface{keel!57!black}\hullmove{-4.40}{-7.57}\hullline{-4.03}{-7.57}\hullline{-5.63}{-6.77}\hullline{-6.07}{-6.77}\hullfill
\hullface{keel!57!black}\hullmove{7.47}{-5.86}\hullline{7.02}{-5.86}\hullline{5.63}{-6.77}\hullline{6.07}{-6.77}\hullfill
\hullface{keel!57!black}\hullmove{-6.07}{-6.77}\hullline{-5.63}{-6.77}\hullline{-7.02}{-5.86}\hullline{-7.47}{-5.86}\hullfill
\hullface{keel!57!black}\hullmove{7.97}{-5.44}\hullline{7.53}{-5.44}\hullline{7.02}{-5.86}\hullline{7.47}{-5.86}\hullfill
\hullface{keel!60!black}\hullmove{-7.47}{-5.86}\hullline{-7.02}{-5.86}\hullline{-7.53}{-5.44}\hullline{-7.97}{-5.44}\hullfill
\hullface{keel!57!black}\hullmove{8.53}{-4.85}\hullline{8.11}{-4.85}\hullline{7.53}{-5.44}\hullline{7.97}{-5.44}\hullfill
\hullface{keel!66!black}\hullmove{-7.97}{-5.44}\hullline{-7.53}{-5.44}\hullline{-8.11}{-4.85}\hullline{-8.53}{-4.85}\hullfill
\hullface{keel!57!black}\hullmove{9.21}{-3.75}\hullline{8.87}{-3.75}\hullline{8.11}{-4.85}\hullline{8.53}{-4.85}\hullfill
\hullface{keel!75!black}\hullmove{-8.53}{-4.85}\hullline{-8.11}{-4.85}\hullline{-8.87}{-3.75}\hullline{-9.21}{-3.75}\hullfill
\hullface{keel!57!black}\hullmove{9.57}{-2.58}\hullline{9.30}{-2.58}\hullline{8.87}{-3.75}\hullline{9.21}{-3.75}\hullfill
\hullface{keel!81!black}\hullmove{-9.21}{-3.75}\hullline{-8.87}{-3.75}\hullline{-9.30}{-2.58}\hullline{-9.57}{-2.58}\hullfill
\hullface{keel!57!black}\hullmove{9.70}{-1.32}\hullline{9.47}{-1.32}\hullline{9.30}{-2.58}\hullline{9.57}{-2.58}\hullfill
\hullface{keel!87!black}\hullmove{-9.57}{-2.58}\hullline{-9.30}{-2.58}\hullline{-9.47}{-1.32}\hullline{-9.70}{-1.32}\hullfill
\hullface{keel!57!black}\hullmove{9.72}{0.00}\hullline{9.50}{0.00}\hullline{9.47}{-1.32}\hullline{9.70}{-1.32}\hullfill
\hullface{keel!90!black}\hullmove{-9.70}{-1.32}\hullline{-9.47}{-1.32}\hullline{-9.50}{0.00}\hullline{-9.72}{0.00}\hullfill
\hullface{topsides!57!black}\hullmove{9.83}{0.78}\hullline{9.63}{0.78}\hullline{9.50}{0.00}\hullline{9.72}{0.00}\hullfill
\hullface{topsides!84!black}\hullmove{-9.72}{0.00}\hullline{-9.50}{0.00}\hullline{-9.63}{0.78}\hullline{-9.83}{0.78}\hullfill
\hullface{topsides!57!black}\hullmove{10.02}{1.56}\hullline{9.85}{1.56}\hullline{9.63}{0.78}\hullline{9.83}{0.78}\hullfill
\hullface{topsides!81!black}\hullmove{-9.83}{0.78}\hullline{-9.63}{0.78}\hullline{-9.85}{1.56}\hullline{-10.02}{1.56}\hullfill
\hullface{topsides!57!black}\hullmove{10.25}{2.34}\hullline{10.12}{2.34}\hullline{9.85}{1.56}\hullline{10.02}{1.56}\hullfill
\hullface{topsides!75!black}\hullmove{-10.02}{1.56}\hullline{-9.85}{1.56}\hullline{-10.12}{2.34}\hullline{-10.25}{2.34}\hullfill
\hullface{topsides!57!black}\hullmove{10.51}{3.12}\hullline{10.44}{3.12}\hullline{10.12}{2.34}\hullline{10.25}{2.34}\hullfill
\hullface{topsides!72!black}\hullmove{-10.25}{2.34}\hullline{-10.12}{2.34}\hullline{-10.44}{3.12}\hullline{-10.51}{3.12}\hullfill
\hullface{topsides!57!black}\hullmove{10.81}{3.90}\hullline{10.78}{3.90}\hullline{10.44}{3.12}\hullline{10.51}{3.12}\hullfill
\hullface{topsides!69!black}\hullmove{-10.51}{3.12}\hullline{-10.44}{3.12}\hullline{-10.78}{3.90}\hullline{-10.81}{3.90}\hullfill
\hullface{keel!57!black}\hullmove{0.90}{-8.77}\hullline{0.83}{-8.77}\hullline{0.00}{-9.06}\hullline{0.00}{-9.06}\hullfill
\hullface{keel!57!black}\hullmove{2.38}{-8.25}\hullline{2.21}{-8.25}\hullline{0.83}{-8.77}\hullline{0.90}{-8.77}\hullfill
\hullface{keel!57!black}\hullmove{-0.90}{-8.77}\hullline{-0.83}{-8.77}\hullline{-2.21}{-8.25}\hullline{-2.38}{-8.25}\hullfill
\hullface{keel!57!black}\hullmove{4.03}{-7.57}\hullline{3.77}{-7.57}\hullline{2.21}{-8.25}\hullline{2.38}{-8.25}\hullfill
\hullface{keel!57!black}\hullmove{-2.38}{-8.25}\hullline{-2.21}{-8.25}\hullline{-3.77}{-7.57}\hullline{-4.03}{-7.57}\hullfill
\hullface{keel!57!black}\hullmove{5.63}{-6.77}\hullline{5.31}{-6.77}\hullline{3.77}{-7.57}\hullline{4.03}{-7.57}\hullfill
\hullface{keel!57!black}\hullmove{-4.03}{-7.57}\hullline{-3.77}{-7.57}\hullline{-5.31}{-6.77}\hullline{-5.63}{-6.77}\hullfill
\hullface{keel!57!black}\hullmove{7.02}{-5.86}\hullline{6.68}{-5.86}\hullline{5.31}{-6.77}\hullline{5.63}{-6.77}\hullfill
\hullface{keel!57!black}\hullmove{-5.63}{-6.77}\hullline{-5.31}{-6.77}\hullline{-6.68}{-5.86}\hullline{-7.02}{-5.86}\hullfill
\hullface{keel!57!black}\hullmove{7.53}{-5.44}\hullline{7.19}{-5.44}\hullline{6.68}{-5.86}\hullline{7.02}{-5.86}\hullfill
\hullface{keel!60!black}\hullmove{-7.02}{-5.86}\hullline{-6.68}{-5.86}\hullline{-7.19}{-5.44}\hullline{-7.53}{-5.44}\hullfill
\hullface{keel!57!black}\hullmove{8.11}{-4.85}\hullline{7.79}{-4.85}\hullline{7.19}{-5.44}\hullline{7.53}{-5.44}\hullfill
\hullface{keel!66!black}\hullmove{-7.53}{-5.44}\hullline{-7.19}{-5.44}\hullline{-7.79}{-4.85}\hullline{-8.11}{-4.85}\hullfill
\hullface{keel!57!black}\hullmove{8.87}{-3.75}\hullline{8.58}{-3.75}\hullline{7.79}{-4.85}\hullline{8.11}{-4.85}\hullfill
\hullface{keel!72!black}\hullmove{-8.11}{-4.85}\hullline{-7.79}{-4.85}\hullline{-8.58}{-3.75}\hullline{-8.87}{-3.75}\hullfill
\hullface{keel!57!black}\hullmove{9.30}{-2.58}\hullline{9.07}{-2.58}\hullline{8.58}{-3.75}\hullline{8.87}{-3.75}\hullfill
\hullface{keel!81!black}\hullmove{-8.87}{-3.75}\hullline{-8.58}{-3.75}\hullline{-9.07}{-2.58}\hullline{-9.30}{-2.58}\hullfill
\hullface{keel!57!black}\hullmove{9.47}{-1.32}\hullline{9.27}{-1.32}\hullline{9.07}{-2.58}\hullline{9.30}{-2.58}\hullfill
\hullface{keel!87!black}\hullmove{-9.30}{-2.58}\hullline{-9.07}{-2.58}\hullline{-9.27}{-1.32}\hullline{-9.47}{-1.32}\hullfill
\hullface{keel!57!black}\hullmove{9.50}{0.00}\hullline{9.31}{0.00}\hullline{9.27}{-1.32}\hullline{9.47}{-1.32}\hullfill
\hullface{keel!90!black}\hullmove{-9.47}{-1.32}\hullline{-9.27}{-1.32}\hullline{-9.31}{0.00}\hullline{-9.50}{0.00}\hullfill
\hullface{topsides!57!black}\hullmove{9.63}{0.78}\hullline{9.46}{0.78}\hullline{9.31}{0.00}\hullline{9.50}{0.00}\hullfill
\hullface{topsides!84!black}\hullmove{-9.50}{0.00}\hullline{-9.31}{0.00}\hullline{-9.46}{0.78}\hullline{-9.63}{0.78}\hullfill
\hullface{topsides!57!black}\hullmove{9.85}{1.56}\hullline{9.71}{1.56}\hullline{9.46}{0.78}\hullline{9.63}{0.78}\hullfill
\hullface{topsides!78!black}\hullmove{-9.63}{0.78}\hullline{-9.46}{0.78}\hullline{-9.71}{1.56}\hullline{-9.85}{1.56}\hullfill
\hullface{topsides!57!black}\hullmove{10.12}{2.34}\hullline{10.02}{2.34}\hullline{9.71}{1.56}\hullline{9.85}{1.56}\hullfill
\hullface{topsides!57!black}\hullmove{10.44}{3.12}\hullline{10.37}{3.12}\hullline{10.02}{2.34}\hullline{10.12}{2.34}\hullfill
\hullface{topsides!57!black}\hullmove{10.78}{3.90}\hullline{10.76}{3.90}\hullline{10.37}{3.12}\hullline{10.44}{3.12}\hullfill
\hullface{topsides!66!black}\hullmove{-10.44}{3.12}\hullline{-10.37}{3.12}\hullline{-10.76}{3.90}\hullline{-10.78}{3.90}\hullfill
\hullface{keel!57!black}\hullmove{0.00}{-9.06}\hullline{0.00}{-9.06}\hullline{-0.83}{-8.77}\hullline{-0.83}{-8.77}\hullfill
\hullface{keel!57!black}\hullmove{2.21}{-8.25}\hullline{2.21}{-8.25}\hullline{0.83}{-8.77}\hullline{0.83}{-8.77}\hullfill
\hullface{keel!57!black}\hullmove{-0.83}{-8.77}\hullline{-0.83}{-8.77}\hullline{-2.21}{-8.25}\hullline{-2.21}{-8.25}\hullfill
\hullface{keel!57!black}\hullmove{3.77}{-7.57}\hullline{3.77}{-7.57}\hullline{2.21}{-8.25}\hullline{2.21}{-8.25}\hullfill
\hullface{keel!57!black}\hullmove{-2.21}{-8.25}\hullline{-2.21}{-8.25}\hullline{-3.77}{-7.57}\hullline{-3.77}{-7.57}\hullfill
\hullface{keel!57!black}\hullmove{5.31}{-6.77}\hullline{5.31}{-6.77}\hullline{3.77}{-7.57}\hullline{3.77}{-7.57}\hullfill
\hullface{keel!57!black}\hullmove{-3.77}{-7.57}\hullline{-3.77}{-7.57}\hullline{-5.31}{-6.77}\hullline{-5.31}{-6.77}\hullfill
\hullface{keel!57!black}\hullmove{6.68}{-5.86}\hullline{6.68}{-5.86}\hullline{5.31}{-6.77}\hullline{5.31}{-6.77}\hullfill
\hullface{keel!57!black}\hullmove{-5.31}{-6.77}\hullline{-5.31}{-6.77}\hullline{-6.68}{-5.86}\hullline{-6.68}{-5.86}\hullfill
\hullface{keel!57!black}\hullmove{7.19}{-5.44}\hullline{7.19}{-5.44}\hullline{6.68}{-5.86}\hullline{6.68}{-5.86}\hullfill
\hullface{keel!60!black}\hullmove{-6.68}{-5.86}\hullline{-6.68}{-5.86}\hullline{-7.19}{-5.44}\hullline{-7.19}{-5.44}\hullfill
\hullface{keel!57!black}\hullmove{7.79}{-4.85}\hullline{7.79}{-4.85}\hullline{7.19}{-5.44}\hullline{7.19}{-5.44}\hullfill
\hullface{keel!66!black}\hullmove{-7.19}{-5.44}\hullline{-7.19}{-5.44}\hullline{-7.79}{-4.85}\hullline{-7.79}{-4.85}\hullfill
\hullface{keel!57!black}\hullmove{8.58}{-3.75}\hullline{8.58}{-3.75}\hullline{7.79}{-4.85}\hullline{7.79}{-4.85}\hullfill
\hullface{keel!72!black}\hullmove{-7.79}{-4.85}\hullline{-7.79}{-4.85}\hullline{-8.58}{-3.75}\hullline{-8.58}{-3.75}\hullfill
\hullface{keel!57!black}\hullmove{9.07}{-2.58}\hullline{9.07}{-2.58}\hullline{8.58}{-3.75}\hullline{8.58}{-3.75}\hullfill
\hullface{keel!81!black}\hullmove{-8.58}{-3.75}\hullline{-8.58}{-3.75}\hullline{-9.07}{-2.58}\hullline{-9.07}{-2.58}\hullfill
\hullface{keel!57!black}\hullmove{9.27}{-1.32}\hullline{9.27}{-1.32}\hullline{9.07}{-2.58}\hullline{9.07}{-2.58}\hullfill
\hullface{keel!87!black}\hullmove{-9.07}{-2.58}\hullline{-9.07}{-2.58}\hullline{-9.27}{-1.32}\hullline{-9.27}{-1.32}\hullfill
\hullface{keel!57!black}\hullmove{9.31}{0.00}\hullline{9.31}{0.00}\hullline{9.27}{-1.32}\hullline{9.27}{-1.32}\hullfill
\hullface{keel!90!black}\hullmove{-9.27}{-1.32}\hullline{-9.27}{-1.32}\hullline{-9.31}{0.00}\hullline{-9.31}{0.00}\hullfill
\hullface{topsides!99!black}\hullmove{9.46}{0.78}\hullline{9.71}{1.56}\hullline{10.02}{2.34}\hullline{10.37}{3.12}\hullline{10.76}{3.90}\hullline{10.76}{7.59}\hullline{10.37}{6.07}\hullline{10.02}{4.55}\hullline{9.71}{3.04}\hullline{9.46}{1.52}\hullline{9.31}{0.00}\hullline{9.31}{0.00}\hullfill
\hullface{topsides!99!black}\hullmove{-9.31}{0.00}\hullline{-9.31}{0.00}\hullline{-9.46}{1.52}\hullline{-9.71}{3.04}\hullline{-10.02}{4.55}\hullline{-10.37}{6.07}\hullline{-10.76}{7.59}\hullline{-10.76}{3.90}\hullline{-10.37}{3.12}\hullline{-10.02}{2.34}\hullline{-9.71}{1.56}\hullline{-9.46}{0.78}\hullfill
\hullface{keel!57!black}\hullmove{0.00}{-9.06}\hullline{0.00}{-9.06}\hullline{-0.81}{-8.77}\hullline{-0.83}{-8.77}\hullfill
\hullface{keel!57!black}\hullmove{2.21}{-8.25}\hullline{2.16}{-8.25}\hullline{0.81}{-8.77}\hullline{0.83}{-8.77}\hullfill
\hullface{keel!57!black}\hullmove{-0.83}{-8.77}\hullline{-0.81}{-8.77}\hullline{-2.16}{-8.25}\hullline{-2.21}{-8.25}\hullfill
\hullface{keel!57!black}\hullmove{3.77}{-7.57}\hullline{3.69}{-7.57}\hullline{2.16}{-8.25}\hullline{2.21}{-8.25}\hullfill
\hullface{keel!57!black}\hullmove{-2.21}{-8.25}\hullline{-2.16}{-8.25}\hullline{-3.69}{-7.57}\hullline{-3.77}{-7.57}\hullfill
\hullface{keel!57!black}\hullmove{5.31}{-6.77}\hullline{5.21}{-6.77}\hullline{3.69}{-7.57}\hullline{3.77}{-7.57}\hullfill
\hullface{keel!57!black}\hullmove{-3.77}{-7.57}\hullline{-3.69}{-7.57}\hullline{-5.21}{-6.77}\hullline{-5.31}{-6.77}\hullfill
\hullface{keel!57!black}\hullmove{6.68}{-5.86}\hullline{6.57}{-5.86}\hullline{5.21}{-6.77}\hullline{5.31}{-6.77}\hullfill
\hullface{keel!57!black}\hullmove{-5.31}{-6.77}\hullline{-5.21}{-6.77}\hullline{-6.57}{-5.86}\hullline{-6.68}{-5.86}\hullfill
\hullface{keel!57!black}\hullmove{7.19}{-5.44}\hullline{7.08}{-5.44}\hullline{6.57}{-5.86}\hullline{6.68}{-5.86}\hullfill
\hullface{keel!60!black}\hullmove{-6.68}{-5.86}\hullline{-6.57}{-5.86}\hullline{-7.08}{-5.44}\hullline{-7.19}{-5.44}\hullfill
\hullface{keel!57!black}\hullmove{7.79}{-4.85}\hullline{7.68}{-4.85}\hullline{7.08}{-5.44}\hullline{7.19}{-5.44}\hullfill
\hullface{keel!66!black}\hullmove{-7.19}{-5.44}\hullline{-7.08}{-5.44}\hullline{-7.68}{-4.85}\hullline{-7.79}{-4.85}\hullfill
\hullface{keel!57!black}\hullmove{8.58}{-3.75}\hullline{8.49}{-3.75}\hullline{7.68}{-4.85}\hullline{7.79}{-4.85}\hullfill
\hullface{keel!72!black}\hullmove{-7.79}{-4.85}\hullline{-7.68}{-4.85}\hullline{-8.49}{-3.75}\hullline{-8.58}{-3.75}\hullfill
\hullface{keel!57!black}\hullmove{9.07}{-2.58}\hullline{8.99}{-2.58}\hullline{8.49}{-3.75}\hullline{8.58}{-3.75}\hullfill
\hullface{keel!81!black}\hullmove{-8.58}{-3.75}\hullline{-8.49}{-3.75}\hullline{-8.99}{-2.58}\hullline{-9.07}{-2.58}\hullfill
\hullface{keel!57!black}\hullmove{9.27}{-1.32}\hullline{9.21}{-1.32}\hullline{8.99}{-2.58}\hullline{9.07}{-2.58}\hullfill
\hullface{keel!87!black}\hullmove{-9.07}{-2.58}\hullline{-8.99}{-2.58}\hullline{-9.21}{-1.32}\hullline{-9.27}{-1.32}\hullfill
\hullface{keel!57!black}\hullmove{9.31}{0.00}\hullline{9.25}{0.00}\hullline{9.21}{-1.32}\hullline{9.27}{-1.32}\hullfill
\hullface{keel!93!black}\hullmove{-9.27}{-1.32}\hullline{-9.21}{-1.32}\hullline{-9.25}{0.00}\hullline{-9.31}{0.00}\hullfill
\hullface{topsides!57!black}\hullmove{9.46}{1.52}\hullline{9.41}{1.52}\hullline{9.25}{0.00}\hullline{9.31}{0.00}\hullfill
\hullface{topsides!87!black}\hullmove{-9.31}{0.00}\hullline{-9.25}{0.00}\hullline{-9.41}{1.52}\hullline{-9.46}{1.52}\hullfill
\hullface{topsides!57!black}\hullmove{9.71}{3.04}\hullline{9.66}{3.04}\hullline{9.41}{1.52}\hullline{9.46}{1.52}\hullfill
\hullface{topsides!84!black}\hullmove{-9.46}{1.52}\hullline{-9.41}{1.52}\hullline{-9.66}{3.04}\hullline{-9.71}{3.04}\hullfill
\hullface{topsides!57!black}\hullmove{10.02}{4.55}\hullline{9.98}{4.55}\hullline{9.66}{3.04}\hullline{9.71}{3.04}\hullfill
\hullface{topsides!81!black}\hullmove{-9.71}{3.04}\hullline{-9.66}{3.04}\hullline{-9.98}{4.55}\hullline{-10.02}{4.55}\hullfill
\hullface{topsides!57!black}\hullmove{10.37}{6.07}\hullline{10.34}{6.07}\hullline{9.98}{4.55}\hullline{10.02}{4.55}\hullfill
\hullface{topsides!78!black}\hullmove{-10.02}{4.55}\hullline{-9.98}{4.55}\hullline{-10.34}{6.07}\hullline{-10.37}{6.07}\hullfill
\hullface{topsides!57!black}\hullmove{10.76}{7.59}\hullline{10.75}{7.59}\hullline{10.34}{6.07}\hullline{10.37}{6.07}\hullfill
\hullface{topsides!75!black}\hullmove{-10.37}{6.07}\hullline{-10.34}{6.07}\hullline{-10.75}{7.59}\hullline{-10.76}{7.59}\hullfill
\hullface{keel!57!black}\hullmove{0.81}{-8.77}\hullline{0.73}{-8.77}\hullline{0.00}{-9.06}\hullline{0.00}{-9.06}\hullfill
\hullface{keel!57!black}\hullmove{0.00}{-9.06}\hullline{0.00}{-9.06}\hullline{-0.73}{-8.77}\hullline{-0.81}{-8.77}\hullfill
\hullface{keel!57!black}\hullmove{2.16}{-8.25}\hullline{1.96}{-8.25}\hullline{0.73}{-8.77}\hullline{0.81}{-8.77}\hullfill
\hullface{keel!57!black}\hullmove{-0.81}{-8.77}\hullline{-0.73}{-8.77}\hullline{-1.96}{-8.25}\hullline{-2.16}{-8.25}\hullfill
\hullface{keel!57!black}\hullmove{3.69}{-7.57}\hullline{3.38}{-7.57}\hullline{1.96}{-8.25}\hullline{2.16}{-8.25}\hullfill
\hullface{keel!57!black}\hullmove{-2.16}{-8.25}\hullline{-1.96}{-8.25}\hullline{-3.38}{-7.57}\hullline{-3.69}{-7.57}\hullfill
\hullface{keel!57!black}\hullmove{5.21}{-6.77}\hullline{4.81}{-6.77}\hullline{3.38}{-7.57}\hullline{3.69}{-7.57}\hullfill
\hullface{keel!57!black}\hullmove{-3.69}{-7.57}\hullline{-3.38}{-7.57}\hullline{-4.81}{-6.77}\hullline{-5.21}{-6.77}\hullfill
\hullface{keel!57!black}\hullmove{6.57}{-5.86}\hullline{6.12}{-5.86}\hullline{4.81}{-6.77}\hullline{5.21}{-6.77}\hullfill
\hullface{keel!57!black}\hullmove{-5.21}{-6.77}\hullline{-4.81}{-6.77}\hullline{-6.12}{-5.86}\hullline{-6.57}{-5.86}\hullfill
\hullface{keel!57!black}\hullmove{7.08}{-5.44}\hullline{6.63}{-5.44}\hullline{6.12}{-5.86}\hullline{6.57}{-5.86}\hullfill
\hullface{keel!60!black}\hullmove{-6.57}{-5.86}\hullline{-6.12}{-5.86}\hullline{-6.63}{-5.44}\hullline{-7.08}{-5.44}\hullfill
\hullface{keel!57!black}\hullmove{7.68}{-4.85}\hullline{7.24}{-4.85}\hullline{6.63}{-5.44}\hullline{7.08}{-5.44}\hullfill
\hullface{keel!66!black}\hullmove{-7.08}{-5.44}\hullline{-6.63}{-5.44}\hullline{-7.24}{-4.85}\hullline{-7.68}{-4.85}\hullfill
\hullface{keel!57!black}\hullmove{8.49}{-3.75}\hullline{8.09}{-3.75}\hullline{7.24}{-4.85}\hullline{7.68}{-4.85}\hullfill
\hullface{keel!72!black}\hullmove{-7.68}{-4.85}\hullline{-7.24}{-4.85}\hullline{-8.09}{-3.75}\hullline{-8.49}{-3.75}\hullfill
\hullface{keel!57!black}\hullmove{8.99}{-2.58}\hullline{8.64}{-2.58}\hullline{8.09}{-3.75}\hullline{8.49}{-3.75}\hullfill
\hullface{keel!81!black}\hullmove{-8.49}{-3.75}\hullline{-8.09}{-3.75}\hullline{-8.64}{-2.58}\hullline{-8.99}{-2.58}\hullfill
\hullface{keel!57!black}\hullmove{9.21}{-1.32}\hullline{8.92}{-1.32}\hullline{8.64}{-2.58}\hullline{8.99}{-2.58}\hullfill
\hullface{keel!87!black}\hullmove{-8.99}{-2.58}\hullline{-8.64}{-2.58}\hullline{-8.92}{-1.32}\hullline{-9.21}{-1.32}\hullfill
\hullface{keel!57!black}\hullmove{9.25}{0.00}\hullline{8.97}{0.00}\hullline{8.92}{-1.32}\hullline{9.21}{-1.32}\hullfill
\hullface{keel!93!black}\hullmove{-9.21}{-1.32}\hullline{-8.92}{-1.32}\hullline{-8.97}{0.00}\hullline{-9.25}{0.00}\hullfill
\hullface{topsides!57!black}\hullmove{9.41}{1.52}\hullline{9.15}{1.52}\hullline{8.97}{0.00}\hullline{9.25}{0.00}\hullfill
\hullface{topsides!90!black}\hullmove{-9.25}{0.00}\hullline{-8.97}{0.00}\hullline{-9.15}{1.52}\hullline{-9.41}{1.52}\hullfill
\hullface{topsides!57!black}\hullmove{9.66}{3.04}\hullline{9.45}{3.04}\hullline{9.15}{1.52}\hullline{9.41}{1.52}\hullfill
\hullface{topsides!84!black}\hullmove{-9.41}{1.52}\hullline{-9.15}{1.52}\hullline{-9.45}{3.04}\hullline{-9.66}{3.04}\hullfill
\hullface{topsides!57!black}\hullmove{9.98}{4.55}\hullline{9.82}{4.55}\hullline{9.45}{3.04}\hullline{9.66}{3.04}\hullfill
\hullface{topsides!81!black}\hullmove{-9.66}{3.04}\hullline{-9.45}{3.04}\hullline{-9.82}{4.55}\hullline{-9.98}{4.55}\hullfill
\hullface{topsides!57!black}\hullmove{10.34}{6.07}\hullline{10.24}{6.07}\hullline{9.82}{4.55}\hullline{9.98}{4.55}\hullfill
\hullface{topsides!78!black}\hullmove{-9.98}{4.55}\hullline{-9.82}{4.55}\hullline{-10.24}{6.07}\hullline{-10.34}{6.07}\hullfill
\hullface{topsides!57!black}\hullmove{10.75}{7.59}\hullline{10.70}{7.59}\hullline{10.24}{6.07}\hullline{10.34}{6.07}\hullfill
\hullface{topsides!75!black}\hullmove{-10.34}{6.07}\hullline{-10.24}{6.07}\hullline{-10.70}{7.59}\hullline{-10.75}{7.59}\hullfill
\hullface{keel!57!black}\hullmove{0.73}{-8.77}\hullline{0.65}{-8.77}\hullline{0.00}{-9.06}\hullline{0.00}{-9.06}\hullfill
\hullface{keel!57!black}\hullmove{0.00}{-9.06}\hullline{0.00}{-9.06}\hullline{-0.65}{-8.77}\hullline{-0.73}{-8.77}\hullfill
\hullface{keel!57!black}\hullmove{1.96}{-8.25}\hullline{1.77}{-8.25}\hullline{0.65}{-8.77}\hullline{0.73}{-8.77}\hullfill
\hullface{keel!57!black}\hullmove{-0.73}{-8.77}\hullline{-0.65}{-8.77}\hullline{-1.77}{-8.25}\hullline{-1.96}{-8.25}\hullfill
\hullface{keel!57!black}\hullmove{3.38}{-7.57}\hullline{3.08}{-7.57}\hullline{1.77}{-8.25}\hullline{1.96}{-8.25}\hullfill
\hullface{keel!57!black}\hullmove{-1.96}{-8.25}\hullline{-1.77}{-8.25}\hullline{-3.08}{-7.57}\hullline{-3.38}{-7.57}\hullfill
\hullface{keel!57!black}\hullmove{4.81}{-6.77}\hullline{4.42}{-6.77}\hullline{3.08}{-7.57}\hullline{3.38}{-7.57}\hullfill
\hullface{keel!57!black}\hullmove{-3.38}{-7.57}\hullline{-3.08}{-7.57}\hullline{-4.42}{-6.77}\hullline{-4.81}{-6.77}\hullfill
\hullface{keel!57!black}\hullmove{6.12}{-5.86}\hullline{5.68}{-5.86}\hullline{4.42}{-6.77}\hullline{4.81}{-6.77}\hullfill
\hullface{keel!57!black}\hullmove{-4.81}{-6.77}\hullline{-4.42}{-6.77}\hullline{-5.68}{-5.86}\hullline{-6.12}{-5.86}\hullfill
\hullface{keel!57!black}\hullmove{6.63}{-5.44}\hullline{6.18}{-5.44}\hullline{5.68}{-5.86}\hullline{6.12}{-5.86}\hullfill
\hullface{keel!60!black}\hullmove{-6.12}{-5.86}\hullline{-5.68}{-5.86}\hullline{-6.18}{-5.44}\hullline{-6.63}{-5.44}\hullfill
\hullface{keel!57!black}\hullmove{7.24}{-4.85}\hullline{6.79}{-4.85}\hullline{6.18}{-5.44}\hullline{6.63}{-5.44}\hullfill
\hullface{keel!66!black}\hullmove{-6.63}{-5.44}\hullline{-6.18}{-5.44}\hullline{-6.79}{-4.85}\hullline{-7.24}{-4.85}\hullfill
\hullface{keel!57!black}\hullmove{8.09}{-3.75}\hullline{7.66}{-3.75}\hullline{6.79}{-4.85}\hullline{7.24}{-4.85}\hullfill
\hullface{keel!72!black}\hullmove{-7.24}{-4.85}\hullline{-6.79}{-4.85}\hullline{-7.66}{-3.75}\hullline{-8.09}{-3.75}\hullfill
\hullface{keel!57!black}\hullmove{8.64}{-2.58}\hullline{8.27}{-2.58}\hullline{7.66}{-3.75}\hullline{8.09}{-3.75}\hullfill
\hullface{keel!81!black}\hullmove{-8.09}{-3.75}\hullline{-7.66}{-3.75}\hullline{-8.27}{-2.58}\hullline{-8.64}{-2.58}\hullfill
\hullface{keel!57!black}\hullmove{8.92}{-1.32}\hullline{8.59}{-1.32}\hullline{8.27}{-2.58}\hullline{8.64}{-2.58}\hullfill
\hullface{keel!87!black}\hullmove{-8.64}{-2.58}\hullline{-8.27}{-2.58}\hullline{-8.59}{-1.32}\hullline{-8.92}{-1.32}\hullfill
\hullface{keel!57!black}\hullmove{8.97}{0.00}\hullline{8.67}{0.00}\hullline{8.59}{-1.32}\hullline{8.92}{-1.32}\hullfill
\hullface{keel!93!black}\hullmove{-8.92}{-1.32}\hullline{-8.59}{-1.32}\hullline{-8.67}{0.00}\hullline{-8.97}{0.00}\hullfill
\hullface{topsides!57!black}\hullmove{9.15}{1.52}\hullline{8.88}{1.52}\hullline{8.67}{0.00}\hullline{8.97}{0.00}\hullfill
\hullface{topsides!90!black}\hullmove{-8.97}{0.00}\hullline{-8.67}{0.00}\hullline{-8.88}{1.52}\hullline{-9.15}{1.52}\hullfill
\hullface{topsides!57!black}\hullmove{9.45}{3.04}\hullline{9.22}{3.04}\hullline{8.88}{1.52}\hullline{9.15}{1.52}\hullfill
\hullface{topsides!84!black}\hullmove{-9.15}{1.52}\hullline{-8.88}{1.52}\hullline{-9.22}{3.04}\hullline{-9.45}{3.04}\hullfill
\hullface{topsides!57!black}\hullmove{9.82}{4.55}\hullline{9.63}{4.55}\hullline{9.22}{3.04}\hullline{9.45}{3.04}\hullfill
\hullface{topsides!81!black}\hullmove{-9.45}{3.04}\hullline{-9.22}{3.04}\hullline{-9.63}{4.55}\hullline{-9.82}{4.55}\hullfill
\hullface{topsides!57!black}\hullmove{10.24}{6.07}\hullline{10.11}{6.07}\hullline{9.63}{4.55}\hullline{9.82}{4.55}\hullfill
\hullface{topsides!78!black}\hullmove{-9.82}{4.55}\hullline{-9.63}{4.55}\hullline{-10.11}{6.07}\hullline{-10.24}{6.07}\hullfill
\hullface{topsides!57!black}\hullmove{10.70}{7.59}\hullline{10.64}{7.59}\hullline{10.11}{6.07}\hullline{10.24}{6.07}\hullfill
\hullface{topsides!72!black}\hullmove{-10.24}{6.07}\hullline{-10.11}{6.07}\hullline{-10.64}{7.59}\hullline{-10.70}{7.59}\hullfill
\hullface{keel!57!black}\hullmove{0.00}{-9.06}\hullline{0.00}{-9.06}\hullline{-0.59}{-8.77}\hullline{-0.65}{-8.77}\hullfill
\hullface{keel!57!black}\hullmove{1.77}{-8.25}\hullline{1.60}{-8.25}\hullline{0.59}{-8.77}\hullline{0.65}{-8.77}\hullfill
\hullface{keel!57!black}\hullmove{-0.65}{-8.77}\hullline{-0.59}{-8.77}\hullline{-1.60}{-8.25}\hullline{-1.77}{-8.25}\hullfill
\hullface{keel!57!black}\hullmove{3.08}{-7.57}\hullline{2.80}{-7.57}\hullline{1.60}{-8.25}\hullline{1.77}{-8.25}\hullfill
\hullface{keel!57!black}\hullmove{-1.77}{-8.25}\hullline{-1.60}{-8.25}\hullline{-2.80}{-7.57}\hullline{-3.08}{-7.57}\hullfill
\hullface{keel!57!black}\hullmove{4.42}{-6.77}\hullline{4.05}{-6.77}\hullline{2.80}{-7.57}\hullline{3.08}{-7.57}\hullfill
\hullface{keel!57!black}\hullmove{-3.08}{-7.57}\hullline{-2.80}{-7.57}\hullline{-4.05}{-6.77}\hullline{-4.42}{-6.77}\hullfill
\hullface{keel!57!black}\hullmove{5.68}{-5.86}\hullline{5.25}{-5.86}\hullline{4.05}{-6.77}\hullline{4.42}{-6.77}\hullfill
\hullface{keel!57!black}\hullmove{-4.42}{-6.77}\hullline{-4.05}{-6.77}\hullline{-5.25}{-5.86}\hullline{-5.68}{-5.86}\hullfill
\hullface{keel!57!black}\hullmove{6.18}{-5.44}\hullline{5.73}{-5.44}\hullline{5.25}{-5.86}\hullline{5.68}{-5.86}\hullfill
\hullface{keel!63!black}\hullmove{-5.68}{-5.86}\hullline{-5.25}{-5.86}\hullline{-5.73}{-5.44}\hullline{-6.18}{-5.44}\hullfill
\hullface{keel!57!black}\hullmove{6.79}{-4.85}\hullline{6.33}{-4.85}\hullline{5.73}{-5.44}\hullline{6.18}{-5.44}\hullfill
\hullface{keel!66!black}\hullmove{-6.18}{-5.44}\hullline{-5.73}{-5.44}\hullline{-6.33}{-4.85}\hullline{-6.79}{-4.85}\hullfill
\hullface{keel!57!black}\hullmove{7.66}{-3.75}\hullline{7.22}{-3.75}\hullline{6.33}{-4.85}\hullline{6.79}{-4.85}\hullfill
\hullface{keel!72!black}\hullmove{-6.79}{-4.85}\hullline{-6.33}{-4.85}\hullline{-7.22}{-3.75}\hullline{-7.66}{-3.75}\hullfill
\hullface{keel!57!black}\hullmove{8.27}{-2.58}\hullline{7.86}{-2.58}\hullline{7.22}{-3.75}\hullline{7.66}{-3.75}\hullfill
\hullface{keel!81!black}\hullmove{-7.66}{-3.75}\hullline{-7.22}{-3.75}\hullline{-7.86}{-2.58}\hullline{-8.27}{-2.58}\hullfill
\hullface{keel!57!black}\hullmove{8.59}{-1.32}\hullline{8.23}{-1.32}\hullline{7.86}{-2.58}\hullline{8.27}{-2.58}\hullfill
\hullface{keel!87!black}\hullmove{-8.27}{-2.58}\hullline{-7.86}{-2.58}\hullline{-8.23}{-1.32}\hullline{-8.59}{-1.32}\hullfill
\hullface{keel!57!black}\hullmove{8.67}{0.00}\hullline{8.34}{0.00}\hullline{8.23}{-1.32}\hullline{8.59}{-1.32}\hullfill
\hullface{keel!93!black}\hullmove{-8.59}{-1.32}\hullline{-8.23}{-1.32}\hullline{-8.34}{0.00}\hullline{-8.67}{0.00}\hullfill
\hullface{topsides!57!black}\hullmove{8.88}{1.52}\hullline{8.57}{1.52}\hullline{8.34}{0.00}\hullline{8.67}{0.00}\hullfill
\hullface{topsides!90!black}\hullmove{-8.67}{0.00}\hullline{-8.34}{0.00}\hullline{-8.57}{1.52}\hullline{-8.88}{1.52}\hullfill
\hullface{topsides!57!black}\hullmove{9.22}{3.04}\hullline{8.96}{3.04}\hullline{8.57}{1.52}\hullline{8.88}{1.52}\hullfill
\hullface{topsides!84!black}\hullmove{-8.88}{1.52}\hullline{-8.57}{1.52}\hullline{-8.96}{3.04}\hullline{-9.22}{3.04}\hullfill
\hullface{topsides!57!black}\hullmove{9.63}{4.55}\hullline{9.43}{4.55}\hullline{8.96}{3.04}\hullline{9.22}{3.04}\hullfill
\hullface{topsides!81!black}\hullmove{-9.22}{3.04}\hullline{-8.96}{3.04}\hullline{-9.43}{4.55}\hullline{-9.63}{4.55}\hullfill
\hullface{topsides!57!black}\hullmove{10.11}{6.07}\hullline{9.97}{6.07}\hullline{9.43}{4.55}\hullline{9.63}{4.55}\hullfill
\hullface{topsides!78!black}\hullmove{-9.63}{4.55}\hullline{-9.43}{4.55}\hullline{-9.97}{6.07}\hullline{-10.11}{6.07}\hullfill
\hullface{topsides!57!black}\hullmove{10.64}{7.59}\hullline{10.57}{7.59}\hullline{9.97}{6.07}\hullline{10.11}{6.07}\hullfill
\hullface{topsides!72!black}\hullmove{-10.11}{6.07}\hullline{-9.97}{6.07}\hullline{-10.57}{7.59}\hullline{-10.64}{7.59}\hullfill
\hullface{keel!57!black}\hullmove{0.59}{-8.77}\hullline{0.53}{-8.77}\hullline{0.00}{-9.06}\hullline{0.00}{-9.06}\hullfill
\hullface{keel!57!black}\hullmove{0.00}{-9.06}\hullline{0.00}{-9.06}\hullline{-0.53}{-8.77}\hullline{-0.59}{-8.77}\hullfill
\hullface{keel!57!black}\hullmove{1.60}{-8.25}\hullline{1.44}{-8.25}\hullline{0.53}{-8.77}\hullline{0.59}{-8.77}\hullfill
\hullface{keel!57!black}\hullmove{-0.59}{-8.77}\hullline{-0.53}{-8.77}\hullline{-1.44}{-8.25}\hullline{-1.60}{-8.25}\hullfill
\hullface{keel!57!black}\hullmove{2.80}{-7.57}\hullline{2.53}{-7.57}\hullline{1.44}{-8.25}\hullline{1.60}{-8.25}\hullfill
\hullface{keel!57!black}\hullmove{-1.60}{-8.25}\hullline{-1.44}{-8.25}\hullline{-2.53}{-7.57}\hullline{-2.80}{-7.57}\hullfill
\hullface{keel!57!black}\hullmove{4.05}{-6.77}\hullline{3.69}{-6.77}\hullline{2.53}{-7.57}\hullline{2.80}{-7.57}\hullfill
\hullface{keel!57!black}\hullmove{-2.80}{-7.57}\hullline{-2.53}{-7.57}\hullline{-3.69}{-6.77}\hullline{-4.05}{-6.77}\hullfill
\hullface{keel!57!black}\hullmove{5.25}{-5.86}\hullline{4.83}{-5.86}\hullline{3.69}{-6.77}\hullline{4.05}{-6.77}\hullfill
\hullface{keel!57!black}\hullmove{-4.05}{-6.77}\hullline{-3.69}{-6.77}\hullline{-4.83}{-5.86}\hullline{-5.25}{-5.86}\hullfill
\hullface{keel!57!black}\hullmove{5.73}{-5.44}\hullline{5.29}{-5.44}\hullline{4.83}{-5.86}\hullline{5.25}{-5.86}\hullfill
\hullface{keel!63!black}\hullmove{-5.25}{-5.86}\hullline{-4.83}{-5.86}\hullline{-5.29}{-5.44}\hullline{-5.73}{-5.44}\hullfill
\hullface{keel!57!black}\hullmove{6.33}{-4.85}\hullline{5.87}{-4.85}\hullline{5.29}{-5.44}\hullline{5.73}{-5.44}\hullfill
\hullface{keel!66!black}\hullmove{-5.73}{-5.44}\hullline{-5.29}{-5.44}\hullline{-5.87}{-4.85}\hullline{-6.33}{-4.85}\hullfill
\hullface{keel!57!black}\hullmove{7.22}{-3.75}\hullline{6.75}{-3.75}\hullline{5.87}{-4.85}\hullline{6.33}{-4.85}\hullfill
\hullface{keel!72!black}\hullmove{-6.33}{-4.85}\hullline{-5.87}{-4.85}\hullline{-6.75}{-3.75}\hullline{-7.22}{-3.75}\hullfill
\hullface{keel!57!black}\hullmove{7.86}{-2.58}\hullline{7.42}{-2.58}\hullline{6.75}{-3.75}\hullline{7.22}{-3.75}\hullfill
\hullface{keel!78!black}\hullmove{-7.22}{-3.75}\hullline{-6.75}{-3.75}\hullline{-7.42}{-2.58}\hullline{-7.86}{-2.58}\hullfill
\hullface{keel!57!black}\hullmove{8.23}{-1.32}\hullline{7.84}{-1.32}\hullline{7.42}{-2.58}\hullline{7.86}{-2.58}\hullfill
\hullface{keel!87!black}\hullmove{-7.86}{-2.58}\hullline{-7.42}{-2.58}\hullline{-7.84}{-1.32}\hullline{-8.23}{-1.32}\hullfill
\hullface{keel!57!black}\hullmove{8.34}{0.00}\hullline{7.97}{0.00}\hullline{7.84}{-1.32}\hullline{8.23}{-1.32}\hullfill
\hullface{keel!93!black}\hullmove{-8.23}{-1.32}\hullline{-7.84}{-1.32}\hullline{-7.97}{0.00}\hullline{-8.34}{0.00}\hullfill
\hullface{topsides!57!black}\hullmove{8.57}{1.52}\hullline{8.23}{1.52}\hullline{7.97}{0.00}\hullline{8.34}{0.00}\hullfill
\hullface{topsides!90!black}\hullmove{-8.34}{0.00}\hullline{-7.97}{0.00}\hullline{-8.23}{1.52}\hullline{-8.57}{1.52}\hullfill
\hullface{topsides!57!black}\hullmove{8.96}{3.04}\hullline{8.67}{3.04}\hullline{8.23}{1.52}\hullline{8.57}{1.52}\hullfill
\hullface{topsides!84!black}\hullmove{-8.57}{1.52}\hullline{-8.23}{1.52}\hullline{-8.67}{3.04}\hullline{-8.96}{3.04}\hullfill
\hullface{topsides!57!black}\hullmove{9.43}{4.55}\hullline{9.20}{4.55}\hullline{8.67}{3.04}\hullline{8.96}{3.04}\hullfill
\hullface{topsides!81!black}\hullmove{-8.96}{3.04}\hullline{-8.67}{3.04}\hullline{-9.20}{4.55}\hullline{-9.43}{4.55}\hullfill
\hullface{topsides!57!black}\hullmove{9.97}{6.07}\hullline{9.80}{6.07}\hullline{9.20}{4.55}\hullline{9.43}{4.55}\hullfill
\hullface{topsides!75!black}\hullmove{-9.43}{4.55}\hullline{-9.20}{4.55}\hullline{-9.80}{6.07}\hullline{-9.97}{6.07}\hullfill
\hullface{topsides!57!black}\hullmove{10.57}{7.59}\hullline{10.47}{7.59}\hullline{9.80}{6.07}\hullline{9.97}{6.07}\hullfill
\hullface{topsides!72!black}\hullmove{-9.97}{6.07}\hullline{-9.80}{6.07}\hullline{-10.47}{7.59}\hullline{-10.57}{7.59}\hullfill
\hullface{keel!57!black}\hullmove{0.53}{-8.77}\hullline{0.47}{-8.77}\hullline{0.00}{-9.06}\hullline{0.00}{-9.06}\hullfill
\hullface{keel!57!black}\hullmove{0.00}{-9.06}\hullline{0.00}{-9.06}\hullline{-0.47}{-8.77}\hullline{-0.53}{-8.77}\hullfill
\hullface{keel!57!black}\hullmove{1.44}{-8.25}\hullline{1.29}{-8.25}\hullline{0.47}{-8.77}\hullline{0.53}{-8.77}\hullfill
\hullface{keel!57!black}\hullmove{-0.53}{-8.77}\hullline{-0.47}{-8.77}\hullline{-1.29}{-8.25}\hullline{-1.44}{-8.25}\hullfill
\hullface{keel!57!black}\hullmove{2.53}{-7.57}\hullline{2.28}{-7.57}\hullline{1.29}{-8.25}\hullline{1.44}{-8.25}\hullfill
\hullface{keel!57!black}\hullmove{-1.44}{-8.25}\hullline{-1.29}{-8.25}\hullline{-2.28}{-7.57}\hullline{-2.53}{-7.57}\hullfill
\hullface{keel!57!black}\hullmove{3.69}{-6.77}\hullline{3.35}{-6.77}\hullline{2.28}{-7.57}\hullline{2.53}{-7.57}\hullfill
\hullface{keel!57!black}\hullmove{-2.53}{-7.57}\hullline{-2.28}{-7.57}\hullline{-3.35}{-6.77}\hullline{-3.69}{-6.77}\hullfill
\hullface{keel!57!black}\hullmove{4.83}{-5.86}\hullline{4.41}{-5.86}\hullline{3.35}{-6.77}\hullline{3.69}{-6.77}\hullfill
\hullface{keel!60!black}\hullmove{-3.69}{-6.77}\hullline{-3.35}{-6.77}\hullline{-4.41}{-5.86}\hullline{-4.83}{-5.86}\hullfill
\hullface{keel!57!black}\hullmove{5.29}{-5.44}\hullline{4.85}{-5.44}\hullline{4.41}{-5.86}\hullline{4.83}{-5.86}\hullfill
\hullface{keel!63!black}\hullmove{-4.83}{-5.86}\hullline{-4.41}{-5.86}\hullline{-4.85}{-5.44}\hullline{-5.29}{-5.44}\hullfill
\hullface{keel!57!black}\hullmove{5.87}{-4.85}\hullline{5.40}{-4.85}\hullline{4.85}{-5.44}\hullline{5.29}{-5.44}\hullfill
\hullface{keel!66!black}\hullmove{-5.29}{-5.44}\hullline{-4.85}{-5.44}\hullline{-5.40}{-4.85}\hullline{-5.87}{-4.85}\hullfill
\hullface{keel!57!black}\hullmove{6.75}{-3.75}\hullline{6.27}{-3.75}\hullline{5.40}{-4.85}\hullline{5.87}{-4.85}\hullfill
\hullface{keel!72!black}\hullmove{-5.87}{-4.85}\hullline{-5.40}{-4.85}\hullline{-6.27}{-3.75}\hullline{-6.75}{-3.75}\hullfill
\hullface{keel!57!black}\hullmove{7.42}{-2.58}\hullline{6.96}{-2.58}\hullline{6.27}{-3.75}\hullline{6.75}{-3.75}\hullfill
\hullface{keel!78!black}\hullmove{-6.75}{-3.75}\hullline{-6.27}{-3.75}\hullline{-6.96}{-2.58}\hullline{-7.42}{-2.58}\hullfill
\hullface{keel!57!black}\hullmove{7.84}{-1.32}\hullline{7.41}{-1.32}\hullline{6.96}{-2.58}\hullline{7.42}{-2.58}\hullfill
\hullface{keel!87!black}\hullmove{-7.42}{-2.58}\hullline{-6.96}{-2.58}\hullline{-7.41}{-1.32}\hullline{-7.84}{-1.32}\hullfill
\hullface{keel!57!black}\hullmove{7.97}{0.00}\hullline{7.57}{0.00}\hullline{7.41}{-1.32}\hullline{7.84}{-1.32}\hullfill
\hullface{keel!93!black}\hullmove{-7.84}{-1.32}\hullline{-7.41}{-1.32}\hullline{-7.57}{0.00}\hullline{-7.97}{0.00}\hullfill
\hullface{topsides!57!black}\hullmove{8.23}{1.52}\hullline{7.87}{1.52}\hullline{7.57}{0.00}\hullline{7.97}{0.00}\hullfill
\hullface{topsides!90!black}\hullmove{-7.97}{0.00}\hullline{-7.57}{0.00}\hullline{-7.87}{1.52}\hullline{-8.23}{1.52}\hullfill
\hullface{topsides!57!black}\hullmove{8.67}{3.04}\hullline{8.35}{3.04}\hullline{7.87}{1.52}\hullline{8.23}{1.52}\hullfill
\hullface{topsides!84!black}\hullmove{-8.23}{1.52}\hullline{-7.87}{1.52}\hullline{-8.35}{3.04}\hullline{-8.67}{3.04}\hullfill
\hullface{topsides!57!black}\hullmove{9.20}{4.55}\hullline{8.94}{4.55}\hullline{8.35}{3.04}\hullline{8.67}{3.04}\hullfill
\hullface{topsides!81!black}\hullmove{-8.67}{3.04}\hullline{-8.35}{3.04}\hullline{-8.94}{4.55}\hullline{-9.20}{4.55}\hullfill
\hullface{topsides!57!black}\hullmove{9.80}{6.07}\hullline{9.61}{6.07}\hullline{8.94}{4.55}\hullline{9.20}{4.55}\hullfill
\hullface{topsides!75!black}\hullmove{-9.20}{4.55}\hullline{-8.94}{4.55}\hullline{-9.61}{6.07}\hullline{-9.80}{6.07}\hullfill
\hullface{topsides!57!black}\hullmove{10.47}{7.59}\hullline{10.36}{7.59}\hullline{9.61}{6.07}\hullline{9.80}{6.07}\hullfill
\hullface{topsides!72!black}\hullmove{-9.80}{6.07}\hullline{-9.61}{6.07}\hullline{-10.36}{7.59}\hullline{-10.47}{7.59}\hullfill
\hullface{keel!57!black}\hullmove{0.47}{-8.77}\hullline{0.42}{-8.77}\hullline{0.00}{-9.06}\hullline{0.00}{-9.06}\hullfill
\hullface{keel!57!black}\hullmove{0.00}{-9.06}\hullline{0.00}{-9.06}\hullline{-0.42}{-8.77}\hullline{-0.47}{-8.77}\hullfill
\hullface{keel!57!black}\hullmove{1.29}{-8.25}\hullline{1.15}{-8.25}\hullline{0.42}{-8.77}\hullline{0.47}{-8.77}\hullfill
\hullface{keel!57!black}\hullmove{-0.47}{-8.77}\hullline{-0.42}{-8.77}\hullline{-1.15}{-8.25}\hullline{-1.29}{-8.25}\hullfill
\hullface{keel!57!black}\hullmove{2.28}{-7.57}\hullline{2.05}{-7.57}\hullline{1.15}{-8.25}\hullline{1.29}{-8.25}\hullfill
\hullface{keel!57!black}\hullmove{-1.29}{-8.25}\hullline{-1.15}{-8.25}\hullline{-2.05}{-7.57}\hullline{-2.28}{-7.57}\hullfill
\hullface{keel!57!black}\hullmove{3.35}{-6.77}\hullline{3.02}{-6.77}\hullline{2.05}{-7.57}\hullline{2.28}{-7.57}\hullfill
\hullface{keel!57!black}\hullmove{-2.28}{-7.57}\hullline{-2.05}{-7.57}\hullline{-3.02}{-6.77}\hullline{-3.35}{-6.77}\hullfill
\hullface{keel!57!black}\hullmove{4.41}{-5.86}\hullline{3.89}{-5.96}\hullline{3.02}{-6.77}\hullline{3.35}{-6.77}\hullfill
\hullface{keel!60!black}\hullmove{-3.35}{-6.77}\hullline{-3.02}{-6.77}\hullline{-3.89}{-5.96}\hullline{-4.41}{-5.86}\hullfill
\hullface{keel!57!black}\hullmove{4.85}{-5.44}\hullline{3.96}{-5.89}\hullline{3.89}{-5.96}\hullline{4.41}{-5.86}\hullfill
\hullface{keel!63!black}\hullmove{-4.41}{-5.86}\hullline{-3.89}{-5.96}\hullline{-3.96}{-5.89}\hullline{-4.85}{-5.44}\hullfill
\hullface{keel!57!black}\hullmove{4.85}{-5.44}\hullline{4.41}{-5.44}\hullline{4.15}{-5.70}\hullline{4.85}{-5.44}\hullfill
\hullface{keel!66!black}\hullmove{-4.85}{-5.44}\hullline{-4.15}{-5.70}\hullline{-4.41}{-5.44}\hullline{-4.85}{-5.44}\hullfill
\hullface{keel!57!black}\hullmove{4.85}{-5.44}\hullline{4.66}{-5.17}\hullline{4.41}{-5.44}\hullline{4.85}{-5.44}\hullfill
\hullface{keel!69!black}\hullmove{-4.85}{-5.44}\hullline{-4.41}{-5.44}\hullline{-4.66}{-5.17}\hullline{-4.85}{-5.44}\hullfill
\hullface{keel!57!black}\hullmove{5.40}{-4.85}\hullline{4.94}{-4.85}\hullline{4.89}{-4.91}\hullline{4.85}{-5.44}\hullfill
\hullface{keel!69!black}\hullmove{-4.85}{-5.44}\hullline{-4.89}{-4.91}\hullline{-4.94}{-4.85}\hullline{-5.40}{-4.85}\hullfill
\hullface{keel!57!black}\hullmove{6.27}{-3.75}\hullline{5.78}{-3.75}\hullline{4.94}{-4.85}\hullline{5.40}{-4.85}\hullfill
\hullface{keel!72!black}\hullmove{-5.40}{-4.85}\hullline{-4.94}{-4.85}\hullline{-5.78}{-3.75}\hullline{-6.27}{-3.75}\hullfill
\hullface{keel!57!black}\hullmove{6.96}{-2.58}\hullline{6.47}{-2.58}\hullline{5.78}{-3.75}\hullline{6.27}{-3.75}\hullfill
\hullface{keel!81!black}\hullmove{-6.27}{-3.75}\hullline{-5.78}{-3.75}\hullline{-6.47}{-2.58}\hullline{-6.96}{-2.58}\hullfill
\hullface{keel!57!black}\hullmove{7.41}{-1.32}\hullline{6.95}{-1.32}\hullline{6.47}{-2.58}\hullline{6.96}{-2.58}\hullfill
\hullface{keel!87!black}\hullmove{-6.96}{-2.58}\hullline{-6.47}{-2.58}\hullline{-6.95}{-1.32}\hullline{-7.41}{-1.32}\hullfill
\hullface{keel!57!black}\hullmove{7.57}{0.00}\hullline{7.14}{0.00}\hullline{6.95}{-1.32}\hullline{7.41}{-1.32}\hullfill
\hullface{keel!93!black}\hullmove{-7.41}{-1.32}\hullline{-6.95}{-1.32}\hullline{-7.14}{0.00}\hullline{-7.57}{0.00}\hullfill
\hullface{topsides!57!black}\hullmove{7.87}{1.52}\hullline{7.47}{1.52}\hullline{7.14}{0.00}\hullline{7.57}{0.00}\hullfill
\hullface{topsides!90!black}\hullmove{-7.57}{0.00}\hullline{-7.14}{0.00}\hullline{-7.47}{1.52}\hullline{-7.87}{1.52}\hullfill
\hullface{topsides!57!black}\hullmove{8.35}{3.04}\hullline{7.99}{3.04}\hullline{7.47}{1.52}\hullline{7.87}{1.52}\hullfill
\hullface{topsides!84!black}\hullmove{-7.87}{1.52}\hullline{-7.47}{1.52}\hullline{-7.99}{3.04}\hullline{-8.35}{3.04}\hullfill
\hullface{topsides!57!black}\hullmove{8.94}{4.55}\hullline{8.65}{4.55}\hullline{7.99}{3.04}\hullline{8.35}{3.04}\hullfill
\hullface{topsides!81!black}\hullmove{-8.35}{3.04}\hullline{-7.99}{3.04}\hullline{-8.65}{4.55}\hullline{-8.94}{4.55}\hullfill
\hullface{topsides!57!black}\hullmove{9.61}{6.07}\hullline{9.39}{6.07}\hullline{8.65}{4.55}\hullline{8.94}{4.55}\hullfill
\hullface{topsides!75!black}\hullmove{-8.94}{4.55}\hullline{-8.65}{4.55}\hullline{-9.39}{6.07}\hullline{-9.61}{6.07}\hullfill
\hullface{topsides!57!black}\hullmove{10.36}{7.59}\hullline{10.21}{7.59}\hullline{9.39}{6.07}\hullline{9.61}{6.07}\hullfill
\hullface{topsides!72!black}\hullmove{-9.61}{6.07}\hullline{-9.39}{6.07}\hullline{-10.21}{7.59}\hullline{-10.36}{7.59}\hullfill
\hullface{keel!57!black}\hullmove{4.85}{-5.44}\hullline{4.85}{-5.44}\hullline{4.15}{-5.70}\hullline{4.00}{-5.86}\hullline{3.96}{-5.89}\hullline{4.85}{-5.44}\hullfill
\hullface{keel!66!black}\hullmove{-4.85}{-5.44}\hullline{-3.96}{-5.89}\hullline{-4.00}{-5.86}\hullline{-4.15}{-5.70}\hullline{-4.85}{-5.44}\hullline{-4.85}{-5.44}\hullfill
\hullface{keel!57!black}\hullmove{4.85}{-5.44}\hullline{4.85}{-5.44}\hullline{4.89}{-4.91}\hullline{4.83}{-4.98}\hullline{4.66}{-5.17}\hullline{4.85}{-5.44}\hullfill
\hullface{keel!69!black}\hullmove{-4.85}{-5.44}\hullline{-4.66}{-5.17}\hullline{-4.83}{-4.98}\hullline{-4.89}{-4.91}\hullline{-4.85}{-5.44}\hullline{-4.85}{-5.44}\hullfill
\hullface{keel!57!black}\hullmove{0.00}{-9.06}\hullline{0.00}{-9.06}\hullline{-0.37}{-8.77}\hullline{-0.42}{-8.77}\hullfill
\hullface{keel!57!black}\hullmove{1.15}{-8.25}\hullline{1.02}{-8.25}\hullline{0.37}{-8.77}\hullline{0.42}{-8.77}\hullfill
\hullface{keel!57!black}\hullmove{-0.42}{-8.77}\hullline{-0.37}{-8.77}\hullline{-1.02}{-8.25}\hullline{-1.15}{-8.25}\hullfill
\hullface{keel!57!black}\hullmove{2.05}{-7.57}\hullline{1.82}{-7.57}\hullline{1.02}{-8.25}\hullline{1.15}{-8.25}\hullfill
\hullface{keel!57!black}\hullmove{-1.15}{-8.25}\hullline{-1.02}{-8.25}\hullline{-1.82}{-7.57}\hullline{-2.05}{-7.57}\hullfill
\hullface{keel!57!black}\hullmove{3.02}{-6.77}\hullline{2.70}{-6.77}\hullline{1.82}{-7.57}\hullline{2.05}{-7.57}\hullfill
\hullface{keel!57!black}\hullmove{-2.05}{-7.57}\hullline{-1.82}{-7.57}\hullline{-2.70}{-6.77}\hullline{-3.02}{-6.77}\hullfill
\hullface{keel!57!black}\hullmove{3.89}{-5.96}\hullline{2.96}{-6.51}\hullline{2.70}{-6.77}\hullline{3.02}{-6.77}\hullfill
\hullface{keel!60!black}\hullmove{-3.02}{-6.77}\hullline{-2.70}{-6.77}\hullline{-2.96}{-6.51}\hullline{-3.89}{-5.96}\hullfill
\hullface{keel!57!black}\hullmove{3.96}{-5.89}\hullline{3.11}{-6.37}\hullline{2.96}{-6.51}\hullline{3.89}{-5.96}\hullfill
\hullface{keel!63!black}\hullmove{-3.89}{-5.96}\hullline{-2.96}{-6.51}\hullline{-3.11}{-6.37}\hullline{-3.96}{-5.89}\hullfill
\hullface{keel!57!black}\hullmove{4.00}{-5.86}\hullline{4.15}{-5.70}\hullline{3.60}{-5.86}\hullline{3.49}{-5.97}\hullline{3.11}{-6.37}\hullline{3.96}{-5.89}\hullfill
\hullface{keel!66!black}\hullmove{-3.96}{-5.89}\hullline{-3.11}{-6.37}\hullline{-3.49}{-5.97}\hullline{-3.60}{-5.86}\hullline{-4.15}{-5.70}\hullline{-4.00}{-5.86}\hullfill
\hullface{keel!57!black}\hullmove{4.41}{-5.44}\hullline{3.98}{-5.44}\hullline{3.60}{-5.86}\hullline{4.15}{-5.70}\hullfill
\hullface{keel!66!black}\hullmove{-4.15}{-5.70}\hullline{-3.60}{-5.86}\hullline{-3.98}{-5.44}\hullline{-4.41}{-5.44}\hullfill
\hullface{keel!57!black}\hullmove{4.66}{-5.17}\hullline{4.83}{-4.98}\hullline{4.47}{-4.85}\hullline{4.43}{-4.90}\hullline{3.98}{-5.44}\hullline{4.41}{-5.44}\hullfill
\hullface{keel!69!black}\hullmove{-4.41}{-5.44}\hullline{-3.98}{-5.44}\hullline{-4.43}{-4.90}\hullline{-4.47}{-4.85}\hullline{-4.83}{-4.98}\hullline{-4.66}{-5.17}\hullfill
\hullface{keel!57!black}\hullmove{4.89}{-4.91}\hullline{4.94}{-4.85}\hullline{4.85}{-4.36}\hullline{4.74}{-4.51}\hullline{4.47}{-4.85}\hullline{4.83}{-4.98}\hullfill
\hullface{keel!72!black}\hullmove{-4.83}{-4.98}\hullline{-4.47}{-4.85}\hullline{-4.74}{-4.51}\hullline{-4.85}{-4.36}\hullline{-4.94}{-4.85}\hullline{-4.89}{-4.91}\hullfill
\hullface{keel!57!black}\hullmove{5.78}{-3.75}\hullline{5.28}{-3.75}\hullline{4.85}{-4.36}\hullline{4.94}{-4.85}\hullfill
\hullface{keel!75!black}\hullmove{-4.94}{-4.85}\hullline{-4.85}{-4.36}\hullline{-5.28}{-3.75}\hullline{-5.78}{-3.75}\hullfill
\hullface{keel!57!black}\hullmove{6.47}{-2.58}\hullline{5.96}{-2.58}\hullline{5.28}{-3.75}\hullline{5.78}{-3.75}\hullfill
\hullface{keel!81!black}\hullmove{-5.78}{-3.75}\hullline{-5.28}{-3.75}\hullline{-5.96}{-2.58}\hullline{-6.47}{-2.58}\hullfill
\hullface{keel!57!black}\hullmove{6.95}{-1.32}\hullline{6.46}{-1.32}\hullline{5.96}{-2.58}\hullline{6.47}{-2.58}\hullfill
\hullface{keel!87!black}\hullmove{-6.47}{-2.58}\hullline{-5.96}{-2.58}\hullline{-6.46}{-1.32}\hullline{-6.95}{-1.32}\hullfill
\hullface{keel!57!black}\hullmove{7.14}{0.00}\hullline{6.68}{0.00}\hullline{6.46}{-1.32}\hullline{6.95}{-1.32}\hullfill
\hullface{keel!93!black}\hullmove{-6.95}{-1.32}\hullline{-6.46}{-1.32}\hullline{-6.68}{0.00}\hullline{-7.14}{0.00}\hullfill
\hullface{topsides!57!black}\hullmove{7.47}{1.52}\hullline{7.03}{1.52}\hullline{6.68}{0.00}\hullline{7.14}{0.00}\hullfill
\hullface{topsides!90!black}\hullmove{-7.14}{0.00}\hullline{-6.68}{0.00}\hullline{-7.03}{1.52}\hullline{-7.47}{1.52}\hullfill
\hullface{topsides!57!black}\hullmove{7.99}{3.04}\hullline{7.61}{3.04}\hullline{7.03}{1.52}\hullline{7.47}{1.52}\hullfill
\hullface{topsides!84!black}\hullmove{-7.47}{1.52}\hullline{-7.03}{1.52}\hullline{-7.61}{3.04}\hullline{-7.99}{3.04}\hullfill
\hullface{topsides!57!black}\hullmove{8.65}{4.55}\hullline{8.32}{4.55}\hullline{7.61}{3.04}\hullline{7.99}{3.04}\hullfill
\hullface{topsides!81!black}\hullmove{-7.99}{3.04}\hullline{-7.61}{3.04}\hullline{-8.32}{4.55}\hullline{-8.65}{4.55}\hullfill
\hullface{topsides!57!black}\hullmove{9.39}{6.07}\hullline{9.14}{6.07}\hullline{8.32}{4.55}\hullline{8.65}{4.55}\hullfill
\hullface{topsides!75!black}\hullmove{-8.65}{4.55}\hullline{-8.32}{4.55}\hullline{-9.14}{6.07}\hullline{-9.39}{6.07}\hullfill
\hullface{topsides!57!black}\hullmove{10.21}{7.59}\hullline{10.04}{7.59}\hullline{9.14}{6.07}\hullline{9.39}{6.07}\hullfill
\hullface{topsides!69!black}\hullmove{-9.39}{6.07}\hullline{-9.14}{6.07}\hullline{-10.04}{7.59}\hullline{-10.21}{7.59}\hullfill
\hullface{keel!57!black}\hullmove{0.37}{-8.77}\hullline{0.32}{-8.77}\hullline{0.00}{-9.06}\hullline{0.00}{-9.06}\hullfill
\hullface{keel!57!black}\hullmove{0.00}{-9.06}\hullline{0.00}{-9.06}\hullline{-0.32}{-8.77}\hullline{-0.37}{-8.77}\hullfill
\hullface{keel!57!black}\hullmove{1.02}{-8.25}\hullline{0.90}{-8.25}\hullline{0.32}{-8.77}\hullline{0.37}{-8.77}\hullfill
\hullface{keel!57!black}\hullmove{-0.37}{-8.77}\hullline{-0.32}{-8.77}\hullline{-0.90}{-8.25}\hullline{-1.02}{-8.25}\hullfill
\hullface{keel!57!black}\hullmove{1.82}{-7.57}\hullline{1.60}{-7.57}\hullline{0.90}{-8.25}\hullline{1.02}{-8.25}\hullfill
\hullface{keel!57!black}\hullmove{-1.02}{-8.25}\hullline{-0.90}{-8.25}\hullline{-1.60}{-7.57}\hullline{-1.82}{-7.57}\hullfill
\hullface{keel!57!black}\hullmove{2.70}{-6.77}\hullline{2.34}{-6.81}\hullline{1.60}{-7.57}\hullline{1.82}{-7.57}\hullfill
\hullface{keel!60!black}\hullmove{-1.82}{-7.57}\hullline{-1.60}{-7.57}\hullline{-2.34}{-6.81}\hullline{-2.70}{-6.77}\hullfill
\hullface{keel!57!black}\hullmove{2.96}{-6.51}\hullline{2.39}{-6.77}\hullline{2.34}{-6.81}\hullline{2.70}{-6.77}\hullfill
\hullface{keel!60!black}\hullmove{-2.70}{-6.77}\hullline{-2.34}{-6.81}\hullline{-2.39}{-6.77}\hullline{-2.96}{-6.51}\hullfill
\hullface{keel!57!black}\hullmove{3.11}{-6.37}\hullline{2.52}{-6.63}\hullline{2.39}{-6.77}\hullline{2.96}{-6.51}\hullfill
\hullface{keel!63!black}\hullmove{-2.96}{-6.51}\hullline{-2.39}{-6.77}\hullline{-2.52}{-6.63}\hullline{-3.11}{-6.37}\hullfill
\hullface{keel!57!black}\hullmove{3.49}{-5.97}\hullline{2.97}{-6.12}\hullline{2.52}{-6.63}\hullline{3.11}{-6.37}\hullfill
\hullface{keel!66!black}\hullmove{-3.11}{-6.37}\hullline{-2.52}{-6.63}\hullline{-2.97}{-6.12}\hullline{-3.49}{-5.97}\hullfill
\hullface{keel!57!black}\hullmove{3.60}{-5.86}\hullline{3.20}{-5.86}\hullline{2.97}{-6.12}\hullline{3.49}{-5.97}\hullfill
\hullface{keel!66!black}\hullmove{-3.49}{-5.97}\hullline{-2.97}{-6.12}\hullline{-3.20}{-5.86}\hullline{-3.60}{-5.86}\hullfill
\hullface{keel!57!black}\hullmove{3.98}{-5.44}\hullline{3.55}{-5.44}\hullline{3.20}{-5.86}\hullline{3.60}{-5.86}\hullfill
\hullface{keel!69!black}\hullmove{-3.60}{-5.86}\hullline{-3.20}{-5.86}\hullline{-3.55}{-5.44}\hullline{-3.98}{-5.44}\hullfill
\hullface{keel!57!black}\hullmove{4.43}{-4.90}\hullline{4.01}{-4.85}\hullline{3.55}{-5.44}\hullline{3.98}{-5.44}\hullfill
\hullface{keel!72!black}\hullmove{-3.98}{-5.44}\hullline{-3.55}{-5.44}\hullline{-4.01}{-4.85}\hullline{-4.43}{-4.90}\hullfill
\hullface{keel!57!black}\hullmove{4.47}{-4.85}\hullline{4.08}{-4.75}\hullline{4.01}{-4.85}\hullline{4.43}{-4.90}\hullfill
\hullface{keel!72!black}\hullmove{-4.43}{-4.90}\hullline{-4.01}{-4.85}\hullline{-4.08}{-4.75}\hullline{-4.47}{-4.85}\hullfill
\hullface{keel!57!black}\hullmove{4.74}{-4.51}\hullline{4.44}{-4.24}\hullline{4.08}{-4.75}\hullline{4.47}{-4.85}\hullfill
\hullface{keel!75!black}\hullmove{-4.47}{-4.85}\hullline{-4.08}{-4.75}\hullline{-4.44}{-4.24}\hullline{-4.74}{-4.51}\hullfill
\hullface{keel!57!black}\hullmove{4.85}{-4.36}\hullline{4.57}{-4.06}\hullline{4.44}{-4.24}\hullline{4.74}{-4.51}\hullfill
\hullface{keel!75!black}\hullmove{-4.74}{-4.51}\hullline{-4.44}{-4.24}\hullline{-4.57}{-4.06}\hullline{-4.85}{-4.36}\hullfill
\hullface{keel!57!black}\hullmove{5.28}{-3.75}\hullline{4.76}{-3.75}\hullline{4.57}{-4.06}\hullline{4.85}{-4.36}\hullfill
\hullface{keel!78!black}\hullmove{-4.85}{-4.36}\hullline{-4.57}{-4.06}\hullline{-4.76}{-3.75}\hullline{-5.28}{-3.75}\hullfill
\hullface{keel!57!black}\hullmove{5.96}{-2.58}\hullline{5.42}{-2.58}\hullline{4.76}{-3.75}\hullline{5.28}{-3.75}\hullfill
\hullface{keel!81!black}\hullmove{-5.28}{-3.75}\hullline{-4.76}{-3.75}\hullline{-5.42}{-2.58}\hullline{-5.96}{-2.58}\hullfill
\hullface{keel!57!black}\hullmove{6.46}{-1.32}\hullline{5.93}{-1.32}\hullline{5.42}{-2.58}\hullline{5.96}{-2.58}\hullfill
\hullface{keel!87!black}\hullmove{-5.96}{-2.58}\hullline{-5.42}{-2.58}\hullline{-5.93}{-1.32}\hullline{-6.46}{-1.32}\hullfill
\hullface{keel!57!black}\hullmove{6.68}{0.00}\hullline{6.17}{0.00}\hullline{5.93}{-1.32}\hullline{6.46}{-1.32}\hullfill
\hullface{keel!93!black}\hullmove{-6.46}{-1.32}\hullline{-5.93}{-1.32}\hullline{-6.17}{0.00}\hullline{-6.68}{0.00}\hullfill
\hullface{topsides!57!black}\hullmove{7.03}{1.52}\hullline{6.56}{1.52}\hullline{6.17}{0.00}\hullline{6.68}{0.00}\hullfill
\hullface{topsides!90!black}\hullmove{-6.68}{0.00}\hullline{-6.17}{0.00}\hullline{-6.56}{1.52}\hullline{-7.03}{1.52}\hullfill
\hullface{topsides!57!black}\hullmove{7.61}{3.04}\hullline{7.19}{3.04}\hullline{6.56}{1.52}\hullline{7.03}{1.52}\hullfill
\hullface{topsides!84!black}\hullmove{-7.03}{1.52}\hullline{-6.56}{1.52}\hullline{-7.19}{3.04}\hullline{-7.61}{3.04}\hullfill
\hullface{topsides!57!black}\hullmove{8.32}{4.55}\hullline{7.96}{4.55}\hullline{7.19}{3.04}\hullline{7.61}{3.04}\hullfill
\hullface{topsides!81!black}\hullmove{-7.61}{3.04}\hullline{-7.19}{3.04}\hullline{-7.96}{4.55}\hullline{-8.32}{4.55}\hullfill
\hullface{topsides!57!black}\hullmove{9.14}{6.07}\hullline{8.85}{6.07}\hullline{7.96}{4.55}\hullline{8.32}{4.55}\hullfill
\hullface{topsides!75!black}\hullmove{-8.32}{4.55}\hullline{-7.96}{4.55}\hullline{-8.85}{6.07}\hullline{-9.14}{6.07}\hullfill
\hullface{topsides!57!black}\hullmove{10.04}{7.59}\hullline{9.83}{7.59}\hullline{8.85}{6.07}\hullline{9.14}{6.07}\hullfill
\hullface{topsides!69!black}\hullmove{-9.14}{6.07}\hullline{-8.85}{6.07}\hullline{-9.83}{7.59}\hullline{-10.04}{7.59}\hullfill
\hullface{keel!57!black}\hullmove{0.32}{-8.77}\hullline{0.28}{-8.77}\hullline{0.00}{-9.06}\hullline{0.00}{-9.06}\hullfill
\hullface{keel!57!black}\hullmove{0.90}{-8.25}\hullline{0.78}{-8.25}\hullline{0.28}{-8.77}\hullline{0.32}{-8.77}\hullfill
\hullface{keel!57!black}\hullmove{-0.32}{-8.77}\hullline{-0.28}{-8.77}\hullline{-0.78}{-8.25}\hullline{-0.90}{-8.25}\hullfill
\hullface{keel!57!black}\hullmove{1.60}{-7.57}\hullline{1.40}{-7.57}\hullline{0.78}{-8.25}\hullline{0.90}{-8.25}\hullfill
\hullface{keel!57!black}\hullmove{-0.90}{-8.25}\hullline{-0.78}{-8.25}\hullline{-1.40}{-7.57}\hullline{-1.60}{-7.57}\hullfill
\hullface{keel!57!black}\hullmove{2.34}{-6.81}\hullline{1.87}{-7.02}\hullline{1.40}{-7.57}\hullline{1.60}{-7.57}\hullfill
\hullface{keel!60!black}\hullmove{-1.60}{-7.57}\hullline{-1.40}{-7.57}\hullline{-1.87}{-7.02}\hullline{-2.34}{-6.81}\hullfill
\hullface{keel!57!black}\hullmove{2.39}{-6.77}\hullline{2.05}{-6.81}\hullline{1.87}{-7.02}\hullline{2.34}{-6.81}\hullfill
\hullface{keel!63!black}\hullmove{-2.34}{-6.81}\hullline{-1.87}{-7.02}\hullline{-2.05}{-6.81}\hullline{-2.39}{-6.77}\hullfill
\hullface{keel!57!black}\hullmove{2.52}{-6.63}\hullline{2.09}{-6.77}\hullline{2.05}{-6.81}\hullline{2.39}{-6.77}\hullfill
\hullface{keel!63!black}\hullmove{-2.39}{-6.77}\hullline{-2.05}{-6.81}\hullline{-2.09}{-6.77}\hullline{-2.52}{-6.63}\hullfill
\hullface{keel!57!black}\hullmove{2.97}{-6.12}\hullline{2.53}{-6.23}\hullline{2.09}{-6.77}\hullline{2.52}{-6.63}\hullfill
\hullface{keel!66!black}\hullmove{-2.52}{-6.63}\hullline{-2.09}{-6.77}\hullline{-2.53}{-6.23}\hullline{-2.97}{-6.12}\hullfill
\hullface{keel!57!black}\hullmove{3.20}{-5.86}\hullline{2.82}{-5.86}\hullline{2.53}{-6.23}\hullline{2.97}{-6.12}\hullfill
\hullface{keel!69!black}\hullmove{-2.97}{-6.12}\hullline{-2.53}{-6.23}\hullline{-2.82}{-5.86}\hullline{-3.20}{-5.86}\hullfill
\hullface{keel!57!black}\hullmove{3.55}{-5.44}\hullline{3.13}{-5.44}\hullline{2.82}{-5.86}\hullline{3.20}{-5.86}\hullfill
\hullface{keel!69!black}\hullmove{-3.20}{-5.86}\hullline{-2.82}{-5.86}\hullline{-3.13}{-5.44}\hullline{-3.55}{-5.44}\hullfill
\hullface{keel!57!black}\hullmove{4.01}{-4.85}\hullline{3.54}{-4.85}\hullline{3.13}{-5.44}\hullline{3.55}{-5.44}\hullfill
\hullface{keel!72!black}\hullmove{-3.55}{-5.44}\hullline{-3.13}{-5.44}\hullline{-3.54}{-4.85}\hullline{-4.01}{-4.85}\hullfill
\hullface{keel!57!black}\hullmove{4.08}{-4.75}\hullline{3.68}{-4.64}\hullline{3.54}{-4.85}\hullline{4.01}{-4.85}\hullfill
\hullface{keel!75!black}\hullmove{-4.01}{-4.85}\hullline{-3.54}{-4.85}\hullline{-3.68}{-4.64}\hullline{-4.08}{-4.75}\hullfill
\hullface{keel!57!black}\hullmove{4.44}{-4.24}\hullline{4.06}{-4.06}\hullline{3.68}{-4.64}\hullline{4.08}{-4.75}\hullfill
\hullface{keel!78!black}\hullmove{-4.08}{-4.75}\hullline{-3.68}{-4.64}\hullline{-4.06}{-4.06}\hullline{-4.44}{-4.24}\hullfill
\hullface{keel!57!black}\hullmove{4.57}{-4.06}\hullline{4.76}{-3.75}\hullline{4.24}{-3.75}\hullline{4.19}{-3.85}\hullline{4.06}{-4.06}\hullline{4.44}{-4.24}\hullfill
\hullface{keel!78!black}\hullmove{-4.44}{-4.24}\hullline{-4.06}{-4.06}\hullline{-4.19}{-3.85}\hullline{-4.24}{-3.75}\hullline{-4.76}{-3.75}\hullline{-4.57}{-4.06}\hullfill
\hullface{keel!57!black}\hullmove{5.42}{-2.58}\hullline{4.87}{-2.58}\hullline{4.24}{-3.75}\hullline{4.76}{-3.75}\hullfill
\hullface{keel!84!black}\hullmove{-4.76}{-3.75}\hullline{-4.24}{-3.75}\hullline{-4.87}{-2.58}\hullline{-5.42}{-2.58}\hullfill
\hullface{keel!57!black}\hullmove{5.93}{-1.32}\hullline{5.37}{-1.32}\hullline{4.87}{-2.58}\hullline{5.42}{-2.58}\hullfill
\hullface{keel!87!black}\hullmove{-5.42}{-2.58}\hullline{-4.87}{-2.58}\hullline{-5.37}{-1.32}\hullline{-5.93}{-1.32}\hullfill
\hullface{keel!57!black}\hullmove{6.17}{0.00}\hullline{5.63}{0.00}\hullline{5.37}{-1.32}\hullline{5.93}{-1.32}\hullfill
\hullface{keel!96!black}\hullmove{-5.93}{-1.32}\hullline{-5.37}{-1.32}\hullline{-5.63}{0.00}\hullline{-6.17}{0.00}\hullfill
\hullface{topsides!57!black}\hullmove{6.56}{1.52}\hullline{6.04}{1.52}\hullline{5.63}{0.00}\hullline{6.17}{0.00}\hullfill
\hullface{topsides!93!black}\hullmove{-6.17}{0.00}\hullline{-5.63}{0.00}\hullline{-6.04}{1.52}\hullline{-6.56}{1.52}\hullfill
\hullface{topsides!57!black}\hullmove{7.19}{3.04}\hullline{6.72}{3.04}\hullline{6.04}{1.52}\hullline{6.56}{1.52}\hullfill
\hullface{topsides!84!black}\hullmove{-6.56}{1.52}\hullline{-6.04}{1.52}\hullline{-6.72}{3.04}\hullline{-7.19}{3.04}\hullfill
\hullface{topsides!57!black}\hullmove{7.96}{4.55}\hullline{7.56}{4.55}\hullline{6.72}{3.04}\hullline{7.19}{3.04}\hullfill
\hullface{topsides!81!black}\hullmove{-7.19}{3.04}\hullline{-6.72}{3.04}\hullline{-7.56}{4.55}\hullline{-7.96}{4.55}\hullfill
\hullface{topsides!57!black}\hullmove{8.85}{6.07}\hullline{8.52}{6.07}\hullline{7.56}{4.55}\hullline{7.96}{4.55}\hullfill
\hullface{topsides!75!black}\hullmove{-7.96}{4.55}\hullline{-7.56}{4.55}\hullline{-8.52}{6.07}\hullline{-8.85}{6.07}\hullfill
\hullface{topsides!57!black}\hullmove{9.83}{7.59}\hullline{9.57}{7.59}\hullline{8.52}{6.07}\hullline{8.85}{6.07}\hullfill
\hullface{topsides!69!black}\hullmove{-8.85}{6.07}\hullline{-8.52}{6.07}\hullline{-9.57}{7.59}\hullline{-9.83}{7.59}\hullfill
\hullface{keel!57!black}\hullmove{0.28}{-8.77}\hullline{0.24}{-8.77}\hullline{0.00}{-9.06}\hullline{0.00}{-9.06}\hullfill
\hullface{keel!57!black}\hullmove{0.00}{-9.06}\hullline{0.00}{-9.06}\hullline{-0.24}{-8.77}\hullline{-0.28}{-8.77}\hullfill
\hullface{keel!57!black}\hullmove{0.78}{-8.25}\hullline{0.67}{-8.25}\hullline{0.24}{-8.77}\hullline{0.28}{-8.77}\hullfill
\hullface{keel!57!black}\hullmove{-0.28}{-8.77}\hullline{-0.24}{-8.77}\hullline{-0.67}{-8.25}\hullline{-0.78}{-8.25}\hullfill
\hullface{keel!57!black}\hullmove{1.40}{-7.57}\hullline{1.20}{-7.57}\hullline{0.67}{-8.25}\hullline{0.78}{-8.25}\hullfill
\hullface{keel!60!black}\hullmove{-0.78}{-8.25}\hullline{-0.67}{-8.25}\hullline{-1.20}{-7.57}\hullline{-1.40}{-7.57}\hullfill
\hullface{keel!57!black}\hullmove{1.87}{-7.02}\hullline{1.50}{-7.17}\hullline{1.20}{-7.57}\hullline{1.40}{-7.57}\hullfill
\hullface{keel!63!black}\hullmove{-1.40}{-7.57}\hullline{-1.20}{-7.57}\hullline{-1.50}{-7.17}\hullline{-1.87}{-7.02}\hullfill
\hullface{keel!57!black}\hullmove{2.05}{-6.81}\hullline{1.67}{-6.94}\hullline{1.50}{-7.17}\hullline{1.87}{-7.02}\hullfill
\hullface{keel!66!black}\hullmove{-1.87}{-7.02}\hullline{-1.50}{-7.17}\hullline{-1.67}{-6.94}\hullline{-2.05}{-6.81}\hullfill
\hullface{keel!57!black}\hullmove{2.09}{-6.77}\hullline{1.80}{-6.77}\hullline{1.67}{-6.94}\hullline{2.05}{-6.81}\hullfill
\hullface{keel!66!black}\hullmove{-2.05}{-6.81}\hullline{-1.67}{-6.94}\hullline{-1.80}{-6.77}\hullline{-2.09}{-6.77}\hullfill
\hullface{keel!57!black}\hullmove{2.53}{-6.23}\hullline{2.13}{-6.30}\hullline{1.80}{-6.77}\hullline{2.09}{-6.77}\hullfill
\hullface{keel!69!black}\hullmove{-2.09}{-6.77}\hullline{-1.80}{-6.77}\hullline{-2.13}{-6.30}\hullline{-2.53}{-6.23}\hullfill
\hullface{keel!57!black}\hullmove{2.82}{-5.86}\hullline{2.43}{-5.86}\hullline{2.13}{-6.30}\hullline{2.53}{-6.23}\hullfill
\hullface{keel!72!black}\hullmove{-2.53}{-6.23}\hullline{-2.13}{-6.30}\hullline{-2.43}{-5.86}\hullline{-2.82}{-5.86}\hullfill
\hullface{keel!57!black}\hullmove{3.13}{-5.44}\hullline{2.71}{-5.44}\hullline{2.43}{-5.86}\hullline{2.82}{-5.86}\hullfill
\hullface{keel!72!black}\hullmove{-2.82}{-5.86}\hullline{-2.43}{-5.86}\hullline{-2.71}{-5.44}\hullline{-3.13}{-5.44}\hullfill
\hullface{keel!57!black}\hullmove{3.54}{-4.85}\hullline{3.08}{-4.85}\hullline{2.71}{-5.44}\hullline{3.13}{-5.44}\hullfill
\hullface{keel!75!black}\hullmove{-3.13}{-5.44}\hullline{-2.71}{-5.44}\hullline{-3.08}{-4.85}\hullline{-3.54}{-4.85}\hullfill
\hullface{keel!57!black}\hullmove{3.68}{-4.64}\hullline{3.25}{-4.57}\hullline{3.08}{-4.85}\hullline{3.54}{-4.85}\hullfill
\hullface{keel!78!black}\hullmove{-3.54}{-4.85}\hullline{-3.08}{-4.85}\hullline{-3.25}{-4.57}\hullline{-3.68}{-4.64}\hullfill
\hullface{keel!57!black}\hullmove{4.06}{-4.06}\hullline{3.61}{-3.93}\hullline{3.25}{-4.57}\hullline{3.68}{-4.64}\hullfill
\hullface{keel!78!black}\hullmove{-3.68}{-4.64}\hullline{-3.25}{-4.57}\hullline{-3.61}{-3.93}\hullline{-4.06}{-4.06}\hullfill
\hullface{keel!57!black}\hullmove{4.19}{-3.85}\hullline{4.24}{-3.75}\hullline{3.74}{-3.70}\hullline{3.71}{-3.75}\hullline{3.61}{-3.93}\hullline{4.06}{-4.06}\hullfill
\hullface{keel!81!black}\hullmove{-4.06}{-4.06}\hullline{-3.61}{-3.93}\hullline{-3.71}{-3.75}\hullline{-3.74}{-3.70}\hullline{-4.24}{-3.75}\hullline{-4.19}{-3.85}\hullfill
\hullface{keel!57!black}\hullmove{4.87}{-2.58}\hullline{4.29}{-2.58}\hullline{3.74}{-3.70}\hullline{4.24}{-3.75}\hullfill
\hullface{keel!84!black}\hullmove{-4.24}{-3.75}\hullline{-3.74}{-3.70}\hullline{-4.29}{-2.58}\hullline{-4.87}{-2.58}\hullfill
\hullface{keel!57!black}\hullmove{5.37}{-1.32}\hullline{4.77}{-1.32}\hullline{4.29}{-2.58}\hullline{4.87}{-2.58}\hullfill
\hullface{keel!90!black}\hullmove{-4.87}{-2.58}\hullline{-4.29}{-2.58}\hullline{-4.77}{-1.32}\hullline{-5.37}{-1.32}\hullfill
\hullface{keel!57!black}\hullmove{5.63}{0.00}\hullline{5.05}{0.00}\hullline{4.77}{-1.32}\hullline{5.37}{-1.32}\hullfill
\hullface{keel!96!black}\hullmove{-5.37}{-1.32}\hullline{-4.77}{-1.32}\hullline{-5.05}{0.00}\hullline{-5.63}{0.00}\hullfill
\hullface{topsides!57!black}\hullmove{6.04}{1.52}\hullline{5.49}{1.52}\hullline{5.05}{0.00}\hullline{5.63}{0.00}\hullfill
\hullface{topsides!93!black}\hullmove{-5.63}{0.00}\hullline{-5.05}{0.00}\hullline{-5.49}{1.52}\hullline{-6.04}{1.52}\hullfill
\hullface{topsides!57!black}\hullmove{6.72}{3.04}\hullline{6.22}{3.04}\hullline{5.49}{1.52}\hullline{6.04}{1.52}\hullfill
\hullface{topsides!84!black}\hullmove{-6.04}{1.52}\hullline{-5.49}{1.52}\hullline{-6.22}{3.04}\hullline{-6.72}{3.04}\hullfill
\hullface{topsides!57!black}\hullmove{7.56}{4.55}\hullline{7.11}{4.55}\hullline{6.22}{3.04}\hullline{6.72}{3.04}\hullfill
\hullface{topsides!81!black}\hullmove{-6.72}{3.04}\hullline{-6.22}{3.04}\hullline{-7.11}{4.55}\hullline{-7.56}{4.55}\hullfill
\hullface{topsides!57!black}\hullmove{8.52}{6.07}\hullline{8.14}{6.07}\hullline{7.11}{4.55}\hullline{7.56}{4.55}\hullfill
\hullface{topsides!75!black}\hullmove{-7.56}{4.55}\hullline{-7.11}{4.55}\hullline{-8.14}{6.07}\hullline{-8.52}{6.07}\hullfill
\hullface{topsides!57!black}\hullmove{9.57}{7.59}\hullline{9.27}{7.59}\hullline{8.14}{6.07}\hullline{8.52}{6.07}\hullfill
\hullface{topsides!69!black}\hullmove{-8.52}{6.07}\hullline{-8.14}{6.07}\hullline{-9.27}{7.59}\hullline{-9.57}{7.59}\hullfill
\hullface{keel!57!black}\hullmove{0.24}{-8.77}\hullline{0.21}{-8.67}\hullline{0.00}{-8.95}\hullline{0.00}{-9.06}\hullfill
\hullface{keel!60!black}\hullmove{0.00}{-9.06}\hullline{0.00}{-8.95}\hullline{-0.21}{-8.67}\hullline{-0.24}{-8.77}\hullfill
\hullface{keel!57!black}\hullmove{0.67}{-8.25}\hullline{0.57}{-8.15}\hullline{0.21}{-8.67}\hullline{0.24}{-8.77}\hullfill
\hullface{keel!60!black}\hullmove{-0.24}{-8.77}\hullline{-0.21}{-8.67}\hullline{-0.57}{-8.15}\hullline{-0.67}{-8.25}\hullfill
\hullface{keel!57!black}\hullmove{1.20}{-7.57}\hullline{1.03}{-7.48}\hullline{0.57}{-8.15}\hullline{0.67}{-8.25}\hullfill
\hullface{keel!63!black}\hullmove{-0.67}{-8.25}\hullline{-0.57}{-8.15}\hullline{-1.03}{-7.48}\hullline{-1.20}{-7.57}\hullfill
\hullface{keel!57!black}\hullmove{1.50}{-7.17}\hullline{1.17}{-7.27}\hullline{1.03}{-7.48}\hullline{1.20}{-7.57}\hullfill
\hullface{keel!66!black}\hullmove{-1.20}{-7.57}\hullline{-1.03}{-7.48}\hullline{-1.17}{-7.27}\hullline{-1.50}{-7.17}\hullfill
\hullface{keel!57!black}\hullmove{1.67}{-6.94}\hullline{1.33}{-7.03}\hullline{1.17}{-7.27}\hullline{1.50}{-7.17}\hullfill
\hullface{keel!69!black}\hullmove{-1.50}{-7.17}\hullline{-1.17}{-7.27}\hullline{-1.33}{-7.03}\hullline{-1.67}{-6.94}\hullfill
\hullface{keel!57!black}\hullmove{1.80}{-6.77}\hullline{1.55}{-6.69}\hullline{1.33}{-7.03}\hullline{1.67}{-6.94}\hullfill
\hullface{keel!69!black}\hullmove{-1.67}{-6.94}\hullline{-1.33}{-7.03}\hullline{-1.55}{-6.69}\hullline{-1.80}{-6.77}\hullfill
\hullface{keel!57!black}\hullmove{2.13}{-6.30}\hullline{1.76}{-6.35}\hullline{1.55}{-6.69}\hullline{1.80}{-6.77}\hullfill
\hullface{keel!72!black}\hullmove{-1.80}{-6.77}\hullline{-1.55}{-6.69}\hullline{-1.76}{-6.35}\hullline{-2.13}{-6.30}\hullfill
\hullface{keel!57!black}\hullmove{2.43}{-5.86}\hullline{2.10}{-5.79}\hullline{1.76}{-6.35}\hullline{2.13}{-6.30}\hullfill
\hullface{keel!72!black}\hullmove{-2.13}{-6.30}\hullline{-1.76}{-6.35}\hullline{-2.10}{-5.79}\hullline{-2.43}{-5.86}\hullfill
\hullface{keel!57!black}\hullmove{2.71}{-5.44}\hullline{2.31}{-5.44}\hullline{2.10}{-5.79}\hullline{2.43}{-5.86}\hullfill
\hullface{keel!75!black}\hullmove{-2.43}{-5.86}\hullline{-2.10}{-5.79}\hullline{-2.31}{-5.44}\hullline{-2.71}{-5.44}\hullfill
\hullface{keel!57!black}\hullmove{3.08}{-4.85}\hullline{2.67}{-4.79}\hullline{2.31}{-5.44}\hullline{2.71}{-5.44}\hullfill
\hullface{keel!78!black}\hullmove{-2.71}{-5.44}\hullline{-2.31}{-5.44}\hullline{-2.67}{-4.79}\hullline{-3.08}{-4.85}\hullfill
\hullface{keel!57!black}\hullmove{3.25}{-4.57}\hullline{2.82}{-4.52}\hullline{2.67}{-4.79}\hullline{3.08}{-4.85}\hullfill
\hullface{keel!78!black}\hullmove{-3.08}{-4.85}\hullline{-2.67}{-4.79}\hullline{-2.82}{-4.52}\hullline{-3.25}{-4.57}\hullfill
\hullface{keel!57!black}\hullmove{3.61}{-3.93}\hullline{3.16}{-3.85}\hullline{2.82}{-4.52}\hullline{3.25}{-4.57}\hullfill
\hullface{keel!81!black}\hullmove{-3.25}{-4.57}\hullline{-2.82}{-4.52}\hullline{-3.16}{-3.85}\hullline{-3.61}{-3.93}\hullfill
\hullface{keel!57!black}\hullmove{3.71}{-3.75}\hullline{3.22}{-3.71}\hullline{3.16}{-3.85}\hullline{3.61}{-3.93}\hullfill
\hullface{keel!84!black}\hullmove{-3.61}{-3.93}\hullline{-3.16}{-3.85}\hullline{-3.22}{-3.71}\hullline{-3.71}{-3.75}\hullfill
\hullface{keel!57!black}\hullmove{3.74}{-3.70}\hullline{3.27}{-3.60}\hullline{3.22}{-3.71}\hullline{3.71}{-3.75}\hullfill
\hullface{keel!84!black}\hullmove{-3.71}{-3.75}\hullline{-3.22}{-3.71}\hullline{-3.27}{-3.60}\hullline{-3.74}{-3.70}\hullfill
\hullface{keel!57!black}\hullmove{4.29}{-2.58}\hullline{3.74}{-2.55}\hullline{3.27}{-3.60}\hullline{3.74}{-3.70}\hullfill
\hullface{keel!84!black}\hullmove{-3.74}{-3.70}\hullline{-3.27}{-3.60}\hullline{-3.74}{-2.55}\hullline{-4.29}{-2.58}\hullfill
\hullface{keel!57!black}\hullmove{4.77}{-1.32}\hullline{4.17}{-1.31}\hullline{3.74}{-2.55}\hullline{4.29}{-2.58}\hullfill
\hullface{keel!90!black}\hullmove{-4.29}{-2.58}\hullline{-3.74}{-2.55}\hullline{-4.17}{-1.31}\hullline{-4.77}{-1.32}\hullfill
\hullface{keel!57!black}\hullmove{5.05}{0.00}\hullline{4.43}{0.00}\hullline{4.17}{-1.31}\hullline{4.77}{-1.32}\hullfill
\hullface{keel!96!black}\hullmove{-4.77}{-1.32}\hullline{-4.17}{-1.31}\hullline{-4.43}{0.00}\hullline{-5.05}{0.00}\hullfill
\hullface{topsides!57!black}\hullmove{5.49}{1.52}\hullline{4.90}{1.52}\hullline{4.43}{0.00}\hullline{5.05}{0.00}\hullfill
\hullface{topsides!93!black}\hullmove{-5.05}{0.00}\hullline{-4.43}{0.00}\hullline{-4.90}{1.52}\hullline{-5.49}{1.52}\hullfill
\hullface{topsides!57!black}\hullmove{6.22}{3.04}\hullline{5.67}{3.04}\hullline{4.90}{1.52}\hullline{5.49}{1.52}\hullfill
\hullface{topsides!84!black}\hullmove{-5.49}{1.52}\hullline{-4.90}{1.52}\hullline{-5.67}{3.04}\hullline{-6.22}{3.04}\hullfill
\hullface{topsides!57!black}\hullmove{7.11}{4.55}\hullline{6.62}{4.55}\hullline{5.67}{3.04}\hullline{6.22}{3.04}\hullfill
\hullface{topsides!81!black}\hullmove{-6.22}{3.04}\hullline{-5.67}{3.04}\hullline{-6.62}{4.55}\hullline{-7.11}{4.55}\hullfill
\hullface{topsides!57!black}\hullmove{8.14}{6.07}\hullline{7.71}{6.07}\hullline{6.62}{4.55}\hullline{7.11}{4.55}\hullfill
\hullface{topsides!75!black}\hullmove{-7.11}{4.55}\hullline{-6.62}{4.55}\hullline{-7.71}{6.07}\hullline{-8.14}{6.07}\hullfill
\hullface{topsides!57!black}\hullmove{9.27}{7.59}\hullline{8.91}{7.59}\hullline{7.71}{6.07}\hullline{8.14}{6.07}\hullfill
\hullface{topsides!72!black}\hullmove{-8.14}{6.07}\hullline{-7.71}{6.07}\hullline{-8.91}{7.59}\hullline{-9.27}{7.59}\hullfill
\hullface{keel!57!black}\hullmove{0.21}{-8.67}\hullline{0.18}{-8.38}\hullline{0.00}{-8.65}\hullline{0.00}{-8.95}\hullfill
\hullface{keel!66!black}\hullmove{0.00}{-8.95}\hullline{0.00}{-8.65}\hullline{-0.18}{-8.38}\hullline{-0.21}{-8.67}\hullfill
\hullface{keel!57!black}\hullmove{0.57}{-8.15}\hullline{0.50}{-7.88}\hullline{0.18}{-8.38}\hullline{0.21}{-8.67}\hullfill
\hullface{keel!69!black}\hullmove{-0.21}{-8.67}\hullline{-0.18}{-8.38}\hullline{-0.50}{-7.88}\hullline{-0.57}{-8.15}\hullfill
\hullface{keel!57!black}\hullmove{1.03}{-7.48}\hullline{0.84}{-7.33}\hullline{0.50}{-7.88}\hullline{0.57}{-8.15}\hullfill
\hullface{keel!69!black}\hullmove{-0.57}{-8.15}\hullline{-0.50}{-7.88}\hullline{-0.84}{-7.33}\hullline{-1.03}{-7.48}\hullfill
\hullface{keel!57!black}\hullmove{1.55}{-6.69}\hullline{1.36}{-6.46}\hullline{0.99}{-7.08}\hullline{1.33}{-7.03}\hullfill
\hullface{keel!72!black}\hullmove{-1.33}{-7.03}\hullline{-0.99}{-7.08}\hullline{-1.36}{-6.46}\hullline{-1.55}{-6.69}\hullfill
\hullface{keel!57!black}\hullmove{1.76}{-6.35}\hullline{1.40}{-6.38}\hullline{1.36}{-6.46}\hullline{1.55}{-6.69}\hullfill
\hullface{keel!72!black}\hullmove{-1.55}{-6.69}\hullline{-1.36}{-6.46}\hullline{-1.40}{-6.38}\hullline{-1.76}{-6.35}\hullfill
\hullface{keel!57!black}\hullmove{2.10}{-5.79}\hullline{1.83}{-5.59}\hullline{1.40}{-6.38}\hullline{1.76}{-6.35}\hullfill
\hullface{keel!75!black}\hullmove{-1.76}{-6.35}\hullline{-1.40}{-6.38}\hullline{-1.83}{-5.59}\hullline{-2.10}{-5.79}\hullfill
\hullface{keel!57!black}\hullmove{2.31}{-5.44}\hullline{1.92}{-5.44}\hullline{1.83}{-5.59}\hullline{2.10}{-5.79}\hullfill
\hullface{keel!78!black}\hullmove{-2.10}{-5.79}\hullline{-1.83}{-5.59}\hullline{-1.92}{-5.44}\hullline{-2.31}{-5.44}\hullfill
\hullface{keel!57!black}\hullmove{2.67}{-4.79}\hullline{2.32}{-4.63}\hullline{1.92}{-5.44}\hullline{2.31}{-5.44}\hullfill
\hullface{keel!78!black}\hullmove{-2.31}{-5.44}\hullline{-1.92}{-5.44}\hullline{-2.32}{-4.63}\hullline{-2.67}{-4.79}\hullfill
\hullface{keel!57!black}\hullmove{2.82}{-4.52}\hullline{2.39}{-4.49}\hullline{2.32}{-4.63}\hullline{2.67}{-4.79}\hullfill
\hullface{keel!81!black}\hullmove{-2.67}{-4.79}\hullline{-2.32}{-4.63}\hullline{-2.39}{-4.49}\hullline{-2.82}{-4.52}\hullfill
\hullface{keel!57!black}\hullmove{3.16}{-3.85}\hullline{2.70}{-3.79}\hullline{2.39}{-4.49}\hullline{2.82}{-4.52}\hullfill
\hullface{keel!84!black}\hullmove{-2.82}{-4.52}\hullline{-2.39}{-4.49}\hullline{-2.70}{-3.79}\hullline{-3.16}{-3.85}\hullfill
\hullface{keel!57!black}\hullmove{3.74}{-2.55}\hullline{3.21}{-2.46}\hullline{2.81}{-3.54}\hullline{3.27}{-3.60}\hullfill
\hullface{keel!87!black}\hullmove{-3.27}{-3.60}\hullline{-2.81}{-3.54}\hullline{-3.21}{-2.46}\hullline{-3.74}{-2.55}\hullfill
\hullface{keel!57!black}\hullmove{4.17}{-1.31}\hullline{3.57}{-1.26}\hullline{3.21}{-2.46}\hullline{3.74}{-2.55}\hullfill
\hullface{keel!90!black}\hullmove{-3.74}{-2.55}\hullline{-3.21}{-2.46}\hullline{-3.57}{-1.26}\hullline{-4.17}{-1.31}\hullfill
\hullface{keel!60!black}\hullmove{4.43}{0.00}\hullline{3.76}{0.00}\hullline{3.57}{-1.26}\hullline{4.17}{-1.31}\hullfill
\hullface{keel!96!black}\hullmove{-4.17}{-1.31}\hullline{-3.57}{-1.26}\hullline{-3.76}{0.00}\hullline{-4.43}{0.00}\hullfill
\hullface{topsides!57!black}\hullmove{4.90}{1.52}\hullline{4.26}{1.52}\hullline{3.76}{0.00}\hullline{4.43}{0.00}\hullfill
\hullface{topsides!93!black}\hullmove{-4.43}{0.00}\hullline{-3.76}{0.00}\hullline{-4.26}{1.52}\hullline{-4.90}{1.52}\hullfill
\hullface{topsides!57!black}\hullmove{5.67}{3.04}\hullline{5.07}{3.04}\hullline{4.26}{1.52}\hullline{4.90}{1.52}\hullfill
\hullface{topsides!84!black}\hullmove{-4.90}{1.52}\hullline{-4.26}{1.52}\hullline{-5.07}{3.04}\hullline{-5.67}{3.04}\hullfill
\hullface{topsides!57!black}\hullmove{6.62}{4.55}\hullline{6.07}{4.55}\hullline{5.07}{3.04}\hullline{5.67}{3.04}\hullfill
\hullface{topsides!81!black}\hullmove{-5.67}{3.04}\hullline{-5.07}{3.04}\hullline{-6.07}{4.55}\hullline{-6.62}{4.55}\hullfill
\hullface{topsides!57!black}\hullmove{7.71}{6.07}\hullline{7.22}{6.07}\hullline{6.07}{4.55}\hullline{6.62}{4.55}\hullfill
\hullface{topsides!75!black}\hullmove{-6.62}{4.55}\hullline{-6.07}{4.55}\hullline{-7.22}{6.07}\hullline{-7.71}{6.07}\hullfill
\hullface{topsides!57!black}\hullmove{8.91}{7.59}\hullline{8.49}{7.59}\hullline{7.22}{6.07}\hullline{7.71}{6.07}\hullfill
\hullface{topsides!72!black}\hullmove{-7.71}{6.07}\hullline{-7.22}{6.07}\hullline{-8.49}{7.59}\hullline{-8.91}{7.59}\hullfill
\hullface{keel!57!black}\hullmove{1.17}{-7.27}\hullline{1.33}{-7.03}\hullline{0.99}{-7.08}\hullline{0.90}{-7.23}\hullline{0.84}{-7.33}\hullline{1.03}{-7.48}\hullfill
\hullface{keel!72!black}\hullmove{-1.03}{-7.48}\hullline{-0.84}{-7.33}\hullline{-0.90}{-7.23}\hullline{-0.99}{-7.08}\hullline{-1.33}{-7.03}\hullline{-1.17}{-7.27}\hullfill
\hullface{keel!57!black}\hullmove{3.22}{-3.71}\hullline{3.27}{-3.60}\hullline{2.81}{-3.54}\hullline{2.79}{-3.58}\hullline{2.70}{-3.79}\hullline{3.16}{-3.85}\hullfill
\hullface{keel!84!black}\hullmove{-3.16}{-3.85}\hullline{-2.70}{-3.79}\hullline{-2.79}{-3.58}\hullline{-2.81}{-3.54}\hullline{-3.27}{-3.60}\hullline{-3.22}{-3.71}\hullfill
\hullface{keel!57!black}\hullmove{0.18}{-8.38}\hullline{0.50}{-7.88}\hullline{0.45}{-7.42}\hullline{0.16}{-7.89}\hullline{0.00}{-8.15}\hullline{0.00}{-8.65}\hullfill
\hullface{keel!72!black}\hullmove{0.00}{-8.65}\hullline{0.00}{-8.15}\hullline{-0.16}{-7.89}\hullline{-0.45}{-7.42}\hullline{-0.50}{-7.88}\hullline{-0.18}{-8.38}\hullfill
\hullface{keel!57!black}\hullmove{0.84}{-7.33}\hullline{0.48}{-7.36}\hullline{0.45}{-7.42}\hullline{0.50}{-7.88}\hullfill
\hullface{keel!72!black}\hullmove{-0.50}{-7.88}\hullline{-0.45}{-7.42}\hullline{-0.48}{-7.36}\hullline{-0.84}{-7.33}\hullfill
\hullface{keel!57!black}\hullmove{0.90}{-7.23}\hullline{0.77}{-7.10}\hullline{0.48}{-7.36}\hullline{0.84}{-7.33}\hullfill
\hullface{keel!60!black}\hullmove{-0.84}{-7.33}\hullline{-0.48}{-7.36}\hullline{-0.77}{-7.10}\hullline{-0.90}{-7.23}\hullfill
\hullface{keel!57!black}\hullmove{0.99}{-7.08}\hullline{1.08}{-6.81}\hullline{0.77}{-7.10}\hullline{0.90}{-7.23}\hullfill
\hullface{keel!60!black}\hullmove{-0.90}{-7.23}\hullline{-0.77}{-7.10}\hullline{-1.08}{-6.81}\hullline{-0.99}{-7.08}\hullfill
\hullface{keel!57!black}\hullmove{1.36}{-6.46}\hullline{1.34}{-6.40}\hullline{1.08}{-6.81}\hullline{0.99}{-7.08}\hullfill
\hullface{keel!63!black}\hullmove{-0.99}{-7.08}\hullline{-1.08}{-6.81}\hullline{-1.34}{-6.40}\hullline{-1.36}{-6.46}\hullfill
\hullface{keel!57!black}\hullmove{1.40}{-6.38}\hullline{1.45}{-6.09}\hullline{1.34}{-6.40}\hullline{1.36}{-6.46}\hullfill
\hullface{keel!69!black}\hullmove{-1.36}{-6.46}\hullline{-1.34}{-6.40}\hullline{-1.45}{-6.09}\hullline{-1.40}{-6.38}\hullfill
\hullface{keel!57!black}\hullmove{1.83}{-5.59}\hullline{1.54}{-5.44}\hullline{1.45}{-6.09}\hullline{1.40}{-6.38}\hullfill
\hullface{keel!78!black}\hullmove{-1.40}{-6.38}\hullline{-1.45}{-6.09}\hullline{-1.54}{-5.44}\hullline{-1.83}{-5.59}\hullfill
\hullface{keel!57!black}\hullmove{1.92}{-5.44}\hullline{1.60}{-5.27}\hullline{1.54}{-5.44}\hullline{1.83}{-5.59}\hullfill
\hullface{keel!81!black}\hullmove{-1.83}{-5.59}\hullline{-1.54}{-5.44}\hullline{-1.60}{-5.27}\hullline{-1.92}{-5.44}\hullfill
\hullface{keel!57!black}\hullmove{2.32}{-4.63}\hullline{1.95}{-4.47}\hullline{1.60}{-5.27}\hullline{1.92}{-5.44}\hullfill
\hullface{keel!81!black}\hullmove{-1.92}{-5.44}\hullline{-1.60}{-5.27}\hullline{-1.95}{-4.47}\hullline{-2.32}{-4.63}\hullfill
\hullface{keel!57!black}\hullmove{2.39}{-4.49}\hullline{2.00}{-4.36}\hullline{1.95}{-4.47}\hullline{2.32}{-4.63}\hullfill
\hullface{keel!84!black}\hullmove{-2.32}{-4.63}\hullline{-1.95}{-4.47}\hullline{-2.00}{-4.36}\hullline{-2.39}{-4.49}\hullfill
\hullface{keel!57!black}\hullmove{2.70}{-3.79}\hullline{2.23}{-3.77}\hullline{2.00}{-4.36}\hullline{2.39}{-4.49}\hullfill
\hullface{keel!84!black}\hullmove{-2.39}{-4.49}\hullline{-2.00}{-4.36}\hullline{-2.23}{-3.77}\hullline{-2.70}{-3.79}\hullfill
\hullface{keel!57!black}\hullmove{2.79}{-3.58}\hullline{2.33}{-3.51}\hullline{2.23}{-3.77}\hullline{2.70}{-3.79}\hullfill
\hullface{keel!87!black}\hullmove{-2.70}{-3.79}\hullline{-2.23}{-3.77}\hullline{-2.33}{-3.51}\hullline{-2.79}{-3.58}\hullfill
\hullface{keel!57!black}\hullmove{2.81}{-3.54}\hullline{2.37}{-3.38}\hullline{2.33}{-3.51}\hullline{2.79}{-3.58}\hullfill
\hullface{keel!87!black}\hullmove{-2.79}{-3.58}\hullline{-2.33}{-3.51}\hullline{-2.37}{-3.38}\hullline{-2.81}{-3.54}\hullfill
\hullface{keel!57!black}\hullmove{3.21}{-2.46}\hullline{2.70}{-2.32}\hullline{2.37}{-3.38}\hullline{2.81}{-3.54}\hullfill
\hullface{keel!90!black}\hullmove{-2.81}{-3.54}\hullline{-2.37}{-3.38}\hullline{-2.70}{-2.32}\hullline{-3.21}{-2.46}\hullfill
\hullface{keel!57!black}\hullmove{3.57}{-1.26}\hullline{2.94}{-1.19}\hullline{2.70}{-2.32}\hullline{3.21}{-2.46}\hullfill
\hullface{keel!93!black}\hullmove{-3.21}{-2.46}\hullline{-2.70}{-2.32}\hullline{-2.94}{-1.19}\hullline{-3.57}{-1.26}\hullfill
\hullface{keel!63!black}\hullmove{3.76}{0.00}\hullline{3.06}{0.00}\hullline{2.94}{-1.19}\hullline{3.57}{-1.26}\hullfill
\hullface{keel!99!black}\hullmove{-3.57}{-1.26}\hullline{-2.94}{-1.19}\hullline{-3.06}{0.00}\hullline{-3.76}{0.00}\hullfill
\hullface{topsides!57!black}\hullmove{4.26}{1.52}\hullline{3.58}{1.52}\hullline{3.06}{0.00}\hullline{3.76}{0.00}\hullfill
\hullface{topsides!93!black}\hullmove{-3.76}{0.00}\hullline{-3.06}{0.00}\hullline{-3.58}{1.52}\hullline{-4.26}{1.52}\hullfill
\hullface{topsides!57!black}\hullmove{5.07}{3.04}\hullline{4.42}{3.04}\hullline{3.58}{1.52}\hullline{4.26}{1.52}\hullfill
\hullface{topsides!87!black}\hullmove{-4.26}{1.52}\hullline{-3.58}{1.52}\hullline{-4.42}{3.04}\hullline{-5.07}{3.04}\hullfill
\hullface{topsides!57!black}\hullmove{6.07}{4.55}\hullline{5.47}{4.55}\hullline{4.42}{3.04}\hullline{5.07}{3.04}\hullfill
\hullface{topsides!81!black}\hullmove{-5.07}{3.04}\hullline{-4.42}{3.04}\hullline{-5.47}{4.55}\hullline{-6.07}{4.55}\hullfill
\hullface{topsides!57!black}\hullmove{7.22}{6.07}\hullline{6.67}{6.07}\hullline{5.47}{4.55}\hullline{6.07}{4.55}\hullfill
\hullface{topsides!75!black}\hullmove{-6.07}{4.55}\hullline{-5.47}{4.55}\hullline{-6.67}{6.07}\hullline{-7.22}{6.07}\hullfill
\hullface{topsides!57!black}\hullmove{8.49}{7.59}\hullline{7.99}{7.59}\hullline{6.67}{6.07}\hullline{7.22}{6.07}\hullfill
\hullface{topsides!72!black}\hullmove{-7.22}{6.07}\hullline{-6.67}{6.07}\hullline{-7.99}{7.59}\hullline{-8.49}{7.59}\hullfill
\hullface{keel!57!black}\hullmove{0.16}{-7.89}\hullline{0.06}{-7.36}\hullline{0.00}{-7.45}\hullline{0.00}{-8.15}\hullfill
\hullface{keel!78!black}\hullmove{0.00}{-8.15}\hullline{0.00}{-7.45}\hullline{-0.06}{-7.36}\hullline{-0.16}{-7.89}\hullfill
\hullface{keel!57!black}\hullmove{0.45}{-7.42}\hullline{0.57}{-7.22}\hullline{0.06}{-7.36}\hullline{0.16}{-7.89}\hullfill
\hullface{keel!60!black}\hullmove{-0.16}{-7.89}\hullline{-0.06}{-7.36}\hullline{-0.57}{-7.22}\hullline{-0.45}{-7.42}\hullfill
\hullface{keel!57!black}\hullmove{0.48}{-7.36}\hullline{0.77}{-7.10}\hullline{0.57}{-7.22}\hullline{0.45}{-7.42}\hullfill
\hullface{keel!57!black}\hullmove{-0.45}{-7.42}\hullline{-0.57}{-7.22}\hullline{-0.77}{-7.10}\hullline{-0.48}{-7.36}\hullfill
\hullface{keel!57!black}\hullmove{0.77}{-7.10}\hullline{1.09}{-6.79}\hullline{0.77}{-7.10}\hullline{0.48}{-7.36}\hullfill
\hullface{keel!57!black}\hullmove{-0.48}{-7.36}\hullline{-0.77}{-7.10}\hullline{-1.09}{-6.79}\hullline{-0.77}{-7.10}\hullfill
\hullface{keel!57!black}\hullmove{1.08}{-6.81}\hullline{1.33}{-6.40}\hullline{1.09}{-6.79}\hullline{0.77}{-7.10}\hullfill
\hullface{keel!57!black}\hullmove{1.34}{-6.40}\hullline{1.40}{-6.23}\hullline{1.33}{-6.40}\hullline{1.08}{-6.81}\hullfill
\hullface{keel!60!black}\hullmove{-1.08}{-6.81}\hullline{-1.33}{-6.40}\hullline{-1.40}{-6.23}\hullline{-1.34}{-6.40}\hullfill
\hullface{keel!57!black}\hullmove{1.45}{-6.09}\hullline{1.53}{-5.57}\hullline{1.40}{-6.23}\hullline{1.34}{-6.40}\hullfill
\hullface{keel!57!black}\hullmove{1.60}{-5.27}\hullline{1.46}{-4.82}\hullline{1.54}{-5.44}\hullline{1.54}{-5.44}\hullfill
\hullface{keel!84!black}\hullmove{-1.54}{-5.44}\hullline{-1.54}{-5.44}\hullline{-1.46}{-4.82}\hullline{-1.60}{-5.27}\hullfill
\hullface{keel!57!black}\hullmove{1.95}{-4.47}\hullline{1.50}{-4.47}\hullline{1.46}{-4.82}\hullline{1.60}{-5.27}\hullfill
\hullface{keel!84!black}\hullmove{-1.60}{-5.27}\hullline{-1.46}{-4.82}\hullline{-1.50}{-4.47}\hullline{-1.95}{-4.47}\hullfill
\hullface{keel!57!black}\hullmove{2.00}{-4.36}\hullline{1.67}{-3.99}\hullline{1.50}{-4.47}\hullline{1.95}{-4.47}\hullfill
\hullface{keel!87!black}\hullmove{-1.95}{-4.47}\hullline{-1.50}{-4.47}\hullline{-1.67}{-3.99}\hullline{-2.00}{-4.36}\hullfill
\hullface{keel!57!black}\hullmove{2.23}{-3.77}\hullline{1.74}{-3.77}\hullline{1.67}{-3.99}\hullline{2.00}{-4.36}\hullfill
\hullface{keel!87!black}\hullmove{-2.00}{-4.36}\hullline{-1.67}{-3.99}\hullline{-1.74}{-3.77}\hullline{-2.23}{-3.77}\hullfill
\hullface{keel!57!black}\hullmove{2.33}{-3.51}\hullline{1.82}{-3.51}\hullline{1.74}{-3.77}\hullline{2.23}{-3.77}\hullfill
\hullface{keel!90!black}\hullmove{-2.23}{-3.77}\hullline{-1.74}{-3.77}\hullline{-1.82}{-3.51}\hullline{-2.33}{-3.51}\hullfill
\hullface{keel!57!black}\hullmove{2.37}{-3.38}\hullline{1.93}{-3.09}\hullline{1.82}{-3.51}\hullline{2.33}{-3.51}\hullfill
\hullface{keel!90!black}\hullmove{-2.33}{-3.51}\hullline{-1.82}{-3.51}\hullline{-1.93}{-3.09}\hullline{-2.37}{-3.38}\hullfill
\hullface{keel!57!black}\hullmove{2.70}{-2.32}\hullline{2.14}{-2.12}\hullline{1.93}{-3.09}\hullline{2.37}{-3.38}\hullfill
\hullface{keel!93!black}\hullmove{-2.37}{-3.38}\hullline{-1.93}{-3.09}\hullline{-2.14}{-2.12}\hullline{-2.70}{-2.32}\hullfill
\hullface{keel!60!black}\hullmove{2.94}{-1.19}\hullline{2.26}{-1.09}\hullline{2.14}{-2.12}\hullline{2.70}{-2.32}\hullfill
\hullface{keel!96!black}\hullmove{-2.70}{-2.32}\hullline{-2.14}{-2.12}\hullline{-2.26}{-1.09}\hullline{-2.94}{-1.19}\hullfill
\hullface{keel!66!black}\hullmove{3.06}{0.00}\hullline{2.31}{0.00}\hullline{2.26}{-1.09}\hullline{2.94}{-1.19}\hullfill
\hullface{keel!98!white}\hullmove{-2.94}{-1.19}\hullline{-2.26}{-1.09}\hullline{-2.31}{0.00}\hullline{-3.06}{0.00}\hullfill
\hullface{topsides!60!black}\hullmove{3.58}{1.52}\hullline{2.84}{1.52}\hullline{2.31}{0.00}\hullline{3.06}{0.00}\hullfill
\hullface{topsides!93!black}\hullmove{-3.06}{0.00}\hullline{-2.31}{0.00}\hullline{-2.84}{1.52}\hullline{-3.58}{1.52}\hullfill
\hullface{topsides!57!black}\hullmove{4.42}{3.04}\hullline{3.72}{3.04}\hullline{2.84}{1.52}\hullline{3.58}{1.52}\hullfill
\hullface{topsides!87!black}\hullmove{-3.58}{1.52}\hullline{-2.84}{1.52}\hullline{-3.72}{3.04}\hullline{-4.42}{3.04}\hullfill
\hullface{topsides!57!black}\hullmove{5.47}{4.55}\hullline{4.80}{4.55}\hullline{3.72}{3.04}\hullline{4.42}{3.04}\hullfill
\hullface{topsides!81!black}\hullmove{-4.42}{3.04}\hullline{-3.72}{3.04}\hullline{-4.80}{4.55}\hullline{-5.47}{4.55}\hullfill
\hullface{topsides!57!black}\hullmove{6.67}{6.07}\hullline{6.04}{6.07}\hullline{4.80}{4.55}\hullline{5.47}{4.55}\hullfill
\hullface{topsides!78!black}\hullmove{-5.47}{4.55}\hullline{-4.80}{4.55}\hullline{-6.04}{6.07}\hullline{-6.67}{6.07}\hullfill
\hullface{topsides!57!black}\hullmove{7.99}{7.59}\hullline{7.41}{7.59}\hullline{6.04}{6.07}\hullline{6.67}{6.07}\hullfill
\hullface{topsides!75!black}\hullmove{-6.67}{6.07}\hullline{-6.04}{6.07}\hullline{-7.41}{7.59}\hullline{-7.99}{7.59}\hullfill
\hullface{keel!57!black}\hullmove{0.06}{-7.36}\hullline{0.00}{-7.32}\hullline{0.00}{-7.32}\hullline{0.00}{-7.45}\hullfill
\hullface{keel!63!black}\hullmove{0.00}{-7.45}\hullline{-0.00}{-7.32}\hullline{-0.00}{-7.32}\hullline{-0.06}{-7.36}\hullfill
\hullface{keel!57!black}\hullmove{0.57}{-7.22}\hullline{0.72}{-7.09}\hullline{0.00}{-7.32}\hullline{0.06}{-7.36}\hullfill
\hullface{keel!57!black}\hullmove{-0.06}{-7.36}\hullline{-0.00}{-7.32}\hullline{-0.72}{-7.09}\hullline{-0.57}{-7.22}\hullfill
\hullface{keel!57!black}\hullmove{0.77}{-7.10}\hullline{0.75}{-7.07}\hullline{0.72}{-7.09}\hullline{0.57}{-7.22}\hullfill
\hullface{keel!57!black}\hullmove{-0.57}{-7.22}\hullline{-0.72}{-7.09}\hullline{-0.75}{-7.07}\hullline{-0.77}{-7.10}\hullfill
\hullface{keel!57!black}\hullmove{-0.77}{-7.10}\hullline{-0.75}{-7.07}\hullline{-1.14}{-6.67}\hullline{-1.09}{-6.79}\hullfill
\hullface{keel!57!black}\hullmove{1.33}{-6.40}\hullline{1.31}{-6.38}\hullline{1.14}{-6.67}\hullline{1.09}{-6.79}\hullfill
\hullface{keel!60!black}\hullmove{-1.09}{-6.79}\hullline{-1.14}{-6.67}\hullline{-1.31}{-6.38}\hullline{-1.33}{-6.40}\hullfill
\hullface{keel!57!black}\hullmove{1.40}{-6.23}\hullline{1.41}{-6.12}\hullline{1.31}{-6.38}\hullline{1.33}{-6.40}\hullfill
\hullface{keel!57!black}\hullmove{1.53}{-5.57}\hullline{1.51}{-5.47}\hullline{1.41}{-6.12}\hullline{1.40}{-6.23}\hullfill
\hullface{keel!75!black}\hullmove{-1.40}{-6.23}\hullline{-1.41}{-6.12}\hullline{-1.51}{-5.47}\hullline{-1.53}{-5.57}\hullfill
\hullface{keel!81!black}\hullmove{-1.53}{-5.57}\hullline{-1.51}{-5.47}\hullline{-1.51}{-5.44}\hullline{-1.54}{-5.44}\hullfill
\hullface{keel!57!black}\hullmove{1.46}{-4.82}\hullline{1.40}{-4.73}\hullline{1.51}{-5.44}\hullline{1.54}{-5.44}\hullfill
\hullface{keel!87!black}\hullmove{-1.54}{-5.44}\hullline{-1.51}{-5.44}\hullline{-1.40}{-4.73}\hullline{-1.46}{-4.82}\hullfill
\hullface{keel!57!black}\hullmove{1.50}{-4.47}\hullline{1.31}{-4.49}\hullline{1.40}{-4.73}\hullline{1.46}{-4.82}\hullfill
\hullface{keel!90!black}\hullmove{-1.46}{-4.82}\hullline{-1.40}{-4.73}\hullline{-1.31}{-4.49}\hullline{-1.50}{-4.47}\hullfill
\hullface{keel!57!black}\hullmove{1.67}{-3.99}\hullline{1.19}{-3.92}\hullline{1.31}{-4.49}\hullline{1.50}{-4.47}\hullfill
\hullface{keel!93!black}\hullmove{-1.50}{-4.47}\hullline{-1.31}{-4.49}\hullline{-1.19}{-3.92}\hullline{-1.67}{-3.99}\hullfill
\hullface{keel!57!black}\hullmove{1.74}{-3.77}\hullline{1.22}{-3.80}\hullline{1.19}{-3.92}\hullline{1.67}{-3.99}\hullfill
\hullface{keel!90!black}\hullmove{-1.67}{-3.99}\hullline{-1.19}{-3.92}\hullline{-1.22}{-3.80}\hullline{-1.74}{-3.77}\hullfill
\hullface{keel!57!black}\hullmove{1.82}{-3.51}\hullline{1.27}{-3.55}\hullline{1.22}{-3.80}\hullline{1.74}{-3.77}\hullfill
\hullface{keel!93!black}\hullmove{-1.74}{-3.77}\hullline{-1.22}{-3.80}\hullline{-1.27}{-3.55}\hullline{-1.82}{-3.51}\hullfill
\hullface{keel!57!black}\hullmove{1.93}{-3.09}\hullline{1.36}{-3.03}\hullline{1.27}{-3.55}\hullline{1.82}{-3.51}\hullfill
\hullface{keel!93!black}\hullmove{-1.82}{-3.51}\hullline{-1.27}{-3.55}\hullline{-1.36}{-3.03}\hullline{-1.93}{-3.09}\hullfill
\hullface{keel!60!black}\hullmove{2.14}{-2.12}\hullline{1.46}{-2.08}\hullline{1.36}{-3.03}\hullline{1.93}{-3.09}\hullfill
\hullface{keel!96!black}\hullmove{-1.93}{-3.09}\hullline{-1.36}{-3.03}\hullline{-1.46}{-2.08}\hullline{-2.14}{-2.12}\hullfill
\hullface{keel!66!black}\hullmove{2.26}{-1.09}\hullline{1.50}{-1.07}\hullline{1.46}{-2.08}\hullline{2.14}{-2.12}\hullfill
\hullface{keel!98!white}\hullmove{-2.14}{-2.12}\hullline{-1.46}{-2.08}\hullline{-1.50}{-1.07}\hullline{-2.26}{-1.09}\hullfill
\hullface{keel!69!black}\hullmove{2.31}{0.00}\hullline{1.51}{0.00}\hullline{1.50}{-1.07}\hullline{2.26}{-1.09}\hullfill
\hullface{keel!95!white}\hullmove{-2.26}{-1.09}\hullline{-1.50}{-1.07}\hullline{-1.51}{0.00}\hullline{-2.31}{0.00}\hullfill
\hullface{topsides!60!black}\hullmove{2.84}{1.52}\hullline{2.06}{1.52}\hullline{1.51}{0.00}\hullline{2.31}{0.00}\hullfill
\hullface{topsides!93!black}\hullmove{-2.31}{0.00}\hullline{-1.51}{0.00}\hullline{-2.06}{1.52}\hullline{-2.84}{1.52}\hullfill
\hullface{topsides!57!black}\hullmove{3.72}{3.04}\hullline{2.96}{3.04}\hullline{2.06}{1.52}\hullline{2.84}{1.52}\hullfill
\hullface{topsides!87!black}\hullmove{-2.84}{1.52}\hullline{-2.06}{1.52}\hullline{-2.96}{3.04}\hullline{-3.72}{3.04}\hullfill
\hullface{topsides!57!black}\hullmove{4.80}{4.55}\hullline{4.06}{4.55}\hullline{2.96}{3.04}\hullline{3.72}{3.04}\hullfill
\hullface{topsides!81!black}\hullmove{-3.72}{3.04}\hullline{-2.96}{3.04}\hullline{-4.06}{4.55}\hullline{-4.80}{4.55}\hullfill
\hullface{topsides!57!black}\hullmove{6.04}{6.07}\hullline{5.33}{6.07}\hullline{4.06}{4.55}\hullline{4.80}{4.55}\hullfill
\hullface{topsides!78!black}\hullmove{-4.80}{4.55}\hullline{-4.06}{4.55}\hullline{-5.33}{6.07}\hullline{-6.04}{6.07}\hullfill
\hullface{topsides!57!black}\hullmove{7.41}{7.59}\hullline{6.73}{7.59}\hullline{5.33}{6.07}\hullline{6.04}{6.07}\hullfill
\hullface{topsides!75!black}\hullmove{-6.04}{6.07}\hullline{-5.33}{6.07}\hullline{-6.73}{7.59}\hullline{-7.41}{7.59}\hullfill
\hullface{keel!57!black}\hullmove{0.72}{-7.09}\hullline{0.71}{-7.02}\hullline{0.00}{-7.25}\hullline{0.00}{-7.32}\hullfill
\hullface{keel!57!black}\hullmove{-0.00}{-7.32}\hullline{0.00}{-7.25}\hullline{-0.71}{-7.02}\hullline{-0.72}{-7.09}\hullfill
\hullface{keel!57!black}\hullmove{-0.72}{-7.09}\hullline{-0.71}{-7.02}\hullline{-0.72}{-7.00}\hullline{-0.75}{-7.07}\hullfill
\hullface{keel!57!black}\hullmove{1.14}{-6.67}\hullline{1.11}{-6.60}\hullline{0.72}{-7.00}\hullline{0.75}{-7.07}\hullfill
\hullface{keel!57!black}\hullmove{-0.75}{-7.07}\hullline{-0.72}{-7.00}\hullline{-1.11}{-6.60}\hullline{-1.14}{-6.67}\hullfill
\hullface{keel!57!black}\hullmove{1.31}{-6.38}\hullline{1.25}{-6.34}\hullline{1.11}{-6.60}\hullline{1.14}{-6.67}\hullfill
\hullface{keel!63!black}\hullmove{-1.14}{-6.67}\hullline{-1.11}{-6.60}\hullline{-1.25}{-6.34}\hullline{-1.31}{-6.38}\hullfill
\hullface{keel!57!black}\hullmove{1.41}{-6.12}\hullline{1.36}{-6.06}\hullline{1.25}{-6.34}\hullline{1.31}{-6.38}\hullfill
\hullface{keel!69!black}\hullmove{-1.31}{-6.38}\hullline{-1.25}{-6.34}\hullline{-1.36}{-6.06}\hullline{-1.41}{-6.12}\hullfill
\hullface{keel!57!black}\hullmove{1.51}{-5.47}\hullline{1.45}{-5.44}\hullline{1.36}{-6.06}\hullline{1.41}{-6.12}\hullfill
\hullface{keel!78!black}\hullmove{-1.41}{-6.12}\hullline{-1.36}{-6.06}\hullline{-1.45}{-5.44}\hullline{-1.51}{-5.47}\hullfill
\hullface{keel!57!black}\hullmove{1.51}{-5.44}\hullline{1.45}{-5.41}\hullline{1.45}{-5.44}\hullline{1.51}{-5.47}\hullfill
\hullface{keel!84!black}\hullmove{-1.51}{-5.47}\hullline{-1.45}{-5.44}\hullline{-1.45}{-5.41}\hullline{-1.51}{-5.44}\hullfill
\hullface{keel!57!black}\hullmove{1.40}{-4.73}\hullline{1.32}{-4.68}\hullline{1.45}{-5.41}\hullline{1.51}{-5.44}\hullfill
\hullface{keel!90!black}\hullmove{-1.51}{-5.44}\hullline{-1.45}{-5.41}\hullline{-1.32}{-4.68}\hullline{-1.40}{-4.73}\hullfill
\hullface{keel!57!black}\hullmove{1.31}{-4.49}\hullline{1.25}{-4.53}\hullline{1.32}{-4.68}\hullline{1.40}{-4.73}\hullfill
\hullface{keel!93!black}\hullmove{-1.40}{-4.73}\hullline{-1.32}{-4.68}\hullline{-1.25}{-4.53}\hullline{-1.31}{-4.49}\hullfill
\hullface{keel!63!black}\hullmove{1.19}{-3.92}\hullline{0.74}{-3.88}\hullline{1.25}{-4.53}\hullline{1.31}{-4.49}\hullfill
\hullface{keel!95!white}\hullmove{-1.31}{-4.49}\hullline{-1.25}{-4.53}\hullline{-0.74}{-3.88}\hullline{-1.19}{-3.92}\hullfill
\hullface{keel!57!black}\hullmove{1.22}{-3.80}\hullline{0.72}{-3.87}\hullline{0.74}{-3.88}\hullline{1.19}{-3.92}\hullfill
\hullface{keel!96!black}\hullmove{-1.19}{-3.92}\hullline{-0.74}{-3.88}\hullline{-0.72}{-3.87}\hullline{-1.22}{-3.80}\hullfill
\hullface{keel!63!black}\hullmove{1.27}{-3.55}\hullline{0.64}{-3.62}\hullline{0.72}{-3.87}\hullline{1.22}{-3.80}\hullfill
\hullface{keel!95!white}\hullmove{-1.22}{-3.80}\hullline{-0.72}{-3.87}\hullline{-0.64}{-3.62}\hullline{-1.27}{-3.55}\hullfill
\hullface{keel!63!black}\hullmove{1.36}{-3.03}\hullline{0.66}{-3.00}\hullline{0.64}{-3.62}\hullline{1.27}{-3.55}\hullfill
\hullface{keel!99!black}\hullmove{-1.27}{-3.55}\hullline{-0.64}{-3.62}\hullline{-0.66}{-3.00}\hullline{-1.36}{-3.03}\hullfill
\hullface{keel!66!black}\hullmove{1.46}{-2.08}\hullline{0.67}{-2.06}\hullline{0.66}{-3.00}\hullline{1.36}{-3.03}\hullfill
\hullface{keel!95!white}\hullmove{-1.36}{-3.03}\hullline{-0.66}{-3.00}\hullline{-0.67}{-2.06}\hullline{-1.46}{-2.08}\hullfill
\hullface{keel!69!black}\hullmove{1.50}{-1.07}\hullline{0.67}{-1.06}\hullline{0.67}{-2.06}\hullline{1.46}{-2.08}\hullfill
\hullface{keel!92!white}\hullmove{-1.46}{-2.08}\hullline{-0.67}{-2.06}\hullline{-0.67}{-1.06}\hullline{-1.50}{-1.07}\hullfill
\hullface{keel!72!black}\hullmove{1.51}{0.00}\hullline{0.67}{0.00}\hullline{0.67}{-1.06}\hullline{1.50}{-1.07}\hullfill
\hullface{keel!92!white}\hullmove{-1.50}{-1.07}\hullline{-0.67}{-1.06}\hullline{-0.67}{0.00}\hullline{-1.51}{0.00}\hullfill
\hullface{topsides!63!black}\hullmove{2.06}{1.52}\hullline{1.22}{1.52}\hullline{0.67}{0.00}\hullline{1.51}{0.00}\hullfill
\hullface{topsides!96!black}\hullmove{-1.51}{0.00}\hullline{-0.67}{0.00}\hullline{-1.22}{1.52}\hullline{-2.06}{1.52}\hullfill
\hullface{topsides!57!black}\hullmove{2.96}{3.04}\hullline{2.13}{3.04}\hullline{1.22}{1.52}\hullline{2.06}{1.52}\hullfill
\hullface{topsides!87!black}\hullmove{-2.06}{1.52}\hullline{-1.22}{1.52}\hullline{-2.13}{3.04}\hullline{-2.96}{3.04}\hullfill
\hullface{topsides!57!black}\hullmove{4.06}{4.55}\hullline{3.25}{4.55}\hullline{2.13}{3.04}\hullline{2.96}{3.04}\hullfill
\hullface{topsides!84!black}\hullmove{-2.96}{3.04}\hullline{-2.13}{3.04}\hullline{-3.25}{4.55}\hullline{-4.06}{4.55}\hullfill
\hullface{topsides!57!black}\hullmove{5.33}{6.07}\hullline{4.53}{6.07}\hullline{3.25}{4.55}\hullline{4.06}{4.55}\hullfill
\hullface{topsides!81!black}\hullmove{-4.06}{4.55}\hullline{-3.25}{4.55}\hullline{-4.53}{6.07}\hullline{-5.33}{6.07}\hullfill
\hullface{topsides!57!black}\hullmove{6.73}{7.59}\hullline{5.95}{7.59}\hullline{4.53}{6.07}\hullline{5.33}{6.07}\hullfill
\hullface{topsides!78!black}\hullmove{-5.33}{6.07}\hullline{-4.53}{6.07}\hullline{-5.95}{7.59}\hullline{-6.73}{7.59}\hullfill
\hullface{keel!57!black}\hullmove{0.71}{-7.02}\hullline{0.68}{-6.91}\hullline{0.00}{-7.13}\hullline{0.00}{-7.25}\hullfill
\hullface{keel!57!black}\hullmove{0.00}{-7.25}\hullline{-0.00}{-7.13}\hullline{-0.68}{-6.91}\hullline{-0.71}{-7.02}\hullfill
\hullface{keel!57!black}\hullmove{0.72}{-7.00}\hullline{0.68}{-6.90}\hullline{0.68}{-6.91}\hullline{0.71}{-7.02}\hullfill
\hullface{keel!57!black}\hullmove{-0.71}{-7.02}\hullline{-0.68}{-6.91}\hullline{-0.68}{-6.90}\hullline{-0.72}{-7.00}\hullfill
\hullface{keel!57!black}\hullmove{1.11}{-6.60}\hullline{1.06}{-6.49}\hullline{0.68}{-6.90}\hullline{0.72}{-7.00}\hullfill
\hullface{keel!57!black}\hullmove{-0.72}{-7.00}\hullline{-0.68}{-6.90}\hullline{-1.06}{-6.49}\hullline{-1.11}{-6.60}\hullfill
\hullface{keel!57!black}\hullmove{1.25}{-6.34}\hullline{1.18}{-6.28}\hullline{1.06}{-6.49}\hullline{1.11}{-6.60}\hullfill
\hullface{keel!66!black}\hullmove{-1.11}{-6.60}\hullline{-1.06}{-6.49}\hullline{-1.18}{-6.28}\hullline{-1.25}{-6.34}\hullfill
\hullface{keel!57!black}\hullmove{1.36}{-6.06}\hullline{1.29}{-5.96}\hullline{1.18}{-6.28}\hullline{1.25}{-6.34}\hullfill
\hullface{keel!72!black}\hullmove{-1.25}{-6.34}\hullline{-1.18}{-6.28}\hullline{-1.29}{-5.96}\hullline{-1.36}{-6.06}\hullfill
\hullface{keel!57!black}\hullmove{1.45}{-5.44}\hullline{1.36}{-5.44}\hullline{1.29}{-5.96}\hullline{1.36}{-6.06}\hullfill
\hullface{keel!81!black}\hullmove{-1.36}{-6.06}\hullline{-1.29}{-5.96}\hullline{-1.36}{-5.44}\hullline{-1.45}{-5.44}\hullfill
\hullface{keel!57!black}\hullmove{1.45}{-5.41}\hullline{1.35}{-5.33}\hullline{1.36}{-5.44}\hullline{1.45}{-5.44}\hullfill
\hullface{keel!87!black}\hullmove{-1.45}{-5.44}\hullline{-1.36}{-5.44}\hullline{-1.35}{-5.33}\hullline{-1.45}{-5.41}\hullfill
\hullface{keel!57!black}\hullmove{1.32}{-4.68}\hullline{1.19}{-4.61}\hullline{1.35}{-5.33}\hullline{1.45}{-5.41}\hullfill
\hullface{keel!90!black}\hullmove{-1.45}{-5.41}\hullline{-1.35}{-5.33}\hullline{-1.19}{-4.61}\hullline{-1.32}{-4.68}\hullfill
\hullface{keel!57!black}\hullmove{1.25}{-4.53}\hullline{1.18}{-4.59}\hullline{1.19}{-4.61}\hullline{1.32}{-4.68}\hullfill
\hullface{keel!96!black}\hullmove{-1.32}{-4.68}\hullline{-1.19}{-4.61}\hullline{-1.18}{-4.59}\hullline{-1.25}{-4.53}\hullfill
\hullface{keel!63!black}\hullmove{0.74}{-3.88}\hullline{0.68}{-3.97}\hullline{1.18}{-4.59}\hullline{1.25}{-4.53}\hullfill
\hullface{keel!98!white}\hullmove{-1.25}{-4.53}\hullline{-1.18}{-4.59}\hullline{-0.68}{-3.97}\hullline{-0.74}{-3.88}\hullfill
\hullface{keel!78!black}\hullmove{0.72}{-3.87}\hullline{0.41}{-3.82}\hullline{0.68}{-3.97}\hullline{0.74}{-3.88}\hullfill
\hullface{keel!95!white}\hullmove{-0.74}{-3.88}\hullline{-0.68}{-3.97}\hullline{-0.41}{-3.82}\hullline{-0.72}{-3.87}\hullfill
\hullface{keel!84!black}\hullmove{0.64}{-3.62}\hullline{0.00}{-3.74}\hullline{0.41}{-3.82}\hullline{0.72}{-3.87}\hullfill
\hullface{keel!83!white}\hullmove{-0.72}{-3.87}\hullline{-0.41}{-3.82}\hullline{-0.00}{-3.74}\hullline{-0.64}{-3.62}\hullfill
\hullface{keel!63!black}\hullmove{0.66}{-3.00}\hullline{0.67}{-2.06}\hullline{0.00}{-2.03}\hullline{0.00}{-2.95}\hullline{0.00}{-3.74}\hullline{0.64}{-3.62}\hullfill
\hullface{keel!95!white}\hullmove{-0.64}{-3.62}\hullline{-0.00}{-3.74}\hullline{0.00}{-2.95}\hullline{0.00}{-2.03}\hullline{-0.67}{-2.06}\hullline{-0.66}{-3.00}\hullfill
\hullface{keel!66!black}\hullmove{0.67}{-1.06}\hullline{0.67}{0.00}\hullline{0.67}{0.00}\hullline{0.00}{0.38}\hullline{0.00}{0.00}\hullline{0.00}{-1.04}\hullline{0.00}{-2.03}\hullline{0.67}{-2.06}\hullfill
\hullface{keel!95!white}\hullmove{-0.67}{-2.06}\hullline{0.00}{-2.03}\hullline{0.00}{-1.04}\hullline{0.00}{0.00}\hullline{0.00}{0.38}\hullline{-0.67}{0.00}\hullline{-0.67}{0.00}\hullline{-0.67}{-1.06}\hullfill
\hullface{topsides!60!black}\hullmove{1.22}{1.52}\hullline{0.53}{1.83}\hullline{0.00}{0.38}\hullline{0.67}{0.00}\hullfill
\hullface{topsides!93!black}\hullmove{-0.67}{0.00}\hullline{0.00}{0.38}\hullline{-0.53}{1.83}\hullline{-1.22}{1.52}\hullfill
\hullface{topsides!57!black}\hullmove{2.13}{3.04}\hullline{1.40}{3.27}\hullline{0.53}{1.83}\hullline{1.22}{1.52}\hullfill
\hullface{topsides!87!black}\hullmove{-1.22}{1.52}\hullline{-0.53}{1.83}\hullline{-1.40}{3.27}\hullline{-2.13}{3.04}\hullfill
\hullface{topsides!57!black}\hullmove{3.25}{4.55}\hullline{2.47}{4.71}\hullline{1.40}{3.27}\hullline{2.13}{3.04}\hullfill
\hullface{topsides!84!black}\hullmove{-2.13}{3.04}\hullline{-1.40}{3.27}\hullline{-2.47}{4.71}\hullline{-3.25}{4.55}\hullfill
\hullface{topsides!57!black}\hullmove{4.53}{6.07}\hullline{3.70}{6.15}\hullline{2.47}{4.71}\hullline{3.25}{4.55}\hullfill
\hullface{topsides!81!black}\hullmove{-3.25}{4.55}\hullline{-2.47}{4.71}\hullline{-3.70}{6.15}\hullline{-4.53}{6.07}\hullfill
\hullface{topsides!57!black}\hullmove{5.95}{7.59}\hullline{5.06}{7.59}\hullline{3.70}{6.15}\hullline{4.53}{6.07}\hullfill
\hullface{topsides!81!black}\hullmove{-4.53}{6.07}\hullline{-3.70}{6.15}\hullline{-5.06}{7.59}\hullline{-5.95}{7.59}\hullfill
\hullface{keel!57!black}\hullmove{0.68}{-6.91}\hullline{0.61}{-6.76}\hullline{0.00}{-6.97}\hullline{0.00}{-7.13}\hullfill
\hullface{keel!57!black}\hullmove{-0.00}{-7.13}\hullline{0.00}{-6.97}\hullline{-0.61}{-6.76}\hullline{-0.68}{-6.91}\hullfill
\hullface{keel!57!black}\hullmove{0.68}{-6.90}\hullline{0.63}{-6.75}\hullline{0.61}{-6.76}\hullline{0.68}{-6.91}\hullfill
\hullface{keel!57!black}\hullmove{-0.68}{-6.91}\hullline{-0.61}{-6.76}\hullline{-0.63}{-6.75}\hullline{-0.68}{-6.90}\hullfill
\hullface{keel!57!black}\hullmove{1.06}{-6.49}\hullline{0.99}{-6.34}\hullline{0.63}{-6.75}\hullline{0.68}{-6.90}\hullfill
\hullface{keel!60!black}\hullmove{-0.68}{-6.90}\hullline{-0.63}{-6.75}\hullline{-0.99}{-6.34}\hullline{-1.06}{-6.49}\hullfill
\hullface{keel!57!black}\hullmove{1.18}{-6.28}\hullline{1.06}{-6.20}\hullline{0.99}{-6.34}\hullline{1.06}{-6.49}\hullfill
\hullface{keel!69!black}\hullmove{-1.06}{-6.49}\hullline{-0.99}{-6.34}\hullline{-1.06}{-6.20}\hullline{-1.18}{-6.28}\hullfill
\hullface{keel!57!black}\hullmove{1.29}{-5.96}\hullline{1.19}{-5.82}\hullline{1.06}{-6.20}\hullline{1.18}{-6.28}\hullfill
\hullface{keel!75!black}\hullmove{-1.18}{-6.28}\hullline{-1.06}{-6.20}\hullline{-1.19}{-5.82}\hullline{-1.29}{-5.96}\hullfill
\hullface{keel!57!black}\hullmove{1.36}{-5.44}\hullline{1.23}{-5.44}\hullline{1.19}{-5.82}\hullline{1.29}{-5.96}\hullfill
\hullface{keel!84!black}\hullmove{-1.29}{-5.96}\hullline{-1.19}{-5.82}\hullline{-1.23}{-5.44}\hullline{-1.36}{-5.44}\hullfill
\hullface{keel!57!black}\hullmove{1.35}{-5.33}\hullline{1.21}{-5.21}\hullline{1.23}{-5.44}\hullline{1.36}{-5.44}\hullfill
\hullface{keel!87!black}\hullmove{-1.36}{-5.44}\hullline{-1.23}{-5.44}\hullline{-1.21}{-5.21}\hullline{-1.35}{-5.33}\hullfill
\hullface{keel!57!black}\hullmove{1.19}{-4.61}\hullline{1.06}{-4.67}\hullline{1.21}{-5.21}\hullline{1.35}{-5.33}\hullfill
\hullface{keel!93!black}\hullmove{-1.35}{-5.33}\hullline{-1.21}{-5.21}\hullline{-1.06}{-4.67}\hullline{-1.19}{-4.61}\hullfill
\hullface{keel!57!black}\hullmove{1.18}{-4.59}\hullline{0.97}{-4.50}\hullline{1.06}{-4.67}\hullline{1.19}{-4.61}\hullfill
\hullface{keel!99!black}\hullmove{-1.19}{-4.61}\hullline{-1.06}{-4.67}\hullline{-0.97}{-4.50}\hullline{-1.18}{-4.59}\hullfill
\hullface{keel!66!black}\hullmove{0.68}{-3.97}\hullline{0.61}{-4.11}\hullline{0.97}{-4.50}\hullline{1.18}{-4.59}\hullfill
\hullface{keel!98!white}\hullmove{-1.18}{-4.59}\hullline{-0.97}{-4.50}\hullline{-0.61}{-4.11}\hullline{-0.68}{-3.97}\hullfill
\hullface{keel!87!black}\hullmove{0.41}{-3.82}\hullline{0.00}{-3.90}\hullline{0.61}{-4.11}\hullline{0.68}{-3.97}\hullfill
\hullface{keel!92!white}\hullmove{-0.68}{-3.97}\hullline{-0.61}{-4.11}\hullline{0.00}{-3.90}\hullline{-0.41}{-3.82}\hullfill
\hullface{keel!78!black}\hullmove{0.00}{-3.74}\hullline{0.00}{-3.73}\hullline{0.00}{-3.90}\hullline{0.41}{-3.82}\hullfill
\hullface{keel!92!white}\hullmove{-0.41}{-3.82}\hullline{0.00}{-3.90}\hullline{0.00}{-3.73}\hullline{-0.00}{-3.74}\hullfill
\hullface{topsides!57!black}\hullmove{0.53}{1.83}\hullline{0.42}{3.07}\hullline{0.00}{1.94}\hullline{0.00}{0.38}\hullfill
\hullface{topsides!90!black}\hullmove{0.00}{0.38}\hullline{0.00}{1.94}\hullline{-0.42}{3.07}\hullline{-0.53}{1.83}\hullfill
\hullface{topsides!57!black}\hullmove{1.40}{3.27}\hullline{1.12}{4.20}\hullline{0.42}{3.07}\hullline{0.53}{1.83}\hullfill
\hullface{topsides!87!black}\hullmove{-0.53}{1.83}\hullline{-0.42}{3.07}\hullline{-1.12}{4.20}\hullline{-1.40}{3.27}\hullfill
\hullface{topsides!57!black}\hullmove{2.47}{4.71}\hullline{1.97}{5.33}\hullline{1.12}{4.20}\hullline{1.40}{3.27}\hullfill
\hullface{topsides!87!black}\hullmove{-1.40}{3.27}\hullline{-1.12}{4.20}\hullline{-1.97}{5.33}\hullline{-2.47}{4.71}\hullfill
\hullface{topsides!57!black}\hullmove{3.70}{6.15}\hullline{2.95}{6.46}\hullline{1.97}{5.33}\hullline{2.47}{4.71}\hullfill
\hullface{topsides!84!black}\hullmove{-2.47}{4.71}\hullline{-1.97}{5.33}\hullline{-2.95}{6.46}\hullline{-3.70}{6.15}\hullfill
\hullface{topsides!57!black}\hullmove{5.06}{7.59}\hullline{4.03}{7.59}\hullline{2.95}{6.46}\hullline{3.70}{6.15}\hullfill
\hullface{topsides!81!black}\hullmove{-3.70}{6.15}\hullline{-2.95}{6.46}\hullline{-4.03}{7.59}\hullline{-5.06}{7.59}\hullfill
\hullface{keel!57!black}\hullmove{0.61}{-6.76}\hullline{0.52}{-6.56}\hullline{0.00}{-6.74}\hullline{0.00}{-6.97}\hullfill
\hullface{keel!57!black}\hullmove{0.00}{-6.97}\hullline{0.00}{-6.74}\hullline{-0.52}{-6.56}\hullline{-0.61}{-6.76}\hullfill
\hullface{keel!57!black}\hullmove{0.63}{-6.75}\hullline{0.57}{-6.52}\hullline{0.52}{-6.56}\hullline{0.61}{-6.76}\hullfill
\hullface{keel!57!black}\hullmove{-0.61}{-6.76}\hullline{-0.52}{-6.56}\hullline{-0.57}{-6.52}\hullline{-0.63}{-6.75}\hullfill
\hullface{keel!57!black}\hullmove{0.99}{-6.34}\hullline{0.88}{-6.13}\hullline{0.57}{-6.52}\hullline{0.63}{-6.75}\hullfill
\hullface{keel!63!black}\hullmove{-0.63}{-6.75}\hullline{-0.57}{-6.52}\hullline{-0.88}{-6.13}\hullline{-0.99}{-6.34}\hullfill
\hullface{keel!57!black}\hullmove{1.06}{-6.20}\hullline{0.90}{-6.09}\hullline{0.88}{-6.13}\hullline{0.99}{-6.34}\hullfill
\hullface{keel!72!black}\hullmove{-0.99}{-6.34}\hullline{-0.88}{-6.13}\hullline{-0.90}{-6.09}\hullline{-1.06}{-6.20}\hullfill
\hullface{keel!57!black}\hullmove{1.19}{-5.82}\hullline{1.03}{-5.63}\hullline{0.90}{-6.09}\hullline{1.06}{-6.20}\hullfill
\hullface{keel!78!black}\hullmove{-1.06}{-6.20}\hullline{-0.90}{-6.09}\hullline{-1.03}{-5.63}\hullline{-1.19}{-5.82}\hullfill
\hullface{keel!57!black}\hullmove{1.23}{-5.44}\hullline{1.04}{-5.44}\hullline{1.03}{-5.63}\hullline{1.19}{-5.82}\hullfill
\hullface{keel!87!black}\hullmove{-1.19}{-5.82}\hullline{-1.03}{-5.63}\hullline{-1.04}{-5.44}\hullline{-1.23}{-5.44}\hullfill
\hullface{keel!57!black}\hullmove{1.21}{-5.21}\hullline{0.99}{-5.03}\hullline{1.04}{-5.44}\hullline{1.23}{-5.44}\hullfill
\hullface{keel!93!black}\hullmove{-1.23}{-5.44}\hullline{-1.04}{-5.44}\hullline{-0.99}{-5.03}\hullline{-1.21}{-5.21}\hullfill
\hullface{keel!57!black}\hullmove{1.06}{-4.67}\hullline{0.90}{-4.78}\hullline{0.99}{-5.03}\hullline{1.21}{-5.21}\hullfill
\hullface{keel!96!black}\hullmove{-1.21}{-5.21}\hullline{-0.99}{-5.03}\hullline{-0.90}{-4.78}\hullline{-1.06}{-4.67}\hullfill
\hullface{keel!66!black}\hullmove{0.97}{-4.50}\hullline{0.58}{-4.35}\hullline{0.90}{-4.78}\hullline{1.06}{-4.67}\hullfill
\hullface{keel!92!white}\hullmove{-1.06}{-4.67}\hullline{-0.90}{-4.78}\hullline{-0.58}{-4.35}\hullline{-0.97}{-4.50}\hullfill
\hullface{keel!72!black}\hullmove{0.61}{-4.11}\hullline{0.52}{-4.31}\hullline{0.58}{-4.35}\hullline{0.97}{-4.50}\hullfill
\hullface{keel!89!white}\hullmove{-0.97}{-4.50}\hullline{-0.58}{-4.35}\hullline{-0.52}{-4.31}\hullline{-0.61}{-4.11}\hullfill
\hullface{keel!90!black}\hullmove{0.00}{-3.90}\hullline{0.00}{-4.13}\hullline{0.52}{-4.31}\hullline{0.61}{-4.11}\hullfill
\hullface{keel!92!white}\hullmove{-0.61}{-4.11}\hullline{-0.52}{-4.31}\hullline{0.00}{-4.13}\hullline{0.00}{-3.90}\hullfill
\hullface{topsides!57!black}\hullmove{0.42}{3.07}\hullline{0.30}{4.31}\hullline{0.00}{3.49}\hullline{0.00}{1.94}\hullfill
\hullface{topsides!90!black}\hullmove{0.00}{1.94}\hullline{0.00}{3.49}\hullline{-0.30}{4.31}\hullline{-0.42}{3.07}\hullfill
\hullface{topsides!57!black}\hullmove{1.12}{4.20}\hullline{0.79}{5.13}\hullline{0.30}{4.31}\hullline{0.42}{3.07}\hullfill
\hullface{topsides!87!black}\hullmove{-0.42}{3.07}\hullline{-0.30}{4.31}\hullline{-0.79}{5.13}\hullline{-1.12}{4.20}\hullfill
\hullface{topsides!60!black}\hullmove{1.97}{5.33}\hullline{1.40}{5.95}\hullline{0.79}{5.13}\hullline{1.12}{4.20}\hullfill
\hullface{topsides!87!black}\hullmove{-1.12}{4.20}\hullline{-0.79}{5.13}\hullline{-1.40}{5.95}\hullline{-1.97}{5.33}\hullfill
\hullface{topsides!60!black}\hullmove{2.95}{6.46}\hullline{2.09}{6.77}\hullline{1.40}{5.95}\hullline{1.97}{5.33}\hullfill
\hullface{topsides!84!black}\hullmove{-1.97}{5.33}\hullline{-1.40}{5.95}\hullline{-2.09}{6.77}\hullline{-2.95}{6.46}\hullfill
\hullface{topsides!60!black}\hullmove{4.03}{7.59}\hullline{2.85}{7.59}\hullline{2.09}{6.77}\hullline{2.95}{6.46}\hullfill
\hullface{topsides!84!black}\hullmove{-2.95}{6.46}\hullline{-2.09}{6.77}\hullline{-2.85}{7.59}\hullline{-4.03}{7.59}\hullfill
\hullface{keel!57!black}\hullmove{0.52}{-6.56}\hullline{0.38}{-6.26}\hullline{0.00}{-6.39}\hullline{0.00}{-6.74}\hullfill
\hullface{keel!57!black}\hullmove{0.00}{-6.74}\hullline{-0.00}{-6.39}\hullline{-0.38}{-6.26}\hullline{-0.52}{-6.56}\hullfill
\hullface{keel!57!black}\hullmove{0.57}{-6.52}\hullline{0.47}{-6.19}\hullline{0.38}{-6.26}\hullline{0.52}{-6.56}\hullfill
\hullface{keel!63!black}\hullmove{-0.52}{-6.56}\hullline{-0.38}{-6.26}\hullline{-0.47}{-6.19}\hullline{-0.57}{-6.52}\hullfill
\hullface{keel!57!black}\hullmove{0.88}{-6.13}\hullline{0.66}{-5.91}\hullline{0.47}{-6.19}\hullline{0.57}{-6.52}\hullfill
\hullface{keel!69!black}\hullmove{-0.57}{-6.52}\hullline{-0.47}{-6.19}\hullline{-0.66}{-5.91}\hullline{-0.88}{-6.13}\hullfill
\hullface{keel!57!black}\hullmove{0.90}{-6.09}\hullline{0.70}{-5.82}\hullline{0.66}{-5.91}\hullline{0.88}{-6.13}\hullfill
\hullface{keel!78!black}\hullmove{-0.88}{-6.13}\hullline{-0.66}{-5.91}\hullline{-0.70}{-5.82}\hullline{-0.90}{-6.09}\hullfill
\hullface{keel!57!black}\hullmove{1.03}{-5.63}\hullline{0.76}{-5.44}\hullline{0.70}{-5.82}\hullline{0.90}{-6.09}\hullfill
\hullface{keel!84!black}\hullmove{-0.90}{-6.09}\hullline{-0.70}{-5.82}\hullline{-0.76}{-5.44}\hullline{-1.03}{-5.63}\hullfill
\hullface{keel!57!black}\hullmove{1.04}{-5.44}\hullline{0.76}{-5.34}\hullline{0.76}{-5.44}\hullline{1.03}{-5.63}\hullfill
\hullface{keel!93!black}\hullmove{-1.03}{-5.63}\hullline{-0.76}{-5.44}\hullline{-0.76}{-5.34}\hullline{-1.04}{-5.44}\hullfill
\hullface{keel!57!black}\hullmove{0.99}{-5.03}\hullline{0.66}{-4.96}\hullline{0.76}{-5.34}\hullline{1.04}{-5.44}\hullfill
\hullface{keel!99!black}\hullmove{-1.04}{-5.44}\hullline{-0.76}{-5.34}\hullline{-0.66}{-4.96}\hullline{-0.99}{-5.03}\hullfill
\hullface{keel!63!black}\hullmove{0.90}{-4.78}\hullline{0.55}{-4.77}\hullline{0.66}{-4.96}\hullline{0.99}{-5.03}\hullfill
\hullface{keel!92!white}\hullmove{-0.99}{-5.03}\hullline{-0.66}{-4.96}\hullline{-0.55}{-4.77}\hullline{-0.90}{-4.78}\hullfill
\hullface{keel!72!black}\hullmove{0.58}{-4.35}\hullline{0.38}{-4.61}\hullline{0.55}{-4.77}\hullline{0.90}{-4.78}\hullfill
\hullface{keel!86!white}\hullmove{-0.90}{-4.78}\hullline{-0.55}{-4.77}\hullline{-0.38}{-4.61}\hullline{-0.58}{-4.35}\hullfill
\hullface{keel!93!black}\hullmove{0.52}{-4.31}\hullline{0.00}{-4.48}\hullline{0.38}{-4.61}\hullline{0.58}{-4.35}\hullfill
\hullface{keel!83!white}\hullmove{-0.58}{-4.35}\hullline{-0.38}{-4.61}\hullline{-0.00}{-4.48}\hullline{-0.52}{-4.31}\hullfill
\hullface{keel!75!black}\hullmove{0.00}{-4.13}\hullline{0.00}{-4.13}\hullline{0.00}{-4.48}\hullline{0.52}{-4.31}\hullfill
\hullface{keel!86!white}\hullmove{-0.52}{-4.31}\hullline{-0.00}{-4.48}\hullline{0.00}{-4.13}\hullline{0.00}{-4.13}\hullfill
\hullface{topsides!57!black}\hullmove{0.30}{4.31}\hullline{0.16}{5.55}\hullline{0.00}{5.04}\hullline{0.00}{3.49}\hullfill
\hullface{topsides!90!black}\hullmove{0.00}{3.49}\hullline{0.00}{5.04}\hullline{-0.16}{5.55}\hullline{-0.30}{4.31}\hullfill
\hullface{topsides!57!black}\hullmove{0.79}{5.13}\hullline{0.42}{6.06}\hullline{0.16}{5.55}\hullline{0.30}{4.31}\hullfill
\hullface{topsides!90!black}\hullmove{-0.30}{4.31}\hullline{-0.16}{5.55}\hullline{-0.42}{6.06}\hullline{-0.79}{5.13}\hullfill
\hullface{topsides!60!black}\hullmove{1.40}{5.95}\hullline{0.74}{6.57}\hullline{0.42}{6.06}\hullline{0.79}{5.13}\hullfill
\hullface{topsides!87!black}\hullmove{-0.79}{5.13}\hullline{-0.42}{6.06}\hullline{-0.74}{6.57}\hullline{-1.40}{5.95}\hullfill
\hullface{topsides!63!black}\hullmove{2.09}{6.77}\hullline{1.11}{7.08}\hullline{0.74}{6.57}\hullline{1.40}{5.95}\hullfill
\hullface{topsides!87!black}\hullmove{-1.40}{5.95}\hullline{-0.74}{6.57}\hullline{-1.11}{7.08}\hullline{-2.09}{6.77}\hullfill
\hullface{topsides!63!black}\hullmove{2.85}{7.59}\hullline{1.52}{7.59}\hullline{1.11}{7.08}\hullline{2.09}{6.77}\hullfill
\hullface{topsides!87!black}\hullmove{-2.09}{6.77}\hullline{-1.11}{7.08}\hullline{-1.52}{7.59}\hullline{-2.85}{7.59}\hullfill
\hullface{keel!63!black}\hullmove{0.38}{-6.26}\hullline{0.00}{-5.44}\hullline{0.00}{-5.44}\hullline{0.00}{-6.39}\hullfill
\hullface{keel!72!black}\hullmove{-0.00}{-6.39}\hullline{0.00}{-5.44}\hullline{-0.00}{-5.44}\hullline{-0.38}{-6.26}\hullfill
\hullface{keel!60!black}\hullmove{0.47}{-6.19}\hullline{0.00}{-5.44}\hullline{0.00}{-5.44}\hullline{0.38}{-6.26}\hullfill
\hullface{keel!81!black}\hullmove{-0.38}{-6.26}\hullline{-0.00}{-5.44}\hullline{-0.00}{-5.44}\hullline{-0.47}{-6.19}\hullfill
\hullface{keel!60!black}\hullmove{0.66}{-5.91}\hullline{0.00}{-5.44}\hullline{0.00}{-5.44}\hullline{0.47}{-6.19}\hullfill
\hullface{keel!87!black}\hullmove{-0.47}{-6.19}\hullline{-0.00}{-5.44}\hullline{-0.00}{-5.44}\hullline{-0.66}{-5.91}\hullfill
\hullface{keel!60!black}\hullmove{0.70}{-5.82}\hullline{0.00}{-5.44}\hullline{0.00}{-5.44}\hullline{0.66}{-5.91}\hullfill
\hullface{keel!93!black}\hullmove{-0.66}{-5.91}\hullline{-0.00}{-5.44}\hullline{-0.00}{-5.44}\hullline{-0.70}{-5.82}\hullfill
\hullface{keel!63!black}\hullmove{0.76}{-5.44}\hullline{0.00}{-5.44}\hullline{0.00}{-5.44}\hullline{0.70}{-5.82}\hullfill
\hullface{keel!99!black}\hullmove{-0.70}{-5.82}\hullline{-0.00}{-5.44}\hullline{-0.00}{-5.44}\hullline{-0.76}{-5.44}\hullfill
\hullface{keel!69!black}\hullmove{0.76}{-5.34}\hullline{0.00}{-5.44}\hullline{0.00}{-5.44}\hullline{0.76}{-5.44}\hullfill
\hullface{keel!92!white}\hullmove{-0.76}{-5.44}\hullline{-0.00}{-5.44}\hullline{0.00}{-5.44}\hullline{-0.76}{-5.34}\hullfill
\hullface{keel!72!black}\hullmove{0.66}{-4.96}\hullline{0.00}{-5.26}\hullline{0.00}{-5.44}\hullline{0.76}{-5.34}\hullfill
\hullface{keel!86!white}\hullmove{-0.76}{-5.34}\hullline{0.00}{-5.44}\hullline{0.00}{-5.26}\hullline{-0.66}{-4.96}\hullfill
\hullface{keel!69!black}\hullmove{0.55}{-4.77}\hullline{0.00}{-4.95}\hullline{0.00}{-5.26}\hullline{0.66}{-4.96}\hullfill
\hullface{keel!86!white}\hullmove{-0.66}{-4.96}\hullline{0.00}{-5.26}\hullline{0.00}{-4.95}\hullline{-0.55}{-4.77}\hullfill
\hullface{keel!66!black}\hullmove{0.38}{-4.61}\hullline{0.00}{-4.54}\hullline{0.00}{-4.95}\hullline{0.55}{-4.77}\hullfill
\hullface{keel!89!white}\hullmove{-0.55}{-4.77}\hullline{0.00}{-4.95}\hullline{0.00}{-4.54}\hullline{-0.38}{-4.61}\hullfill
\hullface{keel!57!black}\hullmove{0.00}{-4.48}\hullline{0.00}{-4.06}\hullline{0.00}{-4.54}\hullline{0.38}{-4.61}\hullfill
\hullface{keel!96!black}\hullmove{-0.38}{-4.61}\hullline{0.00}{-4.54}\hullline{0.00}{-4.06}\hullline{-0.00}{-4.48}\hullfill
\hullface{topsides!57!black}\hullmove{0.16}{5.55}\hullline{0.00}{6.79}\hullline{0.00}{6.59}\hullline{0.00}{5.04}\hullfill
\hullface{topsides!90!black}\hullmove{0.00}{5.04}\hullline{0.00}{6.59}\hullline{0.00}{6.79}\hullline{-0.16}{5.55}\hullfill
\hullface{topsides!57!black}\hullmove{0.42}{6.06}\hullline{0.00}{6.99}\hullline{0.00}{6.79}\hullline{0.16}{5.55}\hullfill
\hullface{topsides!93!black}\hullmove{-0.16}{5.55}\hullline{0.00}{6.79}\hullline{0.00}{6.99}\hullline{-0.42}{6.06}\hullfill
\hullface{topsides!63!black}\hullmove{0.74}{6.57}\hullline{0.00}{7.19}\hullline{0.00}{6.99}\hullline{0.42}{6.06}\hullfill
\hullface{topsides!93!black}\hullmove{-0.42}{6.06}\hullline{0.00}{6.99}\hullline{0.00}{7.19}\hullline{-0.74}{6.57}\hullfill
\hullface{topsides!69!black}\hullmove{1.11}{7.08}\hullline{0.00}{7.39}\hullline{0.00}{7.19}\hullline{0.74}{6.57}\hullfill
\hullface{topsides!93!black}\hullmove{-0.74}{6.57}\hullline{0.00}{7.19}\hullline{0.00}{7.39}\hullline{-1.11}{7.08}\hullfill
\hullface{topsides!72!black}\hullmove{1.52}{7.59}\hullline{0.00}{7.59}\hullline{0.00}{7.39}\hullline{1.11}{7.08}\hullfill
\hullface{topsides!96!black}\hullmove{-1.11}{7.08}\hullline{0.00}{7.39}\hullline{0.00}{7.59}\hullline{-1.52}{7.59}\hullfill
\end{scope}
    \begin{scope}[hull face]\colorlet{topsides}{keel}\colorlet{deck}{keel}
      \colorlet{ship}{keel}\clip[rotate=-#1] (-40,-30) rectangle (40,0);
      \hullface{ship!99!black}\hullmove{-9.09}{2.90}\hullline{9.09}{2.90}\hullline{9.09}{16.57}\hullline{-9.09}{16.57}\hullfill
\hullface{ship!99!black}\hullmove{-10.87}{16.57}\hullline{10.87}{16.57}\hullline{10.87}{19.37}\hullline{-10.87}{19.37}\hullfill
\hullface{keel!57!black}\hullmove{10.87}{-4.85}\hullline{10.87}{-3.75}\hullline{10.87}{-2.58}\hullline{10.87}{-1.32}\hullline{10.87}{0.00}\hullline{10.87}{0.00}\hullline{10.87}{-1.32}\hullline{10.87}{-2.58}\hullline{10.87}{-3.75}\hullline{10.87}{-4.85}\hullline{10.87}{-5.44}\hullline{10.87}{-5.44}\hullfill
\hullface{topsides!57!black}\hullmove{10.87}{0.78}\hullline{10.87}{1.56}\hullline{10.87}{2.34}\hullline{10.87}{3.12}\hullline{10.87}{3.90}\hullline{10.87}{3.90}\hullline{10.87}{3.12}\hullline{10.87}{2.34}\hullline{10.87}{1.56}\hullline{10.87}{0.78}\hullline{10.87}{0.00}\hullline{10.87}{0.00}\hullfill
\hullface{keel!57!black}\hullmove{10.42}{-8.77}\hullline{9.69}{-8.77}\hullline{0.00}{-9.06}\hullline{0.00}{-9.06}\hullfill
\hullface{keel!57!black}\hullmove{0.00}{-9.06}\hullline{0.00}{-9.06}\hullline{-9.69}{-8.77}\hullline{-10.42}{-8.77}\hullfill
\hullface{keel!57!black}\hullmove{10.87}{-8.25}\hullline{10.85}{-8.25}\hullline{9.69}{-8.77}\hullline{10.42}{-8.77}\hullfill
\hullface{keel!57!black}\hullmove{-10.42}{-8.77}\hullline{-9.69}{-8.77}\hullline{-10.85}{-8.25}\hullline{-10.87}{-8.25}\hullfill
\hullface{keel!57!black}\hullmove{10.87}{-7.57}\hullline{10.87}{-6.77}\hullline{10.87}{-5.86}\hullline{10.87}{-5.44}\hullline{10.87}{-5.44}\hullline{10.87}{-5.86}\hullline{10.87}{-6.77}\hullline{10.87}{-7.57}\hullline{10.85}{-8.25}\hullline{10.87}{-8.25}\hullfill
\hullface{keel!57!black}\hullmove{9.69}{-8.77}\hullline{8.24}{-8.77}\hullline{0.00}{-9.06}\hullline{0.00}{-9.06}\hullfill
\hullface{keel!57!black}\hullmove{0.00}{-9.06}\hullline{0.00}{-9.06}\hullline{-8.24}{-8.77}\hullline{-9.69}{-8.77}\hullfill
\hullface{keel!57!black}\hullmove{10.85}{-8.25}\hullline{10.69}{-8.25}\hullline{8.24}{-8.77}\hullline{9.69}{-8.77}\hullfill
\hullface{keel!57!black}\hullmove{-9.69}{-8.77}\hullline{-8.24}{-8.77}\hullline{-10.69}{-8.25}\hullline{-10.85}{-8.25}\hullfill
\hullface{keel!57!black}\hullmove{10.87}{-7.57}\hullline{10.86}{-7.57}\hullline{10.69}{-8.25}\hullline{10.85}{-8.25}\hullfill
\hullface{keel!81!black}\hullmove{-10.85}{-8.25}\hullline{-10.69}{-8.25}\hullline{-10.86}{-7.57}\hullline{-10.87}{-7.57}\hullfill
\hullface{keel!57!black}\hullmove{10.87}{-6.77}\hullline{10.87}{-5.86}\hullline{10.87}{-5.44}\hullline{10.87}{-5.44}\hullline{10.87}{-5.86}\hullline{10.87}{-6.77}\hullline{10.86}{-7.57}\hullline{10.87}{-7.57}\hullfill
\hullface{keel!57!black}\hullmove{10.87}{-4.85}\hullline{10.87}{-3.75}\hullline{10.87}{-2.58}\hullline{10.87}{-1.32}\hullline{10.87}{0.00}\hullline{10.87}{0.00}\hullline{10.87}{-1.32}\hullline{10.87}{-2.58}\hullline{10.87}{-3.75}\hullline{10.87}{-4.85}\hullline{10.87}{-5.44}\hullline{10.87}{-5.44}\hullfill
\hullface{topsides!57!black}\hullmove{10.87}{0.78}\hullline{10.87}{1.56}\hullline{10.87}{2.34}\hullline{10.87}{3.12}\hullline{10.87}{3.90}\hullline{10.87}{3.90}\hullline{10.87}{3.12}\hullline{10.87}{2.34}\hullline{10.87}{1.56}\hullline{10.87}{0.78}\hullline{10.87}{0.00}\hullline{10.87}{0.00}\hullfill
\hullface{keel!57!black}\hullmove{8.24}{-8.77}\hullline{6.72}{-8.77}\hullline{0.00}{-9.06}\hullline{0.00}{-9.06}\hullfill
\hullface{keel!57!black}\hullmove{0.00}{-9.06}\hullline{0.00}{-9.06}\hullline{-6.72}{-8.77}\hullline{-8.24}{-8.77}\hullfill
\hullface{keel!57!black}\hullmove{10.69}{-8.25}\hullline{10.21}{-8.25}\hullline{6.72}{-8.77}\hullline{8.24}{-8.77}\hullfill
\hullface{keel!57!black}\hullmove{-8.24}{-8.77}\hullline{-6.72}{-8.77}\hullline{-10.21}{-8.25}\hullline{-10.69}{-8.25}\hullfill
\hullface{keel!57!black}\hullmove{10.86}{-7.57}\hullline{10.81}{-7.57}\hullline{10.21}{-8.25}\hullline{10.69}{-8.25}\hullfill
\hullface{keel!72!black}\hullmove{-10.69}{-8.25}\hullline{-10.21}{-8.25}\hullline{-10.81}{-7.57}\hullline{-10.86}{-7.57}\hullfill
\hullface{keel!81!black}\hullmove{-10.86}{-7.57}\hullline{-10.81}{-7.57}\hullline{-10.86}{-6.77}\hullline{-10.87}{-6.77}\hullfill
\hullface{keel!57!black}\hullmove{10.87}{-5.86}\hullline{10.87}{-5.44}\hullline{10.86}{-5.44}\hullline{10.86}{-5.86}\hullline{10.86}{-6.77}\hullline{10.87}{-6.77}\hullfill
\hullface{topsides!57!black}\hullmove{10.87}{0.78}\hullline{10.87}{1.56}\hullline{10.87}{2.34}\hullline{10.87}{3.12}\hullline{10.87}{3.90}\hullline{10.87}{3.90}\hullline{10.87}{3.12}\hullline{10.87}{2.34}\hullline{10.86}{1.56}\hullline{10.86}{0.78}\hullline{10.86}{0.00}\hullline{10.87}{0.00}\hullfill
\hullface{keel!57!black}\hullmove{6.72}{-8.77}\hullline{5.45}{-8.77}\hullline{0.00}{-9.06}\hullline{0.00}{-9.06}\hullfill
\hullface{keel!57!black}\hullmove{0.00}{-9.06}\hullline{0.00}{-9.06}\hullline{-5.45}{-8.77}\hullline{-6.72}{-8.77}\hullfill
\hullface{keel!57!black}\hullmove{10.21}{-8.25}\hullline{9.43}{-8.25}\hullline{5.45}{-8.77}\hullline{6.72}{-8.77}\hullfill
\hullface{keel!57!black}\hullmove{-6.72}{-8.77}\hullline{-5.45}{-8.77}\hullline{-9.43}{-8.25}\hullline{-10.21}{-8.25}\hullfill
\hullface{keel!57!black}\hullmove{10.81}{-7.57}\hullline{10.63}{-7.57}\hullline{9.43}{-8.25}\hullline{10.21}{-8.25}\hullfill
\hullface{keel!60!black}\hullmove{-10.21}{-8.25}\hullline{-9.43}{-8.25}\hullline{-10.63}{-7.57}\hullline{-10.81}{-7.57}\hullfill
\hullface{keel!57!black}\hullmove{10.86}{-6.77}\hullline{10.83}{-6.77}\hullline{10.63}{-7.57}\hullline{10.81}{-7.57}\hullfill
\hullface{keel!81!black}\hullmove{-10.81}{-7.57}\hullline{-10.63}{-7.57}\hullline{-10.83}{-6.77}\hullline{-10.86}{-6.77}\hullfill
\hullface{keel!57!black}\hullmove{5.45}{-8.77}\hullline{4.46}{-8.77}\hullline{0.00}{-9.06}\hullline{0.00}{-9.06}\hullfill
\hullface{keel!57!black}\hullmove{0.00}{-9.06}\hullline{0.00}{-9.06}\hullline{-4.46}{-8.77}\hullline{-5.45}{-8.77}\hullfill
\hullface{keel!57!black}\hullmove{9.43}{-8.25}\hullline{8.53}{-8.25}\hullline{4.46}{-8.77}\hullline{5.45}{-8.77}\hullfill
\hullface{keel!57!black}\hullmove{-5.45}{-8.77}\hullline{-4.46}{-8.77}\hullline{-8.53}{-8.25}\hullline{-9.43}{-8.25}\hullfill
\hullface{keel!57!black}\hullmove{10.63}{-7.57}\hullline{10.27}{-7.57}\hullline{8.53}{-8.25}\hullline{9.43}{-8.25}\hullfill
\hullface{keel!57!black}\hullmove{-9.43}{-8.25}\hullline{-8.53}{-8.25}\hullline{-10.27}{-7.57}\hullline{-10.63}{-7.57}\hullfill
\hullface{keel!57!black}\hullmove{10.83}{-6.77}\hullline{10.74}{-6.77}\hullline{10.27}{-7.57}\hullline{10.63}{-7.57}\hullfill
\hullface{keel!75!black}\hullmove{-10.63}{-7.57}\hullline{-10.27}{-7.57}\hullline{-10.74}{-6.77}\hullline{-10.83}{-6.77}\hullfill
\hullface{keel!57!black}\hullmove{10.85}{-5.86}\hullline{10.82}{-5.86}\hullline{10.74}{-6.77}\hullline{10.83}{-6.77}\hullfill
\hullface{keel!81!black}\hullmove{-10.83}{-6.77}\hullline{-10.74}{-6.77}\hullline{-10.82}{-5.86}\hullline{-10.85}{-5.86}\hullfill
\hullface{keel!57!black}\hullmove{10.85}{-5.44}\hullline{10.83}{-5.44}\hullline{10.82}{-5.86}\hullline{10.85}{-5.86}\hullfill
\hullface{topsides!57!black}\hullmove{10.87}{3.90}\hullline{10.87}{3.90}\hullline{10.86}{3.12}\hullline{10.86}{3.12}\hullfill
\hullface{keel!57!black}\hullmove{4.46}{-8.77}\hullline{3.70}{-8.77}\hullline{0.00}{-9.06}\hullline{0.00}{-9.06}\hullfill
\hullface{keel!57!black}\hullmove{0.00}{-9.06}\hullline{0.00}{-9.06}\hullline{-3.70}{-8.77}\hullline{-4.46}{-8.77}\hullfill
\hullface{keel!57!black}\hullmove{8.53}{-8.25}\hullline{7.63}{-8.25}\hullline{3.70}{-8.77}\hullline{4.46}{-8.77}\hullfill
\hullface{keel!57!black}\hullmove{-4.46}{-8.77}\hullline{-3.70}{-8.77}\hullline{-7.63}{-8.25}\hullline{-8.53}{-8.25}\hullfill
\hullface{keel!57!black}\hullmove{10.27}{-7.57}\hullline{9.77}{-7.57}\hullline{7.63}{-8.25}\hullline{8.53}{-8.25}\hullfill
\hullface{keel!57!black}\hullmove{-8.53}{-8.25}\hullline{-7.63}{-8.25}\hullline{-9.77}{-7.57}\hullline{-10.27}{-7.57}\hullfill
\hullface{keel!57!black}\hullmove{10.74}{-6.77}\hullline{10.56}{-6.77}\hullline{9.77}{-7.57}\hullline{10.27}{-7.57}\hullfill
\hullface{keel!69!black}\hullmove{-10.27}{-7.57}\hullline{-9.77}{-7.57}\hullline{-10.56}{-6.77}\hullline{-10.74}{-6.77}\hullfill
\hullface{keel!57!black}\hullmove{10.82}{-5.86}\hullline{10.77}{-5.86}\hullline{10.56}{-6.77}\hullline{10.74}{-6.77}\hullfill
\hullface{keel!81!black}\hullmove{-10.74}{-6.77}\hullline{-10.56}{-6.77}\hullline{-10.77}{-5.86}\hullline{-10.82}{-5.86}\hullfill
\hullface{keel!57!black}\hullmove{10.83}{-5.44}\hullline{10.79}{-5.44}\hullline{10.77}{-5.86}\hullline{10.82}{-5.86}\hullfill
\hullface{keel!81!black}\hullmove{-10.82}{-5.86}\hullline{-10.77}{-5.86}\hullline{-10.79}{-5.44}\hullline{-10.83}{-5.44}\hullfill
\hullface{topsides!81!black}\hullmove{-10.83}{0.00}\hullline{-10.80}{0.00}\hullline{-10.81}{0.78}\hullline{-10.82}{1.56}\hullline{-10.84}{2.34}\hullline{-10.85}{2.34}\hullline{-10.84}{1.56}\hullline{-10.83}{0.78}\hullfill
\hullface{topsides!57!black}\hullmove{10.86}{3.12}\hullline{10.87}{3.90}\hullline{10.87}{3.90}\hullline{10.85}{3.12}\hullline{10.84}{2.34}\hullline{10.85}{2.34}\hullfill
\hullface{keel!81!black}\hullmove{-10.83}{-5.44}\hullline{-10.79}{-5.44}\hullline{-10.80}{-4.85}\hullline{-10.80}{-3.75}\hullline{-10.80}{-2.58}\hullline{-10.80}{-1.32}\hullline{-10.80}{0.00}\hullline{-10.83}{0.00}\hullline{-10.83}{-1.32}\hullline{-10.83}{-2.58}\hullline{-10.83}{-3.75}\hullline{-10.83}{-4.85}\hullfill
\hullface{topsides!57!black}\hullmove{10.81}{0.78}\hullline{10.82}{1.56}\hullline{10.79}{1.56}\hullline{10.78}{0.78}\hullline{10.76}{0.00}\hullline{10.80}{0.00}\hullfill
\hullface{keel!57!black}\hullmove{10.80}{-3.75}\hullline{10.80}{-2.58}\hullline{10.80}{-1.32}\hullline{10.80}{0.00}\hullline{10.76}{0.00}\hullline{10.76}{-1.32}\hullline{10.76}{-2.58}\hullline{10.76}{-3.75}\hullline{10.75}{-4.85}\hullline{10.80}{-4.85}\hullfill
\hullface{keel!57!black}\hullmove{3.70}{-8.77}\hullline{3.12}{-8.77}\hullline{0.00}{-9.06}\hullline{0.00}{-9.06}\hullfill
\hullface{keel!57!black}\hullmove{0.00}{-9.06}\hullline{0.00}{-9.06}\hullline{-3.12}{-8.77}\hullline{-3.70}{-8.77}\hullfill
\hullface{keel!57!black}\hullmove{7.63}{-8.25}\hullline{6.80}{-8.25}\hullline{3.12}{-8.77}\hullline{3.70}{-8.77}\hullfill
\hullface{keel!57!black}\hullmove{-3.70}{-8.77}\hullline{-3.12}{-8.77}\hullline{-6.80}{-8.25}\hullline{-7.63}{-8.25}\hullfill
\hullface{keel!57!black}\hullmove{9.77}{-7.57}\hullline{9.17}{-7.57}\hullline{6.80}{-8.25}\hullline{7.63}{-8.25}\hullfill
\hullface{keel!57!black}\hullmove{-7.63}{-8.25}\hullline{-6.80}{-8.25}\hullline{-9.17}{-7.57}\hullline{-9.77}{-7.57}\hullfill
\hullface{keel!57!black}\hullmove{10.56}{-6.77}\hullline{10.28}{-6.77}\hullline{9.17}{-7.57}\hullline{9.77}{-7.57}\hullfill
\hullface{keel!60!black}\hullmove{-9.77}{-7.57}\hullline{-9.17}{-7.57}\hullline{-10.28}{-6.77}\hullline{-10.56}{-6.77}\hullfill
\hullface{keel!57!black}\hullmove{10.77}{-5.86}\hullline{10.66}{-5.86}\hullline{10.28}{-6.77}\hullline{10.56}{-6.77}\hullfill
\hullface{keel!78!black}\hullmove{-10.56}{-6.77}\hullline{-10.28}{-6.77}\hullline{-10.66}{-5.86}\hullline{-10.77}{-5.86}\hullfill
\hullface{keel!57!black}\hullmove{10.79}{-5.44}\hullline{10.72}{-5.44}\hullline{10.66}{-5.86}\hullline{10.77}{-5.86}\hullfill
\hullface{keel!81!black}\hullmove{-10.77}{-5.86}\hullline{-10.66}{-5.86}\hullline{-10.72}{-5.44}\hullline{-10.79}{-5.44}\hullfill
\hullface{keel!57!black}\hullmove{10.80}{-4.85}\hullline{10.75}{-4.85}\hullline{10.72}{-5.44}\hullline{10.79}{-5.44}\hullfill
\hullface{keel!81!black}\hullmove{-10.79}{-5.44}\hullline{-10.72}{-5.44}\hullline{-10.75}{-4.85}\hullline{-10.80}{-4.85}\hullfill
\hullface{topsides!81!black}\hullmove{-10.80}{0.00}\hullline{-10.76}{0.00}\hullline{-10.78}{0.78}\hullline{-10.79}{1.56}\hullline{-10.82}{1.56}\hullline{-10.81}{0.78}\hullfill
\hullface{topsides!81!black}\hullmove{-10.82}{1.56}\hullline{-10.79}{1.56}\hullline{-10.82}{2.34}\hullline{-10.84}{3.12}\hullline{-10.85}{3.12}\hullline{-10.84}{2.34}\hullfill
\hullface{topsides!57!black}\hullmove{10.87}{3.90}\hullline{10.87}{3.90}\hullline{10.84}{3.12}\hullline{10.85}{3.12}\hullfill
\hullface{keel!84!black}\hullmove{-10.76}{-3.75}\hullline{-10.71}{-3.75}\hullline{-10.71}{-2.58}\hullline{-10.71}{-1.32}\hullline{-10.71}{0.00}\hullline{-10.76}{0.00}\hullline{-10.76}{-1.32}\hullline{-10.76}{-2.58}\hullfill
\hullface{topsides!57!black}\hullmove{10.78}{0.78}\hullline{10.79}{1.56}\hullline{10.76}{1.56}\hullline{10.73}{0.78}\hullline{10.71}{0.00}\hullline{10.76}{0.00}\hullfill
\hullface{topsides!81!black}\hullmove{-10.76}{0.00}\hullline{-10.71}{0.00}\hullline{-10.73}{0.78}\hullline{-10.76}{1.56}\hullline{-10.79}{1.56}\hullline{-10.78}{0.78}\hullfill
\hullface{topsides!57!black}\hullmove{10.82}{2.34}\hullline{10.84}{3.12}\hullline{10.83}{3.12}\hullline{10.79}{2.34}\hullline{10.76}{1.56}\hullline{10.79}{1.56}\hullfill
\hullface{topsides!81!black}\hullmove{-10.79}{1.56}\hullline{-10.76}{1.56}\hullline{-10.79}{2.34}\hullline{-10.83}{3.12}\hullline{-10.84}{3.12}\hullline{-10.82}{2.34}\hullfill
\hullface{keel!57!black}\hullmove{3.12}{-8.77}\hullline{2.66}{-8.77}\hullline{0.00}{-9.06}\hullline{0.00}{-9.06}\hullfill
\hullface{keel!57!black}\hullmove{0.00}{-9.06}\hullline{0.00}{-9.06}\hullline{-2.66}{-8.77}\hullline{-3.12}{-8.77}\hullfill
\hullface{keel!57!black}\hullmove{6.80}{-8.25}\hullline{6.06}{-8.25}\hullline{2.66}{-8.77}\hullline{3.12}{-8.77}\hullfill
\hullface{keel!57!black}\hullmove{-3.12}{-8.77}\hullline{-2.66}{-8.77}\hullline{-6.06}{-8.25}\hullline{-6.80}{-8.25}\hullfill
\hullface{keel!57!black}\hullmove{9.17}{-7.57}\hullline{8.54}{-7.57}\hullline{6.06}{-8.25}\hullline{6.80}{-8.25}\hullfill
\hullface{keel!57!black}\hullmove{-6.80}{-8.25}\hullline{-6.06}{-8.25}\hullline{-8.54}{-7.57}\hullline{-9.17}{-7.57}\hullfill
\hullface{keel!57!black}\hullmove{10.28}{-6.77}\hullline{9.91}{-6.77}\hullline{8.54}{-7.57}\hullline{9.17}{-7.57}\hullfill
\hullface{keel!57!black}\hullmove{-9.17}{-7.57}\hullline{-8.54}{-7.57}\hullline{-9.91}{-6.77}\hullline{-10.28}{-6.77}\hullfill
\hullface{keel!57!black}\hullmove{10.66}{-5.86}\hullline{10.49}{-5.86}\hullline{9.91}{-6.77}\hullline{10.28}{-6.77}\hullfill
\hullface{keel!75!black}\hullmove{-10.28}{-6.77}\hullline{-9.91}{-6.77}\hullline{-10.49}{-5.86}\hullline{-10.66}{-5.86}\hullfill
\hullface{keel!57!black}\hullmove{10.72}{-5.44}\hullline{10.60}{-5.44}\hullline{10.49}{-5.86}\hullline{10.66}{-5.86}\hullfill
\hullface{keel!81!black}\hullmove{-10.66}{-5.86}\hullline{-10.49}{-5.86}\hullline{-10.60}{-5.44}\hullline{-10.72}{-5.44}\hullfill
\hullface{keel!57!black}\hullmove{10.75}{-4.85}\hullline{10.67}{-4.85}\hullline{10.60}{-5.44}\hullline{10.72}{-5.44}\hullfill
\hullface{keel!81!black}\hullmove{-10.72}{-5.44}\hullline{-10.60}{-5.44}\hullline{-10.67}{-4.85}\hullline{-10.75}{-4.85}\hullfill
\hullface{keel!57!black}\hullmove{10.76}{-3.75}\hullline{10.71}{-3.75}\hullline{10.67}{-4.85}\hullline{10.75}{-4.85}\hullfill
\hullface{keel!84!black}\hullmove{-10.75}{-4.85}\hullline{-10.67}{-4.85}\hullline{-10.71}{-3.75}\hullline{-10.76}{-3.75}\hullfill
\hullface{topsides!57!black}\hullmove{10.87}{3.90}\hullline{10.87}{3.90}\hullline{10.83}{3.12}\hullline{10.84}{3.12}\hullfill
\hullface{keel!57!black}\hullmove{2.66}{-8.77}\hullline{2.30}{-8.77}\hullline{0.00}{-9.06}\hullline{0.00}{-9.06}\hullfill
\hullface{keel!57!black}\hullmove{0.00}{-9.06}\hullline{0.00}{-9.06}\hullline{-2.30}{-8.77}\hullline{-2.66}{-8.77}\hullfill
\hullface{keel!57!black}\hullmove{6.06}{-8.25}\hullline{5.41}{-8.25}\hullline{2.30}{-8.77}\hullline{2.66}{-8.77}\hullfill
\hullface{keel!57!black}\hullmove{-2.66}{-8.77}\hullline{-2.30}{-8.77}\hullline{-5.41}{-8.25}\hullline{-6.06}{-8.25}\hullfill
\hullface{keel!57!black}\hullmove{8.54}{-7.57}\hullline{7.91}{-7.57}\hullline{5.41}{-8.25}\hullline{6.06}{-8.25}\hullfill
\hullface{keel!57!black}\hullmove{-6.06}{-8.25}\hullline{-5.41}{-8.25}\hullline{-7.91}{-7.57}\hullline{-8.54}{-7.57}\hullfill
\hullface{keel!57!black}\hullmove{9.91}{-6.77}\hullline{9.48}{-6.77}\hullline{7.91}{-7.57}\hullline{8.54}{-7.57}\hullfill
\hullface{keel!57!black}\hullmove{-8.54}{-7.57}\hullline{-7.91}{-7.57}\hullline{-9.48}{-6.77}\hullline{-9.91}{-6.77}\hullfill
\hullface{keel!57!black}\hullmove{10.49}{-5.86}\hullline{10.26}{-5.86}\hullline{9.48}{-6.77}\hullline{9.91}{-6.77}\hullfill
\hullface{keel!69!black}\hullmove{-9.91}{-6.77}\hullline{-9.48}{-6.77}\hullline{-10.26}{-5.86}\hullline{-10.49}{-5.86}\hullfill
\hullface{keel!57!black}\hullmove{10.60}{-5.44}\hullline{10.43}{-5.44}\hullline{10.26}{-5.86}\hullline{10.49}{-5.86}\hullfill
\hullface{keel!78!black}\hullmove{-10.49}{-5.86}\hullline{-10.26}{-5.86}\hullline{-10.43}{-5.44}\hullline{-10.60}{-5.44}\hullfill
\hullface{keel!57!black}\hullmove{10.67}{-4.85}\hullline{10.56}{-4.85}\hullline{10.43}{-5.44}\hullline{10.60}{-5.44}\hullfill
\hullface{keel!81!black}\hullmove{-10.60}{-5.44}\hullline{-10.43}{-5.44}\hullline{-10.56}{-4.85}\hullline{-10.67}{-4.85}\hullfill
\hullface{keel!57!black}\hullmove{10.71}{-3.75}\hullline{10.64}{-3.75}\hullline{10.56}{-4.85}\hullline{10.67}{-4.85}\hullfill
\hullface{keel!84!black}\hullmove{-10.67}{-4.85}\hullline{-10.56}{-4.85}\hullline{-10.64}{-3.75}\hullline{-10.71}{-3.75}\hullfill
\hullface{keel!57!black}\hullmove{10.71}{-2.58}\hullline{10.71}{-1.32}\hullline{10.71}{0.00}\hullline{10.65}{0.00}\hullline{10.65}{-1.32}\hullline{10.65}{-2.58}\hullline{10.64}{-3.75}\hullline{10.71}{-3.75}\hullfill
\hullface{topsides!57!black}\hullmove{10.73}{0.78}\hullline{10.67}{0.78}\hullline{10.65}{0.00}\hullline{10.71}{0.00}\hullfill
\hullface{topsides!81!black}\hullmove{-10.71}{0.00}\hullline{-10.65}{0.00}\hullline{-10.67}{0.78}\hullline{-10.73}{0.78}\hullfill
\hullface{topsides!57!black}\hullmove{10.76}{1.56}\hullline{10.79}{2.34}\hullline{10.76}{2.34}\hullline{10.71}{1.56}\hullline{10.67}{0.78}\hullline{10.73}{0.78}\hullfill
\hullface{topsides!81!black}\hullmove{-10.73}{0.78}\hullline{-10.67}{0.78}\hullline{-10.71}{1.56}\hullline{-10.76}{2.34}\hullline{-10.79}{2.34}\hullline{-10.76}{1.56}\hullfill
\hullface{topsides!57!black}\hullmove{10.83}{3.12}\hullline{10.81}{3.12}\hullline{10.76}{2.34}\hullline{10.79}{2.34}\hullfill
\hullface{topsides!81!black}\hullmove{-10.79}{2.34}\hullline{-10.76}{2.34}\hullline{-10.81}{3.12}\hullline{-10.83}{3.12}\hullfill
\hullface{topsides!57!black}\hullmove{10.87}{3.90}\hullline{10.87}{3.90}\hullline{10.81}{3.12}\hullline{10.83}{3.12}\hullfill
\hullface{topsides!78!black}\hullmove{-10.83}{3.12}\hullline{-10.81}{3.12}\hullline{-10.87}{3.90}\hullline{-10.87}{3.90}\hullfill
\hullface{keel!57!black}\hullmove{2.30}{-8.77}\hullline{2.00}{-8.77}\hullline{0.00}{-9.06}\hullline{0.00}{-9.06}\hullfill
\hullface{keel!57!black}\hullmove{0.00}{-9.06}\hullline{0.00}{-9.06}\hullline{-2.00}{-8.77}\hullline{-2.30}{-8.77}\hullfill
\hullface{keel!57!black}\hullmove{5.41}{-8.25}\hullline{4.84}{-8.25}\hullline{2.00}{-8.77}\hullline{2.30}{-8.77}\hullfill
\hullface{keel!57!black}\hullmove{-2.30}{-8.77}\hullline{-2.00}{-8.77}\hullline{-4.84}{-8.25}\hullline{-5.41}{-8.25}\hullfill
\hullface{keel!57!black}\hullmove{7.91}{-7.57}\hullline{7.30}{-7.57}\hullline{4.84}{-8.25}\hullline{5.41}{-8.25}\hullfill
\hullface{keel!57!black}\hullmove{-5.41}{-8.25}\hullline{-4.84}{-8.25}\hullline{-7.30}{-7.57}\hullline{-7.91}{-7.57}\hullfill
\hullface{keel!57!black}\hullmove{9.48}{-6.77}\hullline{9.00}{-6.77}\hullline{7.30}{-7.57}\hullline{7.91}{-7.57}\hullfill
\hullface{keel!57!black}\hullmove{-7.91}{-7.57}\hullline{-7.30}{-7.57}\hullline{-9.00}{-6.77}\hullline{-9.48}{-6.77}\hullfill
\hullface{keel!57!black}\hullmove{10.26}{-5.86}\hullline{9.96}{-5.86}\hullline{9.00}{-6.77}\hullline{9.48}{-6.77}\hullfill
\hullface{keel!66!black}\hullmove{-9.48}{-6.77}\hullline{-9.00}{-6.77}\hullline{-9.96}{-5.86}\hullline{-10.26}{-5.86}\hullfill
\hullface{keel!57!black}\hullmove{10.43}{-5.44}\hullline{10.20}{-5.44}\hullline{9.96}{-5.86}\hullline{10.26}{-5.86}\hullfill
\hullface{keel!75!black}\hullmove{-10.26}{-5.86}\hullline{-9.96}{-5.86}\hullline{-10.20}{-5.44}\hullline{-10.43}{-5.44}\hullfill
\hullface{keel!57!black}\hullmove{10.56}{-4.85}\hullline{10.40}{-4.85}\hullline{10.20}{-5.44}\hullline{10.43}{-5.44}\hullfill
\hullface{keel!81!black}\hullmove{-10.43}{-5.44}\hullline{-10.20}{-5.44}\hullline{-10.40}{-4.85}\hullline{-10.56}{-4.85}\hullfill
\hullface{keel!57!black}\hullmove{10.64}{-3.75}\hullline{10.54}{-3.75}\hullline{10.40}{-4.85}\hullline{10.56}{-4.85}\hullfill
\hullface{keel!84!black}\hullmove{-10.56}{-4.85}\hullline{-10.40}{-4.85}\hullline{-10.54}{-3.75}\hullline{-10.64}{-3.75}\hullfill
\hullface{keel!57!black}\hullmove{10.65}{-2.58}\hullline{10.57}{-2.58}\hullline{10.54}{-3.75}\hullline{10.64}{-3.75}\hullfill
\hullface{keel!84!black}\hullmove{-10.64}{-3.75}\hullline{-10.54}{-3.75}\hullline{-10.57}{-2.58}\hullline{-10.65}{-2.58}\hullfill
\hullface{topsides!57!black}\hullmove{10.67}{0.78}\hullline{10.60}{0.78}\hullline{10.57}{0.00}\hullline{10.65}{0.00}\hullfill
\hullface{topsides!84!black}\hullmove{-10.65}{0.00}\hullline{-10.57}{0.00}\hullline{-10.60}{0.78}\hullline{-10.67}{0.78}\hullfill
\hullface{topsides!57!black}\hullmove{10.71}{1.56}\hullline{10.66}{1.56}\hullline{10.60}{0.78}\hullline{10.67}{0.78}\hullfill
\hullface{topsides!81!black}\hullmove{-10.67}{0.78}\hullline{-10.60}{0.78}\hullline{-10.66}{1.56}\hullline{-10.71}{1.56}\hullfill
\hullface{topsides!57!black}\hullmove{10.76}{2.34}\hullline{10.72}{2.34}\hullline{10.66}{1.56}\hullline{10.71}{1.56}\hullfill
\hullface{topsides!81!black}\hullmove{-10.71}{1.56}\hullline{-10.66}{1.56}\hullline{-10.72}{2.34}\hullline{-10.76}{2.34}\hullfill
\hullface{topsides!57!black}\hullmove{10.81}{3.12}\hullline{10.79}{3.12}\hullline{10.72}{2.34}\hullline{10.76}{2.34}\hullfill
\hullface{topsides!57!black}\hullmove{10.87}{3.90}\hullline{10.87}{3.90}\hullline{10.79}{3.12}\hullline{10.81}{3.12}\hullfill
\hullface{topsides!78!black}\hullmove{-10.81}{3.12}\hullline{-10.79}{3.12}\hullline{-10.87}{3.90}\hullline{-10.87}{3.90}\hullfill
\hullface{keel!57!black}\hullmove{10.65}{-1.32}\hullline{10.65}{0.00}\hullline{10.57}{0.00}\hullline{10.57}{-1.32}\hullline{10.57}{-2.58}\hullline{10.65}{-2.58}\hullfill
\hullface{keel!84!black}\hullmove{-10.65}{-2.58}\hullline{-10.57}{-2.58}\hullline{-10.57}{-1.32}\hullline{-10.57}{0.00}\hullline{-10.65}{0.00}\hullline{-10.65}{-1.32}\hullfill
\hullface{keel!57!black}\hullmove{10.57}{-1.32}\hullline{10.57}{0.00}\hullline{10.48}{0.00}\hullline{10.48}{-1.32}\hullline{10.47}{-2.58}\hullline{10.57}{-2.58}\hullfill
\hullface{keel!84!black}\hullmove{-10.57}{-2.58}\hullline{-10.47}{-2.58}\hullline{-10.48}{-1.32}\hullline{-10.48}{0.00}\hullline{-10.57}{0.00}\hullline{-10.57}{-1.32}\hullfill
\hullface{keel!57!black}\hullmove{2.00}{-8.77}\hullline{1.76}{-8.77}\hullline{0.00}{-9.06}\hullline{0.00}{-9.06}\hullfill
\hullface{keel!57!black}\hullmove{0.00}{-9.06}\hullline{0.00}{-9.06}\hullline{-1.76}{-8.77}\hullline{-2.00}{-8.77}\hullfill
\hullface{keel!57!black}\hullmove{4.84}{-8.25}\hullline{4.35}{-8.25}\hullline{1.76}{-8.77}\hullline{2.00}{-8.77}\hullfill
\hullface{keel!57!black}\hullmove{-2.00}{-8.77}\hullline{-1.76}{-8.77}\hullline{-4.35}{-8.25}\hullline{-4.84}{-8.25}\hullfill
\hullface{keel!57!black}\hullmove{7.30}{-7.57}\hullline{6.72}{-7.57}\hullline{4.35}{-8.25}\hullline{4.84}{-8.25}\hullfill
\hullface{keel!57!black}\hullmove{-4.84}{-8.25}\hullline{-4.35}{-8.25}\hullline{-6.72}{-7.57}\hullline{-7.30}{-7.57}\hullfill
\hullface{keel!57!black}\hullmove{9.00}{-6.77}\hullline{8.50}{-6.77}\hullline{6.72}{-7.57}\hullline{7.30}{-7.57}\hullfill
\hullface{keel!57!black}\hullmove{-7.30}{-7.57}\hullline{-6.72}{-7.57}\hullline{-8.50}{-6.77}\hullline{-9.00}{-6.77}\hullfill
\hullface{keel!57!black}\hullmove{9.96}{-5.86}\hullline{9.61}{-5.86}\hullline{8.50}{-6.77}\hullline{9.00}{-6.77}\hullfill
\hullface{keel!63!black}\hullmove{-9.00}{-6.77}\hullline{-8.50}{-6.77}\hullline{-9.61}{-5.86}\hullline{-9.96}{-5.86}\hullfill
\hullface{keel!57!black}\hullmove{10.20}{-5.44}\hullline{9.91}{-5.44}\hullline{9.61}{-5.86}\hullline{9.96}{-5.86}\hullfill
\hullface{keel!72!black}\hullmove{-9.96}{-5.86}\hullline{-9.61}{-5.86}\hullline{-9.91}{-5.44}\hullline{-10.20}{-5.44}\hullfill
\hullface{keel!57!black}\hullmove{10.40}{-4.85}\hullline{10.19}{-4.85}\hullline{9.91}{-5.44}\hullline{10.20}{-5.44}\hullfill
\hullface{keel!78!black}\hullmove{-10.20}{-5.44}\hullline{-9.91}{-5.44}\hullline{-10.19}{-4.85}\hullline{-10.40}{-4.85}\hullfill
\hullface{keel!57!black}\hullmove{10.54}{-3.75}\hullline{10.41}{-3.75}\hullline{10.19}{-4.85}\hullline{10.40}{-4.85}\hullfill
\hullface{keel!84!black}\hullmove{-10.40}{-4.85}\hullline{-10.19}{-4.85}\hullline{-10.41}{-3.75}\hullline{-10.54}{-3.75}\hullfill
\hullface{keel!57!black}\hullmove{10.57}{-2.58}\hullline{10.47}{-2.58}\hullline{10.41}{-3.75}\hullline{10.54}{-3.75}\hullfill
\hullface{keel!84!black}\hullmove{-10.54}{-3.75}\hullline{-10.41}{-3.75}\hullline{-10.47}{-2.58}\hullline{-10.57}{-2.58}\hullfill
\hullface{topsides!57!black}\hullmove{10.60}{0.78}\hullline{10.52}{0.78}\hullline{10.48}{0.00}\hullline{10.57}{0.00}\hullfill
\hullface{topsides!84!black}\hullmove{-10.57}{0.00}\hullline{-10.48}{0.00}\hullline{-10.52}{0.78}\hullline{-10.60}{0.78}\hullfill
\hullface{topsides!57!black}\hullmove{10.66}{1.56}\hullline{10.59}{1.56}\hullline{10.52}{0.78}\hullline{10.60}{0.78}\hullfill
\hullface{topsides!81!black}\hullmove{-10.60}{0.78}\hullline{-10.52}{0.78}\hullline{-10.59}{1.56}\hullline{-10.66}{1.56}\hullfill
\hullface{topsides!57!black}\hullmove{10.72}{2.34}\hullline{10.67}{2.34}\hullline{10.59}{1.56}\hullline{10.66}{1.56}\hullfill
\hullface{topsides!81!black}\hullmove{-10.66}{1.56}\hullline{-10.59}{1.56}\hullline{-10.67}{2.34}\hullline{-10.72}{2.34}\hullfill
\hullface{topsides!57!black}\hullmove{10.79}{3.12}\hullline{10.76}{3.12}\hullline{10.67}{2.34}\hullline{10.72}{2.34}\hullfill
\hullface{topsides!78!black}\hullmove{-10.72}{2.34}\hullline{-10.67}{2.34}\hullline{-10.76}{3.12}\hullline{-10.79}{3.12}\hullfill
\hullface{topsides!57!black}\hullmove{10.87}{3.90}\hullline{10.86}{3.90}\hullline{10.76}{3.12}\hullline{10.79}{3.12}\hullfill
\hullface{topsides!78!black}\hullmove{-10.79}{3.12}\hullline{-10.76}{3.12}\hullline{-10.86}{3.90}\hullline{-10.87}{3.90}\hullfill
\hullface{keel!57!black}\hullmove{1.76}{-8.77}\hullline{1.56}{-8.77}\hullline{0.00}{-9.06}\hullline{0.00}{-9.06}\hullfill
\hullface{keel!57!black}\hullmove{0.00}{-9.06}\hullline{0.00}{-9.06}\hullline{-1.56}{-8.77}\hullline{-1.76}{-8.77}\hullfill
\hullface{keel!57!black}\hullmove{4.35}{-8.25}\hullline{3.91}{-8.25}\hullline{1.56}{-8.77}\hullline{1.76}{-8.77}\hullfill
\hullface{keel!57!black}\hullmove{-1.76}{-8.77}\hullline{-1.56}{-8.77}\hullline{-3.91}{-8.25}\hullline{-4.35}{-8.25}\hullfill
\hullface{keel!57!black}\hullmove{6.72}{-7.57}\hullline{6.19}{-7.57}\hullline{3.91}{-8.25}\hullline{4.35}{-8.25}\hullfill
\hullface{keel!57!black}\hullmove{-4.35}{-8.25}\hullline{-3.91}{-8.25}\hullline{-6.19}{-7.57}\hullline{-6.72}{-7.57}\hullfill
\hullface{keel!57!black}\hullmove{8.50}{-6.77}\hullline{8.00}{-6.77}\hullline{6.19}{-7.57}\hullline{6.72}{-7.57}\hullfill
\hullface{keel!57!black}\hullmove{-6.72}{-7.57}\hullline{-6.19}{-7.57}\hullline{-8.00}{-6.77}\hullline{-8.50}{-6.77}\hullfill
\hullface{keel!57!black}\hullmove{9.61}{-5.86}\hullline{9.23}{-5.86}\hullline{8.00}{-6.77}\hullline{8.50}{-6.77}\hullfill
\hullface{keel!60!black}\hullmove{-8.50}{-6.77}\hullline{-8.00}{-6.77}\hullline{-9.23}{-5.86}\hullline{-9.61}{-5.86}\hullfill
\hullface{keel!57!black}\hullmove{9.91}{-5.44}\hullline{9.59}{-5.44}\hullline{9.23}{-5.86}\hullline{9.61}{-5.86}\hullfill
\hullface{keel!69!black}\hullmove{-9.61}{-5.86}\hullline{-9.23}{-5.86}\hullline{-9.59}{-5.44}\hullline{-9.91}{-5.44}\hullfill
\hullface{keel!57!black}\hullmove{10.19}{-4.85}\hullline{9.93}{-4.85}\hullline{9.59}{-5.44}\hullline{9.91}{-5.44}\hullfill
\hullface{keel!75!black}\hullmove{-9.91}{-5.44}\hullline{-9.59}{-5.44}\hullline{-9.93}{-4.85}\hullline{-10.19}{-4.85}\hullfill
\hullface{keel!57!black}\hullmove{10.41}{-3.75}\hullline{10.25}{-3.75}\hullline{9.93}{-4.85}\hullline{10.19}{-4.85}\hullfill
\hullface{keel!81!black}\hullmove{-10.19}{-4.85}\hullline{-9.93}{-4.85}\hullline{-10.25}{-3.75}\hullline{-10.41}{-3.75}\hullfill
\hullface{keel!57!black}\hullmove{10.47}{-2.58}\hullline{10.35}{-2.58}\hullline{10.25}{-3.75}\hullline{10.41}{-3.75}\hullfill
\hullface{keel!84!black}\hullmove{-10.41}{-3.75}\hullline{-10.25}{-3.75}\hullline{-10.35}{-2.58}\hullline{-10.47}{-2.58}\hullfill
\hullface{keel!57!black}\hullmove{10.48}{-1.32}\hullline{10.48}{0.00}\hullline{10.37}{0.00}\hullline{10.37}{-1.32}\hullline{10.35}{-2.58}\hullline{10.47}{-2.58}\hullfill
\hullface{keel!87!black}\hullmove{-10.47}{-2.58}\hullline{-10.35}{-2.58}\hullline{-10.37}{-1.32}\hullline{-10.37}{0.00}\hullline{-10.48}{0.00}\hullline{-10.48}{-1.32}\hullfill
\hullface{topsides!57!black}\hullmove{10.52}{0.78}\hullline{10.42}{0.78}\hullline{10.37}{0.00}\hullline{10.48}{0.00}\hullfill
\hullface{topsides!84!black}\hullmove{-10.48}{0.00}\hullline{-10.37}{0.00}\hullline{-10.42}{0.78}\hullline{-10.52}{0.78}\hullfill
\hullface{topsides!57!black}\hullmove{10.59}{1.56}\hullline{10.50}{1.56}\hullline{10.42}{0.78}\hullline{10.52}{0.78}\hullfill
\hullface{topsides!81!black}\hullmove{-10.52}{0.78}\hullline{-10.42}{0.78}\hullline{-10.50}{1.56}\hullline{-10.59}{1.56}\hullfill
\hullface{topsides!57!black}\hullmove{10.67}{2.34}\hullline{10.61}{2.34}\hullline{10.50}{1.56}\hullline{10.59}{1.56}\hullfill
\hullface{topsides!81!black}\hullmove{-10.59}{1.56}\hullline{-10.50}{1.56}\hullline{-10.61}{2.34}\hullline{-10.67}{2.34}\hullfill
\hullface{topsides!57!black}\hullmove{10.76}{3.12}\hullline{10.73}{3.12}\hullline{10.61}{2.34}\hullline{10.67}{2.34}\hullfill
\hullface{topsides!78!black}\hullmove{-10.67}{2.34}\hullline{-10.61}{2.34}\hullline{-10.73}{3.12}\hullline{-10.76}{3.12}\hullfill
\hullface{topsides!57!black}\hullmove{10.86}{3.90}\hullline{10.86}{3.90}\hullline{10.73}{3.12}\hullline{10.76}{3.12}\hullfill
\hullface{topsides!75!black}\hullmove{-10.76}{3.12}\hullline{-10.73}{3.12}\hullline{-10.86}{3.90}\hullline{-10.86}{3.90}\hullfill
\hullface{keel!57!black}\hullmove{1.56}{-8.77}\hullline{1.38}{-8.77}\hullline{0.00}{-9.06}\hullline{0.00}{-9.06}\hullfill
\hullface{keel!57!black}\hullmove{0.00}{-9.06}\hullline{0.00}{-9.06}\hullline{-1.38}{-8.77}\hullline{-1.56}{-8.77}\hullfill
\hullface{keel!57!black}\hullmove{3.91}{-8.25}\hullline{3.53}{-8.25}\hullline{1.38}{-8.77}\hullline{1.56}{-8.77}\hullfill
\hullface{keel!57!black}\hullmove{-1.56}{-8.77}\hullline{-1.38}{-8.77}\hullline{-3.53}{-8.25}\hullline{-3.91}{-8.25}\hullfill
\hullface{keel!57!black}\hullmove{6.19}{-7.57}\hullline{5.69}{-7.57}\hullline{3.53}{-8.25}\hullline{3.91}{-8.25}\hullfill
\hullface{keel!57!black}\hullmove{-3.91}{-8.25}\hullline{-3.53}{-8.25}\hullline{-5.69}{-7.57}\hullline{-6.19}{-7.57}\hullfill
\hullface{keel!57!black}\hullmove{8.00}{-6.77}\hullline{7.50}{-6.77}\hullline{5.69}{-7.57}\hullline{6.19}{-7.57}\hullfill
\hullface{keel!57!black}\hullmove{-6.19}{-7.57}\hullline{-5.69}{-7.57}\hullline{-7.50}{-6.77}\hullline{-8.00}{-6.77}\hullfill
\hullface{keel!57!black}\hullmove{9.23}{-5.86}\hullline{8.81}{-5.86}\hullline{7.50}{-6.77}\hullline{8.00}{-6.77}\hullfill
\hullface{keel!57!black}\hullmove{-8.00}{-6.77}\hullline{-7.50}{-6.77}\hullline{-8.81}{-5.86}\hullline{-9.23}{-5.86}\hullfill
\hullface{keel!57!black}\hullmove{9.59}{-5.44}\hullline{9.22}{-5.44}\hullline{8.81}{-5.86}\hullline{9.23}{-5.86}\hullfill
\hullface{keel!66!black}\hullmove{-9.23}{-5.86}\hullline{-8.81}{-5.86}\hullline{-9.22}{-5.44}\hullline{-9.59}{-5.44}\hullfill
\hullface{keel!57!black}\hullmove{9.93}{-4.85}\hullline{9.63}{-4.85}\hullline{9.22}{-5.44}\hullline{9.59}{-5.44}\hullfill
\hullface{keel!72!black}\hullmove{-9.59}{-5.44}\hullline{-9.22}{-5.44}\hullline{-9.63}{-4.85}\hullline{-9.93}{-4.85}\hullfill
\hullface{keel!57!black}\hullmove{10.25}{-3.75}\hullline{10.05}{-3.75}\hullline{9.63}{-4.85}\hullline{9.93}{-4.85}\hullfill
\hullface{keel!81!black}\hullmove{-9.93}{-4.85}\hullline{-9.63}{-4.85}\hullline{-10.05}{-3.75}\hullline{-10.25}{-3.75}\hullfill
\hullface{keel!57!black}\hullmove{10.35}{-2.58}\hullline{10.20}{-2.58}\hullline{10.05}{-3.75}\hullline{10.25}{-3.75}\hullfill
\hullface{keel!84!black}\hullmove{-10.25}{-3.75}\hullline{-10.05}{-3.75}\hullline{-10.20}{-2.58}\hullline{-10.35}{-2.58}\hullfill
\hullface{keel!57!black}\hullmove{10.37}{-1.32}\hullline{10.24}{-1.32}\hullline{10.20}{-2.58}\hullline{10.35}{-2.58}\hullfill
\hullface{keel!87!black}\hullmove{-10.35}{-2.58}\hullline{-10.20}{-2.58}\hullline{-10.24}{-1.32}\hullline{-10.37}{-1.32}\hullfill
\hullface{keel!57!black}\hullmove{10.37}{0.00}\hullline{10.24}{0.00}\hullline{10.24}{-1.32}\hullline{10.37}{-1.32}\hullfill
\hullface{keel!87!black}\hullmove{-10.37}{-1.32}\hullline{-10.24}{-1.32}\hullline{-10.24}{0.00}\hullline{-10.37}{0.00}\hullfill
\hullface{topsides!57!black}\hullmove{10.42}{0.78}\hullline{10.30}{0.78}\hullline{10.24}{0.00}\hullline{10.37}{0.00}\hullfill
\hullface{topsides!84!black}\hullmove{-10.37}{0.00}\hullline{-10.24}{0.00}\hullline{-10.30}{0.78}\hullline{-10.42}{0.78}\hullfill
\hullface{topsides!57!black}\hullmove{10.50}{1.56}\hullline{10.41}{1.56}\hullline{10.30}{0.78}\hullline{10.42}{0.78}\hullfill
\hullface{topsides!81!black}\hullmove{-10.42}{0.78}\hullline{-10.30}{0.78}\hullline{-10.41}{1.56}\hullline{-10.50}{1.56}\hullfill
\hullface{topsides!57!black}\hullmove{10.61}{2.34}\hullline{10.54}{2.34}\hullline{10.41}{1.56}\hullline{10.50}{1.56}\hullfill
\hullface{topsides!78!black}\hullmove{-10.50}{1.56}\hullline{-10.41}{1.56}\hullline{-10.54}{2.34}\hullline{-10.61}{2.34}\hullfill
\hullface{topsides!57!black}\hullmove{10.73}{3.12}\hullline{10.69}{3.12}\hullline{10.54}{2.34}\hullline{10.61}{2.34}\hullfill
\hullface{topsides!78!black}\hullmove{-10.61}{2.34}\hullline{-10.54}{2.34}\hullline{-10.69}{3.12}\hullline{-10.73}{3.12}\hullfill
\hullface{topsides!57!black}\hullmove{10.86}{3.90}\hullline{10.85}{3.90}\hullline{10.69}{3.12}\hullline{10.73}{3.12}\hullfill
\hullface{topsides!75!black}\hullmove{-10.73}{3.12}\hullline{-10.69}{3.12}\hullline{-10.85}{3.90}\hullline{-10.86}{3.90}\hullfill
\hullface{keel!57!black}\hullmove{1.38}{-8.77}\hullline{1.24}{-8.77}\hullline{0.00}{-9.06}\hullline{0.00}{-9.06}\hullfill
\hullface{keel!57!black}\hullmove{0.00}{-9.06}\hullline{0.00}{-9.06}\hullline{-1.24}{-8.77}\hullline{-1.38}{-8.77}\hullfill
\hullface{keel!57!black}\hullmove{3.53}{-8.25}\hullline{3.19}{-8.25}\hullline{1.24}{-8.77}\hullline{1.38}{-8.77}\hullfill
\hullface{keel!57!black}\hullmove{-1.38}{-8.77}\hullline{-1.24}{-8.77}\hullline{-3.19}{-8.25}\hullline{-3.53}{-8.25}\hullfill
\hullface{keel!57!black}\hullmove{5.69}{-7.57}\hullline{5.22}{-7.57}\hullline{3.19}{-8.25}\hullline{3.53}{-8.25}\hullfill
\hullface{keel!57!black}\hullmove{-3.53}{-8.25}\hullline{-3.19}{-8.25}\hullline{-5.22}{-7.57}\hullline{-5.69}{-7.57}\hullfill
\hullface{keel!57!black}\hullmove{7.50}{-6.77}\hullline{7.00}{-6.77}\hullline{5.22}{-7.57}\hullline{5.69}{-7.57}\hullfill
\hullface{keel!57!black}\hullmove{-5.69}{-7.57}\hullline{-5.22}{-7.57}\hullline{-7.00}{-6.77}\hullline{-7.50}{-6.77}\hullfill
\hullface{keel!57!black}\hullmove{8.81}{-5.86}\hullline{8.37}{-5.86}\hullline{7.00}{-6.77}\hullline{7.50}{-6.77}\hullfill
\hullface{keel!57!black}\hullmove{-7.50}{-6.77}\hullline{-7.00}{-6.77}\hullline{-8.37}{-5.86}\hullline{-8.81}{-5.86}\hullfill
\hullface{keel!57!black}\hullmove{9.22}{-5.44}\hullline{8.82}{-5.44}\hullline{8.37}{-5.86}\hullline{8.81}{-5.86}\hullfill
\hullface{keel!63!black}\hullmove{-8.81}{-5.86}\hullline{-8.37}{-5.86}\hullline{-8.82}{-5.44}\hullline{-9.22}{-5.44}\hullfill
\hullface{keel!57!black}\hullmove{9.63}{-4.85}\hullline{9.29}{-4.85}\hullline{8.82}{-5.44}\hullline{9.22}{-5.44}\hullfill
\hullface{keel!69!black}\hullmove{-9.22}{-5.44}\hullline{-8.82}{-5.44}\hullline{-9.29}{-4.85}\hullline{-9.63}{-4.85}\hullfill
\hullface{keel!57!black}\hullmove{10.05}{-3.75}\hullline{9.81}{-3.75}\hullline{9.29}{-4.85}\hullline{9.63}{-4.85}\hullfill
\hullface{keel!78!black}\hullmove{-9.63}{-4.85}\hullline{-9.29}{-4.85}\hullline{-9.81}{-3.75}\hullline{-10.05}{-3.75}\hullfill
\hullface{keel!57!black}\hullmove{10.20}{-2.58}\hullline{10.03}{-2.58}\hullline{9.81}{-3.75}\hullline{10.05}{-3.75}\hullfill
\hullface{keel!84!black}\hullmove{-10.05}{-3.75}\hullline{-9.81}{-3.75}\hullline{-10.03}{-2.58}\hullline{-10.20}{-2.58}\hullfill
\hullface{keel!57!black}\hullmove{10.24}{-1.32}\hullline{10.08}{-1.32}\hullline{10.03}{-2.58}\hullline{10.20}{-2.58}\hullfill
\hullface{keel!87!black}\hullmove{-10.20}{-2.58}\hullline{-10.03}{-2.58}\hullline{-10.08}{-1.32}\hullline{-10.24}{-1.32}\hullfill
\hullface{keel!57!black}\hullmove{10.24}{0.00}\hullline{10.09}{0.00}\hullline{10.08}{-1.32}\hullline{10.24}{-1.32}\hullfill
\hullface{keel!87!black}\hullmove{-10.24}{-1.32}\hullline{-10.08}{-1.32}\hullline{-10.09}{0.00}\hullline{-10.24}{0.00}\hullfill
\hullface{topsides!57!black}\hullmove{10.30}{0.78}\hullline{10.17}{0.78}\hullline{10.09}{0.00}\hullline{10.24}{0.00}\hullfill
\hullface{topsides!84!black}\hullmove{-10.24}{0.00}\hullline{-10.09}{0.00}\hullline{-10.17}{0.78}\hullline{-10.30}{0.78}\hullfill
\hullface{topsides!57!black}\hullmove{10.41}{1.56}\hullline{10.30}{1.56}\hullline{10.17}{0.78}\hullline{10.30}{0.78}\hullfill
\hullface{topsides!81!black}\hullmove{-10.30}{0.78}\hullline{-10.17}{0.78}\hullline{-10.30}{1.56}\hullline{-10.41}{1.56}\hullfill
\hullface{topsides!57!black}\hullmove{10.54}{2.34}\hullline{10.46}{2.34}\hullline{10.30}{1.56}\hullline{10.41}{1.56}\hullfill
\hullface{topsides!78!black}\hullmove{-10.41}{1.56}\hullline{-10.30}{1.56}\hullline{-10.46}{2.34}\hullline{-10.54}{2.34}\hullfill
\hullface{topsides!57!black}\hullmove{10.69}{3.12}\hullline{10.64}{3.12}\hullline{10.46}{2.34}\hullline{10.54}{2.34}\hullfill
\hullface{topsides!75!black}\hullmove{-10.54}{2.34}\hullline{-10.46}{2.34}\hullline{-10.64}{3.12}\hullline{-10.69}{3.12}\hullfill
\hullface{topsides!57!black}\hullmove{10.85}{3.90}\hullline{10.84}{3.90}\hullline{10.64}{3.12}\hullline{10.69}{3.12}\hullfill
\hullface{topsides!72!black}\hullmove{-10.69}{3.12}\hullline{-10.64}{3.12}\hullline{-10.84}{3.90}\hullline{-10.85}{3.90}\hullfill
\hullface{keel!57!black}\hullmove{1.24}{-8.77}\hullline{1.11}{-8.77}\hullline{0.00}{-9.06}\hullline{0.00}{-9.06}\hullfill
\hullface{keel!57!black}\hullmove{0.00}{-9.06}\hullline{0.00}{-9.06}\hullline{-1.11}{-8.77}\hullline{-1.24}{-8.77}\hullfill
\hullface{keel!57!black}\hullmove{3.19}{-8.25}\hullline{2.89}{-8.25}\hullline{1.11}{-8.77}\hullline{1.24}{-8.77}\hullfill
\hullface{keel!57!black}\hullmove{-1.24}{-8.77}\hullline{-1.11}{-8.77}\hullline{-2.89}{-8.25}\hullline{-3.19}{-8.25}\hullfill
\hullface{keel!57!black}\hullmove{5.22}{-7.57}\hullline{4.80}{-7.57}\hullline{2.89}{-8.25}\hullline{3.19}{-8.25}\hullfill
\hullface{keel!57!black}\hullmove{-3.19}{-8.25}\hullline{-2.89}{-8.25}\hullline{-4.80}{-7.57}\hullline{-5.22}{-7.57}\hullfill
\hullface{keel!57!black}\hullmove{7.00}{-6.77}\hullline{6.53}{-6.77}\hullline{4.80}{-7.57}\hullline{5.22}{-7.57}\hullfill
\hullface{keel!57!black}\hullmove{-5.22}{-7.57}\hullline{-4.80}{-7.57}\hullline{-6.53}{-6.77}\hullline{-7.00}{-6.77}\hullfill
\hullface{keel!57!black}\hullmove{8.37}{-5.86}\hullline{7.93}{-5.86}\hullline{6.53}{-6.77}\hullline{7.00}{-6.77}\hullfill
\hullface{keel!57!black}\hullmove{-7.00}{-6.77}\hullline{-6.53}{-6.77}\hullline{-7.93}{-5.86}\hullline{-8.37}{-5.86}\hullfill
\hullface{keel!57!black}\hullmove{8.82}{-5.44}\hullline{8.40}{-5.44}\hullline{7.93}{-5.86}\hullline{8.37}{-5.86}\hullfill
\hullface{keel!63!black}\hullmove{-8.37}{-5.86}\hullline{-7.93}{-5.86}\hullline{-8.40}{-5.44}\hullline{-8.82}{-5.44}\hullfill
\hullface{keel!57!black}\hullmove{9.29}{-4.85}\hullline{8.92}{-4.85}\hullline{8.40}{-5.44}\hullline{8.82}{-5.44}\hullfill
\hullface{keel!69!black}\hullmove{-8.82}{-5.44}\hullline{-8.40}{-5.44}\hullline{-8.92}{-4.85}\hullline{-9.29}{-4.85}\hullfill
\hullface{keel!57!black}\hullmove{9.81}{-3.75}\hullline{9.53}{-3.75}\hullline{8.92}{-4.85}\hullline{9.29}{-4.85}\hullfill
\hullface{keel!78!black}\hullmove{-9.29}{-4.85}\hullline{-8.92}{-4.85}\hullline{-9.53}{-3.75}\hullline{-9.81}{-3.75}\hullfill
\hullface{keel!57!black}\hullmove{10.03}{-2.58}\hullline{9.82}{-2.58}\hullline{9.53}{-3.75}\hullline{9.81}{-3.75}\hullfill
\hullface{keel!84!black}\hullmove{-9.81}{-3.75}\hullline{-9.53}{-3.75}\hullline{-9.82}{-2.58}\hullline{-10.03}{-2.58}\hullfill
\hullface{keel!57!black}\hullmove{10.08}{-1.32}\hullline{9.90}{-1.32}\hullline{9.82}{-2.58}\hullline{10.03}{-2.58}\hullfill
\hullface{keel!87!black}\hullmove{-10.03}{-2.58}\hullline{-9.82}{-2.58}\hullline{-9.90}{-1.32}\hullline{-10.08}{-1.32}\hullfill
\hullface{keel!57!black}\hullmove{10.09}{0.00}\hullline{9.91}{0.00}\hullline{9.90}{-1.32}\hullline{10.08}{-1.32}\hullfill
\hullface{keel!90!black}\hullmove{-10.08}{-1.32}\hullline{-9.90}{-1.32}\hullline{-9.91}{0.00}\hullline{-10.09}{0.00}\hullfill
\hullface{topsides!57!black}\hullmove{10.17}{0.78}\hullline{10.01}{0.78}\hullline{9.91}{0.00}\hullline{10.09}{0.00}\hullfill
\hullface{topsides!84!black}\hullmove{-10.09}{0.00}\hullline{-9.91}{0.00}\hullline{-10.01}{0.78}\hullline{-10.17}{0.78}\hullfill
\hullface{topsides!57!black}\hullmove{10.30}{1.56}\hullline{10.17}{1.56}\hullline{10.01}{0.78}\hullline{10.17}{0.78}\hullfill
\hullface{topsides!81!black}\hullmove{-10.17}{0.78}\hullline{-10.01}{0.78}\hullline{-10.17}{1.56}\hullline{-10.30}{1.56}\hullfill
\hullface{topsides!57!black}\hullmove{10.46}{2.34}\hullline{10.36}{2.34}\hullline{10.17}{1.56}\hullline{10.30}{1.56}\hullfill
\hullface{topsides!78!black}\hullmove{-10.30}{1.56}\hullline{-10.17}{1.56}\hullline{-10.36}{2.34}\hullline{-10.46}{2.34}\hullfill
\hullface{topsides!57!black}\hullmove{10.64}{3.12}\hullline{10.58}{3.12}\hullline{10.36}{2.34}\hullline{10.46}{2.34}\hullfill
\hullface{topsides!75!black}\hullmove{-10.46}{2.34}\hullline{-10.36}{2.34}\hullline{-10.58}{3.12}\hullline{-10.64}{3.12}\hullfill
\hullface{topsides!57!black}\hullmove{10.84}{3.90}\hullline{10.83}{3.90}\hullline{10.58}{3.12}\hullline{10.64}{3.12}\hullfill
\hullface{topsides!72!black}\hullmove{-10.64}{3.12}\hullline{-10.58}{3.12}\hullline{-10.83}{3.90}\hullline{-10.84}{3.90}\hullfill
\hullface{keel!57!black}\hullmove{1.11}{-8.77}\hullline{0.99}{-8.77}\hullline{0.00}{-9.06}\hullline{0.00}{-9.06}\hullfill
\hullface{keel!57!black}\hullmove{0.00}{-9.06}\hullline{0.00}{-9.06}\hullline{-0.99}{-8.77}\hullline{-1.11}{-8.77}\hullfill
\hullface{keel!57!black}\hullmove{2.89}{-8.25}\hullline{2.63}{-8.25}\hullline{0.99}{-8.77}\hullline{1.11}{-8.77}\hullfill
\hullface{keel!57!black}\hullmove{-1.11}{-8.77}\hullline{-0.99}{-8.77}\hullline{-2.63}{-8.25}\hullline{-2.89}{-8.25}\hullfill
\hullface{keel!57!black}\hullmove{4.80}{-7.57}\hullline{4.40}{-7.57}\hullline{2.63}{-8.25}\hullline{2.89}{-8.25}\hullfill
\hullface{keel!57!black}\hullmove{-2.89}{-8.25}\hullline{-2.63}{-8.25}\hullline{-4.40}{-7.57}\hullline{-4.80}{-7.57}\hullfill
\hullface{keel!57!black}\hullmove{6.53}{-6.77}\hullline{6.07}{-6.77}\hullline{4.40}{-7.57}\hullline{4.80}{-7.57}\hullfill
\hullface{keel!57!black}\hullmove{-4.80}{-7.57}\hullline{-4.40}{-7.57}\hullline{-6.07}{-6.77}\hullline{-6.53}{-6.77}\hullfill
\hullface{keel!57!black}\hullmove{7.93}{-5.86}\hullline{7.47}{-5.86}\hullline{6.07}{-6.77}\hullline{6.53}{-6.77}\hullfill
\hullface{keel!57!black}\hullmove{-6.53}{-6.77}\hullline{-6.07}{-6.77}\hullline{-7.47}{-5.86}\hullline{-7.93}{-5.86}\hullfill
\hullface{keel!57!black}\hullmove{8.40}{-5.44}\hullline{7.97}{-5.44}\hullline{7.47}{-5.86}\hullline{7.93}{-5.86}\hullfill
\hullface{keel!63!black}\hullmove{-7.93}{-5.86}\hullline{-7.47}{-5.86}\hullline{-7.97}{-5.44}\hullline{-8.40}{-5.44}\hullfill
\hullface{keel!57!black}\hullmove{8.92}{-4.85}\hullline{8.53}{-4.85}\hullline{7.97}{-5.44}\hullline{8.40}{-5.44}\hullfill
\hullface{keel!66!black}\hullmove{-8.40}{-5.44}\hullline{-7.97}{-5.44}\hullline{-8.53}{-4.85}\hullline{-8.92}{-4.85}\hullfill
\hullface{keel!57!black}\hullmove{9.53}{-3.75}\hullline{9.21}{-3.75}\hullline{8.53}{-4.85}\hullline{8.92}{-4.85}\hullfill
\hullface{keel!75!black}\hullmove{-8.92}{-4.85}\hullline{-8.53}{-4.85}\hullline{-9.21}{-3.75}\hullline{-9.53}{-3.75}\hullfill
\hullface{keel!57!black}\hullmove{9.82}{-2.58}\hullline{9.57}{-2.58}\hullline{9.21}{-3.75}\hullline{9.53}{-3.75}\hullfill
\hullface{keel!84!black}\hullmove{-9.53}{-3.75}\hullline{-9.21}{-3.75}\hullline{-9.57}{-2.58}\hullline{-9.82}{-2.58}\hullfill
\hullface{keel!57!black}\hullmove{9.90}{-1.32}\hullline{9.70}{-1.32}\hullline{9.57}{-2.58}\hullline{9.82}{-2.58}\hullfill
\hullface{keel!87!black}\hullmove{-9.82}{-2.58}\hullline{-9.57}{-2.58}\hullline{-9.70}{-1.32}\hullline{-9.90}{-1.32}\hullfill
\hullface{keel!57!black}\hullmove{9.91}{0.00}\hullline{9.72}{0.00}\hullline{9.70}{-1.32}\hullline{9.90}{-1.32}\hullfill
\hullface{keel!90!black}\hullmove{-9.90}{-1.32}\hullline{-9.70}{-1.32}\hullline{-9.72}{0.00}\hullline{-9.91}{0.00}\hullfill
\hullface{topsides!57!black}\hullmove{10.01}{0.78}\hullline{9.83}{0.78}\hullline{9.72}{0.00}\hullline{9.91}{0.00}\hullfill
\hullface{topsides!84!black}\hullmove{-9.91}{0.00}\hullline{-9.72}{0.00}\hullline{-9.83}{0.78}\hullline{-10.01}{0.78}\hullfill
\hullface{topsides!57!black}\hullmove{10.17}{1.56}\hullline{10.02}{1.56}\hullline{9.83}{0.78}\hullline{10.01}{0.78}\hullfill
\hullface{topsides!81!black}\hullmove{-10.01}{0.78}\hullline{-9.83}{0.78}\hullline{-10.02}{1.56}\hullline{-10.17}{1.56}\hullfill
\hullface{topsides!57!black}\hullmove{10.36}{2.34}\hullline{10.25}{2.34}\hullline{10.02}{1.56}\hullline{10.17}{1.56}\hullfill
\hullface{topsides!78!black}\hullmove{-10.17}{1.56}\hullline{-10.02}{1.56}\hullline{-10.25}{2.34}\hullline{-10.36}{2.34}\hullfill
\hullface{topsides!57!black}\hullmove{10.58}{3.12}\hullline{10.51}{3.12}\hullline{10.25}{2.34}\hullline{10.36}{2.34}\hullfill
\hullface{topsides!75!black}\hullmove{-10.36}{2.34}\hullline{-10.25}{2.34}\hullline{-10.51}{3.12}\hullline{-10.58}{3.12}\hullfill
\hullface{topsides!57!black}\hullmove{10.83}{3.90}\hullline{10.81}{3.90}\hullline{10.51}{3.12}\hullline{10.58}{3.12}\hullfill
\hullface{topsides!69!black}\hullmove{-10.58}{3.12}\hullline{-10.51}{3.12}\hullline{-10.81}{3.90}\hullline{-10.83}{3.90}\hullfill
\hullface{keel!57!black}\hullmove{0.99}{-8.77}\hullline{0.90}{-8.77}\hullline{0.00}{-9.06}\hullline{0.00}{-9.06}\hullfill
\hullface{keel!57!black}\hullmove{0.00}{-9.06}\hullline{0.00}{-9.06}\hullline{-0.90}{-8.77}\hullline{-0.99}{-8.77}\hullfill
\hullface{keel!57!black}\hullmove{2.63}{-8.25}\hullline{2.38}{-8.25}\hullline{0.90}{-8.77}\hullline{0.99}{-8.77}\hullfill
\hullface{keel!57!black}\hullmove{-0.99}{-8.77}\hullline{-0.90}{-8.77}\hullline{-2.38}{-8.25}\hullline{-2.63}{-8.25}\hullfill
\hullface{keel!57!black}\hullmove{4.40}{-7.57}\hullline{4.03}{-7.57}\hullline{2.38}{-8.25}\hullline{2.63}{-8.25}\hullfill
\hullface{keel!57!black}\hullmove{-2.63}{-8.25}\hullline{-2.38}{-8.25}\hullline{-4.03}{-7.57}\hullline{-4.40}{-7.57}\hullfill
\hullface{keel!57!black}\hullmove{6.07}{-6.77}\hullline{5.63}{-6.77}\hullline{4.03}{-7.57}\hullline{4.40}{-7.57}\hullfill
\hullface{keel!57!black}\hullmove{-4.40}{-7.57}\hullline{-4.03}{-7.57}\hullline{-5.63}{-6.77}\hullline{-6.07}{-6.77}\hullfill
\hullface{keel!57!black}\hullmove{7.47}{-5.86}\hullline{7.02}{-5.86}\hullline{5.63}{-6.77}\hullline{6.07}{-6.77}\hullfill
\hullface{keel!57!black}\hullmove{-6.07}{-6.77}\hullline{-5.63}{-6.77}\hullline{-7.02}{-5.86}\hullline{-7.47}{-5.86}\hullfill
\hullface{keel!57!black}\hullmove{7.97}{-5.44}\hullline{7.53}{-5.44}\hullline{7.02}{-5.86}\hullline{7.47}{-5.86}\hullfill
\hullface{keel!60!black}\hullmove{-7.47}{-5.86}\hullline{-7.02}{-5.86}\hullline{-7.53}{-5.44}\hullline{-7.97}{-5.44}\hullfill
\hullface{keel!57!black}\hullmove{8.53}{-4.85}\hullline{8.11}{-4.85}\hullline{7.53}{-5.44}\hullline{7.97}{-5.44}\hullfill
\hullface{keel!66!black}\hullmove{-7.97}{-5.44}\hullline{-7.53}{-5.44}\hullline{-8.11}{-4.85}\hullline{-8.53}{-4.85}\hullfill
\hullface{keel!57!black}\hullmove{9.21}{-3.75}\hullline{8.87}{-3.75}\hullline{8.11}{-4.85}\hullline{8.53}{-4.85}\hullfill
\hullface{keel!75!black}\hullmove{-8.53}{-4.85}\hullline{-8.11}{-4.85}\hullline{-8.87}{-3.75}\hullline{-9.21}{-3.75}\hullfill
\hullface{keel!57!black}\hullmove{9.57}{-2.58}\hullline{9.30}{-2.58}\hullline{8.87}{-3.75}\hullline{9.21}{-3.75}\hullfill
\hullface{keel!81!black}\hullmove{-9.21}{-3.75}\hullline{-8.87}{-3.75}\hullline{-9.30}{-2.58}\hullline{-9.57}{-2.58}\hullfill
\hullface{keel!57!black}\hullmove{9.70}{-1.32}\hullline{9.47}{-1.32}\hullline{9.30}{-2.58}\hullline{9.57}{-2.58}\hullfill
\hullface{keel!87!black}\hullmove{-9.57}{-2.58}\hullline{-9.30}{-2.58}\hullline{-9.47}{-1.32}\hullline{-9.70}{-1.32}\hullfill
\hullface{keel!57!black}\hullmove{9.72}{0.00}\hullline{9.50}{0.00}\hullline{9.47}{-1.32}\hullline{9.70}{-1.32}\hullfill
\hullface{keel!90!black}\hullmove{-9.70}{-1.32}\hullline{-9.47}{-1.32}\hullline{-9.50}{0.00}\hullline{-9.72}{0.00}\hullfill
\hullface{topsides!57!black}\hullmove{9.83}{0.78}\hullline{9.63}{0.78}\hullline{9.50}{0.00}\hullline{9.72}{0.00}\hullfill
\hullface{topsides!84!black}\hullmove{-9.72}{0.00}\hullline{-9.50}{0.00}\hullline{-9.63}{0.78}\hullline{-9.83}{0.78}\hullfill
\hullface{topsides!57!black}\hullmove{10.02}{1.56}\hullline{9.85}{1.56}\hullline{9.63}{0.78}\hullline{9.83}{0.78}\hullfill
\hullface{topsides!81!black}\hullmove{-9.83}{0.78}\hullline{-9.63}{0.78}\hullline{-9.85}{1.56}\hullline{-10.02}{1.56}\hullfill
\hullface{topsides!57!black}\hullmove{10.25}{2.34}\hullline{10.12}{2.34}\hullline{9.85}{1.56}\hullline{10.02}{1.56}\hullfill
\hullface{topsides!75!black}\hullmove{-10.02}{1.56}\hullline{-9.85}{1.56}\hullline{-10.12}{2.34}\hullline{-10.25}{2.34}\hullfill
\hullface{topsides!57!black}\hullmove{10.51}{3.12}\hullline{10.44}{3.12}\hullline{10.12}{2.34}\hullline{10.25}{2.34}\hullfill
\hullface{topsides!72!black}\hullmove{-10.25}{2.34}\hullline{-10.12}{2.34}\hullline{-10.44}{3.12}\hullline{-10.51}{3.12}\hullfill
\hullface{topsides!57!black}\hullmove{10.81}{3.90}\hullline{10.78}{3.90}\hullline{10.44}{3.12}\hullline{10.51}{3.12}\hullfill
\hullface{topsides!69!black}\hullmove{-10.51}{3.12}\hullline{-10.44}{3.12}\hullline{-10.78}{3.90}\hullline{-10.81}{3.90}\hullfill
\hullface{keel!57!black}\hullmove{0.90}{-8.77}\hullline{0.83}{-8.77}\hullline{0.00}{-9.06}\hullline{0.00}{-9.06}\hullfill
\hullface{keel!57!black}\hullmove{2.38}{-8.25}\hullline{2.21}{-8.25}\hullline{0.83}{-8.77}\hullline{0.90}{-8.77}\hullfill
\hullface{keel!57!black}\hullmove{-0.90}{-8.77}\hullline{-0.83}{-8.77}\hullline{-2.21}{-8.25}\hullline{-2.38}{-8.25}\hullfill
\hullface{keel!57!black}\hullmove{4.03}{-7.57}\hullline{3.77}{-7.57}\hullline{2.21}{-8.25}\hullline{2.38}{-8.25}\hullfill
\hullface{keel!57!black}\hullmove{-2.38}{-8.25}\hullline{-2.21}{-8.25}\hullline{-3.77}{-7.57}\hullline{-4.03}{-7.57}\hullfill
\hullface{keel!57!black}\hullmove{5.63}{-6.77}\hullline{5.31}{-6.77}\hullline{3.77}{-7.57}\hullline{4.03}{-7.57}\hullfill
\hullface{keel!57!black}\hullmove{-4.03}{-7.57}\hullline{-3.77}{-7.57}\hullline{-5.31}{-6.77}\hullline{-5.63}{-6.77}\hullfill
\hullface{keel!57!black}\hullmove{7.02}{-5.86}\hullline{6.68}{-5.86}\hullline{5.31}{-6.77}\hullline{5.63}{-6.77}\hullfill
\hullface{keel!57!black}\hullmove{-5.63}{-6.77}\hullline{-5.31}{-6.77}\hullline{-6.68}{-5.86}\hullline{-7.02}{-5.86}\hullfill
\hullface{keel!57!black}\hullmove{7.53}{-5.44}\hullline{7.19}{-5.44}\hullline{6.68}{-5.86}\hullline{7.02}{-5.86}\hullfill
\hullface{keel!60!black}\hullmove{-7.02}{-5.86}\hullline{-6.68}{-5.86}\hullline{-7.19}{-5.44}\hullline{-7.53}{-5.44}\hullfill
\hullface{keel!57!black}\hullmove{8.11}{-4.85}\hullline{7.79}{-4.85}\hullline{7.19}{-5.44}\hullline{7.53}{-5.44}\hullfill
\hullface{keel!66!black}\hullmove{-7.53}{-5.44}\hullline{-7.19}{-5.44}\hullline{-7.79}{-4.85}\hullline{-8.11}{-4.85}\hullfill
\hullface{keel!57!black}\hullmove{8.87}{-3.75}\hullline{8.58}{-3.75}\hullline{7.79}{-4.85}\hullline{8.11}{-4.85}\hullfill
\hullface{keel!72!black}\hullmove{-8.11}{-4.85}\hullline{-7.79}{-4.85}\hullline{-8.58}{-3.75}\hullline{-8.87}{-3.75}\hullfill
\hullface{keel!57!black}\hullmove{9.30}{-2.58}\hullline{9.07}{-2.58}\hullline{8.58}{-3.75}\hullline{8.87}{-3.75}\hullfill
\hullface{keel!81!black}\hullmove{-8.87}{-3.75}\hullline{-8.58}{-3.75}\hullline{-9.07}{-2.58}\hullline{-9.30}{-2.58}\hullfill
\hullface{keel!57!black}\hullmove{9.47}{-1.32}\hullline{9.27}{-1.32}\hullline{9.07}{-2.58}\hullline{9.30}{-2.58}\hullfill
\hullface{keel!87!black}\hullmove{-9.30}{-2.58}\hullline{-9.07}{-2.58}\hullline{-9.27}{-1.32}\hullline{-9.47}{-1.32}\hullfill
\hullface{keel!57!black}\hullmove{9.50}{0.00}\hullline{9.31}{0.00}\hullline{9.27}{-1.32}\hullline{9.47}{-1.32}\hullfill
\hullface{keel!90!black}\hullmove{-9.47}{-1.32}\hullline{-9.27}{-1.32}\hullline{-9.31}{0.00}\hullline{-9.50}{0.00}\hullfill
\hullface{topsides!57!black}\hullmove{9.63}{0.78}\hullline{9.46}{0.78}\hullline{9.31}{0.00}\hullline{9.50}{0.00}\hullfill
\hullface{topsides!84!black}\hullmove{-9.50}{0.00}\hullline{-9.31}{0.00}\hullline{-9.46}{0.78}\hullline{-9.63}{0.78}\hullfill
\hullface{topsides!57!black}\hullmove{9.85}{1.56}\hullline{9.71}{1.56}\hullline{9.46}{0.78}\hullline{9.63}{0.78}\hullfill
\hullface{topsides!78!black}\hullmove{-9.63}{0.78}\hullline{-9.46}{0.78}\hullline{-9.71}{1.56}\hullline{-9.85}{1.56}\hullfill
\hullface{topsides!57!black}\hullmove{10.12}{2.34}\hullline{10.02}{2.34}\hullline{9.71}{1.56}\hullline{9.85}{1.56}\hullfill
\hullface{topsides!57!black}\hullmove{10.44}{3.12}\hullline{10.37}{3.12}\hullline{10.02}{2.34}\hullline{10.12}{2.34}\hullfill
\hullface{topsides!57!black}\hullmove{10.78}{3.90}\hullline{10.76}{3.90}\hullline{10.37}{3.12}\hullline{10.44}{3.12}\hullfill
\hullface{topsides!66!black}\hullmove{-10.44}{3.12}\hullline{-10.37}{3.12}\hullline{-10.76}{3.90}\hullline{-10.78}{3.90}\hullfill
\hullface{keel!57!black}\hullmove{0.00}{-9.06}\hullline{0.00}{-9.06}\hullline{-0.83}{-8.77}\hullline{-0.83}{-8.77}\hullfill
\hullface{keel!57!black}\hullmove{2.21}{-8.25}\hullline{2.21}{-8.25}\hullline{0.83}{-8.77}\hullline{0.83}{-8.77}\hullfill
\hullface{keel!57!black}\hullmove{-0.83}{-8.77}\hullline{-0.83}{-8.77}\hullline{-2.21}{-8.25}\hullline{-2.21}{-8.25}\hullfill
\hullface{keel!57!black}\hullmove{3.77}{-7.57}\hullline{3.77}{-7.57}\hullline{2.21}{-8.25}\hullline{2.21}{-8.25}\hullfill
\hullface{keel!57!black}\hullmove{-2.21}{-8.25}\hullline{-2.21}{-8.25}\hullline{-3.77}{-7.57}\hullline{-3.77}{-7.57}\hullfill
\hullface{keel!57!black}\hullmove{5.31}{-6.77}\hullline{5.31}{-6.77}\hullline{3.77}{-7.57}\hullline{3.77}{-7.57}\hullfill
\hullface{keel!57!black}\hullmove{-3.77}{-7.57}\hullline{-3.77}{-7.57}\hullline{-5.31}{-6.77}\hullline{-5.31}{-6.77}\hullfill
\hullface{keel!57!black}\hullmove{6.68}{-5.86}\hullline{6.68}{-5.86}\hullline{5.31}{-6.77}\hullline{5.31}{-6.77}\hullfill
\hullface{keel!57!black}\hullmove{-5.31}{-6.77}\hullline{-5.31}{-6.77}\hullline{-6.68}{-5.86}\hullline{-6.68}{-5.86}\hullfill
\hullface{keel!57!black}\hullmove{7.19}{-5.44}\hullline{7.19}{-5.44}\hullline{6.68}{-5.86}\hullline{6.68}{-5.86}\hullfill
\hullface{keel!60!black}\hullmove{-6.68}{-5.86}\hullline{-6.68}{-5.86}\hullline{-7.19}{-5.44}\hullline{-7.19}{-5.44}\hullfill
\hullface{keel!57!black}\hullmove{7.79}{-4.85}\hullline{7.79}{-4.85}\hullline{7.19}{-5.44}\hullline{7.19}{-5.44}\hullfill
\hullface{keel!66!black}\hullmove{-7.19}{-5.44}\hullline{-7.19}{-5.44}\hullline{-7.79}{-4.85}\hullline{-7.79}{-4.85}\hullfill
\hullface{keel!57!black}\hullmove{8.58}{-3.75}\hullline{8.58}{-3.75}\hullline{7.79}{-4.85}\hullline{7.79}{-4.85}\hullfill
\hullface{keel!72!black}\hullmove{-7.79}{-4.85}\hullline{-7.79}{-4.85}\hullline{-8.58}{-3.75}\hullline{-8.58}{-3.75}\hullfill
\hullface{keel!57!black}\hullmove{9.07}{-2.58}\hullline{9.07}{-2.58}\hullline{8.58}{-3.75}\hullline{8.58}{-3.75}\hullfill
\hullface{keel!81!black}\hullmove{-8.58}{-3.75}\hullline{-8.58}{-3.75}\hullline{-9.07}{-2.58}\hullline{-9.07}{-2.58}\hullfill
\hullface{keel!57!black}\hullmove{9.27}{-1.32}\hullline{9.27}{-1.32}\hullline{9.07}{-2.58}\hullline{9.07}{-2.58}\hullfill
\hullface{keel!87!black}\hullmove{-9.07}{-2.58}\hullline{-9.07}{-2.58}\hullline{-9.27}{-1.32}\hullline{-9.27}{-1.32}\hullfill
\hullface{keel!57!black}\hullmove{9.31}{0.00}\hullline{9.31}{0.00}\hullline{9.27}{-1.32}\hullline{9.27}{-1.32}\hullfill
\hullface{keel!90!black}\hullmove{-9.27}{-1.32}\hullline{-9.27}{-1.32}\hullline{-9.31}{0.00}\hullline{-9.31}{0.00}\hullfill
\hullface{topsides!99!black}\hullmove{9.46}{0.78}\hullline{9.71}{1.56}\hullline{10.02}{2.34}\hullline{10.37}{3.12}\hullline{10.76}{3.90}\hullline{10.76}{7.59}\hullline{10.37}{6.07}\hullline{10.02}{4.55}\hullline{9.71}{3.04}\hullline{9.46}{1.52}\hullline{9.31}{0.00}\hullline{9.31}{0.00}\hullfill
\hullface{topsides!99!black}\hullmove{-9.31}{0.00}\hullline{-9.31}{0.00}\hullline{-9.46}{1.52}\hullline{-9.71}{3.04}\hullline{-10.02}{4.55}\hullline{-10.37}{6.07}\hullline{-10.76}{7.59}\hullline{-10.76}{3.90}\hullline{-10.37}{3.12}\hullline{-10.02}{2.34}\hullline{-9.71}{1.56}\hullline{-9.46}{0.78}\hullfill
\hullface{keel!57!black}\hullmove{0.00}{-9.06}\hullline{0.00}{-9.06}\hullline{-0.81}{-8.77}\hullline{-0.83}{-8.77}\hullfill
\hullface{keel!57!black}\hullmove{2.21}{-8.25}\hullline{2.16}{-8.25}\hullline{0.81}{-8.77}\hullline{0.83}{-8.77}\hullfill
\hullface{keel!57!black}\hullmove{-0.83}{-8.77}\hullline{-0.81}{-8.77}\hullline{-2.16}{-8.25}\hullline{-2.21}{-8.25}\hullfill
\hullface{keel!57!black}\hullmove{3.77}{-7.57}\hullline{3.69}{-7.57}\hullline{2.16}{-8.25}\hullline{2.21}{-8.25}\hullfill
\hullface{keel!57!black}\hullmove{-2.21}{-8.25}\hullline{-2.16}{-8.25}\hullline{-3.69}{-7.57}\hullline{-3.77}{-7.57}\hullfill
\hullface{keel!57!black}\hullmove{5.31}{-6.77}\hullline{5.21}{-6.77}\hullline{3.69}{-7.57}\hullline{3.77}{-7.57}\hullfill
\hullface{keel!57!black}\hullmove{-3.77}{-7.57}\hullline{-3.69}{-7.57}\hullline{-5.21}{-6.77}\hullline{-5.31}{-6.77}\hullfill
\hullface{keel!57!black}\hullmove{6.68}{-5.86}\hullline{6.57}{-5.86}\hullline{5.21}{-6.77}\hullline{5.31}{-6.77}\hullfill
\hullface{keel!57!black}\hullmove{-5.31}{-6.77}\hullline{-5.21}{-6.77}\hullline{-6.57}{-5.86}\hullline{-6.68}{-5.86}\hullfill
\hullface{keel!57!black}\hullmove{7.19}{-5.44}\hullline{7.08}{-5.44}\hullline{6.57}{-5.86}\hullline{6.68}{-5.86}\hullfill
\hullface{keel!60!black}\hullmove{-6.68}{-5.86}\hullline{-6.57}{-5.86}\hullline{-7.08}{-5.44}\hullline{-7.19}{-5.44}\hullfill
\hullface{keel!57!black}\hullmove{7.79}{-4.85}\hullline{7.68}{-4.85}\hullline{7.08}{-5.44}\hullline{7.19}{-5.44}\hullfill
\hullface{keel!66!black}\hullmove{-7.19}{-5.44}\hullline{-7.08}{-5.44}\hullline{-7.68}{-4.85}\hullline{-7.79}{-4.85}\hullfill
\hullface{keel!57!black}\hullmove{8.58}{-3.75}\hullline{8.49}{-3.75}\hullline{7.68}{-4.85}\hullline{7.79}{-4.85}\hullfill
\hullface{keel!72!black}\hullmove{-7.79}{-4.85}\hullline{-7.68}{-4.85}\hullline{-8.49}{-3.75}\hullline{-8.58}{-3.75}\hullfill
\hullface{keel!57!black}\hullmove{9.07}{-2.58}\hullline{8.99}{-2.58}\hullline{8.49}{-3.75}\hullline{8.58}{-3.75}\hullfill
\hullface{keel!81!black}\hullmove{-8.58}{-3.75}\hullline{-8.49}{-3.75}\hullline{-8.99}{-2.58}\hullline{-9.07}{-2.58}\hullfill
\hullface{keel!57!black}\hullmove{9.27}{-1.32}\hullline{9.21}{-1.32}\hullline{8.99}{-2.58}\hullline{9.07}{-2.58}\hullfill
\hullface{keel!87!black}\hullmove{-9.07}{-2.58}\hullline{-8.99}{-2.58}\hullline{-9.21}{-1.32}\hullline{-9.27}{-1.32}\hullfill
\hullface{keel!57!black}\hullmove{9.31}{0.00}\hullline{9.25}{0.00}\hullline{9.21}{-1.32}\hullline{9.27}{-1.32}\hullfill
\hullface{keel!93!black}\hullmove{-9.27}{-1.32}\hullline{-9.21}{-1.32}\hullline{-9.25}{0.00}\hullline{-9.31}{0.00}\hullfill
\hullface{topsides!57!black}\hullmove{9.46}{1.52}\hullline{9.41}{1.52}\hullline{9.25}{0.00}\hullline{9.31}{0.00}\hullfill
\hullface{topsides!87!black}\hullmove{-9.31}{0.00}\hullline{-9.25}{0.00}\hullline{-9.41}{1.52}\hullline{-9.46}{1.52}\hullfill
\hullface{topsides!57!black}\hullmove{9.71}{3.04}\hullline{9.66}{3.04}\hullline{9.41}{1.52}\hullline{9.46}{1.52}\hullfill
\hullface{topsides!84!black}\hullmove{-9.46}{1.52}\hullline{-9.41}{1.52}\hullline{-9.66}{3.04}\hullline{-9.71}{3.04}\hullfill
\hullface{topsides!57!black}\hullmove{10.02}{4.55}\hullline{9.98}{4.55}\hullline{9.66}{3.04}\hullline{9.71}{3.04}\hullfill
\hullface{topsides!81!black}\hullmove{-9.71}{3.04}\hullline{-9.66}{3.04}\hullline{-9.98}{4.55}\hullline{-10.02}{4.55}\hullfill
\hullface{topsides!57!black}\hullmove{10.37}{6.07}\hullline{10.34}{6.07}\hullline{9.98}{4.55}\hullline{10.02}{4.55}\hullfill
\hullface{topsides!78!black}\hullmove{-10.02}{4.55}\hullline{-9.98}{4.55}\hullline{-10.34}{6.07}\hullline{-10.37}{6.07}\hullfill
\hullface{topsides!57!black}\hullmove{10.76}{7.59}\hullline{10.75}{7.59}\hullline{10.34}{6.07}\hullline{10.37}{6.07}\hullfill
\hullface{topsides!75!black}\hullmove{-10.37}{6.07}\hullline{-10.34}{6.07}\hullline{-10.75}{7.59}\hullline{-10.76}{7.59}\hullfill
\hullface{keel!57!black}\hullmove{0.81}{-8.77}\hullline{0.73}{-8.77}\hullline{0.00}{-9.06}\hullline{0.00}{-9.06}\hullfill
\hullface{keel!57!black}\hullmove{0.00}{-9.06}\hullline{0.00}{-9.06}\hullline{-0.73}{-8.77}\hullline{-0.81}{-8.77}\hullfill
\hullface{keel!57!black}\hullmove{2.16}{-8.25}\hullline{1.96}{-8.25}\hullline{0.73}{-8.77}\hullline{0.81}{-8.77}\hullfill
\hullface{keel!57!black}\hullmove{-0.81}{-8.77}\hullline{-0.73}{-8.77}\hullline{-1.96}{-8.25}\hullline{-2.16}{-8.25}\hullfill
\hullface{keel!57!black}\hullmove{3.69}{-7.57}\hullline{3.38}{-7.57}\hullline{1.96}{-8.25}\hullline{2.16}{-8.25}\hullfill
\hullface{keel!57!black}\hullmove{-2.16}{-8.25}\hullline{-1.96}{-8.25}\hullline{-3.38}{-7.57}\hullline{-3.69}{-7.57}\hullfill
\hullface{keel!57!black}\hullmove{5.21}{-6.77}\hullline{4.81}{-6.77}\hullline{3.38}{-7.57}\hullline{3.69}{-7.57}\hullfill
\hullface{keel!57!black}\hullmove{-3.69}{-7.57}\hullline{-3.38}{-7.57}\hullline{-4.81}{-6.77}\hullline{-5.21}{-6.77}\hullfill
\hullface{keel!57!black}\hullmove{6.57}{-5.86}\hullline{6.12}{-5.86}\hullline{4.81}{-6.77}\hullline{5.21}{-6.77}\hullfill
\hullface{keel!57!black}\hullmove{-5.21}{-6.77}\hullline{-4.81}{-6.77}\hullline{-6.12}{-5.86}\hullline{-6.57}{-5.86}\hullfill
\hullface{keel!57!black}\hullmove{7.08}{-5.44}\hullline{6.63}{-5.44}\hullline{6.12}{-5.86}\hullline{6.57}{-5.86}\hullfill
\hullface{keel!60!black}\hullmove{-6.57}{-5.86}\hullline{-6.12}{-5.86}\hullline{-6.63}{-5.44}\hullline{-7.08}{-5.44}\hullfill
\hullface{keel!57!black}\hullmove{7.68}{-4.85}\hullline{7.24}{-4.85}\hullline{6.63}{-5.44}\hullline{7.08}{-5.44}\hullfill
\hullface{keel!66!black}\hullmove{-7.08}{-5.44}\hullline{-6.63}{-5.44}\hullline{-7.24}{-4.85}\hullline{-7.68}{-4.85}\hullfill
\hullface{keel!57!black}\hullmove{8.49}{-3.75}\hullline{8.09}{-3.75}\hullline{7.24}{-4.85}\hullline{7.68}{-4.85}\hullfill
\hullface{keel!72!black}\hullmove{-7.68}{-4.85}\hullline{-7.24}{-4.85}\hullline{-8.09}{-3.75}\hullline{-8.49}{-3.75}\hullfill
\hullface{keel!57!black}\hullmove{8.99}{-2.58}\hullline{8.64}{-2.58}\hullline{8.09}{-3.75}\hullline{8.49}{-3.75}\hullfill
\hullface{keel!81!black}\hullmove{-8.49}{-3.75}\hullline{-8.09}{-3.75}\hullline{-8.64}{-2.58}\hullline{-8.99}{-2.58}\hullfill
\hullface{keel!57!black}\hullmove{9.21}{-1.32}\hullline{8.92}{-1.32}\hullline{8.64}{-2.58}\hullline{8.99}{-2.58}\hullfill
\hullface{keel!87!black}\hullmove{-8.99}{-2.58}\hullline{-8.64}{-2.58}\hullline{-8.92}{-1.32}\hullline{-9.21}{-1.32}\hullfill
\hullface{keel!57!black}\hullmove{9.25}{0.00}\hullline{8.97}{0.00}\hullline{8.92}{-1.32}\hullline{9.21}{-1.32}\hullfill
\hullface{keel!93!black}\hullmove{-9.21}{-1.32}\hullline{-8.92}{-1.32}\hullline{-8.97}{0.00}\hullline{-9.25}{0.00}\hullfill
\hullface{topsides!57!black}\hullmove{9.41}{1.52}\hullline{9.15}{1.52}\hullline{8.97}{0.00}\hullline{9.25}{0.00}\hullfill
\hullface{topsides!90!black}\hullmove{-9.25}{0.00}\hullline{-8.97}{0.00}\hullline{-9.15}{1.52}\hullline{-9.41}{1.52}\hullfill
\hullface{topsides!57!black}\hullmove{9.66}{3.04}\hullline{9.45}{3.04}\hullline{9.15}{1.52}\hullline{9.41}{1.52}\hullfill
\hullface{topsides!84!black}\hullmove{-9.41}{1.52}\hullline{-9.15}{1.52}\hullline{-9.45}{3.04}\hullline{-9.66}{3.04}\hullfill
\hullface{topsides!57!black}\hullmove{9.98}{4.55}\hullline{9.82}{4.55}\hullline{9.45}{3.04}\hullline{9.66}{3.04}\hullfill
\hullface{topsides!81!black}\hullmove{-9.66}{3.04}\hullline{-9.45}{3.04}\hullline{-9.82}{4.55}\hullline{-9.98}{4.55}\hullfill
\hullface{topsides!57!black}\hullmove{10.34}{6.07}\hullline{10.24}{6.07}\hullline{9.82}{4.55}\hullline{9.98}{4.55}\hullfill
\hullface{topsides!78!black}\hullmove{-9.98}{4.55}\hullline{-9.82}{4.55}\hullline{-10.24}{6.07}\hullline{-10.34}{6.07}\hullfill
\hullface{topsides!57!black}\hullmove{10.75}{7.59}\hullline{10.70}{7.59}\hullline{10.24}{6.07}\hullline{10.34}{6.07}\hullfill
\hullface{topsides!75!black}\hullmove{-10.34}{6.07}\hullline{-10.24}{6.07}\hullline{-10.70}{7.59}\hullline{-10.75}{7.59}\hullfill
\hullface{keel!57!black}\hullmove{0.73}{-8.77}\hullline{0.65}{-8.77}\hullline{0.00}{-9.06}\hullline{0.00}{-9.06}\hullfill
\hullface{keel!57!black}\hullmove{0.00}{-9.06}\hullline{0.00}{-9.06}\hullline{-0.65}{-8.77}\hullline{-0.73}{-8.77}\hullfill
\hullface{keel!57!black}\hullmove{1.96}{-8.25}\hullline{1.77}{-8.25}\hullline{0.65}{-8.77}\hullline{0.73}{-8.77}\hullfill
\hullface{keel!57!black}\hullmove{-0.73}{-8.77}\hullline{-0.65}{-8.77}\hullline{-1.77}{-8.25}\hullline{-1.96}{-8.25}\hullfill
\hullface{keel!57!black}\hullmove{3.38}{-7.57}\hullline{3.08}{-7.57}\hullline{1.77}{-8.25}\hullline{1.96}{-8.25}\hullfill
\hullface{keel!57!black}\hullmove{-1.96}{-8.25}\hullline{-1.77}{-8.25}\hullline{-3.08}{-7.57}\hullline{-3.38}{-7.57}\hullfill
\hullface{keel!57!black}\hullmove{4.81}{-6.77}\hullline{4.42}{-6.77}\hullline{3.08}{-7.57}\hullline{3.38}{-7.57}\hullfill
\hullface{keel!57!black}\hullmove{-3.38}{-7.57}\hullline{-3.08}{-7.57}\hullline{-4.42}{-6.77}\hullline{-4.81}{-6.77}\hullfill
\hullface{keel!57!black}\hullmove{6.12}{-5.86}\hullline{5.68}{-5.86}\hullline{4.42}{-6.77}\hullline{4.81}{-6.77}\hullfill
\hullface{keel!57!black}\hullmove{-4.81}{-6.77}\hullline{-4.42}{-6.77}\hullline{-5.68}{-5.86}\hullline{-6.12}{-5.86}\hullfill
\hullface{keel!57!black}\hullmove{6.63}{-5.44}\hullline{6.18}{-5.44}\hullline{5.68}{-5.86}\hullline{6.12}{-5.86}\hullfill
\hullface{keel!60!black}\hullmove{-6.12}{-5.86}\hullline{-5.68}{-5.86}\hullline{-6.18}{-5.44}\hullline{-6.63}{-5.44}\hullfill
\hullface{keel!57!black}\hullmove{7.24}{-4.85}\hullline{6.79}{-4.85}\hullline{6.18}{-5.44}\hullline{6.63}{-5.44}\hullfill
\hullface{keel!66!black}\hullmove{-6.63}{-5.44}\hullline{-6.18}{-5.44}\hullline{-6.79}{-4.85}\hullline{-7.24}{-4.85}\hullfill
\hullface{keel!57!black}\hullmove{8.09}{-3.75}\hullline{7.66}{-3.75}\hullline{6.79}{-4.85}\hullline{7.24}{-4.85}\hullfill
\hullface{keel!72!black}\hullmove{-7.24}{-4.85}\hullline{-6.79}{-4.85}\hullline{-7.66}{-3.75}\hullline{-8.09}{-3.75}\hullfill
\hullface{keel!57!black}\hullmove{8.64}{-2.58}\hullline{8.27}{-2.58}\hullline{7.66}{-3.75}\hullline{8.09}{-3.75}\hullfill
\hullface{keel!81!black}\hullmove{-8.09}{-3.75}\hullline{-7.66}{-3.75}\hullline{-8.27}{-2.58}\hullline{-8.64}{-2.58}\hullfill
\hullface{keel!57!black}\hullmove{8.92}{-1.32}\hullline{8.59}{-1.32}\hullline{8.27}{-2.58}\hullline{8.64}{-2.58}\hullfill
\hullface{keel!87!black}\hullmove{-8.64}{-2.58}\hullline{-8.27}{-2.58}\hullline{-8.59}{-1.32}\hullline{-8.92}{-1.32}\hullfill
\hullface{keel!57!black}\hullmove{8.97}{0.00}\hullline{8.67}{0.00}\hullline{8.59}{-1.32}\hullline{8.92}{-1.32}\hullfill
\hullface{keel!93!black}\hullmove{-8.92}{-1.32}\hullline{-8.59}{-1.32}\hullline{-8.67}{0.00}\hullline{-8.97}{0.00}\hullfill
\hullface{topsides!57!black}\hullmove{9.15}{1.52}\hullline{8.88}{1.52}\hullline{8.67}{0.00}\hullline{8.97}{0.00}\hullfill
\hullface{topsides!90!black}\hullmove{-8.97}{0.00}\hullline{-8.67}{0.00}\hullline{-8.88}{1.52}\hullline{-9.15}{1.52}\hullfill
\hullface{topsides!57!black}\hullmove{9.45}{3.04}\hullline{9.22}{3.04}\hullline{8.88}{1.52}\hullline{9.15}{1.52}\hullfill
\hullface{topsides!84!black}\hullmove{-9.15}{1.52}\hullline{-8.88}{1.52}\hullline{-9.22}{3.04}\hullline{-9.45}{3.04}\hullfill
\hullface{topsides!57!black}\hullmove{9.82}{4.55}\hullline{9.63}{4.55}\hullline{9.22}{3.04}\hullline{9.45}{3.04}\hullfill
\hullface{topsides!81!black}\hullmove{-9.45}{3.04}\hullline{-9.22}{3.04}\hullline{-9.63}{4.55}\hullline{-9.82}{4.55}\hullfill
\hullface{topsides!57!black}\hullmove{10.24}{6.07}\hullline{10.11}{6.07}\hullline{9.63}{4.55}\hullline{9.82}{4.55}\hullfill
\hullface{topsides!78!black}\hullmove{-9.82}{4.55}\hullline{-9.63}{4.55}\hullline{-10.11}{6.07}\hullline{-10.24}{6.07}\hullfill
\hullface{topsides!57!black}\hullmove{10.70}{7.59}\hullline{10.64}{7.59}\hullline{10.11}{6.07}\hullline{10.24}{6.07}\hullfill
\hullface{topsides!72!black}\hullmove{-10.24}{6.07}\hullline{-10.11}{6.07}\hullline{-10.64}{7.59}\hullline{-10.70}{7.59}\hullfill
\hullface{keel!57!black}\hullmove{0.00}{-9.06}\hullline{0.00}{-9.06}\hullline{-0.59}{-8.77}\hullline{-0.65}{-8.77}\hullfill
\hullface{keel!57!black}\hullmove{1.77}{-8.25}\hullline{1.60}{-8.25}\hullline{0.59}{-8.77}\hullline{0.65}{-8.77}\hullfill
\hullface{keel!57!black}\hullmove{-0.65}{-8.77}\hullline{-0.59}{-8.77}\hullline{-1.60}{-8.25}\hullline{-1.77}{-8.25}\hullfill
\hullface{keel!57!black}\hullmove{3.08}{-7.57}\hullline{2.80}{-7.57}\hullline{1.60}{-8.25}\hullline{1.77}{-8.25}\hullfill
\hullface{keel!57!black}\hullmove{-1.77}{-8.25}\hullline{-1.60}{-8.25}\hullline{-2.80}{-7.57}\hullline{-3.08}{-7.57}\hullfill
\hullface{keel!57!black}\hullmove{4.42}{-6.77}\hullline{4.05}{-6.77}\hullline{2.80}{-7.57}\hullline{3.08}{-7.57}\hullfill
\hullface{keel!57!black}\hullmove{-3.08}{-7.57}\hullline{-2.80}{-7.57}\hullline{-4.05}{-6.77}\hullline{-4.42}{-6.77}\hullfill
\hullface{keel!57!black}\hullmove{5.68}{-5.86}\hullline{5.25}{-5.86}\hullline{4.05}{-6.77}\hullline{4.42}{-6.77}\hullfill
\hullface{keel!57!black}\hullmove{-4.42}{-6.77}\hullline{-4.05}{-6.77}\hullline{-5.25}{-5.86}\hullline{-5.68}{-5.86}\hullfill
\hullface{keel!57!black}\hullmove{6.18}{-5.44}\hullline{5.73}{-5.44}\hullline{5.25}{-5.86}\hullline{5.68}{-5.86}\hullfill
\hullface{keel!63!black}\hullmove{-5.68}{-5.86}\hullline{-5.25}{-5.86}\hullline{-5.73}{-5.44}\hullline{-6.18}{-5.44}\hullfill
\hullface{keel!57!black}\hullmove{6.79}{-4.85}\hullline{6.33}{-4.85}\hullline{5.73}{-5.44}\hullline{6.18}{-5.44}\hullfill
\hullface{keel!66!black}\hullmove{-6.18}{-5.44}\hullline{-5.73}{-5.44}\hullline{-6.33}{-4.85}\hullline{-6.79}{-4.85}\hullfill
\hullface{keel!57!black}\hullmove{7.66}{-3.75}\hullline{7.22}{-3.75}\hullline{6.33}{-4.85}\hullline{6.79}{-4.85}\hullfill
\hullface{keel!72!black}\hullmove{-6.79}{-4.85}\hullline{-6.33}{-4.85}\hullline{-7.22}{-3.75}\hullline{-7.66}{-3.75}\hullfill
\hullface{keel!57!black}\hullmove{8.27}{-2.58}\hullline{7.86}{-2.58}\hullline{7.22}{-3.75}\hullline{7.66}{-3.75}\hullfill
\hullface{keel!81!black}\hullmove{-7.66}{-3.75}\hullline{-7.22}{-3.75}\hullline{-7.86}{-2.58}\hullline{-8.27}{-2.58}\hullfill
\hullface{keel!57!black}\hullmove{8.59}{-1.32}\hullline{8.23}{-1.32}\hullline{7.86}{-2.58}\hullline{8.27}{-2.58}\hullfill
\hullface{keel!87!black}\hullmove{-8.27}{-2.58}\hullline{-7.86}{-2.58}\hullline{-8.23}{-1.32}\hullline{-8.59}{-1.32}\hullfill
\hullface{keel!57!black}\hullmove{8.67}{0.00}\hullline{8.34}{0.00}\hullline{8.23}{-1.32}\hullline{8.59}{-1.32}\hullfill
\hullface{keel!93!black}\hullmove{-8.59}{-1.32}\hullline{-8.23}{-1.32}\hullline{-8.34}{0.00}\hullline{-8.67}{0.00}\hullfill
\hullface{topsides!57!black}\hullmove{8.88}{1.52}\hullline{8.57}{1.52}\hullline{8.34}{0.00}\hullline{8.67}{0.00}\hullfill
\hullface{topsides!90!black}\hullmove{-8.67}{0.00}\hullline{-8.34}{0.00}\hullline{-8.57}{1.52}\hullline{-8.88}{1.52}\hullfill
\hullface{topsides!57!black}\hullmove{9.22}{3.04}\hullline{8.96}{3.04}\hullline{8.57}{1.52}\hullline{8.88}{1.52}\hullfill
\hullface{topsides!84!black}\hullmove{-8.88}{1.52}\hullline{-8.57}{1.52}\hullline{-8.96}{3.04}\hullline{-9.22}{3.04}\hullfill
\hullface{topsides!57!black}\hullmove{9.63}{4.55}\hullline{9.43}{4.55}\hullline{8.96}{3.04}\hullline{9.22}{3.04}\hullfill
\hullface{topsides!81!black}\hullmove{-9.22}{3.04}\hullline{-8.96}{3.04}\hullline{-9.43}{4.55}\hullline{-9.63}{4.55}\hullfill
\hullface{topsides!57!black}\hullmove{10.11}{6.07}\hullline{9.97}{6.07}\hullline{9.43}{4.55}\hullline{9.63}{4.55}\hullfill
\hullface{topsides!78!black}\hullmove{-9.63}{4.55}\hullline{-9.43}{4.55}\hullline{-9.97}{6.07}\hullline{-10.11}{6.07}\hullfill
\hullface{topsides!57!black}\hullmove{10.64}{7.59}\hullline{10.57}{7.59}\hullline{9.97}{6.07}\hullline{10.11}{6.07}\hullfill
\hullface{topsides!72!black}\hullmove{-10.11}{6.07}\hullline{-9.97}{6.07}\hullline{-10.57}{7.59}\hullline{-10.64}{7.59}\hullfill
\hullface{keel!57!black}\hullmove{0.59}{-8.77}\hullline{0.53}{-8.77}\hullline{0.00}{-9.06}\hullline{0.00}{-9.06}\hullfill
\hullface{keel!57!black}\hullmove{0.00}{-9.06}\hullline{0.00}{-9.06}\hullline{-0.53}{-8.77}\hullline{-0.59}{-8.77}\hullfill
\hullface{keel!57!black}\hullmove{1.60}{-8.25}\hullline{1.44}{-8.25}\hullline{0.53}{-8.77}\hullline{0.59}{-8.77}\hullfill
\hullface{keel!57!black}\hullmove{-0.59}{-8.77}\hullline{-0.53}{-8.77}\hullline{-1.44}{-8.25}\hullline{-1.60}{-8.25}\hullfill
\hullface{keel!57!black}\hullmove{2.80}{-7.57}\hullline{2.53}{-7.57}\hullline{1.44}{-8.25}\hullline{1.60}{-8.25}\hullfill
\hullface{keel!57!black}\hullmove{-1.60}{-8.25}\hullline{-1.44}{-8.25}\hullline{-2.53}{-7.57}\hullline{-2.80}{-7.57}\hullfill
\hullface{keel!57!black}\hullmove{4.05}{-6.77}\hullline{3.69}{-6.77}\hullline{2.53}{-7.57}\hullline{2.80}{-7.57}\hullfill
\hullface{keel!57!black}\hullmove{-2.80}{-7.57}\hullline{-2.53}{-7.57}\hullline{-3.69}{-6.77}\hullline{-4.05}{-6.77}\hullfill
\hullface{keel!57!black}\hullmove{5.25}{-5.86}\hullline{4.83}{-5.86}\hullline{3.69}{-6.77}\hullline{4.05}{-6.77}\hullfill
\hullface{keel!57!black}\hullmove{-4.05}{-6.77}\hullline{-3.69}{-6.77}\hullline{-4.83}{-5.86}\hullline{-5.25}{-5.86}\hullfill
\hullface{keel!57!black}\hullmove{5.73}{-5.44}\hullline{5.29}{-5.44}\hullline{4.83}{-5.86}\hullline{5.25}{-5.86}\hullfill
\hullface{keel!63!black}\hullmove{-5.25}{-5.86}\hullline{-4.83}{-5.86}\hullline{-5.29}{-5.44}\hullline{-5.73}{-5.44}\hullfill
\hullface{keel!57!black}\hullmove{6.33}{-4.85}\hullline{5.87}{-4.85}\hullline{5.29}{-5.44}\hullline{5.73}{-5.44}\hullfill
\hullface{keel!66!black}\hullmove{-5.73}{-5.44}\hullline{-5.29}{-5.44}\hullline{-5.87}{-4.85}\hullline{-6.33}{-4.85}\hullfill
\hullface{keel!57!black}\hullmove{7.22}{-3.75}\hullline{6.75}{-3.75}\hullline{5.87}{-4.85}\hullline{6.33}{-4.85}\hullfill
\hullface{keel!72!black}\hullmove{-6.33}{-4.85}\hullline{-5.87}{-4.85}\hullline{-6.75}{-3.75}\hullline{-7.22}{-3.75}\hullfill
\hullface{keel!57!black}\hullmove{7.86}{-2.58}\hullline{7.42}{-2.58}\hullline{6.75}{-3.75}\hullline{7.22}{-3.75}\hullfill
\hullface{keel!78!black}\hullmove{-7.22}{-3.75}\hullline{-6.75}{-3.75}\hullline{-7.42}{-2.58}\hullline{-7.86}{-2.58}\hullfill
\hullface{keel!57!black}\hullmove{8.23}{-1.32}\hullline{7.84}{-1.32}\hullline{7.42}{-2.58}\hullline{7.86}{-2.58}\hullfill
\hullface{keel!87!black}\hullmove{-7.86}{-2.58}\hullline{-7.42}{-2.58}\hullline{-7.84}{-1.32}\hullline{-8.23}{-1.32}\hullfill
\hullface{keel!57!black}\hullmove{8.34}{0.00}\hullline{7.97}{0.00}\hullline{7.84}{-1.32}\hullline{8.23}{-1.32}\hullfill
\hullface{keel!93!black}\hullmove{-8.23}{-1.32}\hullline{-7.84}{-1.32}\hullline{-7.97}{0.00}\hullline{-8.34}{0.00}\hullfill
\hullface{topsides!57!black}\hullmove{8.57}{1.52}\hullline{8.23}{1.52}\hullline{7.97}{0.00}\hullline{8.34}{0.00}\hullfill
\hullface{topsides!90!black}\hullmove{-8.34}{0.00}\hullline{-7.97}{0.00}\hullline{-8.23}{1.52}\hullline{-8.57}{1.52}\hullfill
\hullface{topsides!57!black}\hullmove{8.96}{3.04}\hullline{8.67}{3.04}\hullline{8.23}{1.52}\hullline{8.57}{1.52}\hullfill
\hullface{topsides!84!black}\hullmove{-8.57}{1.52}\hullline{-8.23}{1.52}\hullline{-8.67}{3.04}\hullline{-8.96}{3.04}\hullfill
\hullface{topsides!57!black}\hullmove{9.43}{4.55}\hullline{9.20}{4.55}\hullline{8.67}{3.04}\hullline{8.96}{3.04}\hullfill
\hullface{topsides!81!black}\hullmove{-8.96}{3.04}\hullline{-8.67}{3.04}\hullline{-9.20}{4.55}\hullline{-9.43}{4.55}\hullfill
\hullface{topsides!57!black}\hullmove{9.97}{6.07}\hullline{9.80}{6.07}\hullline{9.20}{4.55}\hullline{9.43}{4.55}\hullfill
\hullface{topsides!75!black}\hullmove{-9.43}{4.55}\hullline{-9.20}{4.55}\hullline{-9.80}{6.07}\hullline{-9.97}{6.07}\hullfill
\hullface{topsides!57!black}\hullmove{10.57}{7.59}\hullline{10.47}{7.59}\hullline{9.80}{6.07}\hullline{9.97}{6.07}\hullfill
\hullface{topsides!72!black}\hullmove{-9.97}{6.07}\hullline{-9.80}{6.07}\hullline{-10.47}{7.59}\hullline{-10.57}{7.59}\hullfill
\hullface{keel!57!black}\hullmove{0.53}{-8.77}\hullline{0.47}{-8.77}\hullline{0.00}{-9.06}\hullline{0.00}{-9.06}\hullfill
\hullface{keel!57!black}\hullmove{0.00}{-9.06}\hullline{0.00}{-9.06}\hullline{-0.47}{-8.77}\hullline{-0.53}{-8.77}\hullfill
\hullface{keel!57!black}\hullmove{1.44}{-8.25}\hullline{1.29}{-8.25}\hullline{0.47}{-8.77}\hullline{0.53}{-8.77}\hullfill
\hullface{keel!57!black}\hullmove{-0.53}{-8.77}\hullline{-0.47}{-8.77}\hullline{-1.29}{-8.25}\hullline{-1.44}{-8.25}\hullfill
\hullface{keel!57!black}\hullmove{2.53}{-7.57}\hullline{2.28}{-7.57}\hullline{1.29}{-8.25}\hullline{1.44}{-8.25}\hullfill
\hullface{keel!57!black}\hullmove{-1.44}{-8.25}\hullline{-1.29}{-8.25}\hullline{-2.28}{-7.57}\hullline{-2.53}{-7.57}\hullfill
\hullface{keel!57!black}\hullmove{3.69}{-6.77}\hullline{3.35}{-6.77}\hullline{2.28}{-7.57}\hullline{2.53}{-7.57}\hullfill
\hullface{keel!57!black}\hullmove{-2.53}{-7.57}\hullline{-2.28}{-7.57}\hullline{-3.35}{-6.77}\hullline{-3.69}{-6.77}\hullfill
\hullface{keel!57!black}\hullmove{4.83}{-5.86}\hullline{4.41}{-5.86}\hullline{3.35}{-6.77}\hullline{3.69}{-6.77}\hullfill
\hullface{keel!60!black}\hullmove{-3.69}{-6.77}\hullline{-3.35}{-6.77}\hullline{-4.41}{-5.86}\hullline{-4.83}{-5.86}\hullfill
\hullface{keel!57!black}\hullmove{5.29}{-5.44}\hullline{4.85}{-5.44}\hullline{4.41}{-5.86}\hullline{4.83}{-5.86}\hullfill
\hullface{keel!63!black}\hullmove{-4.83}{-5.86}\hullline{-4.41}{-5.86}\hullline{-4.85}{-5.44}\hullline{-5.29}{-5.44}\hullfill
\hullface{keel!57!black}\hullmove{5.87}{-4.85}\hullline{5.40}{-4.85}\hullline{4.85}{-5.44}\hullline{5.29}{-5.44}\hullfill
\hullface{keel!66!black}\hullmove{-5.29}{-5.44}\hullline{-4.85}{-5.44}\hullline{-5.40}{-4.85}\hullline{-5.87}{-4.85}\hullfill
\hullface{keel!57!black}\hullmove{6.75}{-3.75}\hullline{6.27}{-3.75}\hullline{5.40}{-4.85}\hullline{5.87}{-4.85}\hullfill
\hullface{keel!72!black}\hullmove{-5.87}{-4.85}\hullline{-5.40}{-4.85}\hullline{-6.27}{-3.75}\hullline{-6.75}{-3.75}\hullfill
\hullface{keel!57!black}\hullmove{7.42}{-2.58}\hullline{6.96}{-2.58}\hullline{6.27}{-3.75}\hullline{6.75}{-3.75}\hullfill
\hullface{keel!78!black}\hullmove{-6.75}{-3.75}\hullline{-6.27}{-3.75}\hullline{-6.96}{-2.58}\hullline{-7.42}{-2.58}\hullfill
\hullface{keel!57!black}\hullmove{7.84}{-1.32}\hullline{7.41}{-1.32}\hullline{6.96}{-2.58}\hullline{7.42}{-2.58}\hullfill
\hullface{keel!87!black}\hullmove{-7.42}{-2.58}\hullline{-6.96}{-2.58}\hullline{-7.41}{-1.32}\hullline{-7.84}{-1.32}\hullfill
\hullface{keel!57!black}\hullmove{7.97}{0.00}\hullline{7.57}{0.00}\hullline{7.41}{-1.32}\hullline{7.84}{-1.32}\hullfill
\hullface{keel!93!black}\hullmove{-7.84}{-1.32}\hullline{-7.41}{-1.32}\hullline{-7.57}{0.00}\hullline{-7.97}{0.00}\hullfill
\hullface{topsides!57!black}\hullmove{8.23}{1.52}\hullline{7.87}{1.52}\hullline{7.57}{0.00}\hullline{7.97}{0.00}\hullfill
\hullface{topsides!90!black}\hullmove{-7.97}{0.00}\hullline{-7.57}{0.00}\hullline{-7.87}{1.52}\hullline{-8.23}{1.52}\hullfill
\hullface{topsides!57!black}\hullmove{8.67}{3.04}\hullline{8.35}{3.04}\hullline{7.87}{1.52}\hullline{8.23}{1.52}\hullfill
\hullface{topsides!84!black}\hullmove{-8.23}{1.52}\hullline{-7.87}{1.52}\hullline{-8.35}{3.04}\hullline{-8.67}{3.04}\hullfill
\hullface{topsides!57!black}\hullmove{9.20}{4.55}\hullline{8.94}{4.55}\hullline{8.35}{3.04}\hullline{8.67}{3.04}\hullfill
\hullface{topsides!81!black}\hullmove{-8.67}{3.04}\hullline{-8.35}{3.04}\hullline{-8.94}{4.55}\hullline{-9.20}{4.55}\hullfill
\hullface{topsides!57!black}\hullmove{9.80}{6.07}\hullline{9.61}{6.07}\hullline{8.94}{4.55}\hullline{9.20}{4.55}\hullfill
\hullface{topsides!75!black}\hullmove{-9.20}{4.55}\hullline{-8.94}{4.55}\hullline{-9.61}{6.07}\hullline{-9.80}{6.07}\hullfill
\hullface{topsides!57!black}\hullmove{10.47}{7.59}\hullline{10.36}{7.59}\hullline{9.61}{6.07}\hullline{9.80}{6.07}\hullfill
\hullface{topsides!72!black}\hullmove{-9.80}{6.07}\hullline{-9.61}{6.07}\hullline{-10.36}{7.59}\hullline{-10.47}{7.59}\hullfill
\hullface{keel!57!black}\hullmove{0.47}{-8.77}\hullline{0.42}{-8.77}\hullline{0.00}{-9.06}\hullline{0.00}{-9.06}\hullfill
\hullface{keel!57!black}\hullmove{0.00}{-9.06}\hullline{0.00}{-9.06}\hullline{-0.42}{-8.77}\hullline{-0.47}{-8.77}\hullfill
\hullface{keel!57!black}\hullmove{1.29}{-8.25}\hullline{1.15}{-8.25}\hullline{0.42}{-8.77}\hullline{0.47}{-8.77}\hullfill
\hullface{keel!57!black}\hullmove{-0.47}{-8.77}\hullline{-0.42}{-8.77}\hullline{-1.15}{-8.25}\hullline{-1.29}{-8.25}\hullfill
\hullface{keel!57!black}\hullmove{2.28}{-7.57}\hullline{2.05}{-7.57}\hullline{1.15}{-8.25}\hullline{1.29}{-8.25}\hullfill
\hullface{keel!57!black}\hullmove{-1.29}{-8.25}\hullline{-1.15}{-8.25}\hullline{-2.05}{-7.57}\hullline{-2.28}{-7.57}\hullfill
\hullface{keel!57!black}\hullmove{3.35}{-6.77}\hullline{3.02}{-6.77}\hullline{2.05}{-7.57}\hullline{2.28}{-7.57}\hullfill
\hullface{keel!57!black}\hullmove{-2.28}{-7.57}\hullline{-2.05}{-7.57}\hullline{-3.02}{-6.77}\hullline{-3.35}{-6.77}\hullfill
\hullface{keel!57!black}\hullmove{4.41}{-5.86}\hullline{3.89}{-5.96}\hullline{3.02}{-6.77}\hullline{3.35}{-6.77}\hullfill
\hullface{keel!60!black}\hullmove{-3.35}{-6.77}\hullline{-3.02}{-6.77}\hullline{-3.89}{-5.96}\hullline{-4.41}{-5.86}\hullfill
\hullface{keel!57!black}\hullmove{4.85}{-5.44}\hullline{3.96}{-5.89}\hullline{3.89}{-5.96}\hullline{4.41}{-5.86}\hullfill
\hullface{keel!63!black}\hullmove{-4.41}{-5.86}\hullline{-3.89}{-5.96}\hullline{-3.96}{-5.89}\hullline{-4.85}{-5.44}\hullfill
\hullface{keel!57!black}\hullmove{4.85}{-5.44}\hullline{4.41}{-5.44}\hullline{4.15}{-5.70}\hullline{4.85}{-5.44}\hullfill
\hullface{keel!66!black}\hullmove{-4.85}{-5.44}\hullline{-4.15}{-5.70}\hullline{-4.41}{-5.44}\hullline{-4.85}{-5.44}\hullfill
\hullface{keel!57!black}\hullmove{4.85}{-5.44}\hullline{4.66}{-5.17}\hullline{4.41}{-5.44}\hullline{4.85}{-5.44}\hullfill
\hullface{keel!69!black}\hullmove{-4.85}{-5.44}\hullline{-4.41}{-5.44}\hullline{-4.66}{-5.17}\hullline{-4.85}{-5.44}\hullfill
\hullface{keel!57!black}\hullmove{5.40}{-4.85}\hullline{4.94}{-4.85}\hullline{4.89}{-4.91}\hullline{4.85}{-5.44}\hullfill
\hullface{keel!69!black}\hullmove{-4.85}{-5.44}\hullline{-4.89}{-4.91}\hullline{-4.94}{-4.85}\hullline{-5.40}{-4.85}\hullfill
\hullface{keel!57!black}\hullmove{6.27}{-3.75}\hullline{5.78}{-3.75}\hullline{4.94}{-4.85}\hullline{5.40}{-4.85}\hullfill
\hullface{keel!72!black}\hullmove{-5.40}{-4.85}\hullline{-4.94}{-4.85}\hullline{-5.78}{-3.75}\hullline{-6.27}{-3.75}\hullfill
\hullface{keel!57!black}\hullmove{6.96}{-2.58}\hullline{6.47}{-2.58}\hullline{5.78}{-3.75}\hullline{6.27}{-3.75}\hullfill
\hullface{keel!81!black}\hullmove{-6.27}{-3.75}\hullline{-5.78}{-3.75}\hullline{-6.47}{-2.58}\hullline{-6.96}{-2.58}\hullfill
\hullface{keel!57!black}\hullmove{7.41}{-1.32}\hullline{6.95}{-1.32}\hullline{6.47}{-2.58}\hullline{6.96}{-2.58}\hullfill
\hullface{keel!87!black}\hullmove{-6.96}{-2.58}\hullline{-6.47}{-2.58}\hullline{-6.95}{-1.32}\hullline{-7.41}{-1.32}\hullfill
\hullface{keel!57!black}\hullmove{7.57}{0.00}\hullline{7.14}{0.00}\hullline{6.95}{-1.32}\hullline{7.41}{-1.32}\hullfill
\hullface{keel!93!black}\hullmove{-7.41}{-1.32}\hullline{-6.95}{-1.32}\hullline{-7.14}{0.00}\hullline{-7.57}{0.00}\hullfill
\hullface{topsides!57!black}\hullmove{7.87}{1.52}\hullline{7.47}{1.52}\hullline{7.14}{0.00}\hullline{7.57}{0.00}\hullfill
\hullface{topsides!90!black}\hullmove{-7.57}{0.00}\hullline{-7.14}{0.00}\hullline{-7.47}{1.52}\hullline{-7.87}{1.52}\hullfill
\hullface{topsides!57!black}\hullmove{8.35}{3.04}\hullline{7.99}{3.04}\hullline{7.47}{1.52}\hullline{7.87}{1.52}\hullfill
\hullface{topsides!84!black}\hullmove{-7.87}{1.52}\hullline{-7.47}{1.52}\hullline{-7.99}{3.04}\hullline{-8.35}{3.04}\hullfill
\hullface{topsides!57!black}\hullmove{8.94}{4.55}\hullline{8.65}{4.55}\hullline{7.99}{3.04}\hullline{8.35}{3.04}\hullfill
\hullface{topsides!81!black}\hullmove{-8.35}{3.04}\hullline{-7.99}{3.04}\hullline{-8.65}{4.55}\hullline{-8.94}{4.55}\hullfill
\hullface{topsides!57!black}\hullmove{9.61}{6.07}\hullline{9.39}{6.07}\hullline{8.65}{4.55}\hullline{8.94}{4.55}\hullfill
\hullface{topsides!75!black}\hullmove{-8.94}{4.55}\hullline{-8.65}{4.55}\hullline{-9.39}{6.07}\hullline{-9.61}{6.07}\hullfill
\hullface{topsides!57!black}\hullmove{10.36}{7.59}\hullline{10.21}{7.59}\hullline{9.39}{6.07}\hullline{9.61}{6.07}\hullfill
\hullface{topsides!72!black}\hullmove{-9.61}{6.07}\hullline{-9.39}{6.07}\hullline{-10.21}{7.59}\hullline{-10.36}{7.59}\hullfill
\hullface{keel!57!black}\hullmove{4.85}{-5.44}\hullline{4.85}{-5.44}\hullline{4.15}{-5.70}\hullline{4.00}{-5.86}\hullline{3.96}{-5.89}\hullline{4.85}{-5.44}\hullfill
\hullface{keel!66!black}\hullmove{-4.85}{-5.44}\hullline{-3.96}{-5.89}\hullline{-4.00}{-5.86}\hullline{-4.15}{-5.70}\hullline{-4.85}{-5.44}\hullline{-4.85}{-5.44}\hullfill
\hullface{keel!57!black}\hullmove{4.85}{-5.44}\hullline{4.85}{-5.44}\hullline{4.89}{-4.91}\hullline{4.83}{-4.98}\hullline{4.66}{-5.17}\hullline{4.85}{-5.44}\hullfill
\hullface{keel!69!black}\hullmove{-4.85}{-5.44}\hullline{-4.66}{-5.17}\hullline{-4.83}{-4.98}\hullline{-4.89}{-4.91}\hullline{-4.85}{-5.44}\hullline{-4.85}{-5.44}\hullfill
\hullface{keel!57!black}\hullmove{0.00}{-9.06}\hullline{0.00}{-9.06}\hullline{-0.37}{-8.77}\hullline{-0.42}{-8.77}\hullfill
\hullface{keel!57!black}\hullmove{1.15}{-8.25}\hullline{1.02}{-8.25}\hullline{0.37}{-8.77}\hullline{0.42}{-8.77}\hullfill
\hullface{keel!57!black}\hullmove{-0.42}{-8.77}\hullline{-0.37}{-8.77}\hullline{-1.02}{-8.25}\hullline{-1.15}{-8.25}\hullfill
\hullface{keel!57!black}\hullmove{2.05}{-7.57}\hullline{1.82}{-7.57}\hullline{1.02}{-8.25}\hullline{1.15}{-8.25}\hullfill
\hullface{keel!57!black}\hullmove{-1.15}{-8.25}\hullline{-1.02}{-8.25}\hullline{-1.82}{-7.57}\hullline{-2.05}{-7.57}\hullfill
\hullface{keel!57!black}\hullmove{3.02}{-6.77}\hullline{2.70}{-6.77}\hullline{1.82}{-7.57}\hullline{2.05}{-7.57}\hullfill
\hullface{keel!57!black}\hullmove{-2.05}{-7.57}\hullline{-1.82}{-7.57}\hullline{-2.70}{-6.77}\hullline{-3.02}{-6.77}\hullfill
\hullface{keel!57!black}\hullmove{3.89}{-5.96}\hullline{2.96}{-6.51}\hullline{2.70}{-6.77}\hullline{3.02}{-6.77}\hullfill
\hullface{keel!60!black}\hullmove{-3.02}{-6.77}\hullline{-2.70}{-6.77}\hullline{-2.96}{-6.51}\hullline{-3.89}{-5.96}\hullfill
\hullface{keel!57!black}\hullmove{3.96}{-5.89}\hullline{3.11}{-6.37}\hullline{2.96}{-6.51}\hullline{3.89}{-5.96}\hullfill
\hullface{keel!63!black}\hullmove{-3.89}{-5.96}\hullline{-2.96}{-6.51}\hullline{-3.11}{-6.37}\hullline{-3.96}{-5.89}\hullfill
\hullface{keel!57!black}\hullmove{4.00}{-5.86}\hullline{4.15}{-5.70}\hullline{3.60}{-5.86}\hullline{3.49}{-5.97}\hullline{3.11}{-6.37}\hullline{3.96}{-5.89}\hullfill
\hullface{keel!66!black}\hullmove{-3.96}{-5.89}\hullline{-3.11}{-6.37}\hullline{-3.49}{-5.97}\hullline{-3.60}{-5.86}\hullline{-4.15}{-5.70}\hullline{-4.00}{-5.86}\hullfill
\hullface{keel!57!black}\hullmove{4.41}{-5.44}\hullline{3.98}{-5.44}\hullline{3.60}{-5.86}\hullline{4.15}{-5.70}\hullfill
\hullface{keel!66!black}\hullmove{-4.15}{-5.70}\hullline{-3.60}{-5.86}\hullline{-3.98}{-5.44}\hullline{-4.41}{-5.44}\hullfill
\hullface{keel!57!black}\hullmove{4.66}{-5.17}\hullline{4.83}{-4.98}\hullline{4.47}{-4.85}\hullline{4.43}{-4.90}\hullline{3.98}{-5.44}\hullline{4.41}{-5.44}\hullfill
\hullface{keel!69!black}\hullmove{-4.41}{-5.44}\hullline{-3.98}{-5.44}\hullline{-4.43}{-4.90}\hullline{-4.47}{-4.85}\hullline{-4.83}{-4.98}\hullline{-4.66}{-5.17}\hullfill
\hullface{keel!57!black}\hullmove{4.89}{-4.91}\hullline{4.94}{-4.85}\hullline{4.85}{-4.36}\hullline{4.74}{-4.51}\hullline{4.47}{-4.85}\hullline{4.83}{-4.98}\hullfill
\hullface{keel!72!black}\hullmove{-4.83}{-4.98}\hullline{-4.47}{-4.85}\hullline{-4.74}{-4.51}\hullline{-4.85}{-4.36}\hullline{-4.94}{-4.85}\hullline{-4.89}{-4.91}\hullfill
\hullface{keel!57!black}\hullmove{5.78}{-3.75}\hullline{5.28}{-3.75}\hullline{4.85}{-4.36}\hullline{4.94}{-4.85}\hullfill
\hullface{keel!75!black}\hullmove{-4.94}{-4.85}\hullline{-4.85}{-4.36}\hullline{-5.28}{-3.75}\hullline{-5.78}{-3.75}\hullfill
\hullface{keel!57!black}\hullmove{6.47}{-2.58}\hullline{5.96}{-2.58}\hullline{5.28}{-3.75}\hullline{5.78}{-3.75}\hullfill
\hullface{keel!81!black}\hullmove{-5.78}{-3.75}\hullline{-5.28}{-3.75}\hullline{-5.96}{-2.58}\hullline{-6.47}{-2.58}\hullfill
\hullface{keel!57!black}\hullmove{6.95}{-1.32}\hullline{6.46}{-1.32}\hullline{5.96}{-2.58}\hullline{6.47}{-2.58}\hullfill
\hullface{keel!87!black}\hullmove{-6.47}{-2.58}\hullline{-5.96}{-2.58}\hullline{-6.46}{-1.32}\hullline{-6.95}{-1.32}\hullfill
\hullface{keel!57!black}\hullmove{7.14}{0.00}\hullline{6.68}{0.00}\hullline{6.46}{-1.32}\hullline{6.95}{-1.32}\hullfill
\hullface{keel!93!black}\hullmove{-6.95}{-1.32}\hullline{-6.46}{-1.32}\hullline{-6.68}{0.00}\hullline{-7.14}{0.00}\hullfill
\hullface{topsides!57!black}\hullmove{7.47}{1.52}\hullline{7.03}{1.52}\hullline{6.68}{0.00}\hullline{7.14}{0.00}\hullfill
\hullface{topsides!90!black}\hullmove{-7.14}{0.00}\hullline{-6.68}{0.00}\hullline{-7.03}{1.52}\hullline{-7.47}{1.52}\hullfill
\hullface{topsides!57!black}\hullmove{7.99}{3.04}\hullline{7.61}{3.04}\hullline{7.03}{1.52}\hullline{7.47}{1.52}\hullfill
\hullface{topsides!84!black}\hullmove{-7.47}{1.52}\hullline{-7.03}{1.52}\hullline{-7.61}{3.04}\hullline{-7.99}{3.04}\hullfill
\hullface{topsides!57!black}\hullmove{8.65}{4.55}\hullline{8.32}{4.55}\hullline{7.61}{3.04}\hullline{7.99}{3.04}\hullfill
\hullface{topsides!81!black}\hullmove{-7.99}{3.04}\hullline{-7.61}{3.04}\hullline{-8.32}{4.55}\hullline{-8.65}{4.55}\hullfill
\hullface{topsides!57!black}\hullmove{9.39}{6.07}\hullline{9.14}{6.07}\hullline{8.32}{4.55}\hullline{8.65}{4.55}\hullfill
\hullface{topsides!75!black}\hullmove{-8.65}{4.55}\hullline{-8.32}{4.55}\hullline{-9.14}{6.07}\hullline{-9.39}{6.07}\hullfill
\hullface{topsides!57!black}\hullmove{10.21}{7.59}\hullline{10.04}{7.59}\hullline{9.14}{6.07}\hullline{9.39}{6.07}\hullfill
\hullface{topsides!69!black}\hullmove{-9.39}{6.07}\hullline{-9.14}{6.07}\hullline{-10.04}{7.59}\hullline{-10.21}{7.59}\hullfill
\hullface{keel!57!black}\hullmove{0.37}{-8.77}\hullline{0.32}{-8.77}\hullline{0.00}{-9.06}\hullline{0.00}{-9.06}\hullfill
\hullface{keel!57!black}\hullmove{0.00}{-9.06}\hullline{0.00}{-9.06}\hullline{-0.32}{-8.77}\hullline{-0.37}{-8.77}\hullfill
\hullface{keel!57!black}\hullmove{1.02}{-8.25}\hullline{0.90}{-8.25}\hullline{0.32}{-8.77}\hullline{0.37}{-8.77}\hullfill
\hullface{keel!57!black}\hullmove{-0.37}{-8.77}\hullline{-0.32}{-8.77}\hullline{-0.90}{-8.25}\hullline{-1.02}{-8.25}\hullfill
\hullface{keel!57!black}\hullmove{1.82}{-7.57}\hullline{1.60}{-7.57}\hullline{0.90}{-8.25}\hullline{1.02}{-8.25}\hullfill
\hullface{keel!57!black}\hullmove{-1.02}{-8.25}\hullline{-0.90}{-8.25}\hullline{-1.60}{-7.57}\hullline{-1.82}{-7.57}\hullfill
\hullface{keel!57!black}\hullmove{2.70}{-6.77}\hullline{2.34}{-6.81}\hullline{1.60}{-7.57}\hullline{1.82}{-7.57}\hullfill
\hullface{keel!60!black}\hullmove{-1.82}{-7.57}\hullline{-1.60}{-7.57}\hullline{-2.34}{-6.81}\hullline{-2.70}{-6.77}\hullfill
\hullface{keel!57!black}\hullmove{2.96}{-6.51}\hullline{2.39}{-6.77}\hullline{2.34}{-6.81}\hullline{2.70}{-6.77}\hullfill
\hullface{keel!60!black}\hullmove{-2.70}{-6.77}\hullline{-2.34}{-6.81}\hullline{-2.39}{-6.77}\hullline{-2.96}{-6.51}\hullfill
\hullface{keel!57!black}\hullmove{3.11}{-6.37}\hullline{2.52}{-6.63}\hullline{2.39}{-6.77}\hullline{2.96}{-6.51}\hullfill
\hullface{keel!63!black}\hullmove{-2.96}{-6.51}\hullline{-2.39}{-6.77}\hullline{-2.52}{-6.63}\hullline{-3.11}{-6.37}\hullfill
\hullface{keel!57!black}\hullmove{3.49}{-5.97}\hullline{2.97}{-6.12}\hullline{2.52}{-6.63}\hullline{3.11}{-6.37}\hullfill
\hullface{keel!66!black}\hullmove{-3.11}{-6.37}\hullline{-2.52}{-6.63}\hullline{-2.97}{-6.12}\hullline{-3.49}{-5.97}\hullfill
\hullface{keel!57!black}\hullmove{3.60}{-5.86}\hullline{3.20}{-5.86}\hullline{2.97}{-6.12}\hullline{3.49}{-5.97}\hullfill
\hullface{keel!66!black}\hullmove{-3.49}{-5.97}\hullline{-2.97}{-6.12}\hullline{-3.20}{-5.86}\hullline{-3.60}{-5.86}\hullfill
\hullface{keel!57!black}\hullmove{3.98}{-5.44}\hullline{3.55}{-5.44}\hullline{3.20}{-5.86}\hullline{3.60}{-5.86}\hullfill
\hullface{keel!69!black}\hullmove{-3.60}{-5.86}\hullline{-3.20}{-5.86}\hullline{-3.55}{-5.44}\hullline{-3.98}{-5.44}\hullfill
\hullface{keel!57!black}\hullmove{4.43}{-4.90}\hullline{4.01}{-4.85}\hullline{3.55}{-5.44}\hullline{3.98}{-5.44}\hullfill
\hullface{keel!72!black}\hullmove{-3.98}{-5.44}\hullline{-3.55}{-5.44}\hullline{-4.01}{-4.85}\hullline{-4.43}{-4.90}\hullfill
\hullface{keel!57!black}\hullmove{4.47}{-4.85}\hullline{4.08}{-4.75}\hullline{4.01}{-4.85}\hullline{4.43}{-4.90}\hullfill
\hullface{keel!72!black}\hullmove{-4.43}{-4.90}\hullline{-4.01}{-4.85}\hullline{-4.08}{-4.75}\hullline{-4.47}{-4.85}\hullfill
\hullface{keel!57!black}\hullmove{4.74}{-4.51}\hullline{4.44}{-4.24}\hullline{4.08}{-4.75}\hullline{4.47}{-4.85}\hullfill
\hullface{keel!75!black}\hullmove{-4.47}{-4.85}\hullline{-4.08}{-4.75}\hullline{-4.44}{-4.24}\hullline{-4.74}{-4.51}\hullfill
\hullface{keel!57!black}\hullmove{4.85}{-4.36}\hullline{4.57}{-4.06}\hullline{4.44}{-4.24}\hullline{4.74}{-4.51}\hullfill
\hullface{keel!75!black}\hullmove{-4.74}{-4.51}\hullline{-4.44}{-4.24}\hullline{-4.57}{-4.06}\hullline{-4.85}{-4.36}\hullfill
\hullface{keel!57!black}\hullmove{5.28}{-3.75}\hullline{4.76}{-3.75}\hullline{4.57}{-4.06}\hullline{4.85}{-4.36}\hullfill
\hullface{keel!78!black}\hullmove{-4.85}{-4.36}\hullline{-4.57}{-4.06}\hullline{-4.76}{-3.75}\hullline{-5.28}{-3.75}\hullfill
\hullface{keel!57!black}\hullmove{5.96}{-2.58}\hullline{5.42}{-2.58}\hullline{4.76}{-3.75}\hullline{5.28}{-3.75}\hullfill
\hullface{keel!81!black}\hullmove{-5.28}{-3.75}\hullline{-4.76}{-3.75}\hullline{-5.42}{-2.58}\hullline{-5.96}{-2.58}\hullfill
\hullface{keel!57!black}\hullmove{6.46}{-1.32}\hullline{5.93}{-1.32}\hullline{5.42}{-2.58}\hullline{5.96}{-2.58}\hullfill
\hullface{keel!87!black}\hullmove{-5.96}{-2.58}\hullline{-5.42}{-2.58}\hullline{-5.93}{-1.32}\hullline{-6.46}{-1.32}\hullfill
\hullface{keel!57!black}\hullmove{6.68}{0.00}\hullline{6.17}{0.00}\hullline{5.93}{-1.32}\hullline{6.46}{-1.32}\hullfill
\hullface{keel!93!black}\hullmove{-6.46}{-1.32}\hullline{-5.93}{-1.32}\hullline{-6.17}{0.00}\hullline{-6.68}{0.00}\hullfill
\hullface{topsides!57!black}\hullmove{7.03}{1.52}\hullline{6.56}{1.52}\hullline{6.17}{0.00}\hullline{6.68}{0.00}\hullfill
\hullface{topsides!90!black}\hullmove{-6.68}{0.00}\hullline{-6.17}{0.00}\hullline{-6.56}{1.52}\hullline{-7.03}{1.52}\hullfill
\hullface{topsides!57!black}\hullmove{7.61}{3.04}\hullline{7.19}{3.04}\hullline{6.56}{1.52}\hullline{7.03}{1.52}\hullfill
\hullface{topsides!84!black}\hullmove{-7.03}{1.52}\hullline{-6.56}{1.52}\hullline{-7.19}{3.04}\hullline{-7.61}{3.04}\hullfill
\hullface{topsides!57!black}\hullmove{8.32}{4.55}\hullline{7.96}{4.55}\hullline{7.19}{3.04}\hullline{7.61}{3.04}\hullfill
\hullface{topsides!81!black}\hullmove{-7.61}{3.04}\hullline{-7.19}{3.04}\hullline{-7.96}{4.55}\hullline{-8.32}{4.55}\hullfill
\hullface{topsides!57!black}\hullmove{9.14}{6.07}\hullline{8.85}{6.07}\hullline{7.96}{4.55}\hullline{8.32}{4.55}\hullfill
\hullface{topsides!75!black}\hullmove{-8.32}{4.55}\hullline{-7.96}{4.55}\hullline{-8.85}{6.07}\hullline{-9.14}{6.07}\hullfill
\hullface{topsides!57!black}\hullmove{10.04}{7.59}\hullline{9.83}{7.59}\hullline{8.85}{6.07}\hullline{9.14}{6.07}\hullfill
\hullface{topsides!69!black}\hullmove{-9.14}{6.07}\hullline{-8.85}{6.07}\hullline{-9.83}{7.59}\hullline{-10.04}{7.59}\hullfill
\hullface{keel!57!black}\hullmove{0.32}{-8.77}\hullline{0.28}{-8.77}\hullline{0.00}{-9.06}\hullline{0.00}{-9.06}\hullfill
\hullface{keel!57!black}\hullmove{0.90}{-8.25}\hullline{0.78}{-8.25}\hullline{0.28}{-8.77}\hullline{0.32}{-8.77}\hullfill
\hullface{keel!57!black}\hullmove{-0.32}{-8.77}\hullline{-0.28}{-8.77}\hullline{-0.78}{-8.25}\hullline{-0.90}{-8.25}\hullfill
\hullface{keel!57!black}\hullmove{1.60}{-7.57}\hullline{1.40}{-7.57}\hullline{0.78}{-8.25}\hullline{0.90}{-8.25}\hullfill
\hullface{keel!57!black}\hullmove{-0.90}{-8.25}\hullline{-0.78}{-8.25}\hullline{-1.40}{-7.57}\hullline{-1.60}{-7.57}\hullfill
\hullface{keel!57!black}\hullmove{2.34}{-6.81}\hullline{1.87}{-7.02}\hullline{1.40}{-7.57}\hullline{1.60}{-7.57}\hullfill
\hullface{keel!60!black}\hullmove{-1.60}{-7.57}\hullline{-1.40}{-7.57}\hullline{-1.87}{-7.02}\hullline{-2.34}{-6.81}\hullfill
\hullface{keel!57!black}\hullmove{2.39}{-6.77}\hullline{2.05}{-6.81}\hullline{1.87}{-7.02}\hullline{2.34}{-6.81}\hullfill
\hullface{keel!63!black}\hullmove{-2.34}{-6.81}\hullline{-1.87}{-7.02}\hullline{-2.05}{-6.81}\hullline{-2.39}{-6.77}\hullfill
\hullface{keel!57!black}\hullmove{2.52}{-6.63}\hullline{2.09}{-6.77}\hullline{2.05}{-6.81}\hullline{2.39}{-6.77}\hullfill
\hullface{keel!63!black}\hullmove{-2.39}{-6.77}\hullline{-2.05}{-6.81}\hullline{-2.09}{-6.77}\hullline{-2.52}{-6.63}\hullfill
\hullface{keel!57!black}\hullmove{2.97}{-6.12}\hullline{2.53}{-6.23}\hullline{2.09}{-6.77}\hullline{2.52}{-6.63}\hullfill
\hullface{keel!66!black}\hullmove{-2.52}{-6.63}\hullline{-2.09}{-6.77}\hullline{-2.53}{-6.23}\hullline{-2.97}{-6.12}\hullfill
\hullface{keel!57!black}\hullmove{3.20}{-5.86}\hullline{2.82}{-5.86}\hullline{2.53}{-6.23}\hullline{2.97}{-6.12}\hullfill
\hullface{keel!69!black}\hullmove{-2.97}{-6.12}\hullline{-2.53}{-6.23}\hullline{-2.82}{-5.86}\hullline{-3.20}{-5.86}\hullfill
\hullface{keel!57!black}\hullmove{3.55}{-5.44}\hullline{3.13}{-5.44}\hullline{2.82}{-5.86}\hullline{3.20}{-5.86}\hullfill
\hullface{keel!69!black}\hullmove{-3.20}{-5.86}\hullline{-2.82}{-5.86}\hullline{-3.13}{-5.44}\hullline{-3.55}{-5.44}\hullfill
\hullface{keel!57!black}\hullmove{4.01}{-4.85}\hullline{3.54}{-4.85}\hullline{3.13}{-5.44}\hullline{3.55}{-5.44}\hullfill
\hullface{keel!72!black}\hullmove{-3.55}{-5.44}\hullline{-3.13}{-5.44}\hullline{-3.54}{-4.85}\hullline{-4.01}{-4.85}\hullfill
\hullface{keel!57!black}\hullmove{4.08}{-4.75}\hullline{3.68}{-4.64}\hullline{3.54}{-4.85}\hullline{4.01}{-4.85}\hullfill
\hullface{keel!75!black}\hullmove{-4.01}{-4.85}\hullline{-3.54}{-4.85}\hullline{-3.68}{-4.64}\hullline{-4.08}{-4.75}\hullfill
\hullface{keel!57!black}\hullmove{4.44}{-4.24}\hullline{4.06}{-4.06}\hullline{3.68}{-4.64}\hullline{4.08}{-4.75}\hullfill
\hullface{keel!78!black}\hullmove{-4.08}{-4.75}\hullline{-3.68}{-4.64}\hullline{-4.06}{-4.06}\hullline{-4.44}{-4.24}\hullfill
\hullface{keel!57!black}\hullmove{4.57}{-4.06}\hullline{4.76}{-3.75}\hullline{4.24}{-3.75}\hullline{4.19}{-3.85}\hullline{4.06}{-4.06}\hullline{4.44}{-4.24}\hullfill
\hullface{keel!78!black}\hullmove{-4.44}{-4.24}\hullline{-4.06}{-4.06}\hullline{-4.19}{-3.85}\hullline{-4.24}{-3.75}\hullline{-4.76}{-3.75}\hullline{-4.57}{-4.06}\hullfill
\hullface{keel!57!black}\hullmove{5.42}{-2.58}\hullline{4.87}{-2.58}\hullline{4.24}{-3.75}\hullline{4.76}{-3.75}\hullfill
\hullface{keel!84!black}\hullmove{-4.76}{-3.75}\hullline{-4.24}{-3.75}\hullline{-4.87}{-2.58}\hullline{-5.42}{-2.58}\hullfill
\hullface{keel!57!black}\hullmove{5.93}{-1.32}\hullline{5.37}{-1.32}\hullline{4.87}{-2.58}\hullline{5.42}{-2.58}\hullfill
\hullface{keel!87!black}\hullmove{-5.42}{-2.58}\hullline{-4.87}{-2.58}\hullline{-5.37}{-1.32}\hullline{-5.93}{-1.32}\hullfill
\hullface{keel!57!black}\hullmove{6.17}{0.00}\hullline{5.63}{0.00}\hullline{5.37}{-1.32}\hullline{5.93}{-1.32}\hullfill
\hullface{keel!96!black}\hullmove{-5.93}{-1.32}\hullline{-5.37}{-1.32}\hullline{-5.63}{0.00}\hullline{-6.17}{0.00}\hullfill
\hullface{topsides!57!black}\hullmove{6.56}{1.52}\hullline{6.04}{1.52}\hullline{5.63}{0.00}\hullline{6.17}{0.00}\hullfill
\hullface{topsides!93!black}\hullmove{-6.17}{0.00}\hullline{-5.63}{0.00}\hullline{-6.04}{1.52}\hullline{-6.56}{1.52}\hullfill
\hullface{topsides!57!black}\hullmove{7.19}{3.04}\hullline{6.72}{3.04}\hullline{6.04}{1.52}\hullline{6.56}{1.52}\hullfill
\hullface{topsides!84!black}\hullmove{-6.56}{1.52}\hullline{-6.04}{1.52}\hullline{-6.72}{3.04}\hullline{-7.19}{3.04}\hullfill
\hullface{topsides!57!black}\hullmove{7.96}{4.55}\hullline{7.56}{4.55}\hullline{6.72}{3.04}\hullline{7.19}{3.04}\hullfill
\hullface{topsides!81!black}\hullmove{-7.19}{3.04}\hullline{-6.72}{3.04}\hullline{-7.56}{4.55}\hullline{-7.96}{4.55}\hullfill
\hullface{topsides!57!black}\hullmove{8.85}{6.07}\hullline{8.52}{6.07}\hullline{7.56}{4.55}\hullline{7.96}{4.55}\hullfill
\hullface{topsides!75!black}\hullmove{-7.96}{4.55}\hullline{-7.56}{4.55}\hullline{-8.52}{6.07}\hullline{-8.85}{6.07}\hullfill
\hullface{topsides!57!black}\hullmove{9.83}{7.59}\hullline{9.57}{7.59}\hullline{8.52}{6.07}\hullline{8.85}{6.07}\hullfill
\hullface{topsides!69!black}\hullmove{-8.85}{6.07}\hullline{-8.52}{6.07}\hullline{-9.57}{7.59}\hullline{-9.83}{7.59}\hullfill
\hullface{keel!57!black}\hullmove{0.28}{-8.77}\hullline{0.24}{-8.77}\hullline{0.00}{-9.06}\hullline{0.00}{-9.06}\hullfill
\hullface{keel!57!black}\hullmove{0.00}{-9.06}\hullline{0.00}{-9.06}\hullline{-0.24}{-8.77}\hullline{-0.28}{-8.77}\hullfill
\hullface{keel!57!black}\hullmove{0.78}{-8.25}\hullline{0.67}{-8.25}\hullline{0.24}{-8.77}\hullline{0.28}{-8.77}\hullfill
\hullface{keel!57!black}\hullmove{-0.28}{-8.77}\hullline{-0.24}{-8.77}\hullline{-0.67}{-8.25}\hullline{-0.78}{-8.25}\hullfill
\hullface{keel!57!black}\hullmove{1.40}{-7.57}\hullline{1.20}{-7.57}\hullline{0.67}{-8.25}\hullline{0.78}{-8.25}\hullfill
\hullface{keel!60!black}\hullmove{-0.78}{-8.25}\hullline{-0.67}{-8.25}\hullline{-1.20}{-7.57}\hullline{-1.40}{-7.57}\hullfill
\hullface{keel!57!black}\hullmove{1.87}{-7.02}\hullline{1.50}{-7.17}\hullline{1.20}{-7.57}\hullline{1.40}{-7.57}\hullfill
\hullface{keel!63!black}\hullmove{-1.40}{-7.57}\hullline{-1.20}{-7.57}\hullline{-1.50}{-7.17}\hullline{-1.87}{-7.02}\hullfill
\hullface{keel!57!black}\hullmove{2.05}{-6.81}\hullline{1.67}{-6.94}\hullline{1.50}{-7.17}\hullline{1.87}{-7.02}\hullfill
\hullface{keel!66!black}\hullmove{-1.87}{-7.02}\hullline{-1.50}{-7.17}\hullline{-1.67}{-6.94}\hullline{-2.05}{-6.81}\hullfill
\hullface{keel!57!black}\hullmove{2.09}{-6.77}\hullline{1.80}{-6.77}\hullline{1.67}{-6.94}\hullline{2.05}{-6.81}\hullfill
\hullface{keel!66!black}\hullmove{-2.05}{-6.81}\hullline{-1.67}{-6.94}\hullline{-1.80}{-6.77}\hullline{-2.09}{-6.77}\hullfill
\hullface{keel!57!black}\hullmove{2.53}{-6.23}\hullline{2.13}{-6.30}\hullline{1.80}{-6.77}\hullline{2.09}{-6.77}\hullfill
\hullface{keel!69!black}\hullmove{-2.09}{-6.77}\hullline{-1.80}{-6.77}\hullline{-2.13}{-6.30}\hullline{-2.53}{-6.23}\hullfill
\hullface{keel!57!black}\hullmove{2.82}{-5.86}\hullline{2.43}{-5.86}\hullline{2.13}{-6.30}\hullline{2.53}{-6.23}\hullfill
\hullface{keel!72!black}\hullmove{-2.53}{-6.23}\hullline{-2.13}{-6.30}\hullline{-2.43}{-5.86}\hullline{-2.82}{-5.86}\hullfill
\hullface{keel!57!black}\hullmove{3.13}{-5.44}\hullline{2.71}{-5.44}\hullline{2.43}{-5.86}\hullline{2.82}{-5.86}\hullfill
\hullface{keel!72!black}\hullmove{-2.82}{-5.86}\hullline{-2.43}{-5.86}\hullline{-2.71}{-5.44}\hullline{-3.13}{-5.44}\hullfill
\hullface{keel!57!black}\hullmove{3.54}{-4.85}\hullline{3.08}{-4.85}\hullline{2.71}{-5.44}\hullline{3.13}{-5.44}\hullfill
\hullface{keel!75!black}\hullmove{-3.13}{-5.44}\hullline{-2.71}{-5.44}\hullline{-3.08}{-4.85}\hullline{-3.54}{-4.85}\hullfill
\hullface{keel!57!black}\hullmove{3.68}{-4.64}\hullline{3.25}{-4.57}\hullline{3.08}{-4.85}\hullline{3.54}{-4.85}\hullfill
\hullface{keel!78!black}\hullmove{-3.54}{-4.85}\hullline{-3.08}{-4.85}\hullline{-3.25}{-4.57}\hullline{-3.68}{-4.64}\hullfill
\hullface{keel!57!black}\hullmove{4.06}{-4.06}\hullline{3.61}{-3.93}\hullline{3.25}{-4.57}\hullline{3.68}{-4.64}\hullfill
\hullface{keel!78!black}\hullmove{-3.68}{-4.64}\hullline{-3.25}{-4.57}\hullline{-3.61}{-3.93}\hullline{-4.06}{-4.06}\hullfill
\hullface{keel!57!black}\hullmove{4.19}{-3.85}\hullline{4.24}{-3.75}\hullline{3.74}{-3.70}\hullline{3.71}{-3.75}\hullline{3.61}{-3.93}\hullline{4.06}{-4.06}\hullfill
\hullface{keel!81!black}\hullmove{-4.06}{-4.06}\hullline{-3.61}{-3.93}\hullline{-3.71}{-3.75}\hullline{-3.74}{-3.70}\hullline{-4.24}{-3.75}\hullline{-4.19}{-3.85}\hullfill
\hullface{keel!57!black}\hullmove{4.87}{-2.58}\hullline{4.29}{-2.58}\hullline{3.74}{-3.70}\hullline{4.24}{-3.75}\hullfill
\hullface{keel!84!black}\hullmove{-4.24}{-3.75}\hullline{-3.74}{-3.70}\hullline{-4.29}{-2.58}\hullline{-4.87}{-2.58}\hullfill
\hullface{keel!57!black}\hullmove{5.37}{-1.32}\hullline{4.77}{-1.32}\hullline{4.29}{-2.58}\hullline{4.87}{-2.58}\hullfill
\hullface{keel!90!black}\hullmove{-4.87}{-2.58}\hullline{-4.29}{-2.58}\hullline{-4.77}{-1.32}\hullline{-5.37}{-1.32}\hullfill
\hullface{keel!57!black}\hullmove{5.63}{0.00}\hullline{5.05}{0.00}\hullline{4.77}{-1.32}\hullline{5.37}{-1.32}\hullfill
\hullface{keel!96!black}\hullmove{-5.37}{-1.32}\hullline{-4.77}{-1.32}\hullline{-5.05}{0.00}\hullline{-5.63}{0.00}\hullfill
\hullface{topsides!57!black}\hullmove{6.04}{1.52}\hullline{5.49}{1.52}\hullline{5.05}{0.00}\hullline{5.63}{0.00}\hullfill
\hullface{topsides!93!black}\hullmove{-5.63}{0.00}\hullline{-5.05}{0.00}\hullline{-5.49}{1.52}\hullline{-6.04}{1.52}\hullfill
\hullface{topsides!57!black}\hullmove{6.72}{3.04}\hullline{6.22}{3.04}\hullline{5.49}{1.52}\hullline{6.04}{1.52}\hullfill
\hullface{topsides!84!black}\hullmove{-6.04}{1.52}\hullline{-5.49}{1.52}\hullline{-6.22}{3.04}\hullline{-6.72}{3.04}\hullfill
\hullface{topsides!57!black}\hullmove{7.56}{4.55}\hullline{7.11}{4.55}\hullline{6.22}{3.04}\hullline{6.72}{3.04}\hullfill
\hullface{topsides!81!black}\hullmove{-6.72}{3.04}\hullline{-6.22}{3.04}\hullline{-7.11}{4.55}\hullline{-7.56}{4.55}\hullfill
\hullface{topsides!57!black}\hullmove{8.52}{6.07}\hullline{8.14}{6.07}\hullline{7.11}{4.55}\hullline{7.56}{4.55}\hullfill
\hullface{topsides!75!black}\hullmove{-7.56}{4.55}\hullline{-7.11}{4.55}\hullline{-8.14}{6.07}\hullline{-8.52}{6.07}\hullfill
\hullface{topsides!57!black}\hullmove{9.57}{7.59}\hullline{9.27}{7.59}\hullline{8.14}{6.07}\hullline{8.52}{6.07}\hullfill
\hullface{topsides!69!black}\hullmove{-8.52}{6.07}\hullline{-8.14}{6.07}\hullline{-9.27}{7.59}\hullline{-9.57}{7.59}\hullfill
\hullface{keel!57!black}\hullmove{0.24}{-8.77}\hullline{0.21}{-8.67}\hullline{0.00}{-8.95}\hullline{0.00}{-9.06}\hullfill
\hullface{keel!60!black}\hullmove{0.00}{-9.06}\hullline{0.00}{-8.95}\hullline{-0.21}{-8.67}\hullline{-0.24}{-8.77}\hullfill
\hullface{keel!57!black}\hullmove{0.67}{-8.25}\hullline{0.57}{-8.15}\hullline{0.21}{-8.67}\hullline{0.24}{-8.77}\hullfill
\hullface{keel!60!black}\hullmove{-0.24}{-8.77}\hullline{-0.21}{-8.67}\hullline{-0.57}{-8.15}\hullline{-0.67}{-8.25}\hullfill
\hullface{keel!57!black}\hullmove{1.20}{-7.57}\hullline{1.03}{-7.48}\hullline{0.57}{-8.15}\hullline{0.67}{-8.25}\hullfill
\hullface{keel!63!black}\hullmove{-0.67}{-8.25}\hullline{-0.57}{-8.15}\hullline{-1.03}{-7.48}\hullline{-1.20}{-7.57}\hullfill
\hullface{keel!57!black}\hullmove{1.50}{-7.17}\hullline{1.17}{-7.27}\hullline{1.03}{-7.48}\hullline{1.20}{-7.57}\hullfill
\hullface{keel!66!black}\hullmove{-1.20}{-7.57}\hullline{-1.03}{-7.48}\hullline{-1.17}{-7.27}\hullline{-1.50}{-7.17}\hullfill
\hullface{keel!57!black}\hullmove{1.67}{-6.94}\hullline{1.33}{-7.03}\hullline{1.17}{-7.27}\hullline{1.50}{-7.17}\hullfill
\hullface{keel!69!black}\hullmove{-1.50}{-7.17}\hullline{-1.17}{-7.27}\hullline{-1.33}{-7.03}\hullline{-1.67}{-6.94}\hullfill
\hullface{keel!57!black}\hullmove{1.80}{-6.77}\hullline{1.55}{-6.69}\hullline{1.33}{-7.03}\hullline{1.67}{-6.94}\hullfill
\hullface{keel!69!black}\hullmove{-1.67}{-6.94}\hullline{-1.33}{-7.03}\hullline{-1.55}{-6.69}\hullline{-1.80}{-6.77}\hullfill
\hullface{keel!57!black}\hullmove{2.13}{-6.30}\hullline{1.76}{-6.35}\hullline{1.55}{-6.69}\hullline{1.80}{-6.77}\hullfill
\hullface{keel!72!black}\hullmove{-1.80}{-6.77}\hullline{-1.55}{-6.69}\hullline{-1.76}{-6.35}\hullline{-2.13}{-6.30}\hullfill
\hullface{keel!57!black}\hullmove{2.43}{-5.86}\hullline{2.10}{-5.79}\hullline{1.76}{-6.35}\hullline{2.13}{-6.30}\hullfill
\hullface{keel!72!black}\hullmove{-2.13}{-6.30}\hullline{-1.76}{-6.35}\hullline{-2.10}{-5.79}\hullline{-2.43}{-5.86}\hullfill
\hullface{keel!57!black}\hullmove{2.71}{-5.44}\hullline{2.31}{-5.44}\hullline{2.10}{-5.79}\hullline{2.43}{-5.86}\hullfill
\hullface{keel!75!black}\hullmove{-2.43}{-5.86}\hullline{-2.10}{-5.79}\hullline{-2.31}{-5.44}\hullline{-2.71}{-5.44}\hullfill
\hullface{keel!57!black}\hullmove{3.08}{-4.85}\hullline{2.67}{-4.79}\hullline{2.31}{-5.44}\hullline{2.71}{-5.44}\hullfill
\hullface{keel!78!black}\hullmove{-2.71}{-5.44}\hullline{-2.31}{-5.44}\hullline{-2.67}{-4.79}\hullline{-3.08}{-4.85}\hullfill
\hullface{keel!57!black}\hullmove{3.25}{-4.57}\hullline{2.82}{-4.52}\hullline{2.67}{-4.79}\hullline{3.08}{-4.85}\hullfill
\hullface{keel!78!black}\hullmove{-3.08}{-4.85}\hullline{-2.67}{-4.79}\hullline{-2.82}{-4.52}\hullline{-3.25}{-4.57}\hullfill
\hullface{keel!57!black}\hullmove{3.61}{-3.93}\hullline{3.16}{-3.85}\hullline{2.82}{-4.52}\hullline{3.25}{-4.57}\hullfill
\hullface{keel!81!black}\hullmove{-3.25}{-4.57}\hullline{-2.82}{-4.52}\hullline{-3.16}{-3.85}\hullline{-3.61}{-3.93}\hullfill
\hullface{keel!57!black}\hullmove{3.71}{-3.75}\hullline{3.22}{-3.71}\hullline{3.16}{-3.85}\hullline{3.61}{-3.93}\hullfill
\hullface{keel!84!black}\hullmove{-3.61}{-3.93}\hullline{-3.16}{-3.85}\hullline{-3.22}{-3.71}\hullline{-3.71}{-3.75}\hullfill
\hullface{keel!57!black}\hullmove{3.74}{-3.70}\hullline{3.27}{-3.60}\hullline{3.22}{-3.71}\hullline{3.71}{-3.75}\hullfill
\hullface{keel!84!black}\hullmove{-3.71}{-3.75}\hullline{-3.22}{-3.71}\hullline{-3.27}{-3.60}\hullline{-3.74}{-3.70}\hullfill
\hullface{keel!57!black}\hullmove{4.29}{-2.58}\hullline{3.74}{-2.55}\hullline{3.27}{-3.60}\hullline{3.74}{-3.70}\hullfill
\hullface{keel!84!black}\hullmove{-3.74}{-3.70}\hullline{-3.27}{-3.60}\hullline{-3.74}{-2.55}\hullline{-4.29}{-2.58}\hullfill
\hullface{keel!57!black}\hullmove{4.77}{-1.32}\hullline{4.17}{-1.31}\hullline{3.74}{-2.55}\hullline{4.29}{-2.58}\hullfill
\hullface{keel!90!black}\hullmove{-4.29}{-2.58}\hullline{-3.74}{-2.55}\hullline{-4.17}{-1.31}\hullline{-4.77}{-1.32}\hullfill
\hullface{keel!57!black}\hullmove{5.05}{0.00}\hullline{4.43}{0.00}\hullline{4.17}{-1.31}\hullline{4.77}{-1.32}\hullfill
\hullface{keel!96!black}\hullmove{-4.77}{-1.32}\hullline{-4.17}{-1.31}\hullline{-4.43}{0.00}\hullline{-5.05}{0.00}\hullfill
\hullface{topsides!57!black}\hullmove{5.49}{1.52}\hullline{4.90}{1.52}\hullline{4.43}{0.00}\hullline{5.05}{0.00}\hullfill
\hullface{topsides!93!black}\hullmove{-5.05}{0.00}\hullline{-4.43}{0.00}\hullline{-4.90}{1.52}\hullline{-5.49}{1.52}\hullfill
\hullface{topsides!57!black}\hullmove{6.22}{3.04}\hullline{5.67}{3.04}\hullline{4.90}{1.52}\hullline{5.49}{1.52}\hullfill
\hullface{topsides!84!black}\hullmove{-5.49}{1.52}\hullline{-4.90}{1.52}\hullline{-5.67}{3.04}\hullline{-6.22}{3.04}\hullfill
\hullface{topsides!57!black}\hullmove{7.11}{4.55}\hullline{6.62}{4.55}\hullline{5.67}{3.04}\hullline{6.22}{3.04}\hullfill
\hullface{topsides!81!black}\hullmove{-6.22}{3.04}\hullline{-5.67}{3.04}\hullline{-6.62}{4.55}\hullline{-7.11}{4.55}\hullfill
\hullface{topsides!57!black}\hullmove{8.14}{6.07}\hullline{7.71}{6.07}\hullline{6.62}{4.55}\hullline{7.11}{4.55}\hullfill
\hullface{topsides!75!black}\hullmove{-7.11}{4.55}\hullline{-6.62}{4.55}\hullline{-7.71}{6.07}\hullline{-8.14}{6.07}\hullfill
\hullface{topsides!57!black}\hullmove{9.27}{7.59}\hullline{8.91}{7.59}\hullline{7.71}{6.07}\hullline{8.14}{6.07}\hullfill
\hullface{topsides!72!black}\hullmove{-8.14}{6.07}\hullline{-7.71}{6.07}\hullline{-8.91}{7.59}\hullline{-9.27}{7.59}\hullfill
\hullface{keel!57!black}\hullmove{0.21}{-8.67}\hullline{0.18}{-8.38}\hullline{0.00}{-8.65}\hullline{0.00}{-8.95}\hullfill
\hullface{keel!66!black}\hullmove{0.00}{-8.95}\hullline{0.00}{-8.65}\hullline{-0.18}{-8.38}\hullline{-0.21}{-8.67}\hullfill
\hullface{keel!57!black}\hullmove{0.57}{-8.15}\hullline{0.50}{-7.88}\hullline{0.18}{-8.38}\hullline{0.21}{-8.67}\hullfill
\hullface{keel!69!black}\hullmove{-0.21}{-8.67}\hullline{-0.18}{-8.38}\hullline{-0.50}{-7.88}\hullline{-0.57}{-8.15}\hullfill
\hullface{keel!57!black}\hullmove{1.03}{-7.48}\hullline{0.84}{-7.33}\hullline{0.50}{-7.88}\hullline{0.57}{-8.15}\hullfill
\hullface{keel!69!black}\hullmove{-0.57}{-8.15}\hullline{-0.50}{-7.88}\hullline{-0.84}{-7.33}\hullline{-1.03}{-7.48}\hullfill
\hullface{keel!57!black}\hullmove{1.55}{-6.69}\hullline{1.36}{-6.46}\hullline{0.99}{-7.08}\hullline{1.33}{-7.03}\hullfill
\hullface{keel!72!black}\hullmove{-1.33}{-7.03}\hullline{-0.99}{-7.08}\hullline{-1.36}{-6.46}\hullline{-1.55}{-6.69}\hullfill
\hullface{keel!57!black}\hullmove{1.76}{-6.35}\hullline{1.40}{-6.38}\hullline{1.36}{-6.46}\hullline{1.55}{-6.69}\hullfill
\hullface{keel!72!black}\hullmove{-1.55}{-6.69}\hullline{-1.36}{-6.46}\hullline{-1.40}{-6.38}\hullline{-1.76}{-6.35}\hullfill
\hullface{keel!57!black}\hullmove{2.10}{-5.79}\hullline{1.83}{-5.59}\hullline{1.40}{-6.38}\hullline{1.76}{-6.35}\hullfill
\hullface{keel!75!black}\hullmove{-1.76}{-6.35}\hullline{-1.40}{-6.38}\hullline{-1.83}{-5.59}\hullline{-2.10}{-5.79}\hullfill
\hullface{keel!57!black}\hullmove{2.31}{-5.44}\hullline{1.92}{-5.44}\hullline{1.83}{-5.59}\hullline{2.10}{-5.79}\hullfill
\hullface{keel!78!black}\hullmove{-2.10}{-5.79}\hullline{-1.83}{-5.59}\hullline{-1.92}{-5.44}\hullline{-2.31}{-5.44}\hullfill
\hullface{keel!57!black}\hullmove{2.67}{-4.79}\hullline{2.32}{-4.63}\hullline{1.92}{-5.44}\hullline{2.31}{-5.44}\hullfill
\hullface{keel!78!black}\hullmove{-2.31}{-5.44}\hullline{-1.92}{-5.44}\hullline{-2.32}{-4.63}\hullline{-2.67}{-4.79}\hullfill
\hullface{keel!57!black}\hullmove{2.82}{-4.52}\hullline{2.39}{-4.49}\hullline{2.32}{-4.63}\hullline{2.67}{-4.79}\hullfill
\hullface{keel!81!black}\hullmove{-2.67}{-4.79}\hullline{-2.32}{-4.63}\hullline{-2.39}{-4.49}\hullline{-2.82}{-4.52}\hullfill
\hullface{keel!57!black}\hullmove{3.16}{-3.85}\hullline{2.70}{-3.79}\hullline{2.39}{-4.49}\hullline{2.82}{-4.52}\hullfill
\hullface{keel!84!black}\hullmove{-2.82}{-4.52}\hullline{-2.39}{-4.49}\hullline{-2.70}{-3.79}\hullline{-3.16}{-3.85}\hullfill
\hullface{keel!57!black}\hullmove{3.74}{-2.55}\hullline{3.21}{-2.46}\hullline{2.81}{-3.54}\hullline{3.27}{-3.60}\hullfill
\hullface{keel!87!black}\hullmove{-3.27}{-3.60}\hullline{-2.81}{-3.54}\hullline{-3.21}{-2.46}\hullline{-3.74}{-2.55}\hullfill
\hullface{keel!57!black}\hullmove{4.17}{-1.31}\hullline{3.57}{-1.26}\hullline{3.21}{-2.46}\hullline{3.74}{-2.55}\hullfill
\hullface{keel!90!black}\hullmove{-3.74}{-2.55}\hullline{-3.21}{-2.46}\hullline{-3.57}{-1.26}\hullline{-4.17}{-1.31}\hullfill
\hullface{keel!60!black}\hullmove{4.43}{0.00}\hullline{3.76}{0.00}\hullline{3.57}{-1.26}\hullline{4.17}{-1.31}\hullfill
\hullface{keel!96!black}\hullmove{-4.17}{-1.31}\hullline{-3.57}{-1.26}\hullline{-3.76}{0.00}\hullline{-4.43}{0.00}\hullfill
\hullface{topsides!57!black}\hullmove{4.90}{1.52}\hullline{4.26}{1.52}\hullline{3.76}{0.00}\hullline{4.43}{0.00}\hullfill
\hullface{topsides!93!black}\hullmove{-4.43}{0.00}\hullline{-3.76}{0.00}\hullline{-4.26}{1.52}\hullline{-4.90}{1.52}\hullfill
\hullface{topsides!57!black}\hullmove{5.67}{3.04}\hullline{5.07}{3.04}\hullline{4.26}{1.52}\hullline{4.90}{1.52}\hullfill
\hullface{topsides!84!black}\hullmove{-4.90}{1.52}\hullline{-4.26}{1.52}\hullline{-5.07}{3.04}\hullline{-5.67}{3.04}\hullfill
\hullface{topsides!57!black}\hullmove{6.62}{4.55}\hullline{6.07}{4.55}\hullline{5.07}{3.04}\hullline{5.67}{3.04}\hullfill
\hullface{topsides!81!black}\hullmove{-5.67}{3.04}\hullline{-5.07}{3.04}\hullline{-6.07}{4.55}\hullline{-6.62}{4.55}\hullfill
\hullface{topsides!57!black}\hullmove{7.71}{6.07}\hullline{7.22}{6.07}\hullline{6.07}{4.55}\hullline{6.62}{4.55}\hullfill
\hullface{topsides!75!black}\hullmove{-6.62}{4.55}\hullline{-6.07}{4.55}\hullline{-7.22}{6.07}\hullline{-7.71}{6.07}\hullfill
\hullface{topsides!57!black}\hullmove{8.91}{7.59}\hullline{8.49}{7.59}\hullline{7.22}{6.07}\hullline{7.71}{6.07}\hullfill
\hullface{topsides!72!black}\hullmove{-7.71}{6.07}\hullline{-7.22}{6.07}\hullline{-8.49}{7.59}\hullline{-8.91}{7.59}\hullfill
\hullface{keel!57!black}\hullmove{1.17}{-7.27}\hullline{1.33}{-7.03}\hullline{0.99}{-7.08}\hullline{0.90}{-7.23}\hullline{0.84}{-7.33}\hullline{1.03}{-7.48}\hullfill
\hullface{keel!72!black}\hullmove{-1.03}{-7.48}\hullline{-0.84}{-7.33}\hullline{-0.90}{-7.23}\hullline{-0.99}{-7.08}\hullline{-1.33}{-7.03}\hullline{-1.17}{-7.27}\hullfill
\hullface{keel!57!black}\hullmove{3.22}{-3.71}\hullline{3.27}{-3.60}\hullline{2.81}{-3.54}\hullline{2.79}{-3.58}\hullline{2.70}{-3.79}\hullline{3.16}{-3.85}\hullfill
\hullface{keel!84!black}\hullmove{-3.16}{-3.85}\hullline{-2.70}{-3.79}\hullline{-2.79}{-3.58}\hullline{-2.81}{-3.54}\hullline{-3.27}{-3.60}\hullline{-3.22}{-3.71}\hullfill
\hullface{keel!57!black}\hullmove{0.18}{-8.38}\hullline{0.50}{-7.88}\hullline{0.45}{-7.42}\hullline{0.16}{-7.89}\hullline{0.00}{-8.15}\hullline{0.00}{-8.65}\hullfill
\hullface{keel!72!black}\hullmove{0.00}{-8.65}\hullline{0.00}{-8.15}\hullline{-0.16}{-7.89}\hullline{-0.45}{-7.42}\hullline{-0.50}{-7.88}\hullline{-0.18}{-8.38}\hullfill
\hullface{keel!57!black}\hullmove{0.84}{-7.33}\hullline{0.48}{-7.36}\hullline{0.45}{-7.42}\hullline{0.50}{-7.88}\hullfill
\hullface{keel!72!black}\hullmove{-0.50}{-7.88}\hullline{-0.45}{-7.42}\hullline{-0.48}{-7.36}\hullline{-0.84}{-7.33}\hullfill
\hullface{keel!57!black}\hullmove{0.90}{-7.23}\hullline{0.77}{-7.10}\hullline{0.48}{-7.36}\hullline{0.84}{-7.33}\hullfill
\hullface{keel!60!black}\hullmove{-0.84}{-7.33}\hullline{-0.48}{-7.36}\hullline{-0.77}{-7.10}\hullline{-0.90}{-7.23}\hullfill
\hullface{keel!57!black}\hullmove{0.99}{-7.08}\hullline{1.08}{-6.81}\hullline{0.77}{-7.10}\hullline{0.90}{-7.23}\hullfill
\hullface{keel!60!black}\hullmove{-0.90}{-7.23}\hullline{-0.77}{-7.10}\hullline{-1.08}{-6.81}\hullline{-0.99}{-7.08}\hullfill
\hullface{keel!57!black}\hullmove{1.36}{-6.46}\hullline{1.34}{-6.40}\hullline{1.08}{-6.81}\hullline{0.99}{-7.08}\hullfill
\hullface{keel!63!black}\hullmove{-0.99}{-7.08}\hullline{-1.08}{-6.81}\hullline{-1.34}{-6.40}\hullline{-1.36}{-6.46}\hullfill
\hullface{keel!57!black}\hullmove{1.40}{-6.38}\hullline{1.45}{-6.09}\hullline{1.34}{-6.40}\hullline{1.36}{-6.46}\hullfill
\hullface{keel!69!black}\hullmove{-1.36}{-6.46}\hullline{-1.34}{-6.40}\hullline{-1.45}{-6.09}\hullline{-1.40}{-6.38}\hullfill
\hullface{keel!57!black}\hullmove{1.83}{-5.59}\hullline{1.54}{-5.44}\hullline{1.45}{-6.09}\hullline{1.40}{-6.38}\hullfill
\hullface{keel!78!black}\hullmove{-1.40}{-6.38}\hullline{-1.45}{-6.09}\hullline{-1.54}{-5.44}\hullline{-1.83}{-5.59}\hullfill
\hullface{keel!57!black}\hullmove{1.92}{-5.44}\hullline{1.60}{-5.27}\hullline{1.54}{-5.44}\hullline{1.83}{-5.59}\hullfill
\hullface{keel!81!black}\hullmove{-1.83}{-5.59}\hullline{-1.54}{-5.44}\hullline{-1.60}{-5.27}\hullline{-1.92}{-5.44}\hullfill
\hullface{keel!57!black}\hullmove{2.32}{-4.63}\hullline{1.95}{-4.47}\hullline{1.60}{-5.27}\hullline{1.92}{-5.44}\hullfill
\hullface{keel!81!black}\hullmove{-1.92}{-5.44}\hullline{-1.60}{-5.27}\hullline{-1.95}{-4.47}\hullline{-2.32}{-4.63}\hullfill
\hullface{keel!57!black}\hullmove{2.39}{-4.49}\hullline{2.00}{-4.36}\hullline{1.95}{-4.47}\hullline{2.32}{-4.63}\hullfill
\hullface{keel!84!black}\hullmove{-2.32}{-4.63}\hullline{-1.95}{-4.47}\hullline{-2.00}{-4.36}\hullline{-2.39}{-4.49}\hullfill
\hullface{keel!57!black}\hullmove{2.70}{-3.79}\hullline{2.23}{-3.77}\hullline{2.00}{-4.36}\hullline{2.39}{-4.49}\hullfill
\hullface{keel!84!black}\hullmove{-2.39}{-4.49}\hullline{-2.00}{-4.36}\hullline{-2.23}{-3.77}\hullline{-2.70}{-3.79}\hullfill
\hullface{keel!57!black}\hullmove{2.79}{-3.58}\hullline{2.33}{-3.51}\hullline{2.23}{-3.77}\hullline{2.70}{-3.79}\hullfill
\hullface{keel!87!black}\hullmove{-2.70}{-3.79}\hullline{-2.23}{-3.77}\hullline{-2.33}{-3.51}\hullline{-2.79}{-3.58}\hullfill
\hullface{keel!57!black}\hullmove{2.81}{-3.54}\hullline{2.37}{-3.38}\hullline{2.33}{-3.51}\hullline{2.79}{-3.58}\hullfill
\hullface{keel!87!black}\hullmove{-2.79}{-3.58}\hullline{-2.33}{-3.51}\hullline{-2.37}{-3.38}\hullline{-2.81}{-3.54}\hullfill
\hullface{keel!57!black}\hullmove{3.21}{-2.46}\hullline{2.70}{-2.32}\hullline{2.37}{-3.38}\hullline{2.81}{-3.54}\hullfill
\hullface{keel!90!black}\hullmove{-2.81}{-3.54}\hullline{-2.37}{-3.38}\hullline{-2.70}{-2.32}\hullline{-3.21}{-2.46}\hullfill
\hullface{keel!57!black}\hullmove{3.57}{-1.26}\hullline{2.94}{-1.19}\hullline{2.70}{-2.32}\hullline{3.21}{-2.46}\hullfill
\hullface{keel!93!black}\hullmove{-3.21}{-2.46}\hullline{-2.70}{-2.32}\hullline{-2.94}{-1.19}\hullline{-3.57}{-1.26}\hullfill
\hullface{keel!63!black}\hullmove{3.76}{0.00}\hullline{3.06}{0.00}\hullline{2.94}{-1.19}\hullline{3.57}{-1.26}\hullfill
\hullface{keel!99!black}\hullmove{-3.57}{-1.26}\hullline{-2.94}{-1.19}\hullline{-3.06}{0.00}\hullline{-3.76}{0.00}\hullfill
\hullface{topsides!57!black}\hullmove{4.26}{1.52}\hullline{3.58}{1.52}\hullline{3.06}{0.00}\hullline{3.76}{0.00}\hullfill
\hullface{topsides!93!black}\hullmove{-3.76}{0.00}\hullline{-3.06}{0.00}\hullline{-3.58}{1.52}\hullline{-4.26}{1.52}\hullfill
\hullface{topsides!57!black}\hullmove{5.07}{3.04}\hullline{4.42}{3.04}\hullline{3.58}{1.52}\hullline{4.26}{1.52}\hullfill
\hullface{topsides!87!black}\hullmove{-4.26}{1.52}\hullline{-3.58}{1.52}\hullline{-4.42}{3.04}\hullline{-5.07}{3.04}\hullfill
\hullface{topsides!57!black}\hullmove{6.07}{4.55}\hullline{5.47}{4.55}\hullline{4.42}{3.04}\hullline{5.07}{3.04}\hullfill
\hullface{topsides!81!black}\hullmove{-5.07}{3.04}\hullline{-4.42}{3.04}\hullline{-5.47}{4.55}\hullline{-6.07}{4.55}\hullfill
\hullface{topsides!57!black}\hullmove{7.22}{6.07}\hullline{6.67}{6.07}\hullline{5.47}{4.55}\hullline{6.07}{4.55}\hullfill
\hullface{topsides!75!black}\hullmove{-6.07}{4.55}\hullline{-5.47}{4.55}\hullline{-6.67}{6.07}\hullline{-7.22}{6.07}\hullfill
\hullface{topsides!57!black}\hullmove{8.49}{7.59}\hullline{7.99}{7.59}\hullline{6.67}{6.07}\hullline{7.22}{6.07}\hullfill
\hullface{topsides!72!black}\hullmove{-7.22}{6.07}\hullline{-6.67}{6.07}\hullline{-7.99}{7.59}\hullline{-8.49}{7.59}\hullfill
\hullface{keel!57!black}\hullmove{0.16}{-7.89}\hullline{0.06}{-7.36}\hullline{0.00}{-7.45}\hullline{0.00}{-8.15}\hullfill
\hullface{keel!78!black}\hullmove{0.00}{-8.15}\hullline{0.00}{-7.45}\hullline{-0.06}{-7.36}\hullline{-0.16}{-7.89}\hullfill
\hullface{keel!57!black}\hullmove{0.45}{-7.42}\hullline{0.57}{-7.22}\hullline{0.06}{-7.36}\hullline{0.16}{-7.89}\hullfill
\hullface{keel!60!black}\hullmove{-0.16}{-7.89}\hullline{-0.06}{-7.36}\hullline{-0.57}{-7.22}\hullline{-0.45}{-7.42}\hullfill
\hullface{keel!57!black}\hullmove{0.48}{-7.36}\hullline{0.77}{-7.10}\hullline{0.57}{-7.22}\hullline{0.45}{-7.42}\hullfill
\hullface{keel!57!black}\hullmove{-0.45}{-7.42}\hullline{-0.57}{-7.22}\hullline{-0.77}{-7.10}\hullline{-0.48}{-7.36}\hullfill
\hullface{keel!57!black}\hullmove{0.77}{-7.10}\hullline{1.09}{-6.79}\hullline{0.77}{-7.10}\hullline{0.48}{-7.36}\hullfill
\hullface{keel!57!black}\hullmove{-0.48}{-7.36}\hullline{-0.77}{-7.10}\hullline{-1.09}{-6.79}\hullline{-0.77}{-7.10}\hullfill
\hullface{keel!57!black}\hullmove{1.08}{-6.81}\hullline{1.33}{-6.40}\hullline{1.09}{-6.79}\hullline{0.77}{-7.10}\hullfill
\hullface{keel!57!black}\hullmove{1.34}{-6.40}\hullline{1.40}{-6.23}\hullline{1.33}{-6.40}\hullline{1.08}{-6.81}\hullfill
\hullface{keel!60!black}\hullmove{-1.08}{-6.81}\hullline{-1.33}{-6.40}\hullline{-1.40}{-6.23}\hullline{-1.34}{-6.40}\hullfill
\hullface{keel!57!black}\hullmove{1.45}{-6.09}\hullline{1.53}{-5.57}\hullline{1.40}{-6.23}\hullline{1.34}{-6.40}\hullfill
\hullface{keel!57!black}\hullmove{1.60}{-5.27}\hullline{1.46}{-4.82}\hullline{1.54}{-5.44}\hullline{1.54}{-5.44}\hullfill
\hullface{keel!84!black}\hullmove{-1.54}{-5.44}\hullline{-1.54}{-5.44}\hullline{-1.46}{-4.82}\hullline{-1.60}{-5.27}\hullfill
\hullface{keel!57!black}\hullmove{1.95}{-4.47}\hullline{1.50}{-4.47}\hullline{1.46}{-4.82}\hullline{1.60}{-5.27}\hullfill
\hullface{keel!84!black}\hullmove{-1.60}{-5.27}\hullline{-1.46}{-4.82}\hullline{-1.50}{-4.47}\hullline{-1.95}{-4.47}\hullfill
\hullface{keel!57!black}\hullmove{2.00}{-4.36}\hullline{1.67}{-3.99}\hullline{1.50}{-4.47}\hullline{1.95}{-4.47}\hullfill
\hullface{keel!87!black}\hullmove{-1.95}{-4.47}\hullline{-1.50}{-4.47}\hullline{-1.67}{-3.99}\hullline{-2.00}{-4.36}\hullfill
\hullface{keel!57!black}\hullmove{2.23}{-3.77}\hullline{1.74}{-3.77}\hullline{1.67}{-3.99}\hullline{2.00}{-4.36}\hullfill
\hullface{keel!87!black}\hullmove{-2.00}{-4.36}\hullline{-1.67}{-3.99}\hullline{-1.74}{-3.77}\hullline{-2.23}{-3.77}\hullfill
\hullface{keel!57!black}\hullmove{2.33}{-3.51}\hullline{1.82}{-3.51}\hullline{1.74}{-3.77}\hullline{2.23}{-3.77}\hullfill
\hullface{keel!90!black}\hullmove{-2.23}{-3.77}\hullline{-1.74}{-3.77}\hullline{-1.82}{-3.51}\hullline{-2.33}{-3.51}\hullfill
\hullface{keel!57!black}\hullmove{2.37}{-3.38}\hullline{1.93}{-3.09}\hullline{1.82}{-3.51}\hullline{2.33}{-3.51}\hullfill
\hullface{keel!90!black}\hullmove{-2.33}{-3.51}\hullline{-1.82}{-3.51}\hullline{-1.93}{-3.09}\hullline{-2.37}{-3.38}\hullfill
\hullface{keel!57!black}\hullmove{2.70}{-2.32}\hullline{2.14}{-2.12}\hullline{1.93}{-3.09}\hullline{2.37}{-3.38}\hullfill
\hullface{keel!93!black}\hullmove{-2.37}{-3.38}\hullline{-1.93}{-3.09}\hullline{-2.14}{-2.12}\hullline{-2.70}{-2.32}\hullfill
\hullface{keel!60!black}\hullmove{2.94}{-1.19}\hullline{2.26}{-1.09}\hullline{2.14}{-2.12}\hullline{2.70}{-2.32}\hullfill
\hullface{keel!96!black}\hullmove{-2.70}{-2.32}\hullline{-2.14}{-2.12}\hullline{-2.26}{-1.09}\hullline{-2.94}{-1.19}\hullfill
\hullface{keel!66!black}\hullmove{3.06}{0.00}\hullline{2.31}{0.00}\hullline{2.26}{-1.09}\hullline{2.94}{-1.19}\hullfill
\hullface{keel!98!white}\hullmove{-2.94}{-1.19}\hullline{-2.26}{-1.09}\hullline{-2.31}{0.00}\hullline{-3.06}{0.00}\hullfill
\hullface{topsides!60!black}\hullmove{3.58}{1.52}\hullline{2.84}{1.52}\hullline{2.31}{0.00}\hullline{3.06}{0.00}\hullfill
\hullface{topsides!93!black}\hullmove{-3.06}{0.00}\hullline{-2.31}{0.00}\hullline{-2.84}{1.52}\hullline{-3.58}{1.52}\hullfill
\hullface{topsides!57!black}\hullmove{4.42}{3.04}\hullline{3.72}{3.04}\hullline{2.84}{1.52}\hullline{3.58}{1.52}\hullfill
\hullface{topsides!87!black}\hullmove{-3.58}{1.52}\hullline{-2.84}{1.52}\hullline{-3.72}{3.04}\hullline{-4.42}{3.04}\hullfill
\hullface{topsides!57!black}\hullmove{5.47}{4.55}\hullline{4.80}{4.55}\hullline{3.72}{3.04}\hullline{4.42}{3.04}\hullfill
\hullface{topsides!81!black}\hullmove{-4.42}{3.04}\hullline{-3.72}{3.04}\hullline{-4.80}{4.55}\hullline{-5.47}{4.55}\hullfill
\hullface{topsides!57!black}\hullmove{6.67}{6.07}\hullline{6.04}{6.07}\hullline{4.80}{4.55}\hullline{5.47}{4.55}\hullfill
\hullface{topsides!78!black}\hullmove{-5.47}{4.55}\hullline{-4.80}{4.55}\hullline{-6.04}{6.07}\hullline{-6.67}{6.07}\hullfill
\hullface{topsides!57!black}\hullmove{7.99}{7.59}\hullline{7.41}{7.59}\hullline{6.04}{6.07}\hullline{6.67}{6.07}\hullfill
\hullface{topsides!75!black}\hullmove{-6.67}{6.07}\hullline{-6.04}{6.07}\hullline{-7.41}{7.59}\hullline{-7.99}{7.59}\hullfill
\hullface{keel!57!black}\hullmove{0.06}{-7.36}\hullline{0.00}{-7.32}\hullline{0.00}{-7.32}\hullline{0.00}{-7.45}\hullfill
\hullface{keel!63!black}\hullmove{0.00}{-7.45}\hullline{-0.00}{-7.32}\hullline{-0.00}{-7.32}\hullline{-0.06}{-7.36}\hullfill
\hullface{keel!57!black}\hullmove{0.57}{-7.22}\hullline{0.72}{-7.09}\hullline{0.00}{-7.32}\hullline{0.06}{-7.36}\hullfill
\hullface{keel!57!black}\hullmove{-0.06}{-7.36}\hullline{-0.00}{-7.32}\hullline{-0.72}{-7.09}\hullline{-0.57}{-7.22}\hullfill
\hullface{keel!57!black}\hullmove{0.77}{-7.10}\hullline{0.75}{-7.07}\hullline{0.72}{-7.09}\hullline{0.57}{-7.22}\hullfill
\hullface{keel!57!black}\hullmove{-0.57}{-7.22}\hullline{-0.72}{-7.09}\hullline{-0.75}{-7.07}\hullline{-0.77}{-7.10}\hullfill
\hullface{keel!57!black}\hullmove{-0.77}{-7.10}\hullline{-0.75}{-7.07}\hullline{-1.14}{-6.67}\hullline{-1.09}{-6.79}\hullfill
\hullface{keel!57!black}\hullmove{1.33}{-6.40}\hullline{1.31}{-6.38}\hullline{1.14}{-6.67}\hullline{1.09}{-6.79}\hullfill
\hullface{keel!60!black}\hullmove{-1.09}{-6.79}\hullline{-1.14}{-6.67}\hullline{-1.31}{-6.38}\hullline{-1.33}{-6.40}\hullfill
\hullface{keel!57!black}\hullmove{1.40}{-6.23}\hullline{1.41}{-6.12}\hullline{1.31}{-6.38}\hullline{1.33}{-6.40}\hullfill
\hullface{keel!57!black}\hullmove{1.53}{-5.57}\hullline{1.51}{-5.47}\hullline{1.41}{-6.12}\hullline{1.40}{-6.23}\hullfill
\hullface{keel!75!black}\hullmove{-1.40}{-6.23}\hullline{-1.41}{-6.12}\hullline{-1.51}{-5.47}\hullline{-1.53}{-5.57}\hullfill
\hullface{keel!81!black}\hullmove{-1.53}{-5.57}\hullline{-1.51}{-5.47}\hullline{-1.51}{-5.44}\hullline{-1.54}{-5.44}\hullfill
\hullface{keel!57!black}\hullmove{1.46}{-4.82}\hullline{1.40}{-4.73}\hullline{1.51}{-5.44}\hullline{1.54}{-5.44}\hullfill
\hullface{keel!87!black}\hullmove{-1.54}{-5.44}\hullline{-1.51}{-5.44}\hullline{-1.40}{-4.73}\hullline{-1.46}{-4.82}\hullfill
\hullface{keel!57!black}\hullmove{1.50}{-4.47}\hullline{1.31}{-4.49}\hullline{1.40}{-4.73}\hullline{1.46}{-4.82}\hullfill
\hullface{keel!90!black}\hullmove{-1.46}{-4.82}\hullline{-1.40}{-4.73}\hullline{-1.31}{-4.49}\hullline{-1.50}{-4.47}\hullfill
\hullface{keel!57!black}\hullmove{1.67}{-3.99}\hullline{1.19}{-3.92}\hullline{1.31}{-4.49}\hullline{1.50}{-4.47}\hullfill
\hullface{keel!93!black}\hullmove{-1.50}{-4.47}\hullline{-1.31}{-4.49}\hullline{-1.19}{-3.92}\hullline{-1.67}{-3.99}\hullfill
\hullface{keel!57!black}\hullmove{1.74}{-3.77}\hullline{1.22}{-3.80}\hullline{1.19}{-3.92}\hullline{1.67}{-3.99}\hullfill
\hullface{keel!90!black}\hullmove{-1.67}{-3.99}\hullline{-1.19}{-3.92}\hullline{-1.22}{-3.80}\hullline{-1.74}{-3.77}\hullfill
\hullface{keel!57!black}\hullmove{1.82}{-3.51}\hullline{1.27}{-3.55}\hullline{1.22}{-3.80}\hullline{1.74}{-3.77}\hullfill
\hullface{keel!93!black}\hullmove{-1.74}{-3.77}\hullline{-1.22}{-3.80}\hullline{-1.27}{-3.55}\hullline{-1.82}{-3.51}\hullfill
\hullface{keel!57!black}\hullmove{1.93}{-3.09}\hullline{1.36}{-3.03}\hullline{1.27}{-3.55}\hullline{1.82}{-3.51}\hullfill
\hullface{keel!93!black}\hullmove{-1.82}{-3.51}\hullline{-1.27}{-3.55}\hullline{-1.36}{-3.03}\hullline{-1.93}{-3.09}\hullfill
\hullface{keel!60!black}\hullmove{2.14}{-2.12}\hullline{1.46}{-2.08}\hullline{1.36}{-3.03}\hullline{1.93}{-3.09}\hullfill
\hullface{keel!96!black}\hullmove{-1.93}{-3.09}\hullline{-1.36}{-3.03}\hullline{-1.46}{-2.08}\hullline{-2.14}{-2.12}\hullfill
\hullface{keel!66!black}\hullmove{2.26}{-1.09}\hullline{1.50}{-1.07}\hullline{1.46}{-2.08}\hullline{2.14}{-2.12}\hullfill
\hullface{keel!98!white}\hullmove{-2.14}{-2.12}\hullline{-1.46}{-2.08}\hullline{-1.50}{-1.07}\hullline{-2.26}{-1.09}\hullfill
\hullface{keel!69!black}\hullmove{2.31}{0.00}\hullline{1.51}{0.00}\hullline{1.50}{-1.07}\hullline{2.26}{-1.09}\hullfill
\hullface{keel!95!white}\hullmove{-2.26}{-1.09}\hullline{-1.50}{-1.07}\hullline{-1.51}{0.00}\hullline{-2.31}{0.00}\hullfill
\hullface{topsides!60!black}\hullmove{2.84}{1.52}\hullline{2.06}{1.52}\hullline{1.51}{0.00}\hullline{2.31}{0.00}\hullfill
\hullface{topsides!93!black}\hullmove{-2.31}{0.00}\hullline{-1.51}{0.00}\hullline{-2.06}{1.52}\hullline{-2.84}{1.52}\hullfill
\hullface{topsides!57!black}\hullmove{3.72}{3.04}\hullline{2.96}{3.04}\hullline{2.06}{1.52}\hullline{2.84}{1.52}\hullfill
\hullface{topsides!87!black}\hullmove{-2.84}{1.52}\hullline{-2.06}{1.52}\hullline{-2.96}{3.04}\hullline{-3.72}{3.04}\hullfill
\hullface{topsides!57!black}\hullmove{4.80}{4.55}\hullline{4.06}{4.55}\hullline{2.96}{3.04}\hullline{3.72}{3.04}\hullfill
\hullface{topsides!81!black}\hullmove{-3.72}{3.04}\hullline{-2.96}{3.04}\hullline{-4.06}{4.55}\hullline{-4.80}{4.55}\hullfill
\hullface{topsides!57!black}\hullmove{6.04}{6.07}\hullline{5.33}{6.07}\hullline{4.06}{4.55}\hullline{4.80}{4.55}\hullfill
\hullface{topsides!78!black}\hullmove{-4.80}{4.55}\hullline{-4.06}{4.55}\hullline{-5.33}{6.07}\hullline{-6.04}{6.07}\hullfill
\hullface{topsides!57!black}\hullmove{7.41}{7.59}\hullline{6.73}{7.59}\hullline{5.33}{6.07}\hullline{6.04}{6.07}\hullfill
\hullface{topsides!75!black}\hullmove{-6.04}{6.07}\hullline{-5.33}{6.07}\hullline{-6.73}{7.59}\hullline{-7.41}{7.59}\hullfill
\hullface{keel!57!black}\hullmove{0.72}{-7.09}\hullline{0.71}{-7.02}\hullline{0.00}{-7.25}\hullline{0.00}{-7.32}\hullfill
\hullface{keel!57!black}\hullmove{-0.00}{-7.32}\hullline{0.00}{-7.25}\hullline{-0.71}{-7.02}\hullline{-0.72}{-7.09}\hullfill
\hullface{keel!57!black}\hullmove{-0.72}{-7.09}\hullline{-0.71}{-7.02}\hullline{-0.72}{-7.00}\hullline{-0.75}{-7.07}\hullfill
\hullface{keel!57!black}\hullmove{1.14}{-6.67}\hullline{1.11}{-6.60}\hullline{0.72}{-7.00}\hullline{0.75}{-7.07}\hullfill
\hullface{keel!57!black}\hullmove{-0.75}{-7.07}\hullline{-0.72}{-7.00}\hullline{-1.11}{-6.60}\hullline{-1.14}{-6.67}\hullfill
\hullface{keel!57!black}\hullmove{1.31}{-6.38}\hullline{1.25}{-6.34}\hullline{1.11}{-6.60}\hullline{1.14}{-6.67}\hullfill
\hullface{keel!63!black}\hullmove{-1.14}{-6.67}\hullline{-1.11}{-6.60}\hullline{-1.25}{-6.34}\hullline{-1.31}{-6.38}\hullfill
\hullface{keel!57!black}\hullmove{1.41}{-6.12}\hullline{1.36}{-6.06}\hullline{1.25}{-6.34}\hullline{1.31}{-6.38}\hullfill
\hullface{keel!69!black}\hullmove{-1.31}{-6.38}\hullline{-1.25}{-6.34}\hullline{-1.36}{-6.06}\hullline{-1.41}{-6.12}\hullfill
\hullface{keel!57!black}\hullmove{1.51}{-5.47}\hullline{1.45}{-5.44}\hullline{1.36}{-6.06}\hullline{1.41}{-6.12}\hullfill
\hullface{keel!78!black}\hullmove{-1.41}{-6.12}\hullline{-1.36}{-6.06}\hullline{-1.45}{-5.44}\hullline{-1.51}{-5.47}\hullfill
\hullface{keel!57!black}\hullmove{1.51}{-5.44}\hullline{1.45}{-5.41}\hullline{1.45}{-5.44}\hullline{1.51}{-5.47}\hullfill
\hullface{keel!84!black}\hullmove{-1.51}{-5.47}\hullline{-1.45}{-5.44}\hullline{-1.45}{-5.41}\hullline{-1.51}{-5.44}\hullfill
\hullface{keel!57!black}\hullmove{1.40}{-4.73}\hullline{1.32}{-4.68}\hullline{1.45}{-5.41}\hullline{1.51}{-5.44}\hullfill
\hullface{keel!90!black}\hullmove{-1.51}{-5.44}\hullline{-1.45}{-5.41}\hullline{-1.32}{-4.68}\hullline{-1.40}{-4.73}\hullfill
\hullface{keel!57!black}\hullmove{1.31}{-4.49}\hullline{1.25}{-4.53}\hullline{1.32}{-4.68}\hullline{1.40}{-4.73}\hullfill
\hullface{keel!93!black}\hullmove{-1.40}{-4.73}\hullline{-1.32}{-4.68}\hullline{-1.25}{-4.53}\hullline{-1.31}{-4.49}\hullfill
\hullface{keel!63!black}\hullmove{1.19}{-3.92}\hullline{0.74}{-3.88}\hullline{1.25}{-4.53}\hullline{1.31}{-4.49}\hullfill
\hullface{keel!95!white}\hullmove{-1.31}{-4.49}\hullline{-1.25}{-4.53}\hullline{-0.74}{-3.88}\hullline{-1.19}{-3.92}\hullfill
\hullface{keel!57!black}\hullmove{1.22}{-3.80}\hullline{0.72}{-3.87}\hullline{0.74}{-3.88}\hullline{1.19}{-3.92}\hullfill
\hullface{keel!96!black}\hullmove{-1.19}{-3.92}\hullline{-0.74}{-3.88}\hullline{-0.72}{-3.87}\hullline{-1.22}{-3.80}\hullfill
\hullface{keel!63!black}\hullmove{1.27}{-3.55}\hullline{0.64}{-3.62}\hullline{0.72}{-3.87}\hullline{1.22}{-3.80}\hullfill
\hullface{keel!95!white}\hullmove{-1.22}{-3.80}\hullline{-0.72}{-3.87}\hullline{-0.64}{-3.62}\hullline{-1.27}{-3.55}\hullfill
\hullface{keel!63!black}\hullmove{1.36}{-3.03}\hullline{0.66}{-3.00}\hullline{0.64}{-3.62}\hullline{1.27}{-3.55}\hullfill
\hullface{keel!99!black}\hullmove{-1.27}{-3.55}\hullline{-0.64}{-3.62}\hullline{-0.66}{-3.00}\hullline{-1.36}{-3.03}\hullfill
\hullface{keel!66!black}\hullmove{1.46}{-2.08}\hullline{0.67}{-2.06}\hullline{0.66}{-3.00}\hullline{1.36}{-3.03}\hullfill
\hullface{keel!95!white}\hullmove{-1.36}{-3.03}\hullline{-0.66}{-3.00}\hullline{-0.67}{-2.06}\hullline{-1.46}{-2.08}\hullfill
\hullface{keel!69!black}\hullmove{1.50}{-1.07}\hullline{0.67}{-1.06}\hullline{0.67}{-2.06}\hullline{1.46}{-2.08}\hullfill
\hullface{keel!92!white}\hullmove{-1.46}{-2.08}\hullline{-0.67}{-2.06}\hullline{-0.67}{-1.06}\hullline{-1.50}{-1.07}\hullfill
\hullface{keel!72!black}\hullmove{1.51}{0.00}\hullline{0.67}{0.00}\hullline{0.67}{-1.06}\hullline{1.50}{-1.07}\hullfill
\hullface{keel!92!white}\hullmove{-1.50}{-1.07}\hullline{-0.67}{-1.06}\hullline{-0.67}{0.00}\hullline{-1.51}{0.00}\hullfill
\hullface{topsides!63!black}\hullmove{2.06}{1.52}\hullline{1.22}{1.52}\hullline{0.67}{0.00}\hullline{1.51}{0.00}\hullfill
\hullface{topsides!96!black}\hullmove{-1.51}{0.00}\hullline{-0.67}{0.00}\hullline{-1.22}{1.52}\hullline{-2.06}{1.52}\hullfill
\hullface{topsides!57!black}\hullmove{2.96}{3.04}\hullline{2.13}{3.04}\hullline{1.22}{1.52}\hullline{2.06}{1.52}\hullfill
\hullface{topsides!87!black}\hullmove{-2.06}{1.52}\hullline{-1.22}{1.52}\hullline{-2.13}{3.04}\hullline{-2.96}{3.04}\hullfill
\hullface{topsides!57!black}\hullmove{4.06}{4.55}\hullline{3.25}{4.55}\hullline{2.13}{3.04}\hullline{2.96}{3.04}\hullfill
\hullface{topsides!84!black}\hullmove{-2.96}{3.04}\hullline{-2.13}{3.04}\hullline{-3.25}{4.55}\hullline{-4.06}{4.55}\hullfill
\hullface{topsides!57!black}\hullmove{5.33}{6.07}\hullline{4.53}{6.07}\hullline{3.25}{4.55}\hullline{4.06}{4.55}\hullfill
\hullface{topsides!81!black}\hullmove{-4.06}{4.55}\hullline{-3.25}{4.55}\hullline{-4.53}{6.07}\hullline{-5.33}{6.07}\hullfill
\hullface{topsides!57!black}\hullmove{6.73}{7.59}\hullline{5.95}{7.59}\hullline{4.53}{6.07}\hullline{5.33}{6.07}\hullfill
\hullface{topsides!78!black}\hullmove{-5.33}{6.07}\hullline{-4.53}{6.07}\hullline{-5.95}{7.59}\hullline{-6.73}{7.59}\hullfill
\hullface{keel!57!black}\hullmove{0.71}{-7.02}\hullline{0.68}{-6.91}\hullline{0.00}{-7.13}\hullline{0.00}{-7.25}\hullfill
\hullface{keel!57!black}\hullmove{0.00}{-7.25}\hullline{-0.00}{-7.13}\hullline{-0.68}{-6.91}\hullline{-0.71}{-7.02}\hullfill
\hullface{keel!57!black}\hullmove{0.72}{-7.00}\hullline{0.68}{-6.90}\hullline{0.68}{-6.91}\hullline{0.71}{-7.02}\hullfill
\hullface{keel!57!black}\hullmove{-0.71}{-7.02}\hullline{-0.68}{-6.91}\hullline{-0.68}{-6.90}\hullline{-0.72}{-7.00}\hullfill
\hullface{keel!57!black}\hullmove{1.11}{-6.60}\hullline{1.06}{-6.49}\hullline{0.68}{-6.90}\hullline{0.72}{-7.00}\hullfill
\hullface{keel!57!black}\hullmove{-0.72}{-7.00}\hullline{-0.68}{-6.90}\hullline{-1.06}{-6.49}\hullline{-1.11}{-6.60}\hullfill
\hullface{keel!57!black}\hullmove{1.25}{-6.34}\hullline{1.18}{-6.28}\hullline{1.06}{-6.49}\hullline{1.11}{-6.60}\hullfill
\hullface{keel!66!black}\hullmove{-1.11}{-6.60}\hullline{-1.06}{-6.49}\hullline{-1.18}{-6.28}\hullline{-1.25}{-6.34}\hullfill
\hullface{keel!57!black}\hullmove{1.36}{-6.06}\hullline{1.29}{-5.96}\hullline{1.18}{-6.28}\hullline{1.25}{-6.34}\hullfill
\hullface{keel!72!black}\hullmove{-1.25}{-6.34}\hullline{-1.18}{-6.28}\hullline{-1.29}{-5.96}\hullline{-1.36}{-6.06}\hullfill
\hullface{keel!57!black}\hullmove{1.45}{-5.44}\hullline{1.36}{-5.44}\hullline{1.29}{-5.96}\hullline{1.36}{-6.06}\hullfill
\hullface{keel!81!black}\hullmove{-1.36}{-6.06}\hullline{-1.29}{-5.96}\hullline{-1.36}{-5.44}\hullline{-1.45}{-5.44}\hullfill
\hullface{keel!57!black}\hullmove{1.45}{-5.41}\hullline{1.35}{-5.33}\hullline{1.36}{-5.44}\hullline{1.45}{-5.44}\hullfill
\hullface{keel!87!black}\hullmove{-1.45}{-5.44}\hullline{-1.36}{-5.44}\hullline{-1.35}{-5.33}\hullline{-1.45}{-5.41}\hullfill
\hullface{keel!57!black}\hullmove{1.32}{-4.68}\hullline{1.19}{-4.61}\hullline{1.35}{-5.33}\hullline{1.45}{-5.41}\hullfill
\hullface{keel!90!black}\hullmove{-1.45}{-5.41}\hullline{-1.35}{-5.33}\hullline{-1.19}{-4.61}\hullline{-1.32}{-4.68}\hullfill
\hullface{keel!57!black}\hullmove{1.25}{-4.53}\hullline{1.18}{-4.59}\hullline{1.19}{-4.61}\hullline{1.32}{-4.68}\hullfill
\hullface{keel!96!black}\hullmove{-1.32}{-4.68}\hullline{-1.19}{-4.61}\hullline{-1.18}{-4.59}\hullline{-1.25}{-4.53}\hullfill
\hullface{keel!63!black}\hullmove{0.74}{-3.88}\hullline{0.68}{-3.97}\hullline{1.18}{-4.59}\hullline{1.25}{-4.53}\hullfill
\hullface{keel!98!white}\hullmove{-1.25}{-4.53}\hullline{-1.18}{-4.59}\hullline{-0.68}{-3.97}\hullline{-0.74}{-3.88}\hullfill
\hullface{keel!78!black}\hullmove{0.72}{-3.87}\hullline{0.41}{-3.82}\hullline{0.68}{-3.97}\hullline{0.74}{-3.88}\hullfill
\hullface{keel!95!white}\hullmove{-0.74}{-3.88}\hullline{-0.68}{-3.97}\hullline{-0.41}{-3.82}\hullline{-0.72}{-3.87}\hullfill
\hullface{keel!84!black}\hullmove{0.64}{-3.62}\hullline{0.00}{-3.74}\hullline{0.41}{-3.82}\hullline{0.72}{-3.87}\hullfill
\hullface{keel!83!white}\hullmove{-0.72}{-3.87}\hullline{-0.41}{-3.82}\hullline{-0.00}{-3.74}\hullline{-0.64}{-3.62}\hullfill
\hullface{keel!63!black}\hullmove{0.66}{-3.00}\hullline{0.67}{-2.06}\hullline{0.00}{-2.03}\hullline{0.00}{-2.95}\hullline{0.00}{-3.74}\hullline{0.64}{-3.62}\hullfill
\hullface{keel!95!white}\hullmove{-0.64}{-3.62}\hullline{-0.00}{-3.74}\hullline{0.00}{-2.95}\hullline{0.00}{-2.03}\hullline{-0.67}{-2.06}\hullline{-0.66}{-3.00}\hullfill
\hullface{keel!66!black}\hullmove{0.67}{-1.06}\hullline{0.67}{0.00}\hullline{0.67}{0.00}\hullline{0.00}{0.38}\hullline{0.00}{0.00}\hullline{0.00}{-1.04}\hullline{0.00}{-2.03}\hullline{0.67}{-2.06}\hullfill
\hullface{keel!95!white}\hullmove{-0.67}{-2.06}\hullline{0.00}{-2.03}\hullline{0.00}{-1.04}\hullline{0.00}{0.00}\hullline{0.00}{0.38}\hullline{-0.67}{0.00}\hullline{-0.67}{0.00}\hullline{-0.67}{-1.06}\hullfill
\hullface{topsides!60!black}\hullmove{1.22}{1.52}\hullline{0.53}{1.83}\hullline{0.00}{0.38}\hullline{0.67}{0.00}\hullfill
\hullface{topsides!93!black}\hullmove{-0.67}{0.00}\hullline{0.00}{0.38}\hullline{-0.53}{1.83}\hullline{-1.22}{1.52}\hullfill
\hullface{topsides!57!black}\hullmove{2.13}{3.04}\hullline{1.40}{3.27}\hullline{0.53}{1.83}\hullline{1.22}{1.52}\hullfill
\hullface{topsides!87!black}\hullmove{-1.22}{1.52}\hullline{-0.53}{1.83}\hullline{-1.40}{3.27}\hullline{-2.13}{3.04}\hullfill
\hullface{topsides!57!black}\hullmove{3.25}{4.55}\hullline{2.47}{4.71}\hullline{1.40}{3.27}\hullline{2.13}{3.04}\hullfill
\hullface{topsides!84!black}\hullmove{-2.13}{3.04}\hullline{-1.40}{3.27}\hullline{-2.47}{4.71}\hullline{-3.25}{4.55}\hullfill
\hullface{topsides!57!black}\hullmove{4.53}{6.07}\hullline{3.70}{6.15}\hullline{2.47}{4.71}\hullline{3.25}{4.55}\hullfill
\hullface{topsides!81!black}\hullmove{-3.25}{4.55}\hullline{-2.47}{4.71}\hullline{-3.70}{6.15}\hullline{-4.53}{6.07}\hullfill
\hullface{topsides!57!black}\hullmove{5.95}{7.59}\hullline{5.06}{7.59}\hullline{3.70}{6.15}\hullline{4.53}{6.07}\hullfill
\hullface{topsides!81!black}\hullmove{-4.53}{6.07}\hullline{-3.70}{6.15}\hullline{-5.06}{7.59}\hullline{-5.95}{7.59}\hullfill
\hullface{keel!57!black}\hullmove{0.68}{-6.91}\hullline{0.61}{-6.76}\hullline{0.00}{-6.97}\hullline{0.00}{-7.13}\hullfill
\hullface{keel!57!black}\hullmove{-0.00}{-7.13}\hullline{0.00}{-6.97}\hullline{-0.61}{-6.76}\hullline{-0.68}{-6.91}\hullfill
\hullface{keel!57!black}\hullmove{0.68}{-6.90}\hullline{0.63}{-6.75}\hullline{0.61}{-6.76}\hullline{0.68}{-6.91}\hullfill
\hullface{keel!57!black}\hullmove{-0.68}{-6.91}\hullline{-0.61}{-6.76}\hullline{-0.63}{-6.75}\hullline{-0.68}{-6.90}\hullfill
\hullface{keel!57!black}\hullmove{1.06}{-6.49}\hullline{0.99}{-6.34}\hullline{0.63}{-6.75}\hullline{0.68}{-6.90}\hullfill
\hullface{keel!60!black}\hullmove{-0.68}{-6.90}\hullline{-0.63}{-6.75}\hullline{-0.99}{-6.34}\hullline{-1.06}{-6.49}\hullfill
\hullface{keel!57!black}\hullmove{1.18}{-6.28}\hullline{1.06}{-6.20}\hullline{0.99}{-6.34}\hullline{1.06}{-6.49}\hullfill
\hullface{keel!69!black}\hullmove{-1.06}{-6.49}\hullline{-0.99}{-6.34}\hullline{-1.06}{-6.20}\hullline{-1.18}{-6.28}\hullfill
\hullface{keel!57!black}\hullmove{1.29}{-5.96}\hullline{1.19}{-5.82}\hullline{1.06}{-6.20}\hullline{1.18}{-6.28}\hullfill
\hullface{keel!75!black}\hullmove{-1.18}{-6.28}\hullline{-1.06}{-6.20}\hullline{-1.19}{-5.82}\hullline{-1.29}{-5.96}\hullfill
\hullface{keel!57!black}\hullmove{1.36}{-5.44}\hullline{1.23}{-5.44}\hullline{1.19}{-5.82}\hullline{1.29}{-5.96}\hullfill
\hullface{keel!84!black}\hullmove{-1.29}{-5.96}\hullline{-1.19}{-5.82}\hullline{-1.23}{-5.44}\hullline{-1.36}{-5.44}\hullfill
\hullface{keel!57!black}\hullmove{1.35}{-5.33}\hullline{1.21}{-5.21}\hullline{1.23}{-5.44}\hullline{1.36}{-5.44}\hullfill
\hullface{keel!87!black}\hullmove{-1.36}{-5.44}\hullline{-1.23}{-5.44}\hullline{-1.21}{-5.21}\hullline{-1.35}{-5.33}\hullfill
\hullface{keel!57!black}\hullmove{1.19}{-4.61}\hullline{1.06}{-4.67}\hullline{1.21}{-5.21}\hullline{1.35}{-5.33}\hullfill
\hullface{keel!93!black}\hullmove{-1.35}{-5.33}\hullline{-1.21}{-5.21}\hullline{-1.06}{-4.67}\hullline{-1.19}{-4.61}\hullfill
\hullface{keel!57!black}\hullmove{1.18}{-4.59}\hullline{0.97}{-4.50}\hullline{1.06}{-4.67}\hullline{1.19}{-4.61}\hullfill
\hullface{keel!99!black}\hullmove{-1.19}{-4.61}\hullline{-1.06}{-4.67}\hullline{-0.97}{-4.50}\hullline{-1.18}{-4.59}\hullfill
\hullface{keel!66!black}\hullmove{0.68}{-3.97}\hullline{0.61}{-4.11}\hullline{0.97}{-4.50}\hullline{1.18}{-4.59}\hullfill
\hullface{keel!98!white}\hullmove{-1.18}{-4.59}\hullline{-0.97}{-4.50}\hullline{-0.61}{-4.11}\hullline{-0.68}{-3.97}\hullfill
\hullface{keel!87!black}\hullmove{0.41}{-3.82}\hullline{0.00}{-3.90}\hullline{0.61}{-4.11}\hullline{0.68}{-3.97}\hullfill
\hullface{keel!92!white}\hullmove{-0.68}{-3.97}\hullline{-0.61}{-4.11}\hullline{0.00}{-3.90}\hullline{-0.41}{-3.82}\hullfill
\hullface{keel!78!black}\hullmove{0.00}{-3.74}\hullline{0.00}{-3.73}\hullline{0.00}{-3.90}\hullline{0.41}{-3.82}\hullfill
\hullface{keel!92!white}\hullmove{-0.41}{-3.82}\hullline{0.00}{-3.90}\hullline{0.00}{-3.73}\hullline{-0.00}{-3.74}\hullfill
\hullface{topsides!57!black}\hullmove{0.53}{1.83}\hullline{0.42}{3.07}\hullline{0.00}{1.94}\hullline{0.00}{0.38}\hullfill
\hullface{topsides!90!black}\hullmove{0.00}{0.38}\hullline{0.00}{1.94}\hullline{-0.42}{3.07}\hullline{-0.53}{1.83}\hullfill
\hullface{topsides!57!black}\hullmove{1.40}{3.27}\hullline{1.12}{4.20}\hullline{0.42}{3.07}\hullline{0.53}{1.83}\hullfill
\hullface{topsides!87!black}\hullmove{-0.53}{1.83}\hullline{-0.42}{3.07}\hullline{-1.12}{4.20}\hullline{-1.40}{3.27}\hullfill
\hullface{topsides!57!black}\hullmove{2.47}{4.71}\hullline{1.97}{5.33}\hullline{1.12}{4.20}\hullline{1.40}{3.27}\hullfill
\hullface{topsides!87!black}\hullmove{-1.40}{3.27}\hullline{-1.12}{4.20}\hullline{-1.97}{5.33}\hullline{-2.47}{4.71}\hullfill
\hullface{topsides!57!black}\hullmove{3.70}{6.15}\hullline{2.95}{6.46}\hullline{1.97}{5.33}\hullline{2.47}{4.71}\hullfill
\hullface{topsides!84!black}\hullmove{-2.47}{4.71}\hullline{-1.97}{5.33}\hullline{-2.95}{6.46}\hullline{-3.70}{6.15}\hullfill
\hullface{topsides!57!black}\hullmove{5.06}{7.59}\hullline{4.03}{7.59}\hullline{2.95}{6.46}\hullline{3.70}{6.15}\hullfill
\hullface{topsides!81!black}\hullmove{-3.70}{6.15}\hullline{-2.95}{6.46}\hullline{-4.03}{7.59}\hullline{-5.06}{7.59}\hullfill
\hullface{keel!57!black}\hullmove{0.61}{-6.76}\hullline{0.52}{-6.56}\hullline{0.00}{-6.74}\hullline{0.00}{-6.97}\hullfill
\hullface{keel!57!black}\hullmove{0.00}{-6.97}\hullline{0.00}{-6.74}\hullline{-0.52}{-6.56}\hullline{-0.61}{-6.76}\hullfill
\hullface{keel!57!black}\hullmove{0.63}{-6.75}\hullline{0.57}{-6.52}\hullline{0.52}{-6.56}\hullline{0.61}{-6.76}\hullfill
\hullface{keel!57!black}\hullmove{-0.61}{-6.76}\hullline{-0.52}{-6.56}\hullline{-0.57}{-6.52}\hullline{-0.63}{-6.75}\hullfill
\hullface{keel!57!black}\hullmove{0.99}{-6.34}\hullline{0.88}{-6.13}\hullline{0.57}{-6.52}\hullline{0.63}{-6.75}\hullfill
\hullface{keel!63!black}\hullmove{-0.63}{-6.75}\hullline{-0.57}{-6.52}\hullline{-0.88}{-6.13}\hullline{-0.99}{-6.34}\hullfill
\hullface{keel!57!black}\hullmove{1.06}{-6.20}\hullline{0.90}{-6.09}\hullline{0.88}{-6.13}\hullline{0.99}{-6.34}\hullfill
\hullface{keel!72!black}\hullmove{-0.99}{-6.34}\hullline{-0.88}{-6.13}\hullline{-0.90}{-6.09}\hullline{-1.06}{-6.20}\hullfill
\hullface{keel!57!black}\hullmove{1.19}{-5.82}\hullline{1.03}{-5.63}\hullline{0.90}{-6.09}\hullline{1.06}{-6.20}\hullfill
\hullface{keel!78!black}\hullmove{-1.06}{-6.20}\hullline{-0.90}{-6.09}\hullline{-1.03}{-5.63}\hullline{-1.19}{-5.82}\hullfill
\hullface{keel!57!black}\hullmove{1.23}{-5.44}\hullline{1.04}{-5.44}\hullline{1.03}{-5.63}\hullline{1.19}{-5.82}\hullfill
\hullface{keel!87!black}\hullmove{-1.19}{-5.82}\hullline{-1.03}{-5.63}\hullline{-1.04}{-5.44}\hullline{-1.23}{-5.44}\hullfill
\hullface{keel!57!black}\hullmove{1.21}{-5.21}\hullline{0.99}{-5.03}\hullline{1.04}{-5.44}\hullline{1.23}{-5.44}\hullfill
\hullface{keel!93!black}\hullmove{-1.23}{-5.44}\hullline{-1.04}{-5.44}\hullline{-0.99}{-5.03}\hullline{-1.21}{-5.21}\hullfill
\hullface{keel!57!black}\hullmove{1.06}{-4.67}\hullline{0.90}{-4.78}\hullline{0.99}{-5.03}\hullline{1.21}{-5.21}\hullfill
\hullface{keel!96!black}\hullmove{-1.21}{-5.21}\hullline{-0.99}{-5.03}\hullline{-0.90}{-4.78}\hullline{-1.06}{-4.67}\hullfill
\hullface{keel!66!black}\hullmove{0.97}{-4.50}\hullline{0.58}{-4.35}\hullline{0.90}{-4.78}\hullline{1.06}{-4.67}\hullfill
\hullface{keel!92!white}\hullmove{-1.06}{-4.67}\hullline{-0.90}{-4.78}\hullline{-0.58}{-4.35}\hullline{-0.97}{-4.50}\hullfill
\hullface{keel!72!black}\hullmove{0.61}{-4.11}\hullline{0.52}{-4.31}\hullline{0.58}{-4.35}\hullline{0.97}{-4.50}\hullfill
\hullface{keel!89!white}\hullmove{-0.97}{-4.50}\hullline{-0.58}{-4.35}\hullline{-0.52}{-4.31}\hullline{-0.61}{-4.11}\hullfill
\hullface{keel!90!black}\hullmove{0.00}{-3.90}\hullline{0.00}{-4.13}\hullline{0.52}{-4.31}\hullline{0.61}{-4.11}\hullfill
\hullface{keel!92!white}\hullmove{-0.61}{-4.11}\hullline{-0.52}{-4.31}\hullline{0.00}{-4.13}\hullline{0.00}{-3.90}\hullfill
\hullface{topsides!57!black}\hullmove{0.42}{3.07}\hullline{0.30}{4.31}\hullline{0.00}{3.49}\hullline{0.00}{1.94}\hullfill
\hullface{topsides!90!black}\hullmove{0.00}{1.94}\hullline{0.00}{3.49}\hullline{-0.30}{4.31}\hullline{-0.42}{3.07}\hullfill
\hullface{topsides!57!black}\hullmove{1.12}{4.20}\hullline{0.79}{5.13}\hullline{0.30}{4.31}\hullline{0.42}{3.07}\hullfill
\hullface{topsides!87!black}\hullmove{-0.42}{3.07}\hullline{-0.30}{4.31}\hullline{-0.79}{5.13}\hullline{-1.12}{4.20}\hullfill
\hullface{topsides!60!black}\hullmove{1.97}{5.33}\hullline{1.40}{5.95}\hullline{0.79}{5.13}\hullline{1.12}{4.20}\hullfill
\hullface{topsides!87!black}\hullmove{-1.12}{4.20}\hullline{-0.79}{5.13}\hullline{-1.40}{5.95}\hullline{-1.97}{5.33}\hullfill
\hullface{topsides!60!black}\hullmove{2.95}{6.46}\hullline{2.09}{6.77}\hullline{1.40}{5.95}\hullline{1.97}{5.33}\hullfill
\hullface{topsides!84!black}\hullmove{-1.97}{5.33}\hullline{-1.40}{5.95}\hullline{-2.09}{6.77}\hullline{-2.95}{6.46}\hullfill
\hullface{topsides!60!black}\hullmove{4.03}{7.59}\hullline{2.85}{7.59}\hullline{2.09}{6.77}\hullline{2.95}{6.46}\hullfill
\hullface{topsides!84!black}\hullmove{-2.95}{6.46}\hullline{-2.09}{6.77}\hullline{-2.85}{7.59}\hullline{-4.03}{7.59}\hullfill
\hullface{keel!57!black}\hullmove{0.52}{-6.56}\hullline{0.38}{-6.26}\hullline{0.00}{-6.39}\hullline{0.00}{-6.74}\hullfill
\hullface{keel!57!black}\hullmove{0.00}{-6.74}\hullline{-0.00}{-6.39}\hullline{-0.38}{-6.26}\hullline{-0.52}{-6.56}\hullfill
\hullface{keel!57!black}\hullmove{0.57}{-6.52}\hullline{0.47}{-6.19}\hullline{0.38}{-6.26}\hullline{0.52}{-6.56}\hullfill
\hullface{keel!63!black}\hullmove{-0.52}{-6.56}\hullline{-0.38}{-6.26}\hullline{-0.47}{-6.19}\hullline{-0.57}{-6.52}\hullfill
\hullface{keel!57!black}\hullmove{0.88}{-6.13}\hullline{0.66}{-5.91}\hullline{0.47}{-6.19}\hullline{0.57}{-6.52}\hullfill
\hullface{keel!69!black}\hullmove{-0.57}{-6.52}\hullline{-0.47}{-6.19}\hullline{-0.66}{-5.91}\hullline{-0.88}{-6.13}\hullfill
\hullface{keel!57!black}\hullmove{0.90}{-6.09}\hullline{0.70}{-5.82}\hullline{0.66}{-5.91}\hullline{0.88}{-6.13}\hullfill
\hullface{keel!78!black}\hullmove{-0.88}{-6.13}\hullline{-0.66}{-5.91}\hullline{-0.70}{-5.82}\hullline{-0.90}{-6.09}\hullfill
\hullface{keel!57!black}\hullmove{1.03}{-5.63}\hullline{0.76}{-5.44}\hullline{0.70}{-5.82}\hullline{0.90}{-6.09}\hullfill
\hullface{keel!84!black}\hullmove{-0.90}{-6.09}\hullline{-0.70}{-5.82}\hullline{-0.76}{-5.44}\hullline{-1.03}{-5.63}\hullfill
\hullface{keel!57!black}\hullmove{1.04}{-5.44}\hullline{0.76}{-5.34}\hullline{0.76}{-5.44}\hullline{1.03}{-5.63}\hullfill
\hullface{keel!93!black}\hullmove{-1.03}{-5.63}\hullline{-0.76}{-5.44}\hullline{-0.76}{-5.34}\hullline{-1.04}{-5.44}\hullfill
\hullface{keel!57!black}\hullmove{0.99}{-5.03}\hullline{0.66}{-4.96}\hullline{0.76}{-5.34}\hullline{1.04}{-5.44}\hullfill
\hullface{keel!99!black}\hullmove{-1.04}{-5.44}\hullline{-0.76}{-5.34}\hullline{-0.66}{-4.96}\hullline{-0.99}{-5.03}\hullfill
\hullface{keel!63!black}\hullmove{0.90}{-4.78}\hullline{0.55}{-4.77}\hullline{0.66}{-4.96}\hullline{0.99}{-5.03}\hullfill
\hullface{keel!92!white}\hullmove{-0.99}{-5.03}\hullline{-0.66}{-4.96}\hullline{-0.55}{-4.77}\hullline{-0.90}{-4.78}\hullfill
\hullface{keel!72!black}\hullmove{0.58}{-4.35}\hullline{0.38}{-4.61}\hullline{0.55}{-4.77}\hullline{0.90}{-4.78}\hullfill
\hullface{keel!86!white}\hullmove{-0.90}{-4.78}\hullline{-0.55}{-4.77}\hullline{-0.38}{-4.61}\hullline{-0.58}{-4.35}\hullfill
\hullface{keel!93!black}\hullmove{0.52}{-4.31}\hullline{0.00}{-4.48}\hullline{0.38}{-4.61}\hullline{0.58}{-4.35}\hullfill
\hullface{keel!83!white}\hullmove{-0.58}{-4.35}\hullline{-0.38}{-4.61}\hullline{-0.00}{-4.48}\hullline{-0.52}{-4.31}\hullfill
\hullface{keel!75!black}\hullmove{0.00}{-4.13}\hullline{0.00}{-4.13}\hullline{0.00}{-4.48}\hullline{0.52}{-4.31}\hullfill
\hullface{keel!86!white}\hullmove{-0.52}{-4.31}\hullline{-0.00}{-4.48}\hullline{0.00}{-4.13}\hullline{0.00}{-4.13}\hullfill
\hullface{topsides!57!black}\hullmove{0.30}{4.31}\hullline{0.16}{5.55}\hullline{0.00}{5.04}\hullline{0.00}{3.49}\hullfill
\hullface{topsides!90!black}\hullmove{0.00}{3.49}\hullline{0.00}{5.04}\hullline{-0.16}{5.55}\hullline{-0.30}{4.31}\hullfill
\hullface{topsides!57!black}\hullmove{0.79}{5.13}\hullline{0.42}{6.06}\hullline{0.16}{5.55}\hullline{0.30}{4.31}\hullfill
\hullface{topsides!90!black}\hullmove{-0.30}{4.31}\hullline{-0.16}{5.55}\hullline{-0.42}{6.06}\hullline{-0.79}{5.13}\hullfill
\hullface{topsides!60!black}\hullmove{1.40}{5.95}\hullline{0.74}{6.57}\hullline{0.42}{6.06}\hullline{0.79}{5.13}\hullfill
\hullface{topsides!87!black}\hullmove{-0.79}{5.13}\hullline{-0.42}{6.06}\hullline{-0.74}{6.57}\hullline{-1.40}{5.95}\hullfill
\hullface{topsides!63!black}\hullmove{2.09}{6.77}\hullline{1.11}{7.08}\hullline{0.74}{6.57}\hullline{1.40}{5.95}\hullfill
\hullface{topsides!87!black}\hullmove{-1.40}{5.95}\hullline{-0.74}{6.57}\hullline{-1.11}{7.08}\hullline{-2.09}{6.77}\hullfill
\hullface{topsides!63!black}\hullmove{2.85}{7.59}\hullline{1.52}{7.59}\hullline{1.11}{7.08}\hullline{2.09}{6.77}\hullfill
\hullface{topsides!87!black}\hullmove{-2.09}{6.77}\hullline{-1.11}{7.08}\hullline{-1.52}{7.59}\hullline{-2.85}{7.59}\hullfill
\hullface{keel!63!black}\hullmove{0.38}{-6.26}\hullline{0.00}{-5.44}\hullline{0.00}{-5.44}\hullline{0.00}{-6.39}\hullfill
\hullface{keel!72!black}\hullmove{-0.00}{-6.39}\hullline{0.00}{-5.44}\hullline{-0.00}{-5.44}\hullline{-0.38}{-6.26}\hullfill
\hullface{keel!60!black}\hullmove{0.47}{-6.19}\hullline{0.00}{-5.44}\hullline{0.00}{-5.44}\hullline{0.38}{-6.26}\hullfill
\hullface{keel!81!black}\hullmove{-0.38}{-6.26}\hullline{-0.00}{-5.44}\hullline{-0.00}{-5.44}\hullline{-0.47}{-6.19}\hullfill
\hullface{keel!60!black}\hullmove{0.66}{-5.91}\hullline{0.00}{-5.44}\hullline{0.00}{-5.44}\hullline{0.47}{-6.19}\hullfill
\hullface{keel!87!black}\hullmove{-0.47}{-6.19}\hullline{-0.00}{-5.44}\hullline{-0.00}{-5.44}\hullline{-0.66}{-5.91}\hullfill
\hullface{keel!60!black}\hullmove{0.70}{-5.82}\hullline{0.00}{-5.44}\hullline{0.00}{-5.44}\hullline{0.66}{-5.91}\hullfill
\hullface{keel!93!black}\hullmove{-0.66}{-5.91}\hullline{-0.00}{-5.44}\hullline{-0.00}{-5.44}\hullline{-0.70}{-5.82}\hullfill
\hullface{keel!63!black}\hullmove{0.76}{-5.44}\hullline{0.00}{-5.44}\hullline{0.00}{-5.44}\hullline{0.70}{-5.82}\hullfill
\hullface{keel!99!black}\hullmove{-0.70}{-5.82}\hullline{-0.00}{-5.44}\hullline{-0.00}{-5.44}\hullline{-0.76}{-5.44}\hullfill
\hullface{keel!69!black}\hullmove{0.76}{-5.34}\hullline{0.00}{-5.44}\hullline{0.00}{-5.44}\hullline{0.76}{-5.44}\hullfill
\hullface{keel!92!white}\hullmove{-0.76}{-5.44}\hullline{-0.00}{-5.44}\hullline{0.00}{-5.44}\hullline{-0.76}{-5.34}\hullfill
\hullface{keel!72!black}\hullmove{0.66}{-4.96}\hullline{0.00}{-5.26}\hullline{0.00}{-5.44}\hullline{0.76}{-5.34}\hullfill
\hullface{keel!86!white}\hullmove{-0.76}{-5.34}\hullline{0.00}{-5.44}\hullline{0.00}{-5.26}\hullline{-0.66}{-4.96}\hullfill
\hullface{keel!69!black}\hullmove{0.55}{-4.77}\hullline{0.00}{-4.95}\hullline{0.00}{-5.26}\hullline{0.66}{-4.96}\hullfill
\hullface{keel!86!white}\hullmove{-0.66}{-4.96}\hullline{0.00}{-5.26}\hullline{0.00}{-4.95}\hullline{-0.55}{-4.77}\hullfill
\hullface{keel!66!black}\hullmove{0.38}{-4.61}\hullline{0.00}{-4.54}\hullline{0.00}{-4.95}\hullline{0.55}{-4.77}\hullfill
\hullface{keel!89!white}\hullmove{-0.55}{-4.77}\hullline{0.00}{-4.95}\hullline{0.00}{-4.54}\hullline{-0.38}{-4.61}\hullfill
\hullface{keel!57!black}\hullmove{0.00}{-4.48}\hullline{0.00}{-4.06}\hullline{0.00}{-4.54}\hullline{0.38}{-4.61}\hullfill
\hullface{keel!96!black}\hullmove{-0.38}{-4.61}\hullline{0.00}{-4.54}\hullline{0.00}{-4.06}\hullline{-0.00}{-4.48}\hullfill
\hullface{topsides!57!black}\hullmove{0.16}{5.55}\hullline{0.00}{6.79}\hullline{0.00}{6.59}\hullline{0.00}{5.04}\hullfill
\hullface{topsides!90!black}\hullmove{0.00}{5.04}\hullline{0.00}{6.59}\hullline{0.00}{6.79}\hullline{-0.16}{5.55}\hullfill
\hullface{topsides!57!black}\hullmove{0.42}{6.06}\hullline{0.00}{6.99}\hullline{0.00}{6.79}\hullline{0.16}{5.55}\hullfill
\hullface{topsides!93!black}\hullmove{-0.16}{5.55}\hullline{0.00}{6.79}\hullline{0.00}{6.99}\hullline{-0.42}{6.06}\hullfill
\hullface{topsides!63!black}\hullmove{0.74}{6.57}\hullline{0.00}{7.19}\hullline{0.00}{6.99}\hullline{0.42}{6.06}\hullfill
\hullface{topsides!93!black}\hullmove{-0.42}{6.06}\hullline{0.00}{6.99}\hullline{0.00}{7.19}\hullline{-0.74}{6.57}\hullfill
\hullface{topsides!69!black}\hullmove{1.11}{7.08}\hullline{0.00}{7.39}\hullline{0.00}{7.19}\hullline{0.74}{6.57}\hullfill
\hullface{topsides!93!black}\hullmove{-0.74}{6.57}\hullline{0.00}{7.19}\hullline{0.00}{7.39}\hullline{-1.11}{7.08}\hullfill
\hullface{topsides!72!black}\hullmove{1.52}{7.59}\hullline{0.00}{7.59}\hullline{0.00}{7.39}\hullline{1.11}{7.08}\hullfill
\hullface{topsides!96!black}\hullmove{-1.11}{7.08}\hullline{0.00}{7.39}\hullline{0.00}{7.59}\hullline{-1.52}{7.59}\hullfill
\end{scope}},
}

\tikzset{
  % the rudder from above, its axis 70 m aft of midship
  rudder above/.pic={\fill[keel] (-71,-0.5) rectangle (-69,0.5);},
  % a Flettner rotor from the side or ahead: its cylinder between two end plates
  rotor/.pic={
    \fill[rotors] (-2.89,0.4) rectangle (2.89,34.2);
    \fill[rotors] (0,34.2) ellipse (4.6 and 0.9);
    \fill[rotors] (0,0.9) ellipse (4.6 and 0.9);},
  % a Flettner rotor from above: its cylinder within its larger top end plate
  rotor above/.pic={
    \fill[rotors!45] (0,0) circle (4.6);
    \fill[rotors] (0,0) circle (2.89);},
  % a wing sail from the side, 40 m tall: one wing tapering to its rounded tip, its leading
  % edge forward and its flap aft of the hinge line
  wing/.pic={
    \path[save path=\wingoutline] (-7.05,0.4) -- (7.05,0.4) -- (6.4,35.4)
      .. controls (6.3,39) and (5,40.4) .. (2.6,40.4) -- (-3.8,40.4)
      .. controls (-4.9,40.4) and (-5.5,40) .. (-5.6,39.2) -- cycle;
    \fill[wings, use path=\wingoutline];
    \begin{scope}
      \clip[use path=\wingoutline];
      \fill[wings!55] (-8,0) -- (-3.2,0) -- (-2.3,41) -- (-8,41) -- cycle;
    \end{scope}
    \draw[wingstext] (-3.2,0.4) -- (-2.32,40.4);},
  % a wing sail from above: the aerofoil section with its flap, leading edge to the right
  wing above/.pic={
    \fill[wings] (7.05,0) .. controls (7.05,1.4) and (4,1.9) .. (1,1.6) -- (-2.8,0.9)
      -- (-2.8,-0.9) -- (1,-1.6) .. controls (4,-1.9) and (7.05,-1.4) .. cycle;
    \fill[wings!55] (-3.5,0.7) -- (-7.05,0.15) -- (-7.05,-0.15) -- (-3.5,-0.7) -- cycle;},
  % a soft kite: its canopy's arc, centred on the origin, and its bridle lines, hiding what
  % lies behind them, meeting at the control unit, from whose foot 15 m below the origin the
  % tether runs
  kite/.pic={
    \fill[white] (0,-10) -- ++(40:12) arc[start angle=40, end angle=140, radius=12]
      -- cycle;
    \fill[kite] (0,-10) ++(150:14) arc[start angle=150, end angle=30, radius=14]
      -- ++(30:-2) arc[start angle=30, end angle=150, radius=12] -- cycle;
    \draw[kite] (0,-10) ++(140:12) -- (0,-10) -- ++(40:12) (0,-10) ++(110:12) -- (0,-10)
      -- ++(70:12);
    \fill[keel] (-2,-15) rectangle (2,-10);},
  % the kite's telescopic launch mast, 35 m up from the deck at its foot, the origin, dashed
  % as it stands only for launch and landing, with the docking head that grips the kite on top
  launch mast/.pic={
    \draw[keel, densely dashed] (0,0) -- (0,33.4);
    \fill[keel] (-2,33.4) -- (2,33.4) -- (3.4,35.2) -- (-3.4,35.2) -- cycle;},
  % the ship from the side, 8 mm long, the size of an icon
  small ship/.pic={\pic[scale=0.055] {ship side};},
}
